% Complete inclusion-ready proof body; no omitted inputs.
% Requires amsmath, amssymb, amsthm, booktabs, hyperref, mathrsfs and proposition/theorem/lemma/corollary/definition/remark environments.
% Bibliography keys are in complete_bibliography.bib.
% Our prose and proofs only; external primary-source TeX is not embedded.
\providecommand{\GL}{\operatorname{GL}}
\providecommand{\End}{\operatorname{End}}
\section{Scope and the directed construction}
The starting objects are the original polynomial map and its retained
inverse-root cubic, followed by its two-sheet conductor and arithmetic
operator extension. The geometric-complexity route intended here is
Mulmuley's nonstandard Riemann-hypothesis programme over finite fields.
Its source-stated proposed steps are recorded at their actual scope.
All the algebraic maps proved in this reader are explicit. No off-critical
zero of the classical Riemann zeta function and no equality or separation
of Boolean complexity classes is established here.

The following complete calculations retain the constants in the original
map, including the coordinate relation $a=2s=-r^2$. The term
\emph{normalization ring} below denotes a specified integral-closure
construction together with its injection, conductor and cokernel;
the original singular ring remains in every diagram.

% Inclusion-ready fragment. Required packages: amsmath, amssymb, amsthm, hyperref.
% Bibliography keys are supplied by nonstandard_rh_sources.bib.
\section{The nonstandard Riemann-hypothesis bridge in the original GCT sources}
\label{sec:gct-nrh-source-bridge}

The bridge named by the user is explicitly present in Mulmuley's
\emph{On P vs. NP, Geometric Complexity Theory, and the Riemann Hypothesis},
arXiv:0908.1936v2~\cite{mulmuley2009gctrh}.  Its final section, headed ``How to prove PH1 and why
should it hold?'', proposes extending the geometric ingredients surrounding
the Riemann hypothesis over finite fields to nonstandard quantum groups.
The stated destination is a characteristic-zero complexity separation.
This is an actual research programme in the source.  In that source the
nonstandard Riemann hypotheses themselves are described as still unknown;
no nonstandard zeta function, Frobenius space, weight bound, or numbered
nonstandard-RH conjecture is specified there.  Consequently the reconstruction
below retains the source's proved constructions and its conjectured arrows
with their original status.  It does not replace the proposed finite-field
bridge by the classical hypothesis for the complex zeros of
$\zeta_{\mathbb Q}(s)$.

\subsection{Objects and arrows before any complexity claim}

In GCT VII, arXiv:0709.0749v2~\cite{mulmuley2008gct7}, Section~1, $H$ is a complex connected
classical reductive group, $\mu$ is a highest weight,
$X=V_\mu(H)$, $G=GL(X)$, and
\[
 \rho:H\longrightarrow G
\]
is the representation homomorphism.  The plethysm multiplicity is
\[
 a_{\lambda,\mu}^{\pi}
 =\dim_{\mathbb C}\operatorname{Hom}_{H}
       \bigl(V_\pi(H),\operatorname{Res}_{\rho}V_\lambda(G)\bigr).
\]
Here $\lambda$ is a highest weight of $G$ and $\pi$ a highest weight of
$H$.  This notation is that of GCT VII; the survey uses a different order
of indices in its plethysm discussion, which is not silently substituted
in these formulas.

The quantum parameter $q$ is independent of any prime-power cardinality.
GCT VII's Section~6, \texttt{srefined}, explicitly warns that a splitting
extension $K/\mathbb Q(q)$ may be needed.  Its general reciprocity statement
is not asserted over $\mathbb Q(q)$ without that extension, and Section~7.1
gives a failure over $\mathbb Q(q)$.  The following description therefore
retains the chosen coefficient field $K$, rather than erasing it.

Let $X_q=V_{q,\mu}(H_q)$ and let $\widehat R_{X,X}^{H}$ be the braiding
operator attached to the standard quantum group $H_q$.  GCT VII,
Section~2, splits $X_q\otimes_K X_q$ according to the signs of the
eigenvalues $\pm q^{a/2}$ and writes
\[
 I=P_{X,X}^{+,H}+P_{X,X}^{-,H}.
\]
With $u^i_j$ the entries of a variable $\dim(X)\times\dim(X)$ matrix,
its matrix bialgebra is the explicit quotient
\begin{equation}\label{eq:gct-nrh-frt-source}
 \mathcal O(M_q^H(X))
 =K\langle u^i_j\rangle\Big/
 \left\langle
 P_{X,X}^{+,H}(\mathbf u\otimes\mathbf u)
       -(\mathbf u\otimes\mathbf u)P_{X,X}^{+,H}
 \right\rangle.
\end{equation}
Replacing $P^{+,H}$ by $P^{-,H}=I-P^{+,H}$ gives the same ideal because
the corresponding matrix of relations changes by the scalar $-1$.
The source defines braided symmetric and exterior algebras with the
relations $P^{-,H}\mathbf x_1\mathbf x_2=0$ and
$P^{+,H}\mathbf x_1\mathbf x_2=0$, respectively.  Nonstandard minors
are the matrix coefficients of the resulting exterior-algebra coactions.
The top exterior power can vanish.  Thus the determinant localization
and an antipode cannot be presumed to exist in the general plethysm case.
The source's term ``quantum group'' includes this possible semigroup case.

Theorems labelled \texttt{tnonstd} and \texttt{tnonstdformal} in GCT VII
state the construction of the quantum map $H_q\to G_q^H$, its matrix
coordinate bialgebra map, cosemisimplicity, and a Peter--Weyl decomposition.
The coordinate-ring direction is
\[
 \mathcal O(M_q^H(X))\longrightarrow\mathcal O(M_q(V))
 \qquad(H=GL(V)),
\]
which is opposite to the displayed direction $H_q\to G_q^H$ of virtual
quantum spaces.  We use the concrete coordinate map, whose elementary
quotient construction is proved below, without importing an unproved
canonical basis from the existence of the bialgebra.

GCT VII's conjecture \texttt{cduality} proposes a decomposition
\begin{equation}\label{eq:gct-nrh-duality-source}
 X_q^{\otimes r}
   \cong\bigoplus_\alpha W_{q,\alpha}\otimes_K T_{q,\alpha}
\end{equation}
for the commuting nonstandard quantum-group and
$\mathcal B_r^H(q)$ actions.  Its conjecture \texttt{creciprocity}, refined
by \texttt{crecirefined}, proposes
\begin{equation}\label{eq:gct-nrh-reciprocity-source}
 V_{q,\lambda}^{H}=\bigoplus_\alpha
                 (W_{q,\alpha})^{\oplus m_\lambda^\alpha},
 \qquad
 V_{q,\lambda}^{H}\big|_{q=1}
       \cong\operatorname{Res}_{\rho}V_\lambda(G).
\end{equation}
The multiplicity $m_\lambda^\alpha$ counts the factors isomorphic to the
Specht module $S_\lambda$ in a composition series of
$T_{q,\alpha}(1)$ as a $\mathcal B_r^H(1)$-module.  A semisimple generic
algebra does not license replacing that composition series by an assumed
semisimple special fibre.  The source itself makes this point in
Section~6.  If $n_\pi^\alpha$ denotes the multiplicity of $V_{q,\pi}$ in
$W_{q,\alpha}$, the proposed reciprocity yields the source's equation
\begin{equation}\label{eq:gct-nrh-multiplicities-source}
 a_{\lambda,\mu}^{\pi}
       =\sum_\alpha m_\lambda^\alpha n_\pi^\alpha.
\end{equation}
This displayed equality is reported here as the consequence asserted by
that reciprocity conjecture, not as a newly proved general-plethysm result.
The elementary dimension calculation involved is explicit: on a direct
sum as in \eqref{eq:gct-nrh-reciprocity-source}, an $H_q$-linear map from
$V_{q,\pi}$ is exactly one such map to each summand, and projection and
inclusion supply inverse linear maps between the corresponding Hom spaces.
Each of the $m_\lambda^\alpha$ copies contributes
$n_\pi^\alpha$ dimensions.  What this calculation does not establish is
the asserted specialization isomorphism in
\eqref{eq:gct-nrh-reciprocity-source}.

\subsection{The positivity and effectivity steps retained separately}

GCT VIII, arXiv:0709.0751v2~\cite{mulmuley2008gct8}, presents algorithms whose correctness and
required properties are conjectured.  Its source labels include
\texttt{clocalcrystalb} (local crystal-basis existence),
\texttt{cuniquemindegleft} and \texttt{cunimindegright} (minimum degree),
\texttt{ccelldecompleft} and \texttt{ccelldecompright} (cell decomposition),
and \texttt{cposkroneckerleft}, \texttt{cposkroneckerright} (positivity).
The cells must recover the irreducible modules, and the subcells must
recover the restricted $H_q$-modules.  This is structural data beyond a
list of positive numbers.

In the Kronecker setting
\[
 H=GL(V)\times GL(W),\qquad X=V\otimes W,
\]
the left positivity conjecture in Section~2.4.2 does not assert that
every raw multiplication constant is a nonnegative polynomial.  It allows
the exact form
\begin{equation}\label{eq:gct-nrh-signed-positivity}
 \pm(q-q^{-1})^a f(q),\qquad a\in\mathbb Z_{\ge0},
\end{equation}
where $f$ is invariant under $q\leftrightarrow q^{-1}$ and has
nonnegative rational coefficients, together with the stated unimodality
requirement.  The factor $(q-q^{-1})^a$ and its sign are part of the
conjectured statement and cannot be suppressed at $q=1$.

The general-plethysm proposal is weaker and more complicated.  The
coefficients can be algebraic over $\mathbb Q(q)$, and Section~2.4.3
discusses coefficients of their minimal polynomials.  The conjecture
\texttt{cnonstdposleft0} uses positivity of a nonzero residual polynomial
at $q=1$ after explicitly retaining signs and powers of $q-q^{-1}$.
The conjecture \texttt{cnonstdposleft1} proposes a decomposition
$c(q)=c^0(q)+c^1(q)$ with a dominant positive unimodal term and a small
correction.  The source includes asymptotic language and a polynomial
bitlength bound; it does not provide a fully specified Frobenius complex
whose weights prove that decomposition.

Section~6 of GCT VIII then adds computational requirements on the labels
of the basis and its cells.  In the standard setting, efficiently
computable tableau/path labels and crystal operations yield counting
rules even when constructing all basis vectors is expensive.  The
nonstandard analogues are proposed there as additional properties whose
full specification is deferred.  Thus coefficient positivity alone
does not provide a uniform polynomial-time algorithm or a $\#P$ formula.
For example, from a predicate $R(x,y)$ and a bound $|y|\le p(|x|)$, both
explicitly supplied, one obtains a $\#P$ count by nondeterministically
choosing a padded witness and checking $R$.  Without a witness encoding,
polynomial bound, and polynomial-time checking procedure, the mere
nonnegative integer value of a multiplicity supplies no such machine.

The survey's final section next discusses Plethysm PH0, the existence
of compatible positive bases for $\rho:H\to GL(V_\mu(H))$, and
Plethysm PH/PH1, positive polyhedral formulas and their effective
versions.  The intended further step is to quantize the triple
\[
 \overline H_{\widehat h}\longrightarrow
 \overline H=SL(\overline X)\longrightarrow L=GL(\overline W)
\]
associated with a class variety.  It is not a proved automatic passage
from a single plethysm pair to all class-variety multiplicities.
The survey's general PH1, labelled \texttt{hph1restatement}, requires
bounds polynomial in $\langle\lambda\rangle,n,\langle m\rangle$,
including the bitlength of $m$, for the padded polynomial
$\overline h(\overline X)=z^{m-n}h(X)$.

Finally the displayed source theorem \texttt{tgct6refined} also requires
an obstruction hypothesis OH.  It takes
$m=2^{\log^a n}$, with $a>1$ fixed, and asks for labels $\lambda$ with
$P_{\lambda,n,m}(k)\ne\varnothing$ for all sufficiently large $k$
but $Q_\lambda=\varnothing$.  This separates a positive formula from
the additional proof that the required obstruction actually exists.
Its displayed target is the general $\#P$ versus $NC$ problem in
characteristic zero; the source routes the analogous $P$ versus $NP$
formulation to GCT VI.  This record does not substitute a Boolean uniform
$P$ versus $NP$ theorem for that stated algebraic target.

The primary-source chain therefore has the following exact status:
\begin{enumerate}
\item The standard geometric positivity story uses the finite-field
purity results surrounding Deligne's work and the canonical bases of
standard quantum groups.  This is the source's established model.
\item The matrix bialgebra quantization is explicitly constructed.
\item General nonstandard canonical bases, their required cells,
positivity, reciprocity and effective labels are proposed additional
statements.  They are not all consequences of the existence of the
matrix bialgebra.
\item The nonstandard finite-field purity/RH theory is a proposed way to
prove those statements; the cited works do not state it as a complete
zeta-function hypothesis.
\item Effective positive formulas must still be connected to the class
varieties and to actual obstruction existence, with the declared size
parameters, before the intended complexity separation follows.
\end{enumerate}
These labels describe the literature's proposed programme, and introduce
no conditional theorem as a replacement for a missing proof in the
present work.

\subsection{A fully explicit FRT coordinate morphism}

Here is the elementary algebraic construction that underlies
\eqref{eq:gct-nrh-frt-source}.  It is included with a full proof so that
the quantum-space arrow has a concrete coordinate-ring meaning.
Let $K$ be a field, $E=K^d$ with its ordered basis $e_1,\ldots,e_d$, and
$P\in\operatorname{End}_K(E\otimes E)$.  Put
\[
 T=K\langle u_{ij}:1\le i,j\le d\rangle,
 \quad W_{(i,k),(j,l)}=u_{ij}u_{kl}.
\]
\[
 \quad R=PW-WP,
 \quad A_P=T/(R_{ab}:a,b\in\{1,\ldots,d\}^2).
\]
Here $(R_{ab})$ denotes the two-sided ideal generated by the displayed
entries; no commutation between different $u_{ij}$ is presumed.
On the free algebra define
\[
 \Delta(u_{ij})=\sum_{h=1}^{d}u_{ih}\otimes u_{hj},
 \qquad \epsilon(u_{ij})=\delta_{ij}.
\]
Both extend uniquely as algebra homomorphisms.  Their coassociativity
and counit identities hold on a generator because both iterated
coproducts equal
$\sum_{h,k}u_{ih}\otimes u_{hk}\otimes u_{kj}$, and contracting either
factor with $\epsilon$ gives $u_{ij}$.  The same identities then hold on
each word by multiplicativity and on all of $T$ by linearity.

For indices $a,b,c$ in $\{1,\ldots,d\}^2$,
\[
 \Delta(W_{ab})=\sum_c W_{ac}\otimes W_{cb}.
\]
Consequently, with each occurrence of the scalar matrix $P$ acting on
these pair-indices,
\[
 \Delta(R_{ab})
 =\sum_c R_{ac}\otimes W_{cb}
       +\sum_c W_{ac}\otimes R_{cb},
 \qquad \epsilon(R_{ab})=0.
\]
The two middle terms cancel entry by entry.  The ideal of relations is
therefore a coideal and lies in the kernel of the counit.  It follows
that these formulas define a bialgebra structure on $A_P$.

Now let $B$ be a bialgebra with a right comodule $E$ written
\[
 \delta_B(e_j)=\sum_i e_i\otimes y_{ij}.
\]
The comodule identities are exactly
$\Delta_B(y_{ij})=\sum_h y_{ih}\otimes y_{hj}$ and
$\epsilon_B(y_{ij})=\delta_{ij}$.  When $P$ is a comodule endomorphism
of $E\otimes E$, applying $\delta_B P=(P\otimes1)\delta_B$ to every
basis vector gives
$P(y\otimes y)=(y\otimes y)P$ entry by entry.  Thus the explicit map
\begin{equation}\label{eq:gct-nrh-universal-map}
 \Phi_P:A_P\longrightarrow B,\qquad
       [u_{ij}]\longmapsto y_{ij}
\end{equation}
is well-defined.  It preserves coproduct and counit on generators and
hence everywhere.  It is the unique bialgebra map with those specified
generator images.  Its image is the subbialgebra generated by the
$y_{ij}$; its exact kernel is
\[
 \{[F(u)]\in A_P:F(y)=0\text{ in }B\}.
\]
This does not claim that the map is injective, nor that its image is all
of $B$.

For comparison at the classical symmetric projector, assume
$\operatorname{char}(K)\ne2$ and take $P=(I+\tau)/2$, where
$\tau(e_i\otimes e_k)=e_k\otimes e_i$.  Then
\[
 2R_{(i,k),(j,l)}
   =u_{kj}u_{il}-u_{il}u_{kj}.
\]
As the four indices vary these are exactly all pairwise commutators of
the matrix variables.  Sending $[u_{ij}]$ to the corresponding commuting
variable defines an isomorphism
\[
 A_{(I+\tau)/2}\longrightarrow K[u_{ij}:1\le i,j\le d],
\]
whose inverse sends a commuting variable to $[u_{ij}]$.  The commutator
relations prove that the inverse is defined, and the two compositions
fix every generator.  This proves the comparison map with its field,
size and characteristic restriction retained.

\subsection{Specialization and the information in its kernel}

Let $D=\mathbb Q[q,q^{-1}]_{(q-1)}$, $t=q-1$, $L$ a finite free
$D$-module, $J\subseteq L$ a submodule, and $A=L/J$.  These are explicit
algebraic data rather than a claim that a particular GCT conjecture
holds.  Let $A[t]=\{a\in A:ta=0\}$.  There is an exact sequence
\begin{equation}\label{eq:gct-nrh-specialization-exact}
 0\longrightarrow A[t]\xrightarrow{\eta}J/tJ
 \longrightarrow L/tL\longrightarrow A/tA\longrightarrow0,
 \qquad \eta([v])=[tv].
\end{equation}
Indeed $ta=0$ means that a lift $v\in L$ satisfies $tv\in J$.
Changing the lift by $j\in J$ changes $tv$ by $tj$, so $\eta$ is
well-defined.  If $tv\in tJ$, then $tv=tj$ for some $j\in J$; multiplication
by $t$ is injective on the free module $L$, so $v=j$ and $[v]=0$.
Thus $\eta$ is injective.  An element $j\in J$ maps to zero in $L/tL$
exactly when $j=tv$ for some $v\in L$, which is exactly the image of
$\eta$.  The kernel of $L/tL\to A/tA$ consists of the classes of
$J+tL$, and the last map is surjective by lifting a representative from
$A$ to $L$.  These observations prove every assertion of exactness.

Put
\[
 J^{\mathrm{sat}}=
  \{v\in L:t^e v\in J\text{ for some }e\in\mathbb Z_{\ge0}\}.
\]
The map $A\to A\otimes_D\mathbb Q(q)$ has kernel
$J^{\mathrm{sat}}/J$.  To prove this, localization sends the class of
$v$ to zero precisely when some nonzero element of $D$ multiplies $v$
into $J$.  Every nonzero element of $D$ is a power of $t$ times a unit:
factor its numerator by its exact order of vanishing at $q=1$ and note
that its denominator and remaining numerator are units in this local
ring.  Multiplying or dividing by that unit does not change membership
in $J$.  This proves the kernel formula.  Passing to a generic fibre
therefore removes an explicitly identified module, not an unspecified
``small error''.  Formula \eqref{eq:gct-nrh-specialization-exact} records
the first special-fibre layer that a proposed reciprocity construction
must preserve in the data to which it is applied.

\subsection{A source-specific coefficient and its retained double zero}

The final GCT IV, by Blasiak, Mulmuley and Sohoni,
arXiv:cs/0703110v4~\cite{blasiakmulmuleysohoni2013gct4}, proves nonstandard Schur--Weyl duality in Section~12
(source label \texttt{t nonstandard schur-weyl duality}).  Its displayed
statement uses $\check X^{\otimes r}$, and the immediately following
remark explicitly records the dual convention needed by its proof.
The source then constructs a global crystal-basis rule for the
two-row problem, $\dim V=\dim W=2$, with
\[
 g_{\lambda\mu\nu}
  =\#\{\text{highest-weight elements of }{+}\operatorname{HNSTC}(\nu)
          \text{ of weight }(\lambda,\mu)\}.
\]
This is a proved bounded case in the source, not a claim about all
dimensions.  Theorem \texttt{t positivity} gives sign-coherent
coefficients for individual Chevalley generators.  Example
\texttt{ex not positive} displays a divided-power expression with mixed
signs. The literal printed target tableau has a weight inconsistency.
Section~\ref{sec:gct-generator-interval} proves that obstruction, constructs
the corrected target $Q$ from the actual action rules, and proves that
its divided-power matrix entry is the displayed Laurent polynomial.
The printed and corrected targets remain separately specified objects.

We now independently prove all algebraic identities used from that
example.  The example takes column-count shape
$\nu=[0,n_3,n_2,2]$ with $n_3=3,n_2=2$; in ordinary row-length notation
this is the same partition $(7,5,3)$, since there are three columns of
height three, two of height two and two of height one.  Its size is
$3\cdot3+2\cdot2+2=15$.  We retain the original column-count shape.
For $n\ge1$ set
\[
 [n]_q=\sum_{j=0}^{n-1}q^{n-1-2j}
 \quad\text{in }\mathbb Z[q,q^{-1}].
\]
Multiplication gives
$(q-q^{-1})[n]_q=q^n-q^{-n}$ by cancellation of explicitly adjacent
terms; the finite sum also defines its values at $q=1$ and $q=-1$.
The source displays the following rational expression and Laurent polynomial:
\begin{equation}\label{eq:gct-nrh-coefficient-original}
 C(q)=\frac{[6]_q\bigl([7]_q-[3]_q\bigr)-[3]_q[8]_q}{[2]_q}
      =q^{10}-q^4-q^{-4}+q^{-10}.
\end{equation}
The first expression lives in $\mathbb Q(q)$ and the second is its
Laurent-polynomial representative; multiplying the equality by
$[2]_q=q+q^{-1}$ verifies it without imposing any evaluation at a zero
of $[2]_q$.

The original Laurent polynomial, including all four signs, satisfies
\begin{align}
 C(q)
  &=(q^7-q^{-7})(q^3-q^{-3})\notag\\
  &=(q-q^{-1})^2[7]_q[3]_q.\label{eq:gct-nrh-coefficient-factors}
\end{align}
The first equality follows by multiplying its two binomials, giving
the four exponents $10,4,-4,-10$ with coefficients $1,-1,-1,1$.
The second follows from the preceding finite-sum identity for $n=7,3$.
Moreover
\[
 [7]_q[3]_q
 =q^8+2q^6+3q^4+3q^2+3+3q^{-2}+3q^{-4}+2q^{-6}+q^{-8}.
\]
The displayed coefficients count the ordered pairs of summands with
each exponent.  This proves nonnegativity of this retained factor,
while $C$ itself has the mixed signs displayed in
\eqref{eq:gct-nrh-coefficient-original}.  No assertion about every
positivity requirement of GCT VIII follows from this single example.

Since
\[
 q-q^{-1}=\frac{(q-1)(q+1)}q,
\]
equation \eqref{eq:gct-nrh-coefficient-factors} gives the equality in
$D$ with all factors retained:
\[
 C(q)=(q-1)^2\frac{(q+1)^2}{q^2}[7]_q[3]_q,
 \qquad
 \left.\frac{(q+1)^2}{q^2}[7]_q[3]_q\right|_{q=1}=4\cdot7\cdot3=84.
\]
Thus $C(1)=C'(1)=0$, $C''(1)=168$, and its exact order at $q=1$ is two.
Specializing $C$ to its value $0$ would lose the nonzero retained second
jet $84(q-1)^2$ in $D/(q-1)^3$.

There is also an explicit two-sheet presentation.  Here $p$ is a new trace
coordinate, distinct from the source polynomial map's retained coordinate
$w$. Define
\[
 B=\mathbb Q[p,Q]/(Q^2-pQ+1),\qquad
 \theta:B\longrightarrow\mathbb Q[q,q^{-1}],
 \quad Q\longmapsto q,\quad p\longmapsto q+q^{-1}.
\]
In $B$, $Q(p-Q)=1$.  Hence the inverse homomorphism sends $q$ to $Q$
and $q^{-1}$ to $p-Q$.  Both compositions fix the displayed generators,
proving that $\theta$ is an isomorphism.  The ring $B$ is free of rank
two over $\mathbb Q[p]$, with ordered basis $(1,Q)$: division by the
monic relation produces a unique representative $a(p)+b(p)Q$, and a
nonzero multiple of a monic degree-two polynomial cannot have smaller
$Q$-degree.  The deck involution sends $Q$ to $p-Q$; its image under
$\theta$ is $q\mapsto q^{-1}$.

Put $r_{\mathrm{GCT}}=2Q-p$ in $B$.  Then
\[
 r_{\mathrm{GCT}}^2=p^2-4.
\]
The associated conic presentation is the isomorphism
\[
 \mathbb Q[p,r_{\mathrm{GCT}}]/(p^2-r_{\mathrm{GCT}}^2-4)
  \longrightarrow\mathbb Q[q,q^{-1}],
 \quad (p,r_{\mathrm{GCT}})\longmapsto
        (q+q^{-1},q-q^{-1}),
\]
whose inverse sends $q$ to $(p+r_{\mathrm{GCT}})/2$ and $q^{-1}$ to
$(p-r_{\mathrm{GCT}})/2$.  Their product is $1$ by the conic equation.
The use of $1/2$ is explicit; these displayed inverses are over
$\mathbb Q$ and are not asserted in characteristic two.

Writing $p=q+q^{-1}$ and applying the recurrence
$[n+1]_q=p[n]_q-[n-1]_q$ gives
\[
 [3]_q=p^2-1,\qquad [7]_q=p^6-5p^4+6p^2-1.
\]
The recurrence is verified directly by multiplying the two finite
sums and cancelling the shared interior terms.  Thus the same retained
coefficient is
\begin{align*}
 C(q)
  &=(p^2-4)(p^6-5p^4+6p^2-1)(p^2-1)\\
  &=p^{10}-10p^8+35p^6-51p^4+29p^2-4.
\end{align*}
At $p=2$ the factor multiplying $p-2$ has value
$4\cdot7\cdot3=84$, whereas
\[
 p-2=\frac{(q-1)^2}{q}.
\]
This exact ramification equation explains why a first zero in the trace
coordinate $p$ is the retained second zero in the coordinate $q$.
The discriminant of $Q^2-pQ+1$ is $p^2-4$; the two excluded branch
values for the etale double cover are $p=2$ and $p=-2$.
Indeed the relative differential module is
$\Omega_{B/\mathbb Q[p]}=B/(2Q-p)\,dQ$, by differentiating the single
monic relation with $p$ held fixed.  Since $(2Q-p)^2=p^2-4$, this module
vanishes after inverting $p^2-4$.  The algebra is already free of rank
two over the base, so this localization is finite etale.  At $p=2$ the
fibre is $\mathbb Q[Q]/(Q-1)^2$, and at $p=-2$ it is
$\mathbb Q[Q]/(Q+1)^2$, proving that both exceptional fibres retain a
nonzero square-zero element.  Further, if $a(p)+b(p)Q$ is invariant under
$Q\mapsto p-Q$, then $b(p)(2Q-p)=0$.  The proved embedding into the
Laurent-polynomial domain gives $b=0$, so the invariant subring is
exactly $\mathbb Q[p]$.
These identities give a concrete, fully calculated nonstandard-GCT
coefficient and its rank-two coordinate algebra for any subsequent
arithmetic transport to act on.  They prove no general canonical-basis
conjecture and no complexity separation.

\section{A source-correct GCT-IV generator interval and its integral endpoint}
\label{sec:gct-generator-interval}

Put $A=\mathbb Z[q,q^{-1}]$, $K=\mathbb Q(q)$ and
$[m]=q^{m-1}+q^{m-3}+\cdots+q^{1-m}$ for $m\geq1$, with $[0]=0$.
Write $p=[2]=q+q^{-1}$ and $r=q-q^{-1}$ throughout this section.
These retain the source parameter (called $u$ in its TeX), rather than
specializing it in the generator calculation. A column is denoted by
its entries in parentheses; concatenation means the source's projected
upper based product, not an ordinary commutative monomial product.
In particular $(23)$ and $(32)$ are distinct columns.

The source is the April 2013 version of GCT IV,
\emph{Geometric complexity theory IV: nonstandard quantum group for the
Kronecker problem}, arXiv:cs/0703110v4. Its literal example labelled
\texttt{ex not positive} and equations labelled \texttt{e aFV computation}
must be read together with the column dictionary, the contextual
nonintegral straightening rule, and the Chevalley action rule. The
precise source bytes and line intervals are recorded in
\texttt{source\_receipts.json}; this proof retains the printed target
as a separate object.

\subsection{Objects, weights, and the printed endpoint}

Let $X=\check X_{(7,5,3)}$ be the original source quotient
$\check Y_{(3,3,3,2,2,1,1)}/\check Y_{\triangleright(3,3,3,2,2,1,1)}$.
Its source basis is the set of positive honest nonstandard tabloid
classes. Its integral form is the $A$-span of those classes. We use
the following actual vectors, including the source's rescaling:
\begin{align*}
 T&=(123)(123)(124)(12)(12)(1)(1),\\
 T_1&=(123)(123)(124)(23)(12)(1)(1),\\
 T_2&=-p^{-1}(123)(123)(123)(12)(12)(4)(1),\\
 P_{\rm printed}&=-p^{-1}(123)(123)(124)(23)(12)(4)(1),\\
 Q&=-p^{-1}(123)(123)(124)(13)(12)(4)(1).
\end{align*}
Thus $\nu=[0,3,2,2]$ in the source's column-count notation and
$\nu=(7,5,3)$ in row-length notation. There are fifteen boxes.
The alphabet dictionary is
$1=(1,1)$, $2=(1,2)$, $3=(2,1)$, $4=(2,2)$ for the $V,W$ letters.
Counting each letter, without reordering columns, gives
\[
\begin{array}{c|cc}
 &V\text{-weight}&W\text{-weight}\\\hline
T&(12,3)&(9,6)\\
T_1,T_2&(11,4)&(9,6)\\
P_{\rm printed}&(10,5)&(8,7)\\
Q&(10,5)&(9,6)
\end{array}
\]
The Chevalley operator $F_V$ changes the $V$-weight by $(-1,1)$ and
preserves both $W$-weight coordinates. Since the basis consists of
weight vectors and $q$ is an indeterminate, its distinct $W$-weight
spaces form a direct sum. Therefore the coefficient of
$P_{\rm printed}$ in $F_V^{(2)}T$ is zero. This proves a discrepancy
in the literal displayed endpoint. It does not infer a corrected
endpoint from the scalar formula. The direct calculation below
constructs $Q$ and its coefficient.

\subsection{Contextual straightening and every two-step path}

The selected column words in the source dictionary are
\[
\begin{array}{c|cc@{\qquad}c|cc}
\text{column}&V&W&\text{column}&V&W\\\hline
1&1&1&3&2&1\\4&2&2&12&11&21\\
13&21&11&23&21&12\\32&12&21&123&211&121\\
124&211&221&134&212&121\\234&212&221&&&
\end{array}
\]
These are diagram labels: they are not claims that the corresponding
nonstandard column equals just one canonical tensor word. For example
the column $(23)$ includes the source's additional $p^{-1}$ term.
The source's sign convention takes $(-p^{-1})^{\deg}$ times an
honest representative with no $3$--$2$ arc, or with its $3$--$2$
arc on the $W$ side, as positive. The displayed $T,T_1$ have degree
zero, while $T_2,Q$ have degree one. The diagram of $Q$ has its
$3$--$2$ arc on the $W$ side. Thus the displayed rescalings have
the source's positive orientation. The diagram classification places
$T,T_1,T_2,Q$ respectively in the pure $3$--$1$ $W$, $3$--$2$--$1$
$W$, no-external-arc, and pure $3$--$2$ $W$ cases after deleting
invariants; these are honest classes in the source basis.
Pair each $2$ with the next available $1$ to its right, recursively
matching $21$ pairs; equivalently scan left to right using a stack
of unmatched $2$ positions. The source's generator rule is
\[
 F_VD=\sum_{j=1}^{\varphi_V(D)}[j]\mathcal F_{(j)V}(D),
\]
where $\varphi_V$ is the number of unpaired $V$-ones. The $j$th term
changes the column containing that one by its column Kashiwara
operator, unless the changed word has an extra internal arc, in
which case the term is zero. The changes needed below are
$123\to134$, $124\to234$, $12\to32$, $1\to3$.
For an unpaired $(12)$ column, choosing the first of its two
$V$-ones creates an extra internal arc and gives zero; choosing
the second gives $(32)$.

For clarity we specify the contextual nonintegral rule used here.
For a consecutive pair of equal-height columns in one of the source
nonintegral strings, write the entire based product as $BCD$.
Let $I$ be the corresponding equal-height invariant, and let
$\omega$ be the unpaired word of $C$ on the indicated side. The
source's contextual projection identity, after taking the quotient,
is
\[
 BCD=-p^{-1}\begin{cases}B'ID&\omega=11,\\
 BID&\omega=12,\\BID'&\omega=22.\end{cases}
\]
In the $11$ instances below, $B'$ is obtained by changing the
unpaired $2$ in $B$ paired to the first local $1$ to a $1$; if
there is no such $2$, $B'=0$. The $12$ instance leaves both outside
words unchanged. These identities follow by projecting the source
formula $B\mathbin{\heartsuit_g}\widetilde C\mathbin{\heartsuit_g}D
=BCD+p^{-1}B'ID$ (and its $12$ version) to the quotient in which
the left-hand side is zero. This is the source corollary labelled
\texttt{c projecting to heartproj basis Kronecker}; the distinction
between the two based products in that formula is essential.

The pairs used here are $(134)(123)$, with $W$-word $11$ and
height-three invariant $(234)(123)$; $(134)(124)$ with $W$-word
$12$ and the same invariant; $(32)(12)$ with $V$-word $11$ and
height-two invariant $(34)(12)$; and $(3)(1)$ with $W$-word $11$
and height-one invariant $(4)(1)$. No isolated relation such as
$(3)(1)=0$ is substituted inside a longer based product.

The integral relations needed in the same quotient are
\begin{gather*}
 (234)(12)=-(124)(23),\qquad (134)(12)=-(124)(13),\\
 (32)(1)=(23)(1),\qquad (12)(23)=(23)(12).
\end{gather*}
They are the source's height-$(3,2)$, $(2,1)$ and $(2,2)$
degree-preserving relations. Their contextual products remain
relations: for integral relations the source's two based products
agree. Thus there is no suppressed nonintegral correction in these
four replacements.

In $T$, the prefix before its last two columns has exactly one
unpaired $W$-two, in column $(124)$. It is paired with the first
of the two suffix $W$-ones. The contextual rule therefore gives
\[
 (123)(123)(124)(12)(12)(3)(1)
 =-p^{-1}(123)(123)(123)(12)(12)(4)(1)=T_2.
\]
For $T_1$ the analogous unpaired $W$-two is in column $(23)$,
and changing it to a $1$ changes that column to $(13)$. Hence
\[
 (123)(123)(124)(23)(12)(3)(1)
 =-p^{-1}(123)(123)(124)(13)(12)(4)(1)=Q.
\]
This constructs the endpoint and its minus sign directly. The second
calculation is source case (F.1-1.5): its hypothesis holds because
$T_1$ has a $3$--$2$--$1$ $W$-arc and no pure $3$--$1$ $W$-arc.

We record how invariant counts exclude paths. An equal-height
invariant pair has all of its letters paired internally or within
the pair on both diagrams, with no exiting arc. An operator
$\mathcal F_{(j)V}$ changes an unpaired letter, so it cannot change
such a pair. Degree-preserving straightening moves preserve the
invariant record. In a nonintegral correction, the added pair is
an invariant and the contextual change is at an outside unpaired
letter; it therefore leaves pre-existing pairs intact. Integral
graded alternatives give zero. These facts prove, step by step,
that nonzero generator coefficients cannot decrease any of the
height-specific invariant counts. They also prove the direction
of containment used here independently of the reversed containment
sentence in source line 5007. In particular a vector with a
height-three or height-two invariant cannot reach $Q$, whose only
invariant is its final $(4)(1)$.

For $T$ the ordered unpaired $V$-one positions (column, word position)
are
\[
 (1,3),(2,3),(3,3),(4,1),(4,2),(5,1),(5,2),(6,1),(7,1).
\]
Put
\[
 H=-p^{-1}(123)(234)(123)(12)(12)(1)(1),\qquad
 R_0=(123)(123)(124)(12)(12)(1)(3).
\]
The following table lists all nine terms before multiplication by
their quantum integer $[j]$:
\[
\begin{array}{c|c|l}
j&\mathcal F_{(j)V}T&\text{verification}\\\hline
1&0&(134)(123),\ W:11,\ B'=0\\
2&H&(134)(124),\ W:12,\ \text{context unchanged}\\
3&-T_1&(234)(12)=-(124)(23)\\
4&0&\text{extra internal arc in column }4\\
5&0&(32)(12),\ V:11,\ B'=0\\
6&0&\text{extra internal arc in column }5\\
7&T_1&(32)(1)=(23)(1),\ \text{then commute }12,23\\
8&T_2&\text{the first contextual suffix calculation above}\\
9&R_0&\text{last unpaired one, the crystal term}
\end{array}
\]
All displayed zero-prefix assertions follow by the stack pairing:
for $j=1$ the prefix is empty; for $j=5$ the prefix has no unpaired
$V$-two. The vector $H$ has a height-three invariant, so none of its
nonzero successors can equal $Q$. This proves
$F_VT=([7]-[3])T_1+[8]T_2+[2]H+[9]R_0$, and only the last term
besides $T_1,T_2$ still requires a possible-return check.

There are eight unpaired $V$-ones in $R_0$, the first eight
positions in the preceding list; its last $3$ supplies an unpaired
$V$-two. Its $W$-diagram equals that of $T$. Put
\[
 H_R=-p^{-1}(123)(234)(123)(12)(12)(1)(3),\qquad
 R_1=(123)(123)(124)(23)(12)(1)(3),
\]
and $R_2=(123)(123)(124)(12)(12)(3)(3)$. Every term is
\[
\begin{array}{c|rrrrrrrr}
j&1&2&3&4&5&6&7&8\\\hline
\mathcal F_{(j)V}R_0&0&H_R&-R_1&0&0&0&R_1&R_2.
\end{array}
\]
The same local relations and zero-prefix calculations apply since
the last column lies outside them. The class $H_R$ has a height-three
invariant. The honest class $R_1$ has no invariant and one unpaired
$V$-two; $R_2$ has no invariant and two unpaired $V$-twos. The source
diagram classification identifies these honest classes, so neither
is $Q$, which has invariant record $(0,0,0,1)$ and no unpaired
$V$-two. This direct enumeration proves that the coefficient of
$Q$ in $F_VR_0$ is zero; no general monotonicity assertion about
unpaired twos is needed.

The ordered positions for $T_1$ are
$(1,3),(2,3),(3,3),(5,1),(5,2),(6,1),(7,1)$.
Its seven terms, with the source cases made explicit, are
\[
\begin{array}{c|l}
j&\text{outcome and reason it can or cannot equal }Q\\\hline
1&0:\ (134)(123),\ W:11,\ \text{empty prefix}\\
2&\text{a height-three invariant, case (F.3-3.1)}\\
3&\text{a height-two invariant, case (F.3-2)}\\
4&0:\ \text{extra internal arc in column }5\\
5&\text{a height-two invariant, case (F.2-1)}\\
6&Q:\ \text{the second contextual suffix calculation above}\\
7&R_1:\ \text{the crystal term, invariant count }i_1=0.
\end{array}
\]
The stated invariant cases can also be checked with the finite
relations: for $j=2$, apply the $W:12$ replacement
$(134)(124)\mapsto-p^{-1}(234)(123)$. For $j=3$, first use
$(234)(23)=-(234)(32)$ and then the contextual $V:11$ replacement
of $(32)(12)$, whose preceding unpaired $V$-two is in $(234)$;
it raises that column to $(124)$ and creates $(34)(12)$.
For $j=5$, use $(32)(1)=(23)(1)$ and then the $W:12$ replacement
$(23)(23)\mapsto-p^{-1}(24)(13)$, which creates a height-two
invariant. The extra integral relation here is the source's
$(3,2)$ bottom-weight relation, and the latter contextual relation
is its height-two nonintegral string. These computations account
for all indices and prove $a^{F_V}_{Q,T_1}=[6]$.

For $T_2$ the seven positions are
$(1,3),(2,3),(3,3),(4,1),(4,2),(5,1),(5,2)$; its existing
$(4)(1)$ invariant contributes none. Terms $j=1,2$ vanish because
their $(134)(123)$ $W:11$ replacements have no preceding unpaired
$W$-two. Terms $j=4,6$ vanish by extra internal arcs. Term $j=5$
vanishes by the $V:11$ $(32)(12)$ relation with no preceding
unpaired $V$-two. Term $j=7$ is the crystal vector
$-p^{-1}(123)(123)(123)(12)(32)(4)(1)$, with one unpaired
$V$-two; after removing its invariant its diagram has no external
arc, whereas $Q$ has a pure $3$--$2$ $W$-arc and no unpaired
$V$-two. Finally the exact $j=3$ calculation is
\[
 -p^{-1}(123)(123)(134)(12)(12)(4)(1)
 =p^{-1}(123)(123)(124)(13)(12)(4)(1)=-Q.
\]
Thus $a^{F_V}_{Q,T_2}=-[3]$. Every possible first step and every
remaining second step has now been accounted for. Therefore precisely
$T_1$ and $T_2$ contribute to the two-step path from $T$ to $Q$,
with four edge coefficients $[7]-[3]$, $[8]$, $[6]$, and $-[3]$.


\subsection{An actual negative-Borel subquotient}

Here $B^-_K$ denotes the algebra generated by $F_V,F_W$, invertible
$K_V,K_W$, with commuting Cartan generators,
$K_VF_VK_V^{-1}=q^{-2}F_V$,
$K_WF_WK_W^{-1}=q^{-2}F_W$, commuting $F_V,F_W$, and trivial cross
Cartan commutators. These are the two $\mathfrak{sl}_2$ Cartan
ratios in the original $U_q(\mathfrak{gl}_2\oplus\mathfrak{gl}_2)$;
the two central total-weight generators act by $q^{15}$.

Take the finite directed support graph on the actual source basis of
$X$: an arrow $b\to b'$ of colour $V$ or $W$ means the corresponding
coefficient in $F_Vb$ or $F_Wb$ is nonzero in $K$. Let $\mathcal D$
be all descendants of $T$, including $T$, and let $\mathcal J$ be
the descendants which cannot reach $Q$, with length-zero paths allowed.
Both sets are closed under outgoing arrows. Indeed a descendant's
successor is a descendant; and a successor which could reach $Q$
would give a path to $Q$ from its predecessor. Consequently
\[
 X_{\mathcal D}=\bigoplus_{b\in\mathcal D}Kb,
 \qquad X_{\mathcal J}=\bigoplus_{b\in\mathcal J}Kb
\]
are $B^-_K$-submodules, since Cartan acts diagonally. Their quotient
$M=X_{\mathcal D}/X_{\mathcal J}$ is an actual source subquotient.
This construction uses the span of support descendants; it makes no
claim that the span equals the cyclic submodule generated by $T$.

Every path from $T$ to $Q$ has two $V$ arrows and no $W$ arrows:
each $W$ arrow lowers the first $W$ coordinate by one, which cannot
be restored, and each $V$ arrow lowers the first $V$ coordinate by
one. The exhaustive calculation above therefore proves
$\mathcal D\setminus\mathcal J=\{T,T_1,T_2,Q\}$.
In this ordered quotient basis, denote the vectors by $e_0,e_1,e_2,e_3$.
Set
\[
 \alpha=[7]-[3],\quad \beta=[8],\quad d=[3],\quad
 z=q^3+q^{-3},\quad\gamma=[6]=dz,\quad\delta=-d.
\]
Then every entry of the induced operators is
\[
 N=[F_V]=\begin{pmatrix}0&0&0&0\\
 \alpha&0&0&0\\ \beta&0&0&0\\0&\gamma&\delta&0\end{pmatrix},
 \qquad [F_W]=0,\qquad
 [K_V]=\operatorname{diag}(q^9,q^7,q^7,q^5),\quad[K_W]=q^3I_4.
\]
Direct multiplication verifies every negative-Borel relation, including
$K_VN=q^{-2}NK_V$ and the cross relations. In particular the four
vectors are not asserted to span an invariant subspace inside $X$:
the discarded vectors in the first-step table are genuinely present
before taking this specified quotient.

The Laurent identity calculated from the two paths is
\begin{align}
 C&=\frac{[6]([7]-[3])-[3][8]}{[2]}
   =q^{10}-q^4-q^{-4}+q^{-10}
   =r^2[7][3].\label{eq:gct-block-C}
\end{align}
For a direct verification, $[6]=[3](q^3+q^{-3})$ and
$(q^3+q^{-3})([7]-[3])-[8]=[2]r^2[7]$ follow by multiplying
their displayed finite Laurent sums; substituting gives the last
equality, and expanding gives the four displayed terms.
Thus
\[
 N^2=pD,\qquad D=CE_{30},\qquad N^3=ND=DN=D^2=0.
\]
Here $E_{30}e_0=e_3$ and it kills the other three vectors. This is
the divided square $F_V^{(2)}=F_V^2/[2]$, with its denominator retained
until the Laurent divisibility has been proved.

Define $M_A=\bigoplus_{i=0}^3Ae_i$. It is precisely the quotient of
the $A$-spans of $\mathcal D$ and $\mathcal J$ in the source integral
form: the source Chevalley coefficients are in $A$, so both spans
are stable and their quotient has this basis. The displayed matrices
preserve $M_A$. The induced divided powers preserve it as well:
$F_V^{(0)}=I$, $F_V^{(1)}=N$, $F_V^{(2)}=D$, and higher $V$ divided
powers and all positive $W$ divided powers vanish. The only product
identity with a nonzero product of positive divided powers is
$F_V^{(1)}F_V^{(1)}=[2]F_V^{(2)}$, already verified. All other
divided-power product identities have both sides zero. This specifies
the divided-power integral negative-Borel action on this subquotient.

There is no extension of these four matrices to the full original
quantum group. Such an extension would require an endomorphism $E_V$
satisfying $E_VN-NE_V=(K_V-K_V^{-1})/(q-q^{-1})$. The left side has
trace zero because $\operatorname{tr}(AB)=\operatorname{tr}(BA)$,
whereas the right side has trace $[9]+2[7]+[5]$, a nonzero Laurent
polynomial (its $q^8$ coefficient is one). This proves the exact
obstruction to that extension. The full source module $X$, with
its additional weights and all four Chevalley generators, is constructed
explicitly in Section~\ref{sec:full-source-generators}. Its original
$1260$ basis vectors, all divided powers, and the $420$-dimensional
domain with $416$-dimensional kernel give the exact quotient
morphism to $M$. The commutator obstruction on $M$ itself remains
the one just proved.

\subsection{The retained integral kernel, cokernel, and specialization}

We first retain the original base $A=\mathbb Z[q,q^{-1}]$ before
inverting $2$. Define
\[
 a_*=q^5+q^{-5},\qquad b_*=q^6+q^2+q^{-2}+q^{-6},\qquad
 I=(a_*,b_*)\subset A.
\]
Then $I=(4,q^2+1)=(4,p)$. Here is an explicit proof of the ideal
equality. Put $x=q^2$, $a_x=x^5+1=q^5a_*$ and
$b_x=x^6+x^4+x^2+1=q^6b_*$. Direct multiplication gives
\[
 (x^4+x^3+x^2+x+1)a_x-x^2(x+1)b_x=x+1.
\]
The polynomials displayed below in the identity
$u_*a_*+v_*b_*=1$ have denominators only $4$, so multiplication by
$4$ gives $4\in I$ already over $A$. Conversely $a_x$ is divisible
by $x+1$, and division of $b_x$ by $x+1$ has remainder $4$ because
$b_x(-1)=4$. This proves both containments. Since $q$ is invertible
and $q^2+1=qp$, the second equality follows. In particular
$A/I\simeq(\mathbb Z/4)[q]/(q^2+1)$, with inverse of $q$ equal to
$-q$; the monic relation makes this free of rank two over $\mathbb Z/4$.

The equation $Nx=0$ gives, over the domain $A$,
\[
 \ker_AN=A(e_1+ze_2)\oplus Ae_3.
\]
For the cokernel let $P=A^2/pA(a_*,b_*)$. Since the first nonzero
column of $N$ is $p(a_*,b_*)$ and the last row generates $dAe_3$,
\[
 \operatorname{coker}_AN=Ae_0\oplus P\oplus A/(d).
\]
There is the exact sequence
\[
 0\longrightarrow A/(p)\xrightarrow{\bar f\mapsto[f(a_*,b_*)]}
 P\xrightarrow{[(x_1,x_2)]\mapsto-b_*x_1+a_*x_2} I
 \longrightarrow0.
\]
Well-definedness and surjectivity of the last map follow from its
formula. The two elements $a_*,b_*$ have no common nonunit divisor:
a common divisor would divide both $4$ and $q^2+1$, which is
impossible in the unique factorization domain $A$. Thus
$a_*x_2=b_*x_1$ implies $x_1=a_*f$, $x_2=b_*f$ in $A$ by
prime-factor divisibility. Quotienting this kernel by $pA(a_*,b_*)$
gives exactly $A/(p)$, proving injectivity of the first map and
exactness in the middle. This records the original $2$-primary
ideal defect which is removed by the base change made next.

At $q=1$ over $\mathbb Z$, the image in the middle two coordinates
is generated by $(4,8)$ and the bottom-coordinate image is $3\mathbb Z$.
The integral basis $(e_1+2e_2,e_2)$ therefore proves
$\operatorname{coker}_{\mathbb Z}N_0\simeq
\mathbb Z^2\oplus\mathbb Z/4\oplus\mathbb Z/3$.

For the node comparison make the explicit base change
$A\to R=\mathbb Z[1/2][q,q^{-1}]$, and put $O=\mathbb Z[1/2]$.
No quantum integer, $3$, $7$, or $q-1$ is inverted.
The finite-sum identities give $\alpha=pa_*$, $\beta=pb_*$ and
$za_*-b_*=r^2[7]$. For an explicit Bezout identity put
\begin{align*}
 u_*&=\frac{q^5}{4}(-q^{10}+q^8+2q^6+2q^4+q^2+3),\\
 v_*&=\frac{q^6}{4}(q^8-q^6-3q^4-q^2+1).
\end{align*}
Multiplication gives $u_*a_*+v_*b_*=1$. Hence
$f=a_*e_1+b_*e_2$ and $g=-v_*e_1+u_*e_2$ are an $R$-basis of
$Re_1\oplus Re_2$: the determinant of their two columns is one
and their inverse sends $xe_1+ye_2$ to
$(u_*x+v_*y)f+(-b_*x+a_*y)g$.

For $x=\sum x_ie_i$, the equation $Nx=0$ first gives $x_0=0$,
because $\alpha$ is nonzero in the domain $R$. The last coordinate
then gives $d(zx_1-x_2)=0$, hence $x_2=zx_1$. Therefore
\[
 \ker_RN=R(e_1+ze_2)\oplus Re_3.
\]
Also $Ne_0=pf$, $Ne_1=dz e_3$, $Ne_2=-d e_3$, $Ne_3=0$,
so $\operatorname{im}N=Rpf\oplus Rd e_3$ and the explicit quotient
map is
\[
 \sum x_ie_i\longmapsto
 (x_0,\ -b_*x_1+a_*x_2,\ 
  \overline{u_*x_1+v_*x_2},\ \overline{x_3})
 \in R\oplus R\oplus R/(p)\oplus R/(d).
\]
Its kernel is exactly the displayed image, by the inverse basis
formula, and its surjectivity follows by lifting the four coordinates.
Thus this is a complete cokernel description, including integral torsion.

Specialize by $R\to O$, $q\mapsto1$, and denote reduction by a subscript
zero. Since $\alpha_0=4$, $\beta_0=8$, $\gamma_0=6$, $\delta_0=-3$,
\[
 N_0=\begin{pmatrix}0&0&0&0\\4&0&0&0\\8&0&0&0\\0&6&-3&0\end{pmatrix},
 \quad N_0^2=0,
 \quad\ker_ON_0=O(e_1+2e_2)\oplus Oe_3,
 \quad\operatorname{coker}_ON_0\simeq O^2\oplus O/(3).
\]
The same coordinate proof establishes the kernel and cokernel after
specialization; $p_0=2$ is a unit of $O$ but $d_0=3$ is not.
In particular the integral torsion cannot be inferred from rational
rank alone. The kernel base-change map takes the two displayed
generators to the two displayed generators and is an isomorphism.

Over $K$, $N^2\ne0$, $N^3=0$, and $\ker N$ has dimension two.
For completeness the columns
$e_0,Ne_0,N^2e_0,e_1+ze_2$ have determinant
$-p^2r^4[7]^2[3]\ne0$. Thus they give the exact chain decomposition
with Jordan block lengths $(3,1)$. Over $\mathbb Q$ at $q=1$, the
columns $e_0,N_0e_0,e_2,N_0e_2$ have determinant $-12\ne0$,
giving block lengths $(2,2)$. These are proved changes of basis over
their stated fields, not identifications of integral lattices.

The unrescaled tableau span has ordered basis
\[
 (e_0,e_1,-pe_2,-pe_3).
\]
Its inclusion into $M_R$ has matrix
$\operatorname{diag}(1,1,-p,-p)$.
Its inclusion cokernel is $(R/(p))^2$, and its kernel is zero because
$R$ is a domain. It is stable under $N$: the four images in that
ordered basis are respectively
$\alpha e_1-(\beta/p)(-pe_2)$, $-(\gamma/p)(-pe_3)$,
$\delta(-pe_3)$, and zero. Here $\beta/p=b_*$ and
$\gamma/p=[3](q^2-1+q^{-2})$ lie in $R$.
It need not be stable under divided powers: $De_0=-(C/p)(-pe_3)$,
and $p$ does not divide $C$. Indeed evaluation at $q=i$ over
$\mathbb C$ gives $p(i)=0$, $r(i)^2=-4$, $[7](i)=-1$,
$[3](i)=-1$, so $C(i)=-4\ne0$. This retains the distinction between
the actual source-rescaled lattice and the raw tableau lattice.

\subsection{An equivariant endpoint map and the original node}

The flag $Re_3\subset Re_1\oplus Re_2\oplus Re_3\subset M_R$
is negative-Borel stable, as follows directly from the matrices.
Let $L_{\rm top}=M_R/(Re_1\oplus Re_2\oplus Re_3)$ and
$L_{\rm bot}=Re_3$. The first line has Cartan ratios $(q^9,q^3)$
and the second $(q^5,q^3)$; both lowering operators kill both lines.
Let $\chi$ be the rank-one negative-Borel module with $V$-weight
$(-2,2)$ and $W$-weight $(0,0)$, hence Cartan ratios
$(q^{-4},1)$ and zero lowering operators. With the usual tensor-product
action, define
\[
 \Phi:L_{\rm top}\otimes\chi\longrightarrow L_{\rm bot},\qquad
 \overline e_0\otimes1\longmapsto Ce_3.
\]
This is the actual map induced by $F_V^{(2)}$ on $M_R$: it kills
the submodule divided out, and its image is $CR e_3$ inside the displayed
bottom line. Its Cartan actions agree after the specified twist, and both
lowering actions vanish on source and target, proving equivariance
on every generator. The map is injective because $C\ne0$ in $R$,
and its exact cokernel is $(R/(C))e_3$.

Here is the exact comparison with the original two-sheet node, including
all coordinates. Over $O[a]$ set
\[
 B=O[a,r_1,r_2]/(r_1^2+a,r_2^2+a),\qquad
 N_{\rm node}=O[y]\oplus O[y],\qquad a=(-y^2,-y^2).
\]
Restriction sends $r_1\mapsto(y,y)$ and $r_2\mapsto(y,-y)$.
In the ordered bases
$(1,r_1,r_2,r_1r_2)$ and
$((1,1),(y,y),(y,-y),(1,-1))$, respectively, it is
$\operatorname{diag}(1,1,1,-a)$: the last image is
$(y^2,-y^2)=-a(1,-1)$. Monic division proves that the first is a
basis; the even/odd decomposition and invertibility of $2$ prove
that the second is a basis. Substitute the original base map
$a=-r^2$ by $O[a]\to R$. The resulting restriction map is
\[
 V:B_R\longrightarrow (N_{\rm node})_R,
 \qquad [V]=\operatorname{diag}(1,1,1,r^2).
\]
Write $h=[7][3]$. Stabilizing the actual endpoint map $\Phi$ by
three identity maps gives a rank-four map $\widehat\Phi$ with matrix
$\operatorname{diag}(1,1,1,C)$. Give both endpoint coordinates
their already specified negative-Borel character $(q^5,q^3)$ and
the three added coordinates any fixed equal characters on source
and target. Give the intermediate and node modules these same
diagonal actions. These explicit module actions make the following
maps equivariant as module maps; no action by algebra automorphisms
of the node ring is asserted:
\[
 \widehat\Phi=VH,\qquad H=\operatorname{diag}(1,1,1,h).
\]
Both maps are injective, and, inside the common target, the image
lattices satisfy $L_C=R^3\oplus CR\subset L_r=R^3\oplus r^2R$.
Consequently
\[
 0\longrightarrow R/(h)\xrightarrow{\;\bar f\mapsto\overline{r^2f}\;}
 R/(C)\xrightarrow{\;\bar f\mapsto\bar f\;}R/(r^2)
 \longrightarrow0.
\]
The first map is well defined because $r^2h=C$; its kernel is zero
by cancellation of $r^2$. The kernel of the second map is precisely
the classes of $r^2f$, and surjectivity follows by lifting a
representative. This proves the full sequence and connects the actual
GCT generator endpoint to the retained conductor lattice.

Finally put $t=q-1$, $u=(q+1)/q$, so $r=tu$, $u(1)=2$ and $h(1)=21$.
The first nonzero $t$-coefficient of $C=t^2u^2h$ is $84t^2$;
equivalently $C(1)=C'(1)=0$ and $C''(1)=168$. The two rank-four
maps $V$ and $\widehat\Phi$ specialize to
$\operatorname{diag}(1,1,1,0)$, with kernel and cokernel $O$ on
the fourth coordinate. The induced kernel map from $H$ is $21$,
and the cokernel map is the identity. To retain the full derived
sequence, the free resolution of $O=R/(t)$ by multiplication by $t$
gives
\[
 0\longrightarrow O\xrightarrow{21}O
 \xrightarrow{\mathrm{mod}\ 21}O/(21)
 \xrightarrow{0}O\xrightarrow{1}O\longrightarrow0.
\]
There is no omitted initial Tor term: if $tf=hg$ in $R$, reduction
at $q=1$ gives $21g(1)=0$ in the domain $O$, hence $g=tg_1$.
Cancellation of $t$ then gives $f=hg_1$. Therefore
$\operatorname{Tor}_1^R(R/(h),O)=0$. The resolution has length one,
so all higher Tor groups vanish as well.
Indeed the two Tor generators are $C/t$ and $r^2/t$; projection takes
the first to $h$ times the second and hence to $21$ after reduction.
Lifting $r^2/t$ to $R/(C)$ and multiplying by $t$ gives $r^2$, the
image of $1\in R/(h)$, proving the positive sign of the connecting
map. The next map is multiplication by $r^2$ and vanishes modulo
$t$; the final map is the identity. Multiplication by $21$ is injective
on $O$, with exactly the displayed cokernel. No passage to a field
has discarded this residual module.

The finite interval, its exact source quotient maps, the divided-power
endpoint, the two integral lattices, and this conductor correspondence
are therefore constructed. They do not constitute the full generator
representation of $\check X_{(7,5,3)}$ or a Frobenius realization of
the nonstandard Riemann-hypothesis programme.

\section{The complete original GCT-IV module and its integral projection}
\label{sec:full-source-generators}

Let $A=\mathbb Z[q,q^{-1}]$, $K=\mathbb Q(q)$, $p=q+q^{-1}$,
and $[j]=q^{j-1}+q^{j-3}+\cdots+q^{1-j}$ for $j>0$.
This section constructs the four original Chevalley matrices on
$\check X_{(7,5,3)}$, rather than extending the four-dimensional
negative-Borel quotient of Section~\ref{sec:gct-generator-interval}.
The column shape is exactly $\alpha=(3,3,3,2,2,1,1)$; the row
shape is $(7,5,3,0)$; the column-count shape is $[0,3,2,2]$.
The source parameter, denoted $u$ in its TeX, remains the
indeterminate $q$ here.

The construction is finite and explicit.  The file
\texttt{agents/}\allowbreak
\texttt{full\_source\_generators/}\allowbreak
\texttt{full\_generators.json}
contains every entry of all four $1260\times1260$ matrices and
the image of every original tabloid. Its coefficient convention is
a list of pairs $(k,a_k)$ denoting the Laurent polynomial
$\sum_k a_kq^k$; its sparse matrix convention lists target indices
in each source column. Indices start at zero. The file
\texttt{source\_basis/basis\_1260.json} gives every original
column label, scalar, invariant record, twelve-tuple, and weight.
The definitions and recurrences below specify these files exactly;
the attached Python programs implement those definitions with
integer arithmetic and exhaustive, finite checks.

\subsection{The original ambient basis and all adjacent relations}

Write $Y=\check Y_\alpha$. Its original NST basis is indexed by
the set $\mathcal W$ of all column tuples in
\[
 \{123,124,134,234\}^3\times
 \{12,13,23,32,24,34\}^2\times\{1,2,3,4\}^2.
\]
The product is the source based product, with the indicated column
composition retained. Thus $|\mathcal W|=4^3 6^2 4^2=36864$.
These are column labels, not ordinary tensor monomials. In
particular, the two columns $23$ and $32$ are distinct.
The alphabet is $1=(1,1),2=(1,2),3=(2,1),4=(2,2)$.
For diagram calculations read each column from bottom to top,
then extract its $V$ and $W$ letters:
\[
\begin{array}{c|cc@{\qquad}c|cc}
c&V&W&c&V&W\\\hline
1&1&1&2&1&2\\3&2&1&4&2&2\\
12&11&21&13&21&11\\23&21&12&32&12&21\\
24&21&22&34&22&21\\
123&211&121&124&211&221\\
134&221&211&234&221&212
\end{array}
\]
Some selected canonical tensor summands of the height-three
columns use other words. The displayed dictionary is the literal
diagram reading rule; it avoids identifying a whole column with
one summand of its canonical expansion.

Let $R_K$ be the original sum of the six adjacent source
submodules $\check Y_{\triangleright^i\alpha}\subset Y$.
The following tables give all of its integral two-column basis
relations. Each row is an equality of the images in $Y/R_K$;
the equality means that the difference belongs to the indicated
local submodule before taking the quotient:
\[
\begin{array}{c|l}
(2,1)&(23)(1)=(32)(1),\quad(23)(3)=(34)(1),\\
 &(24)(3)=(34)(2),\quad(24)(1)=(32)(2)\\\hline
(3,2)&(124)(13)=-(134)(12),\quad(234)(23)=-(234)(32),\\
 &(234)(13)=-(134)(32),\quad(124)(23)=-(234)(12)\\\hline
(2,2)&(32)(24)=(24)(32),\quad(12)(13)=(13)(12),\\
 &(34)(24)=(24)(34),\quad(12)(23)=(23)(12),\\
 &(32)(13)=(13)(32),\quad(34)(13)=(13)(34),\\
 &(34)(23)=(23)(34),\quad(12)(24)=(24)(12),\\
 &(32)(23)=(23)(32),\quad(24)(13)=(34)(12).
\end{array}
\]
The first nine $(2,2)$ relations are its nine-dimensional
degree-preserving component; the last is its one-dimensional
invariant component. For an integral local relation the source
based products for the submodule and ambient module agree, so
these relations hold with every left and right NST context.

There are two further local strings at each equal height. The
active side, its three words $C_{11},C_{12},C_{22}$, and its
invariant $I$ are as follows:
\[
\begin{array}{c|c|ccc|c}
h&\text{side}&C_{11}&C_{12}&C_{22}&I\\\hline
1&V&(2)(1)&(2)(3)&(4)(3)&(4)(1)\\
1&W&(3)(1)&(3)(2)&(4)(2)&(4)(1)\\
2&V&(32)(12)&(32)(32)&(34)(32)&(34)(12)\\
2&W&(23)(13)&(23)(23)&(24)(23)&(24)(13)\\
3&V&(124)(123)&(124)(134)&(234)(134)&(234)(123)\\
3&W&(134)(123)&(134)(124)&(234)(124)&(234)(123)
\end{array}
\]
In the ambient two-column module their projected string vectors
are $\widetilde C_{11}=C_{11}$,
$\widetilde C_{12}=C_{12}+p^{-1}I$,
$\widetilde C_{22}=C_{22}$.  Their contextual products in
$\check Y_{\triangleright^i\alpha}$ give exactly
\begin{equation}\label{eq:full-contextual-relations}
 BC_{11}D=-p^{-1}B'ID,\qquad
 BC_{12}D=-p^{-1}BID,\qquad
 BC_{22}D=-p^{-1}BID'.
\end{equation}
Here is the full meaning of $B'$ and $D'$. Pair every $2$ with
the next available $1$ to its right, using a stack of unmatched
$2$ positions. For $11$, take the leftmost local unpaired $1$.
If it has no partner in $B$, set $B'=0$; otherwise change its
partner from $2$ to $1$ and apply the raising column Kashiwara
operator at that column. For $22$, take the rightmost local
unpaired $2$. If it has no partner in $D$, set $D'=0$;
otherwise change that partner from $1$ to $2$ and apply the
lowering column Kashiwara operator there. In either case an
extra internal arc makes the corrected word zero. The other
context is unchanged. The middle equality changes no context.
These are the source contextual projection formulas, with every
outside letter retained. An isolated equality from this table
cannot be substituted in a longer product while discarding
$B'$ or $D'$.

The integral table and these six strings exhaust the local
submodule bases for the adjacent shapes occurring in $\alpha$:
$(3,3)$ twice, $(3,2)$ once, $(2,2)$ once, $(2,1)$ once and
$(1,1)$ once. Taking all contexts therefore spans exactly $R_K$,
by its defining sum. This proves that the following relation
construction presents the original $\check X_{(7,5,3)}$.

\subsection{A finite exact quotient algorithm}

Order $\mathcal W$ lexicographically using the column orders in
its displayed Cartesian product. There are $22528$ contextual
instances of integral table rows and $47616$ instances of
\eqref{eq:full-contextual-relations}. Each says
$t_i=\epsilon p^kt_j$ with $\epsilon\in\{1,-1\}$ and
$k\in\{0,-1\}$, or $t_i=0$.
Join the vertices by these equalities, keeping for each vertex
an accumulated scalar $\epsilon p^k$ relative to a chosen root.
An equality within one component either agrees with its
accumulated scalar, or yields
$(1-\epsilon p^k)t_{\rm root}=0$. In the latter case set that
component to zero; do the same when a zero relation is present.

This is Gaussian elimination specialized to binomial relations
over $K$. Indeed a spanning tree in each component expresses
every vertex as its root times a nonzero scalar. Every remaining
edge is then either $0=0$ or a nonzero scalar times that root.
Consequently a component contributes exactly one free $K$
coordinate if every cycle agrees and no zero relation is present,
and contributes zero otherwise. For this original presentation
the algorithm leaves exactly $1260$ nonzero components; $18652$
original words are zero and $18212$ original words have nonzero
image. The program checks every equation again through its
computed projection.

This elimination takes place over $K$. Inverting a nonzero
$1-\epsilon p^k$ in a conflicting component does not prove that
an integral presentation obtained merely by clearing denominators
has no torsion. We make no such claim. The integral lattice and
its exact kernel are constructed separately below inside the
original field quotient.

\subsection{The complete original honest basis and its signs}
\label{fullbasis:enumeration}

We specify the original basis at the original parameter.  Put
$K=\mathbb Q(q)$, $p=q+q^{-1}$, $\nu=(7,5,3,0)$, and
$\alpha=\nu'=(3,3,3,2,2,1,1)$.  In this subsection a string such as
$124$ denotes a single column, read from top to bottom; parentheses list
columns from left to right.  In particular $23$ and $32$ are distinct
nonstandard columns.  The source column alphabets, in their source order,
are
\[
\mathcal C_3=(123,124,134,234),\qquad
\mathcal C_2=(12,13,23,32,24,34),\qquad
\mathcal C_1=(1,2,3,4).
\]
All references to source lines below refer to the retained original
\texttt{Apr13KroneckerGCT4.tex}, whose SHA-256 is
\begin{center}\small
\texttt{5c5b990306c3b817d822df9f5088fdd42dda146fc3449155a0b2f6202d2ed49e}
\end{center}
The alphabets and their full canonical tensor expressions occur at
lines~4630--4635; their order occurs at~6064.

The ambient source space is
\[
\check Y_\alpha=
 \bigotimes_{j=1}^{7}\check\Lambda^{\alpha_j}X,
\qquad \dim_K\check Y_\alpha=4^3\,6^2\,4^2=36864.
\]
Its NST basis uses the source projected canonical product: if
$T=(C_1,\ldots,C_7)$, its vector is
$\check p_\alpha(C_1\mathbin\heartsuit\cdots\mathbin\heartsuit C_7)$.
Here $\heartsuit$ concatenates the original upper canonical words in each
of the $V,W$ factors, and $\check p_\alpha$ is the source projector.
This is the definition at~4751--4765 and agrees with the based tensor
product by~4795--4805.  Thus these words do not mean an unqualified pure
tensor in the NSC basis.  The original quotient is
\[
 \check X_\nu=\check Y_\alpha/\sum_{j=1}^{6}\check Y_{>^j\alpha},
 \qquad
 \check Y_{>^j\alpha}=
 \check Y_{(\alpha_1,\ldots,\alpha_{j-1})}\otimes
 \check Y_{>(\alpha_j,\alpha_{j+1})}\otimes
 \check Y_{(\alpha_{j+2},\ldots,\alpha_7)}.
\]
These are exactly the objects in source~4918--4920, with the source local
relation submodules.  When written in the canonical NST product the
nonintegral relations must retain their context, as in~5773--5785.

Here is an exhaustive finite index set for the honest basis.  Choose
$i_3,i_2,i_1\in\{0,1\}$ and nonnegative integers $m_C$ for each column
$C\in\mathcal C_3\sqcup\mathcal C_2\sqcup\mathcal C_1$, subject to
\begin{equation}\label{fullbasis:counts}
 \sum_{C\in\mathcal C_3}m_C=3-2i_3,\qquad
 \sum_{C\in\mathcal C_2}m_C=2-2i_2,\qquad
 \sum_{C\in\mathcal C_1}m_C=2-2i_1.
\end{equation}
Retain precisely the tuples satisfying all the following source
conditions:
\begin{align}
 m_2m_3&=0,&
 m_1>0&\ \Longrightarrow\ m_{34}=m_{24}=m_{32}=0,\nonumber\\
 m_2>0&\ \Longrightarrow\ m_{34}=0,&
 m_3>1&\ \Longrightarrow\ m_{24}=0,\nonumber\\
 m_{23}&\leq1,&m_{32}&\leq1,\nonumber\\
 m_{124}m_{134}&=0,&
 m_{234}>0&\ \Longrightarrow\ m_{12}=m_{13}=m_{32}=0,\nonumber\\
 m_{134}>0&\ \Longrightarrow\ m_{12}=0,&
 m_{124}>1&\ \Longrightarrow\ m_{13}=0,\nonumber\\
 m_{234}m_1&=0.&&\label{fullbasis:predicates}
\end{align}
Within each height the invariant-free columns are in the displayed
alphabet order.  Equations~\eqref{fullbasis:predicates} retain every
condition of source~6114--6130; conditions with identical antecedents
have been grouped without changing their content.

From the tuple form the exact source word
\begin{equation}\label{fullbasis:word}
\begin{aligned}
 T(m,i)={}&123^{m_{123}}124^{m_{124}}134^{m_{134}}
 (234,123)^{i_3}234^{m_{234}}\\[-1mm]
 &12^{m_{12}}13^{m_{13}}23^{m_{23}}32^{m_{32}}
 (24,13)^{i_2}24^{m_{24}}34^{m_{34}}\\[-1mm]
 &1^{m_1}2^{m_2}3^{m_3}(4,1)^{i_1}4^{m_4}.
\end{aligned}
\end{equation}
Exponentiation of a column or pair means that many consecutive copies.
The invariant insertion positions and orientations are those of
source~6169.  The invariant record is $(0,i_3,i_2,i_1)$ and the degree
is $d=i_3+i_2+i_1$.  Indeed the inserted pairs have two arcs connecting
their columns and no arcs exiting the pair.  Removing them gives the
invariant-free word with multiplicities $m_C$.  The source classification
at~6263--6269 identifies these straightened words bijectively with the
honest classes.  The complete quotient computation in this chapter also
checks, independently for this fixed shape, that the words in
\eqref{fullbasis:word} project to distinct nonzero components and exhaust
all nonzero quotient components.  Thus both the representatives and the
module are the original source objects.

For completeness we give an exact diagram algorithm.  Read the columns
bottom to top and then left to right and apply
\[
 1\mapsto(1,1),\quad2\mapsto(1,2),\quad
 3\mapsto(2,1),\quad4\mapsto(2,2).
\]
This produces the original length-$15$ $V$ and $W$ words (source~4723).
For each word scan left to right, pushing every $2$ on a stack; a $1$
pairs with the last unmatched $2$ if there is one.  Otherwise it remains
unpaired.  This is exactly the source parenthesis rule~2448.  Record the
column containing each letter.  An arc is external when its ends lie in
different columns; two arcs between a pair of columns give an invariant
pair (source~4938--4954).  A two-edge component involving one column of
each height $3,2,1$ is a $3$--$2$--$1$ arc.  Its color is the common
edge color, or the color of the longer edge if the colors differ, as in
source~6320.  All other remaining edges are pure.  A height-$2$ column
with no external edge is externally free; its type is $V$ for
$12,32,34$ and $W$ for $13,23,24$ (source~6068--6073).

The exact twelve-entry label is
\begin{equation}\label{fullbasis:zeta}
 \zeta=(i_4,i_3,i_2,i_1,a_{21},a_{32},a_{31V},a_{31W},
                 a_{321V},a_{321W},f_V,f_W).
\end{equation}
Here $a_{21},a_{32}$ count pure arcs, the other $a$'s retain their
specified source colors, and $f_V,f_W$ count externally free columns.
This is the original order in source~7233--7238.  Together with the two
unpaired words, necessarily of the form $1^a2^b$, it determines the
honest class by~6473--6481.  Original torus weights count all letters,
including paired letters: each factor's weight is
$(\#1,\#2)$.  This retains the determinant-weight contribution of every
invariant.

\paragraph{A source sign ambiguity and its exact resolution.}
Source~5489--5490 specifies a positive class using the multiplier
$(-1/p)^d$ and a representative with a $W$-colored $3$--$2$ arc whenever
the class has a $3$--$2$ arc.  It also asserts that having no $3$--$2$
arc for every representative is equivalent to having none for any
representative.  That latter assertion fails; the $3$--$2$--$1$ example
at~6337--6388 already has equivalent representatives of both types.
We preserve the original sign choice through its stated
$W$-arc representative criterion and explicitly record the needed
correction, instead of treating the inaccurate equivalence as a theorem.

Here is the exact integral witness relevant to this shape:
\begin{equation}\label{fullbasis:signwitness}
 \begin{aligned}
 (134,13,23,1,x)
 &=(134,13,32,1,x)\\
 &=(134,32,13,1,x)\\
 &=-(234,13,13,1,x),\qquad x\in\{1,2,3,4\}.
 \end{aligned}
\end{equation}
The first equality is the source integral $(23,1)=(32,1)$ relation
at~5080.  The second is the type-$V$/type-$W$ commutation at~5139.
The last is $(134,32)=-(234,13)$ at~5209.  All are degree-preserving
integral moves, so their insertion into the retained contexts is valid.
The original left word has a $3$--$2$--$1$ $V$-arc and no $3$--$2$
edge; the right word has a $W$-colored $3$--$2$ edge.  Direct application
of the pairing algorithm verifies both assertions for each of the four
values of $x$.  Consequently the positive vector represented by the
left word is $-(-1/p)^d$ times that word.

Within the present shape, the additional words requiring this sign are
exactly
\[
\begin{gathered}
 (123,123,134,13,23,1,x),\quad
 (123,134,134,13,23,1,x),\\
 (134,134,134,13,23,1,x),\quad
 (134,234,123,13,23,1,x),\\
 x\in\{1,2,3,4\}.
\end{gathered}
\]
For the last family first use
$(134,234,123)=(234,123,134)$, the invariant commutation
at~5286--5287, then apply~\eqref{fullbasis:signwitness}.  The first
three families use~\eqref{fullbasis:signwitness} directly in their
height-$3$ prefixes.  Exhaustiveness follows by evaluating
\eqref{fullbasis:predicates} and the pairing algorithm on the finite
index set; every index and its predicate values are computed by
\texttt{enumerate\_basis.py}, and independently by
\texttt{audit/independent\_audit.py}.  All sixteen equalities are also
checked against the full original quotient projection by
\texttt{check\_sign\_witnesses.py}.  In source lexicographic order their
zero-based indices are $262$--$265$, $478$--$481$, $1000$--$1003$, and
$1108$--$1111$.

Thus the final completely explicit scaling is
\begin{equation}\label{fullbasis:sign}
 b_{m,i}=\epsilon(T)(-1/p)^d T,
 \qquad
 \epsilon(T)=
 \begin{cases}
 -1,&\text{$T$ has a $V$-colored $3$--$2$ edge},\\
 -1,&\text{$T$ has no $3$--$2$ edge and has a $3$--$2$--$1$ $V$-arc},\\
 1,&\text{otherwise}.
 \end{cases}
\end{equation}
The sixteen corrections affect neither the index set nor any original
quotient relation.  They give an explicitly specified diagonal sign
change, with its exact inverse equal to itself.  All four retained
labels below have their original established scalings unchanged.

\paragraph{Exhaustive enumeration and a second exact bijection check.}
For each invariant record the numbers of words in
\eqref{fullbasis:word} are as follows:
\[
\begin{array}{c|rrrrrrrr}
(i_3,i_2,i_1)&000&001&010&011&100&101&110&111\\\hline
\text{number}&628&148&108&16&268&56&32&4.
\end{array}
\]
Their sum is $1260$; the degree counts are $628,524,104,4$ for
$d=0,1,2,3$.  All $1260$ keys consisting of~\eqref{fullbasis:zeta}
and the two unpaired words are distinct.  There are $48$ distinct
twelve-tuples.  For each such tuple, if its two unpaired lengths are
$h_V,h_W$, its vectors contain exactly all pairs
\[
 (1^{h_V-a}2^a,1^{h_W-b}2^b),\qquad
 0\leq a\leq h_V,\quad0\leq b\leq h_W.
\]
The complete data are in \texttt{basis\_1260.json} and
\texttt{ssyt\_bijection\_check.json}.

We additionally verify the original map to ordinary semistandard
tableaux, retaining every source case rather than using its cardinality
alone.  For clarity here is that map in multiplicity form.  Let $m'_C$
be the ordinary output multiplicities, initially equal to the
invariant-free $m_C$; set $m'_{32}=0$.  The columns $123,234,1,4$ keep
their invariant-free multiplicities.  The following tables specify
all other multiplicities; entries not changed in a row keep their
invariant-free values.  First, the right part (source~6195--6210):
\[
\begin{array}{c|cc|cc}
\text{condition}&m'_2&m'_3&m'_{24}&m'_{34}\\\hline
m_1>0&i_1+m_2&i_1+m_3&m_{24}&m_{34}\\
m_1=0,\ m_{24}>0,\ m_2=0&0&2i_1+m_3&m_{24}&m_{34}\\
m_1=0,\ m_{24}>0,\ m_2>0&i_1+m_2&i_1&m_{24}&0\\
m_1=m_{24}=0,\ i_1>0&i_1+m_2&i_1+m_3&m_{34}&0\\
m_1=m_{24}=i_1=0&m_2&m_3&0&m_{34}
\end{array}
\]
Next the middle part (source~6217--6232):
\[
\begin{array}{c|cc}
\text{condition}&m'_{14}&m'_{23}\\\hline
m_{234}=0,\ m_1>0&2i_2+m_{23}&0\\
m_{234}>0,\ m_1=0&0&2i_2+m_{23}\\
m_{234}=m_1=m_{23}=0&2i_2+m_{32}&0\\
m_{234}=m_1=0,\ m_{23}>0&0&2i_2+m_{32}+m_{23}
\end{array}
\]
Finally the left part (source~6237--6250):
\[
\begin{array}{c|cc|cc}
\text{condition}&m'_{124}&m'_{134}&m'_{12}&m'_{13}\\\hline
m_{234}>0&i_3+m_{124}&i_3+m_{134}&m_{12}&m_{13}\\
m_{234}=0,\ m_{13}>0,\ m_{134}=0&2i_3+m_{124}&0&m_{12}&m_{13}\\
m_{234}=0,\ m_{13}>0,\ m_{134}>0&i_3&i_3+m_{134}&0&m_{13}\\
m_{234}=m_{13}=0,\ i_3>0&i_3+m_{124}&i_3+m_{134}&0&m_{12}\\
m_{234}=m_{13}=i_3=0&m_{124}&m_{134}&m_{12}&0
\end{array}
\]
Conditions~\eqref{fullbasis:predicates} imply that exactly one row of
each table applies.  The three tables write disjoint groups of output
multiplicities and use only original $m_C$, so they do not interfere.
Arrange the output columns in ordinary column lexicographic order.

The exact finite verification in \texttt{check\_ssyt\_bijection.py}
enumerates every tuple in~\eqref{fullbasis:counts} by combinations with
repetition, tests every predicate~\eqref{fullbasis:predicates}, and
evaluates the three tables.  Independently it enumerates every ordinary
height-$3$ column multiset of size $3$, height-$2$ column multiset of
size $2$, and height-$1$ column multiset of size $2$.  Each ordinary
column is a strictly increasing sequence in $\{1,2,3,4\}$; every common
row coordinate of adjacent columns is tested for weak increase.  These
are exactly the defining semistandard conditions, and the enumeration
has no other exclusion.  The two sets are equal, have $1260$ elements,
and the map's $1260$ outputs are distinct.  The complete inverse is
therefore given without an existence assumption by the reversed
finite correspondence in \texttt{ssyt\_correspondence.json}.
As a further exact check the ordinary Weyl product is
\[
 \prod_{1\leq i<j\leq4}\frac{\nu_i-\nu_j+j-i}{j-i}
 =3\cdot3\cdot\frac{10}{3}\cdot3\cdot\frac72\cdot4=1260.
\]
The numerical identity is a check on the explicitly constructed
bijection, not a replacement for its construction.

\paragraph{Retained original labels.}
In the zero-based basis order defined above,
\begin{align*}
 b_{126}=T&=(123,123,124,12,12,1,1),\\
 b_{146}=T^1&=(123,123,124,23,12,1,1),\\
 b_8=T^2&=-p^{-1}(123,123,123,12,12,4,1),\\
 b_{144}=Q&=-p^{-1}(123,123,124,13,12,4,1).
\end{align*}
For $T^1$ and $Q$, the lexicographic word has the $12$ column before
$23$ and $13$, respectively; their equality with the displayed original
representatives is the integral type-$V$/type-$W$ commutation.  Their
twelve-tuples and unpaired lengths are
\[
\begin{array}{c|c|c}
T&(0,0,0,0,0,0,0,1,0,0,2,0)&(9,3)\\
T^1&(0,0,0,0,0,0,0,0,0,1,1,0)&(7,3)\\
T^2&(0,0,0,1,0,0,0,0,0,0,2,0)&(7,3)\\
Q&(0,0,0,1,0,1,0,0,0,0,1,0)&(5,3).
\end{array}
\]
All these unpaired words consist entirely of $1$'s.  The original
weights $(\mathrm{wt}_V,\mathrm{wt}_W)$ are respectively
$((12,3),(9,6))$, $((11,4),(9,6))$, $((11,4),(9,6))$,
and $((10,5),(9,6))$.

The literal source target at~7333 is retained separately:
\[
 b_{154}=T'_{\mathrm{printed}}
 =-p^{-1}(123,123,124,23,12,4,1).
\]
It has the same twelve-tuple as $Q$ but unpaired words $(11111,112)$
and original weights $((10,5),(8,7))$.  Since $F_V$ commutes with
the $W$ torus, its coefficient in $F_V^2T$ or $F_V^{(2)}T$ is zero.
This is the precise weight obstruction to the printed target; the
corrected $Q$ retains the already proved contextual endpoint and is
not silently identified with the printed object.

\paragraph{Original ambient generator algorithm.}
For $Z=V,W$, let $\varphi_Z(T)$ count its unpaired $1$'s.  Source
4881--4908 gives
\[
 F_ZT=\sum_{j=1}^{\varphi_Z(T)}[j]\,F_{(j)Z}T.
\]
The $j$th operation changes the column containing the $j$th unpaired
$1$, in left-to-right order, by its local Kashiwara arrow; it is zero
if changing that selected letter creates an extra internal arc.  For
$E_Z$, replace the $j$th unpaired $2$ counted from the right, with
coefficient $[j]$, exactly as in source~2508--2517.  The local $F$
strings are
\[
\begin{array}{c|ccccc}
V&1\to3&2\to4&123\to134&124\to234&12\to32\to34\\
W&1\to2&3\to4&123\to124&134\to234&13\to23\to24.
\end{array}
\]
All unlisted arrows are zero, and the local $E$ arrows reverse these
strings.  Choosing the first of the two unpaired $1$'s in a height-$2$
spin-one column creates an extra internal arc; the second gives its
first $F$ arrow.  This zero test is performed for each summand before
the contextual quotient reduction.  Scaling the input by
\eqref{fullbasis:sign}, reducing in the original quotient, and expressing
each output in that same signed basis specifies every original
Chevalley coefficient.  The complete matrices, divided powers, and
their exact checks are supplied by the full quotient construction in
this chapter.


For each of the $1260$ displayed representatives $b_n$, let
$b_n=\sigma_n p^{-d_n}\overline t_{w_n}$, with its fully
specified sign and degree. Each nonzero relation component
contains exactly one such representative. This is checked by
matching all $1260$ roots; hence the vectors $b_n$ are an actual
basis of the original quotient. For every original word $w$
the table records
\begin{equation}\label{eq:full-tab-projection}
 \pi(t_w)=0\quad\hbox{or}\quad
 \pi(t_w)=\epsilon_w p^{k_w}b_{n_w}.
\end{equation}
The algorithm produces only $k_w\in\{0,1,2,3\}$.
The numbers of nonzero words with these four exponents are
$12544,5100,552,16$, respectively. No sign is suppressed in
the table.

\subsection{Every entry of all four original generators}

For a side $Z\in\{V,W\}$ concatenate the displayed diagram
words of an original tabloid $w$. Let its unpaired ones from
left to right be $a_1,\ldots,a_\varphi$, and its unpaired twos
from left to right be $c_1,\ldots,c_\varepsilon$. Define
$\mathcal F_{jZ}(w)$ by changing the column containing $a_j$
with its lowering column Kashiwara operator; define
$\mathcal E_{jZ}(w)$ by raising the column containing
$c_{\varepsilon-j+1}$. An extra internal arc gives zero.
To specify the column operator without an implicit convention:
the other side's weight and internal-arc number stay fixed,
the chosen side's number of twos changes by one, and the chosen
side's original internal-arc number stays fixed. Among the
columns of that same height there is exactly one such column.
If the intermediate changed diagram has more internal arcs
than the old diagram, the result is zero instead.

The original ambient action is exactly
\begin{equation}\label{eq:full-original-action}
 F_Zt_w=\sum_{j=1}^{\varphi_Z(w)}[j]t_{\mathcal F_{jZ}(w)},
 \qquad
 E_Zt_w=\sum_{j=1}^{\varepsilon_Z(w)}[j]t_{\mathcal E_{jZ}(w)}.
\end{equation}
Terms whose changed word is zero are omitted. Apply
\eqref{eq:full-tab-projection} to every summand and multiply by
$\sigma_np^{-d_n}$ to obtain the column for $g b_n$.
Equivalently, for $g=F_Z$ or $E_Z$ its $(m,n)$ entry is the
finite sum
\begin{equation}\label{eq:full-all-matrix-entries}
 (G_g)_{mn}=
 \sum_{\substack{j:\,\pi(t_{g_{jZ}(w_n)})
       =\epsilon_jp^{k_j}b_m}}
       [j]\sigma_n\epsilon_jp^{k_j-d_n}.
\end{equation}
This supplies every matrix entry, including all cancellations,
signs, and possible intermediate denominators.

The nonzero entry counts are
\[
\begin{array}{c|rrrr}
g&F_V&E_V&F_W&E_W\\\hline
\#\operatorname{supp}(G_g)&1998&2083&1998&2083.
\end{array}
\]
Every resulting entry
lies in $A$: finite integral Laurent addition and division by
$p$ with remainder zero verify this directly. The original
projection commutes with the action on every one of the $36864$
ambient words for each of the four generators, including all
$18652$ words with zero image. These $147456$ equalities prove
descent throughout the original relation space, independently
of the choice of representatives.

Let $v_n$ and $w_n'$ be the numbers of $V$ and $W$ ones in
the fifteen-letter diagram of $b_n$. The four diagonal
$\mathfrak{gl}_2$ weight generators have eigenvalues
$q^{v_n},q^{15-v_n},q^{w_n'},q^{15-w_n'}$. Thus the two
Cartan ratios are
$K_Vb_n=q^{2v_n-15}b_n$ and
$K_Wb_n=q^{2w_n'-15}b_n$; their central products are $q^{15}$.
The exact matrices satisfy
\[
 [E_Z,F_Z]=\frac{K_Z-K_Z^{-1}}{q-q^{-1}},\qquad
 [g_V,h_W]=0\quad(g,h\in\{E,F\}),
\]
and $K_ZE_ZK_Z^{-1}=q^2E_Z$,
$K_ZF_ZK_Z^{-1}=q^{-2}F_Z$, with trivial cross Cartan
commutators. These equalities are checked as equality of
integer Laurent coefficient dictionaries in every matrix
column. They also follow by descent from
\eqref{eq:full-original-action}, the source graphical action
on the original tensor product. Thus these are the full
original $U_q(\mathfrak{gl}_2\oplus\mathfrak{gl}_2)$
generators, with their exact quotient identification.

\subsection{Integral lattices, divided powers, and specialization}

Define $L_A=\bigoplus_{n=0}^{1259}Ab_n\subset Y/R_K$.
The exact Laurent integrality just proved makes this an
integral generator lattice. For each $g\in\{E_V,F_V,E_W,F_W\}$
define the matrices recursively by
\[
 G_g^{(0)}=I,\qquad G_g^{(m)}=G_gG_g^{(m-1)}/[m].
\]
This retains the original divided-power denominator until the
division is established. The certificate computes each
coefficient by exact integral Laurent division and verifies
zero remainder for every $m=1,\ldots,10$. Both raising and
lowering matrices have $G_g^{(9)}\ne0$ and $G_g^{(10)}=0$.
Consequently all divided powers preserve $L_A$ and every
higher divided power is zero. The nonzero-entry counts are
\[
\begin{array}{c|rrrrrrrrrrr}
m&0&1&2&3&4&5&6&7&8&9&10\\\hline
F_V,F_W&1260&1998&2314&2100&1489&892&447&185&56&12&0\\
E_V,E_W&1260&2083&2425&2082&1486&907&452&186&56&12&0.
\end{array}
\]
Over $K$ the recurrence is $G_g^m/[m]!$, by induction.
The divided-power identities therefore hold over $K$; both
sides are integral matrices, and the injection $A\hookrightarrow K$
proves the same identities over $A$. This proves the full
divided-power integral action, not only the four ordinary
generator identities.

There is also an exact map from the original ambient integral
word span. Put $Y_A=\bigoplus_{w\in\mathcal W}At_w$ and
use \eqref{eq:full-tab-projection} to define
$\pi_A:Y_A\to L_A$. Its exponents are nonnegative, so this
map is defined over $A$. For each $n$ choose the first word
$w(n)$ in the fixed order whose projection is
$\epsilon(n)b_n$, with exponent zero. The exhaustive table
contains such a word for every $n$. Define
$s_A(b_n)=\epsilon(n)t_{w(n)}$. Then
$\pi_As_A(b_n)=b_n$ for each basis vector, proving
$\pi_As_A=I$.

Here is the entire kernel, with no unspecified torsion. For a
nonpivot word $w$, set
\[
 h_w=\begin{cases}
 t_w,&\pi_A(t_w)=0,\\
 t_w-\epsilon_w\epsilon(n_w)p^{k_w}t_{w(n_w)},
       &\pi_A(t_w)=\epsilon_wp^{k_w}b_{n_w}.
 \end{cases}
\]
The $35604$ vectors $h_w$ form an $A$-basis of $\ker\pi_A$.
Their image is zero by the displayed formula. They are
independent because in a relation among them the coefficient
of each distinct nonpivot original word is its own coefficient.
Subtracting $\sum_w a_wh_w$ from any kernel vector leaves
only pivot words; their images are signed distinct basis
vectors of $L_A$, so all remaining coefficients vanish.
This proves spanning. Therefore
\begin{equation}\label{eq:full-integral-split}
 0\longrightarrow A^{35604}\xrightarrow{h}Y_A
 \xrightarrow{\pi_A}L_A\longrightarrow0
\end{equation}
is split exact as a sequence of $A$-modules, with cokernel
zero. Its section is a specified linear section; no
equivariance of that section is asserted. The projection
itself is equivariant by the $147456$ descent identities.
The ambient graphical module has its original divided-power
action, so its kernel is stable under that action as well.
The section and all kernel vectors are listed in
\texttt{raw\_integral\_projection.json}.

Upon $q=1$ over $\mathbb Z$, $p=2$, and the same formulas
with $p^{k_w}=2^{k_w}$ remain a split exact sequence of free
abelian groups. In particular the natural kernel base-change
map is an isomorphism. Every full generator and divided-power
matrix specializes by replacing $\sum a_kq^k$ by $\sum a_k$;
\texttt{specialized\_generators.json} gives all these exact
integer entries. Thus this specialization retains every
coefficient and has rank $1260$ over $\mathbb Z$.
This is a statement about the actual image lattice inside
the field quotient. It does not replace the unsaturated
integral span of the original binomial relations by its
saturation without identifying the map.

For clarity the corresponding saturation map can also be
stated exactly. Let $P_A\subset Y_A$ be generated by the
integral contextual rows and by every nonintegral row
$pt_w+t_{w'}$ (or $t_w$ when its correction is zero).
Then $P_A\subset\ker\pi_A$ and
$K\otimes_AP_A=K\otimes_A\ker\pi_A=R_K$.
The induced map $Y_A/P_A\to L_A$ is onto and has exactly
the torsion kernel $(\ker\pi_A)/P_A$.
Here ``torsion'' follows because tensoring that kernel with
$K$ is zero: each of its elements is annihilated by some
nonzero element of $A$. It need not be zero. A fully explicit
presentation of this kernel is obtained by expressing each
displayed generator of $P_A$ in the basis $h_w$: its
coefficient at $h_w$ is just its coefficient at that nonpivot
word, as proved in the elimination above. This preserves
the original relation presentation and states the exact
map to its torsion-free image lattice.

There is a sharper exact answer after a specified
localization. Put $B=A[p^{-1}]$, retaining $2$ as a
nonunit. In this ring every tree scalar in the original
relation graph is a unit. Elimination along a spanning
tree is therefore an invertible change of presentation
coordinates over $B$. For each connected component it
leaves one root and relations
$(1-\epsilon p^k)\,\mathrm{root}=0$ for its non-tree
edges, plus $\mathrm{root}=0$ if there is an explicit
zero relation. The exhaustive graph has $1260$ components
with no defect and no zero relation, $4147$ with an
explicit zero relation, and $40$ further components
whose entire nonzero cycle ideal is generated by
$1-(-1)=2$. Consequently the actual localized
presentation has the exact decomposition
\begin{equation}\label{eq:full-localized-two-torsion}
\begin{aligned}
 B^{36864}/P_B&\simeq B^{1260}\oplus(B/(2))^{40},\\
 \ker(B^{36864}/P_B\longrightarrow B\otimes_A L_A)
       &\simeq(B/(2))^{40}.
\end{aligned}
\end{equation}
The map to the image lattice is the previously specified
original-word projection; the displayed isomorphism is
given by the spanning-tree coordinates. This identifies
the kernel of the actual map, with every original
component listed in
\texttt{quotient\_audit/}\allowbreak
\texttt{relation\_graph\_certificate.json}.

One of the forty torsion generators has the explicit
original representative
\[
u=(123)(123)(234)(23)(23)(1)(1).
\]
The three original integral rows give, in succession,
\[
\begin{aligned}
u&=-v,\qquad
 v=(123)(123)(234)(32)(23)(1)(1),\\
v&=w,\qquad
 w=(123)(123)(234)(23)(32)(1)(1),\\
w&=u.
\end{aligned}
\]
The first row uses the adjacent $(3,2)$ relation, the
second commutes $(32)(23)$, and the third uses
$(32)(1)=(23)(1)$. Thus $2u=0$. Its component has
six original vertices, no explicit zero relation, and
cycle ideal exactly $(2)$; the spanning-tree proof shows
that $u\ne0$ and generates one $B/(2)$ summand.
Indeed $B/(2)$ is the localization of
$\mathbb F_2[q,q^{-1}]$ at its nonzero element
$q+q^{-1}$, so it is not the zero ring. This proves
nonvanishing, not merely the absence of a proof that
the torsion class vanishes.

\subsection{The original rational specialization bridge}

We first specify the source's map $s_X$, which is a map of
quotients. For any original NST $t_w$, let $d(w)$ be its
number of invariant column pairs, including pairs of
different heights. Such a pair is detected by two arcs
between its two columns. Define
$Y_A^S=\bigoplus_w A(-p^{-1})^{d(w)}t_w$ inside $Y$.
Let $R_A^S$ be the source's sum of the correspondingly
scaled adjacent submodule lattices: scale each contextual
submodule basis element $\widetilde t$ by
$(-p^{-1})^{\deg(\widetilde t)}$, with its ambient filtered
degree retained. Thus the source has the canonical surjection
\[
 s_X:Y_A^S/R_A^S\longrightarrow\pi(Y_A^S),\qquad
 [y]\longmapsto\pi(y).
\]
The identification of its localized relation lattice
with $P_B$ requires more than equality of their
$K$-spans. We give both containments explicitly.
By the source definition at lines $5620$--$5624$, a
contextual basis vector is
$g=B_0\mathbin{\heartsuit_g}\widetilde C
\mathbin{\heartsuit_g}D_0$, with $B_0,D_0$ ranging
over all original NST bases of the two contexts.
Its actual lattice generator at lines $5881$--$5889$
is $(-p^{-1})^{\deg(g)}g$. For an integral local
$\widetilde C$, the equality at lines $5680$--$5685$
is exact in $Y$: $g$ equals the ambient based product
of the listed binomial with its two contexts. Thus it
is the corresponding integral graph row with its
original coefficients $+1,-1$ or $+1,+1$.
For a nonintegral local $\widetilde C$, the equality
at lines $5773$--$5785$ is also exact in $Y$:
$g=t_w+p^{-1}t_{w'}$, or $g=t_w$ when the contextual
correction is zero. This uses the explicit contextual
corollary, not the weaker congruence for nonintegral
vectors in lines $5680$--$5685$. The cleared graph
row is respectively $pg$ or $g$. Therefore each
source lattice generator is $\pm p^j$ times its
corresponding graph row for a specified integer $j$,
and conversely each graph row is $\pm p^{-j}$ times
that source lattice generator. Every such coefficient
is a unit in $B$; no factor $2$ is inverted in either
containment. The contexts and local rows are
exhaustive in both definitions. It follows that
$B\otimes_A R_A^S=P_B$ as actual submodules of
$B\otimes_A Y_A^S=B^{36864}$.
Hence the localization of this exact source presentation is the presentation
in \eqref{eq:full-localized-two-torsion}; its localized
map $s_X$ has precisely the same nonzero kernel
$(B/(2))^{40}$. In particular the source remark at
lines $5901$--$5908$ asserting equality of its two
integral forms cannot be used as an injectivity proof
for these literal prescribed contextual lattices.
The displayed original cycle gives the exact integral
obstruction, while the rational specialization below
constructs their isomorphism over its stated base.
In this shape $\pi(Y_A^S)=L_A$: applying the explicit
projection to each scaled word gives zero or
$\epsilon(-1)^{d(w)}p^{k_w-d(w)}b_{n_w}$.
All nonzero exponents are $0$ or $1$; their counts are
$17728$ and $484$. Conversely the scaled honest
representatives include every $b_n$. This proves both
containments without assuming that $s_X$ is injective over $A$.

Its rational specialization is injective, by the following
additional argument. Every edge of the original relation
graph has scalar $\epsilon p^k$, and every cycle product
has that form. At $p=2$ all edges remain nonzero, and
$\epsilon 2^k=1$ if and only if $\epsilon=1$ and $k=0$.
These are exactly the conditions for $\epsilon p^k=1$ in
$K$. Hence the $1260$ nonzero components remain nonzero,
and every zero component remains zero. More specifically,
the independently reconstructed graph has $4147$ components
with an explicit zero relation and $40$ further components
killed by a cycle of scalar $-1$. The specialized relation
space has rank $36864-1260=35604$ over $\mathbb Q$.
At $q=1$ the rescalings $(-p^{-1})^d=(-1/2)^d$ are
nonzero rational numbers, so they change no rational span.
Thus $\mathbb Q\otimes_A(Y_A^S/R_A^S)$ has dimension
$1260$. The specialization of $s_X$ is a surjection onto
$\mathbb Q\otimes_A L_A$, also of dimension $1260$ by
\eqref{eq:full-integral-split}; hence it is an isomorphism.
Its formula is the original quotient map, not a dimension
comparison without a specified morphism.

The rational specialization of $L_A$ is the restriction of
the classical Schur module of row shape $(7,5,3,0)$, through
an explicit quotient map, as follows. Realize each original
column at $q=1$ in $(V\otimes W)^{\otimes h}$ by its source
canonical column formula; take the ordinary exterior quotient
$e_h:(V\otimes W)^{\otimes h}\to\bigwedge^h(V\otimes W)$
on each column factor. The restriction to the original
nonstandard exterior subspace is a $\mathrm{GL}(V)\times
\mathrm{GL}(W)$ isomorphism over $\mathbb Q$. To see
injectivity, define $\operatorname{Alt}_h:
\bigwedge^hX\to X^{\otimes h}$ by the unnormalized signed
sum over all permutations. Directly from the alternating
relations, $e_h\operatorname{Alt}_h=h!I$ on $\bigwedge^hX$
and $\operatorname{Alt}_he_h=h!I$ on the alternating
tensor subspace. The source subspace at $q=1$ is exactly
this alternating image, since the two tensor permutation
actions on $V,W$ combine to the permutation action on
$X=V\otimes W$. Thus $e_h$ has inverse
$\operatorname{Alt}_h/h!$ on these subspaces. This proof retains
the factor $h!$ and works over $\mathbb Q$, as stated.

Tensor these seven maps, keeping the original ordered column
factors. Apply this tensor product linearly to the actual
specialized based-product vector $t_w$ and its tensor
expansion. A based product at $q=1$ need not be the ordinary
tensor product of its column labels; replacing each label
independently and concatenating would not specify this map.
The source defines its two-column submodules as
the exact $U^\tau$ summands removed by the classical
two-column Schur quotient: this is its defining prescription
at source lines $5015$--$5021$, not an inference from the
global dimension $1260$. Here the ambient two-column module
is itself multiplicity-free as a $U^\tau$-module; the
involution $\tau$ interchanges $V,W$ and distinguishes
otherwise repeated restrictions. We verify the small cases
to justify uniqueness of the prescribed subspace. Write
$D_V=\det V,D_W=\det W$. First
$Y_{(1,1)}$ has the three distinct $U^\tau$ summands of
dimensions $9,6,1$, obtained from
$(S^2V\oplus D_V)\otimes(S^2W\oplus D_W)$ by pairing
the two cross terms. Next
\[
 \bigwedge^2X=(S^2V\otimes D_W)
                 \oplus(D_V\otimes S^2W).
\]
Tensor with $X$ and use
$S^2V\otimes V=S^3V\oplus D_VV$ and its $W$ version.
The two cross outer terms pair to dimension $16$; the
two copies of $D_VV\otimes D_WW$ split into the two
distinct $\tau$ signs, each of dimension $4$.
Thus $Y_{(2,1)}$ is multiplicity-free with dimensions
$16,4,4$, the two last summands having opposite signs.
For $Y_{(2,2)}$ use
$S^2V\otimes S^2V=S^4V\oplus D_VS^2V\oplus D_V^2$
and its $W$ version in the previous exterior decomposition.
The outer terms pair to dimensions $10$ and $6$;
the two cross terms of dimension $9$ split with opposite
$\tau$ signs, and the two invariant lines also split
with opposite signs. These six summands, of dimensions
$10,6,9,9,1,1$, are pairwise inequivalent as
$U^\tau$-modules. The displayed tensor decompositions
follow by the explicit symmetric product and alternating
contraction maps for two-dimensional $V,W$: in each
case the weights are the disjoint strings with respective
highest weights $(a+b-j,j)$, for $0\le j\le\min(a,b)$;
their dimensions sum to $(a+1)(b+1)$, proving exhaustion.
Finally $\bigwedge^3X\simeq D_VD_W\otimes X$ with
the corresponding overall $\tau$ sign. Hence
$Y_{(3,2)}$ and $Y_{(3,3)}$ are character twists of
$Y_{(1,2)}$ and $Y_{(1,1)}$ and remain
multiplicity-free. This proves the ambient fact needed
here. Selecting the source's specified classical
quotient summands therefore uniquely determines its
kernel as the sum of all remaining ambient summands.
The local bases in the two tables are the explicit bases
given for this prescription in the source's Figures
\texttt{straightening11} through \texttt{straightening33}.
Consequently the columnwise exterior maps take each such
specialized local submodule to the classical adjacent
straightening submodule. This is the exact local
identification used in the source proof at lines
$6013$--$6018$; all its contexts are the unchanged tensor
factors. Quotient by their sum on both sides. The relation
graph argument above identifies the specialized relation
space with the entire specialized kernel, so the resulting
map is
\[
 s_{(7,5,3)}:\mathbb Q\otimes_A L_A
   \longrightarrow
 \left(\bigwedge^3X\right)^{\otimes3}
 \otimes\left(\bigwedge^2X\right)^{\otimes2}
 \otimes X^{\otimes2}\big/\text{classical adjacent relations}.
\]
Its formula takes any representative in $Y_A$ to the class
of its columnwise exterior images. Two representatives
differ by the kernel, whose rational specialization was
proved equal to the specialized relation space; hence the map is
well defined and injective. It is onto because the
columnwise exterior maps are onto. The target is the
classical Schur module $S_{(7,5,3)}X$ by its exterior-column
presentation, and each map is equivariant. This is the
explicit rational specialization correspondence used for
the source character calculation. No literal identification
of a crystal-highest tabloid with an $E$-annihilated vector
is used here.

\subsection{The full source map to the retained interval}

In the full basis the original vectors are exactly
\[
 T=b_{126},\qquad T_1=b_{146},\qquad T_2=b_8,\qquad Q=b_{144}.
\]
These equalities include their original signs and the factors
$-p^{-1}$ in $T_2,Q$. The height-two columns in $T_1,Q$
commute by an integral table row when their order differs
from the representative table; no scalar correction occurs.

Form the directed support graph of the full $F_V,F_W$
matrices. Let $\mathcal D$ be the descendants of $126$,
and let $\mathcal J$ be the descendants which cannot reach
$144$. Exhaustive traversal of these explicit matrices gives
\[
 |\mathcal D|=420,\quad |\mathcal J|=416,\quad
 \mathcal D\setminus\mathcal J=\{126,146,8,144\}.
\]
Both sets are closed under outgoing lowering edges: a
successor remains a descendant, and a successor able to
reach $144$ would provide such a path from its predecessor.
Their $A$-spans are therefore stable under the integral
negative-Borel action and its divided powers. The exact
projection $\bigoplus_{n\in\mathcal D}Ab_n\to M_A$ kills
the $416$ basis vectors in $\mathcal J$ and takes the four
remaining vectors, in the displayed order, to
$e_0,e_1,e_2,e_3$. Its kernel is exactly the displayed
$416$-dimensional span and its cokernel is zero.

Projecting every column of the full matrices gives
\[
 F_V=\begin{pmatrix}0&0&0&0\\
 [7]-[3]&0&0&0\\{}[8]&0&0&0\\0&[6]&-[3]&0
 \end{pmatrix},\qquad F_W=0,
\]
and the original Cartan ratios
$\operatorname{diag}(q^9,q^7,q^7,q^5)$ and $q^3I$.
The full divided square has the actual entry
\[
 (F_V^{(2)})_{144,126}
 =\frac{[6]([7]-[3])-[3][8]}{[2]}
 =q^{10}-q^4-q^{-4}+q^{-10}.
\]
The relation with the earlier restricted raw tableau lattice
is especially explicit. Its four raw words have images
$(e_0,e_1,-pe_2,-pe_3)$, and therefore their inclusion matrix
in $M_A$ is $\operatorname{diag}(1,1,-p,-p)$, with the
previously proved cokernel $(A/(p))^2$. Enlarging the raw
domain to $Y_A$ supplies other original words. The section
chosen above gives the following four preimages, each with
coefficient $+1$:
\[
\begin{array}{c|l}
T&(123)(123)(124)(12)(12)(1)(1)\\
T_1&(123)(123)(124)(12)(23)(1)(1)\\
T_2&(123)(123)(123)(12)(12)(2)(3)\\
Q&(123)(123)(124)(12)(13)(2)(3).
\end{array}
\]
For the last two rows the active $V$-string is $C_{12}=(2)(3)$;
the middle formula in \eqref{eq:full-contextual-relations}
gives $(2)(3)=-p^{-1}(4)(1)$ with these exact contexts.
This proves their stated images directly. The $T_1$ and
$Q$ column-order adjustments use the integral $(2,2)$
commutations already listed.

Put $Y_{\mathcal D,A}=\pi_A^{-1}
(\bigoplus_{n\in\mathcal D}Ab_n)$. Inclusion of the four
old raw words into this module and the composite
$Y_{\mathcal D,A}\xrightarrow{\pi_A}
\bigoplus_{n\in\mathcal D}Ab_n\twoheadrightarrow M_A$
give the commutative square
\[
\begin{array}{ccc}
A^4_{\rm old\ raw}&\longrightarrow&Y_{\mathcal D,A}\\
\big\downarrow{\scriptstyle\operatorname{diag}(1,1,-p,-p)}
 &&\big\downarrow{\scriptstyle\pi_A\ {\rm followed\ by\ quotient}}\\
M_A&\xrightarrow{\quad I\quad}&M_A.
\end{array}
\]
Its right arrow is onto by the four additional preimages.
Its left arrow retains its exact nonzero cokernel; enlarging
the domain does not erase that restricted-lattice defect.
Thus the earlier four-vector construction is recovered by
an exact morphism from the full original generator module.
Its integral endpoint map, conductor factorization and
specialization kernel/cokernel remain the maps already
proved in Section~\ref{sec:gct-generator-interval}.
The extra $1256$ basis vectors here are original source
vectors with their original raising and lowering actions;
they are not chosen to repair the trace of a hypothetical
four-dimensional quantum-group representation.

\section{The full original character and specialization-compatible projectors}
\label{sec:gct-full-character}

The original row shape in this section is $\nu=(7,5,3)$, its column
lengths are $(3,3,3,2,2,1,1)$, and its column-count shape is
$[0,3,2,2]$. None of these data is replaced by a smaller partition.
Set $A=\mathbb Z[q,q^{-1}]$, $K=\mathbb Q(q)$, and
$X_A=A\,\mathrm{pNSTC}(\nu)$ inside the source quotient
\[
 X=\check X_\nu=\check Y_{(3,3,3,2,2,1,1)}/
 \sum_{i=1}^{6}\check Y_{\triangleright^i(3,3,3,2,2,1,1)}.
\]
Write $[t]=(q^t-q^{-t})/(q-q^{-1})$ for every integer $t$,
with $[0]=0$ and $[-t]=-[t]$. The source parameter is retained.

We use the actual two-row theorem of the April 2013 GCT-IV source,
whose TeX labels are \texttt{t main canonical basis},
\texttt{p highest weight straightened NST}, and
\texttt{c basis correct size}. Its hypothesis is
$\dim V=\dim W=2$, with arbitrary $\nu$ of length at most four;
it is not a theorem constructing all conjectured nonstandard tensor
fat cells. At this scope the positive honest classes are an $A$-basis,
the divided Chevalley generators preserve that lattice, and the
crystal cells have the highest weights obtained by the source's
straightened-class rule. These are facts about the quotient just
displayed. Crystal-highest class representatives are not being
identified with simultaneous Chevalley kernels in $X$.

\subsection{An exhaustive count in the original column shape}

An original straightened highest class is specified by its invariant
record $(0,i_3,i_2,i_1)$ and invariant-free columns
\[
 (123)^A(124)^B(134)^C(12)^D(13)^E(23)^G(1)^f.
\]
Here powers count consecutive columns, not commutative products.
All signs and scalings remain those of the positive source class:
an underlying tabloid of degree $h=i_3+i_2+i_1$ has the source
factor $(-1/[2])^h$, followed by its prescribed orientation sign.
Counting a class does not change this scalar.
The inserted height-three invariant is $(234)(123)$,
the prescribed height-two invariant is $(24)(13)$, and the
height-one invariant is $(4)(1)$. They are positioned by the
source's straightening convention; their respective two-sided
weights are $(3,3;3,3)$, $(2,2;2,2)$, and $(1,1;1,1)$.
Thus the complete list of constraints in the retained shape is
\begin{gather}
 i_3,i_2,i_1\in\{0,1\},\quad
 A+B+C=3-2i_3,\quad D+E+G=2-2i_2,\quad f=2-2i_1,
 \label{eq:gct-honest-shape}\\
 A,B,C,D,E,G,f\geq0,\quad BC=0,\quad
 D>0\Longrightarrow C=0,\quad E>0\Longrightarrow B\leq1,
 \nonumber\\
 G\leq1,\quad G=1\Longrightarrow f>0,\quad
 B\leq f+E,\quad C\leq f+D.\label{eq:gct-honest-inequalities}
\end{gather}
The source's highest-class classification and invariant insertion
give one positive honest highest class for each solution, and its
unique straightened representative recovers these integers. Hence
there is neither an omitted class nor multiplicity from choosing
different representatives of one class.

Put $k=3i_3+2i_2+i_1$. The original alphabet
$1=(1,1)$, $2=(1,2)$, $3=(2,1)$, $4=(2,2)$ gives, by counting
each letter in these columns,
\begin{align}
 a&=2A+2B+C+2D+E+G+f+k,\nonumber\\
 b&=2A+B+2C+D+2E+G+f+k,\label{eq:gct-highest-letter-weights}\\
 \lambda&=(a,15-a),\qquad \mu=(b,15-b).\nonumber
\end{align}
The invariants add the displayed $k$ to both coordinates of both
weights. Consequently the total degree stays fifteen on each side.
The inequalities ensure $a,b\geq8$.

Here is the entire finite enumeration. Each row is one solution of
\eqref{eq:gct-honest-shape}--\eqref{eq:gct-honest-inequalities};
the last column is $(2a-14)(2b-14)$.
\begin{center}\small
\begin{tabular}{ccccccr}
\toprule
$(i_3,i_2,i_1)$ & $(A,B,C)$ & $(D,E,G)$ & $f$ & $\lambda$ & $\mu$ & dimension\\
\midrule
$(0,0,0)$ & $(1,0,2)$ & $(0,1,1)$ & $2$ & $(8,7)$ & $(11,4)$ & 16\\
$(0,0,0)$ & $(1,0,2)$ & $(0,2,0)$ & $2$ & $(8,7)$ & $(12,3)$ & 20\\
$(0,0,0)$ & $(1,2,0)$ & $(1,0,1)$ & $2$ & $(11,4)$ & $(8,7)$ & 16\\
$(0,0,0)$ & $(1,2,0)$ & $(2,0,0)$ & $2$ & $(12,3)$ & $(8,7)$ & 20\\
$(0,0,0)$ & $(2,0,1)$ & $(0,1,1)$ & $2$ & $(9,6)$ & $(11,4)$ & 32\\
$(0,0,0)$ & $(2,0,1)$ & $(0,2,0)$ & $2$ & $(9,6)$ & $(12,3)$ & 40\\
$(0,0,0)$ & $(2,1,0)$ & $(0,1,1)$ & $2$ & $(10,5)$ & $(10,5)$ & 36\\
$(0,0,0)$ & $(2,1,0)$ & $(0,2,0)$ & $2$ & $(10,5)$ & $(11,4)$ & 48\\
$(0,0,0)$ & $(2,1,0)$ & $(1,0,1)$ & $2$ & $(11,4)$ & $(9,6)$ & 32\\
$(0,0,0)$ & $(2,1,0)$ & $(1,1,0)$ & $2$ & $(11,4)$ & $(10,5)$ & 48\\
$(0,0,0)$ & $(2,1,0)$ & $(2,0,0)$ & $2$ & $(12,3)$ & $(9,6)$ & 40\\
$(0,0,0)$ & $(3,0,0)$ & $(0,1,1)$ & $2$ & $(10,5)$ & $(11,4)$ & 48\\
$(0,0,0)$ & $(3,0,0)$ & $(0,2,0)$ & $2$ & $(10,5)$ & $(12,3)$ & 60\\
$(0,0,0)$ & $(3,0,0)$ & $(1,0,1)$ & $2$ & $(11,4)$ & $(10,5)$ & 48\\
$(0,0,0)$ & $(3,0,0)$ & $(1,1,0)$ & $2$ & $(11,4)$ & $(11,4)$ & 64\\
$(0,0,0)$ & $(3,0,0)$ & $(2,0,0)$ & $2$ & $(12,3)$ & $(10,5)$ & 60\\
$(0,0,1)$ & $(2,1,0)$ & $(0,2,0)$ & $0$ & $(9,6)$ & $(10,5)$ & 24\\
$(0,0,1)$ & $(2,1,0)$ & $(1,1,0)$ & $0$ & $(10,5)$ & $(9,6)$ & 24\\
$(0,0,1)$ & $(3,0,0)$ & $(0,2,0)$ & $0$ & $(9,6)$ & $(11,4)$ & 32\\
$(0,0,1)$ & $(3,0,0)$ & $(1,1,0)$ & $0$ & $(10,5)$ & $(10,5)$ & 36\\
$(0,0,1)$ & $(3,0,0)$ & $(2,0,0)$ & $0$ & $(11,4)$ & $(9,6)$ & 32\\
$(0,1,0)$ & $(1,0,2)$ & $(0,0,0)$ & $2$ & $(8,7)$ & $(10,5)$ & 12\\
$(0,1,0)$ & $(1,2,0)$ & $(0,0,0)$ & $2$ & $(10,5)$ & $(8,7)$ & 12\\
$(0,1,0)$ & $(2,0,1)$ & $(0,0,0)$ & $2$ & $(9,6)$ & $(10,5)$ & 24\\
\bottomrule
\end{tabular}
\end{center}

\begin{center}\small
\begin{tabular}{ccccccr}
\toprule
$(i_3,i_2,i_1)$ & $(A,B,C)$ & $(D,E,G)$ & $f$ & $\lambda$ & $\mu$ & dimension\\
\midrule
$(0,1,0)$ & $(2,1,0)$ & $(0,0,0)$ & $2$ & $(10,5)$ & $(9,6)$ & 24\\
$(0,1,0)$ & $(3,0,0)$ & $(0,0,0)$ & $2$ & $(10,5)$ & $(10,5)$ & 36\\
$(0,1,1)$ & $(3,0,0)$ & $(0,0,0)$ & $0$ & $(9,6)$ & $(9,6)$ & 16\\
$(1,0,0)$ & $(0,0,1)$ & $(0,1,1)$ & $2$ & $(8,7)$ & $(10,5)$ & 12\\
$(1,0,0)$ & $(0,0,1)$ & $(0,2,0)$ & $2$ & $(8,7)$ & $(11,4)$ & 16\\
$(1,0,0)$ & $(0,1,0)$ & $(0,1,1)$ & $2$ & $(9,6)$ & $(9,6)$ & 16\\
$(1,0,0)$ & $(0,1,0)$ & $(0,2,0)$ & $2$ & $(9,6)$ & $(10,5)$ & 24\\
$(1,0,0)$ & $(0,1,0)$ & $(1,0,1)$ & $2$ & $(10,5)$ & $(8,7)$ & 12\\
$(1,0,0)$ & $(0,1,0)$ & $(1,1,0)$ & $2$ & $(10,5)$ & $(9,6)$ & 24\\
$(1,0,0)$ & $(0,1,0)$ & $(2,0,0)$ & $2$ & $(11,4)$ & $(8,7)$ & 16\\
$(1,0,0)$ & $(1,0,0)$ & $(0,1,1)$ & $2$ & $(9,6)$ & $(10,5)$ & 24\\
$(1,0,0)$ & $(1,0,0)$ & $(0,2,0)$ & $2$ & $(9,6)$ & $(11,4)$ & 32\\
$(1,0,0)$ & $(1,0,0)$ & $(1,0,1)$ & $2$ & $(10,5)$ & $(9,6)$ & 24\\
$(1,0,0)$ & $(1,0,0)$ & $(1,1,0)$ & $2$ & $(10,5)$ & $(10,5)$ & 36\\
$(1,0,0)$ & $(1,0,0)$ & $(2,0,0)$ & $2$ & $(11,4)$ & $(9,6)$ & 32\\
$(1,0,1)$ & $(0,1,0)$ & $(0,2,0)$ & $0$ & $(8,7)$ & $(9,6)$ & 8\\
$(1,0,1)$ & $(0,1,0)$ & $(1,1,0)$ & $0$ & $(9,6)$ & $(8,7)$ & 8\\
$(1,0,1)$ & $(1,0,0)$ & $(0,2,0)$ & $0$ & $(8,7)$ & $(10,5)$ & 12\\
$(1,0,1)$ & $(1,0,0)$ & $(1,1,0)$ & $0$ & $(9,6)$ & $(9,6)$ & 16\\
$(1,0,1)$ & $(1,0,0)$ & $(2,0,0)$ & $0$ & $(10,5)$ & $(8,7)$ & 12\\
$(1,1,0)$ & $(0,0,1)$ & $(0,0,0)$ & $2$ & $(8,7)$ & $(9,6)$ & 8\\
$(1,1,0)$ & $(0,1,0)$ & $(0,0,0)$ & $2$ & $(9,6)$ & $(8,7)$ & 8\\
$(1,1,0)$ & $(1,0,0)$ & $(0,0,0)$ & $2$ & $(9,6)$ & $(9,6)$ & 16\\
$(1,1,1)$ & $(1,0,0)$ & $(0,0,0)$ & $0$ & $(8,7)$ & $(8,7)$ & 4\\
\bottomrule
\end{tabular}
\end{center}


For a fully specified exhaustive procedure, choose the eight invariant
records, then $A=0,\ldots,3-2i_3$,
$B=0,\ldots,3-2i_3-A$, and set $C=3-2i_3-A-B$.
Choose $D=0,\ldots,2-2i_2$ and
$E=0,\ldots,2-2i_2-D$, and set $G=2-2i_2-D-E$.
Reject exactly the failed inequalities above, set $f=2-2i_1$,
and evaluate \eqref{eq:gct-highest-letter-weights}.
Every nonnegative solution is visited exactly once because all its
free coordinates are visited exactly once. This proves the enumeration
without a numerical interpolation or a guess about saturation.

Let $g_{ab}$ be the number of these solutions with highest first
coordinates $(a,b)$. Rows and columns in the following matrix have
the order $12,11,10,9,8$:
\begin{equation}
 (g_{ab})=
 \begin{pmatrix}
 0&0&1&1&1\\0&1&2&3&2\\1&2&4&4&3\\
 1&3&4&4&2\\1&2&3&2&1
 \end{pmatrix}.
 \label{eq:gct-full-multiplicity-matrix}
\end{equation}
All other highest weights have multiplicity zero. The sum of the
entries is $48$, and
\begin{equation}
 \dim_K X=\sum_{a,b=8}^{12}g_{ab}(2a-14)(2b-14)=1260.
 \label{eq:gct-full-dimension}
\end{equation}

\subsection{The full character and the classical specialization map}

Let $x_1,x_2$ and $y_1,y_2$ record the two original torus weights.
The complete character, retaining every coordinate, is
\begin{equation}
 \operatorname{ch}X=
 \sum_{a,b=8}^{12}g_{ab}
 \left(\sum_{i=0}^{2a-15}x_1^{a-i}x_2^{15-a+i}\right)
 \left(\sum_{j=0}^{2b-15}y_1^{b-j}y_2^{15-b+j}\right).
 \label{eq:gct-full-character-polynomial}
\end{equation}
Each crystal component is a rectangle with one basis vector at
every listed weight, so the source highest-class count proves this
formula for the original quotient, not for an unrelated space of
the same dimension.

For reference, the dimensions $w_{uv}$ of the dominant weight
spaces $(u,15-u;v,15-v)$, with rows and columns again in order
$12,11,10,9,8$, are
\begin{equation}
 (w_{uv})=
 \begin{pmatrix}
 0&0&1&2&3\\0&1&4&8&11\\1&4&11&19&25\\
 2&8&19&31&39\\3&11&25&39&48
 \end{pmatrix}.
 \label{eq:gct-full-dominant-weight-matrix}
\end{equation}
The other weights satisfy $w_{uv}=w_{15-u,v}=w_{u,15-v}$,
and vanish outside $3\leq u,v\leq12$. Formula
\eqref{eq:gct-full-character-polynomial} proves these reflection
identities, since each of its two intervals is invariant under the
corresponding reflection. On the dominant chamber the exact inverse
to rectangle summation is
\[
 g_{ab}=w_{ab}-w_{a+1,b}-w_{a,b+1}+w_{a+1,b+1}.
\]
Indeed $w_{ab}=\sum_{c\geq a,d\geq b}g_{cd}$ there, and both
successive finite differences cancel precisely the terms outside
$(c,d)=(a,b)$.

The source specialization is a diagram of actual morphisms
\[
 X\ \longleftarrow\ X_A\ \longrightarrow\
 X_A/(q-1)X_A\ \longrightarrow\
 \mathbb Q\otimes_{\mathbb Z}(X_A/(q-1)X_A)
 \xrightarrow{\ \sigma\ }
 S_{(7,5,3)}(V_{\mathbb Q}\otimes W_{\mathbb Q}).
\]
The first arrow is inclusion followed by scalar extension; the second
is the quotient. We make explicit the source's two integral presentations.
With $Y_A^S$ its rescaled tabloid lattice and $M_A^S$ the sum of its
six specified integral relation lattices, the source first forms
$X'_A=Y_A^S/M_A^S$. The inclusion
$M_A^S\subseteq (K M_A^S)\cap Y_A^S$ gives the actual surjection
\[
 s_X:X'_A\longrightarrow
 Y_A^S/((K M_A^S)\cap Y_A^S)=X_A,\qquad [y]\longmapsto[y].
\]
The source's specialization theorem is stated first for $X'_A$.
Its remark at TeX lines 5901--5908 also asserts an integral
identification, but that stronger assertion is not used here:
an integral relation presentation can retain torsion that the
honest image lattice does not retain. We now prove exactly the
rational specialization map required for the displayed diagram.
Let $\mathbb Q_1$ denote $\mathbb Q$ as an $A$-algebra through
$q\mapsto1$. Right exactness gives a surjective module map
\[
 \overline s_X=\mathbb Q_1\otimes_A s_X:
 \mathbb Q_1\otimes_A X'_A\twoheadrightarrow
 \mathbb Q_1\otimes_A X_A.
\]
The original contextual relation graph has $36864$ tabloid
vertices; its specialization at $q=1$, $p=[2]=2$, has relation
rank $35604$, so its quotient has dimension $1260$ over
$\mathbb Q$. All original rescaling factors are signed powers
of $p$, hence units over $\mathbb Q_1$. They do not change
that rational row span. This is the explicit source relation
specialization computation, recorded in
\texttt{specialization\_bridge.json} with its graph hash and
the original scaled-image check. Equivalently it is the
rational dimension statement in the source's two-row theorem.
Thus the domain of $\overline s_X$ has dimension $1260$.
Its codomain has dimension $1260$ because the original honest
lattice has its proved $1260$-element $A$-basis.

The inverse can be given directly, without an inverse to $s_X$
over $A$. For each original positive honest basis vector $b$,
choose its displayed positive rescaled tabloid representative
$\widetilde b\in Y_A^S$. Define
\[
 j:\mathbb Q_1\otimes_A X_A\longrightarrow
 \mathbb Q_1\otimes_A X'_A,\qquad
 1\otimes b\longmapsto1\otimes[\widetilde b].
\]
This is a well-defined linear map on a basis, and
$\overline s_X j=1$ by the quotient definitions. These $1260$
lifted vectors are independent, since their images are the
honest basis; the dimension calculation proves that they are
also a basis of the domain. It follows that
$j\overline s_X=1$. As the inverse of the surjective module
map $\overline s_X$, $j$ intertwines the specialized original
actions. In particular its value is independent of the choice
of representative. This proves the exact rational isomorphism
and makes no injectivity assertion about $s_X$ over $A$.

The last arrow in the displayed diagram is $\sigma=\tau j$,
where $\tau$ is the source theorem's rational Schur-quotient
map, described intrinsically as follows. The nonstandard exterior
algebra is the quotient of the tensor algebra on the four original
letters by its nonstandard symmetric-square relations. At $q=1$
these relations are the ordinary symmetric-square relations. The
identity on the four letters therefore induces an isomorphism of
the specialized exterior algebra with the ordinary exterior algebra.
Apply its degree maps in the original seven tensor factors. The
source's adjacent-column relation spaces specialize to the adjacent
Schur relation spaces. Thus the tensor map carries the sum of the
six relation spaces onto the sum of the six classical relation
spaces, and it induces $\tau$ on their rational quotients. Applying the
inverse tensor map gives an inverse on these quotients. Naturality
on the four-letter tensor algebra shows that these maps intertwine
the $\mathfrak{gl}_2\oplus\mathfrak{gl}_2$ actions and their
polynomial $\mathrm{GL}_2\times\mathrm{GL}_2$ actions over
$\mathbb Q$. This is the relation proved in the source's
specialization theorem; it does not assert specialization of arbitrary
elements of $K\otimes_A X_A$ that have poles at $q=1$.

In particular
\[
 \operatorname{ch}X=s_{(7,5,3)}
 (x_1y_1,x_1y_2,x_2y_1,x_2y_2).
\]
There is also a separate exact verification of the entire polynomial.
Enumerate all weakly increasing rows of lengths $7,5,3$ on
$\{1,2,3,4\}$, retaining a triple of rows exactly when its columns
strictly increase. This gives all $1260$ semistandard tableaux
exactly once. If their contents are $(c_1,c_2,c_3,c_4)$, their
monomial is
\[
 x_1^{c_1+c_2}x_2^{c_3+c_4}
 y_1^{c_1+c_3}y_2^{c_2+c_4}.
\]
The resulting weight multiplicities equal
\eqref{eq:gct-full-character-polynomial} at every weight.
The separate checker
\texttt{independent\_check/check\_character.py} additionally evaluates the Jacobi--Trudi
determinant in exact integer polynomials. The source's set bijection
between honest classes and ordinary semistandard tableaux is not
being used as a weight-preserving linear map; $\sigma$ is the
module morphism that establishes the classical connection.

\subsection{Explicit intertwiners into the actual source module}

We prove the needed splitting and then give maps whose input
operators are the full original source generators. For
$U_q(\mathfrak{sl}_2)$ use
$KEK^{-1}=q^2E$, $KFK^{-1}=q^{-2}F$ and
$EF-FE=(K-K^{-1})/(q-q^{-1})$.
An integer $h$ denotes a weight on which $K$ acts by $q^h$.

\begin{lemma}\label{lem:gct-explicit-string-splitting}
Every finite-dimensional type-one integrable $K$-module for this
algebra is the direct sum of its highest-weight strings. A string
from a vector $z$ with $Ez=0$ and $Kz=q^n z$ has $n\geq0$,
basis $F^{(i)}z$, $0\leq i\leq n$, and actions
\[
 F F^{(i)}z=[i+1]F^{(i+1)}z,\qquad
 E F^{(i)}z=[n-i+1]F^{(i-1)}z.
\]
\end{lemma}
\begin{proof}
Repeatedly applying the displayed commutator gives
$EF^jz=[j][n-j+1]F^{j-1}z$. This follows by induction:
moving $E$ past one additional $F$ adds the next Cartan term,
and the finite Laurent identity
$[j+1][n-j]-[j][n-j+1]=[n-2j]$ supplies precisely that term.
Let $d$ be the last nonzero power of $F$ on $z$; it exists by
integrability. Applying $E$ to $F^{d+1}z=0$ gives
$[d+1][n-d]F^dz=0$. Over $K$, $[t]=0$ for an integer $t$
only when $t=0$, so $d=n\geq0$. Distinct string vectors have
distinct weights and are independent. Dividing by $[i]!$ yields
the claimed formulas.

For completeness, splitting is proved by induction on dimension.
Choose the largest weight $n$ occurring in a nonzero module.
Its entire weight space consists of highest vectors. The strings
from a basis of that space are independent: a relation at weight
$n-2i$ can be raised by $E^i$, which multiplies each top vector
by the same nonzero scalar. They span a submodule $N$, a sum
of copies of the length-$n$ string. The quotient has smaller
dimension and no weight $n$, so by induction it is a sum of
strings with highest weights $0\leq m<n$.
Lift each quotient highest vector to a vector $z$ of weight $m$.
Then $Ez\in N_{m+2}$. On every copy of the length-$n$ string,
$E:N_m\to N_{m+2}$ is surjective whenever the target occurs:
its coefficient in the formulas above is nonzero. Subtract a
preimage in $N_m$ to obtain a highest lift. Its length-$m$
string maps isomorphically onto the chosen quotient string by
the preceding formulas. All these lifted strings together map
isomorphically to the quotient, so their sum is disjoint from
$N$ and splits the original module. This completes the induction.
\end{proof}

For $a,b\in\{8,9,10,11,12\}$ put $n=2a-15$, $m=2b-15$ and
define the actual source subspace
\[
 H_{ab}=X_{(a,15-a;b,15-b)}\cap\ker E_V\cap\ker E_W.
\]
Let $L_{(a,15-a)}$ be the standard quantum $\mathfrak{gl}_2$
module with basis $e_i$, $0\leq i\leq n$, of weight
$(a-i,15-a+i)$, with $Fe_i=[i+1]e_{i+1}$ and
$Ee_i=[n-i+1]e_{i-1}$; the out-of-range terms are zero.
Define $L_{(b,15-b)}$ in the same way, with basis $f_j$.
In the next tensor product $H_{ab}$ is used as a multiplicity
vector space with trivial quantum-group action, not with an
action inherited from its embedding in $X$. Then the map
\begin{equation}
 \Phi:\bigoplus_{a,b=8}^{12}
 L_{(a,15-a)}\otimes L_{(b,15-b)}\otimes H_{ab}
 \longrightarrow X,\qquad
 e_i\otimes f_j\otimes z\longmapsto F_V^{(i)}F_W^{(j)}z
 \label{eq:gct-actual-string-intertwiner}
\end{equation}
is an isomorphism of the full quantum
$\mathfrak{gl}_2\oplus\mathfrak{gl}_2$ modules.
To prove this, first apply Lemma
\ref{lem:gct-explicit-string-splitting} to the $V$ generators.
The commuting $W$ generators act on each $V$ highest subspace.
Apply the lemma to that action. The two successive string maps
are exactly \eqref{eq:gct-actual-string-intertwiner}.
The lemma verifies both raising and lowering operators, and
the weights verify all four original diagonal generators and
both total degree-fifteen central characters. This proves the
intertwiner and its bijectivity. The character comparison now
gives $\dim_K H_{ab}=g_{ab}$; no choice of purported highest
source basis vectors was made.

Here is an explicit inverse requiring only the original operators.
For either colour, put
\[
 C=(q-q^{-1})^2FE+qK+q^{-1}K^{-1}.
\]
Its commutators with $K,E,F$ are zero. For example its commutator
with $E$ is
\[
 -(q-q^{-1})E(q^2K-q^{-2}K^{-1})
 +q(q^2-1)EK+q^{-1}(q^{-2}-1)EK^{-1}=0;
\]
the $F$ calculation follows by the same defining relations.
On a highest vector of weight $n$ it acts by
$c_n=q^{n+1}+q^{-n-1}$, hence by that scalar on the whole string.
All source weights $h$ are odd. Consequently the operator
\begin{equation}
 B=FE+\sum_h[(h+1)/2]^2\,\mathrm{pr}_h
       =\frac{C-2}{(q-q^{-1})^2}
 \label{eq:gct-integral-casimir}
\end{equation}
preserves $X_A$. The equality holds because
$q^{h+1}+q^{-h-1}-2=(q^{(h+1)/2}-q^{-(h+1)/2})^2$.
Here $\mathrm{pr}_h$ is projection to the original weight
subspace of the source weight basis, so it is $A$-linear.
The apparent quotient in \eqref{eq:gct-integral-casimir} has
already been replaced by its exact integral operator formula;
it is not a limit used to discard any term.

For $n\in\mathcal N=\{1,3,5,7,9\}$ set
$\beta_n=[(n+1)/2]^2$ and
\begin{equation}
 \Pi_n(B)=\prod_{t\in\mathcal N\setminus\{n\}}
 \frac{B-\beta_t}{\beta_n-\beta_t}.
 \label{eq:gct-source-isotypic-projector}
\end{equation}
The denominators are nonzero since at $q=1$ the five eigenvalues
are respectively $1,4,9,16,25$. The string decomposition proves
that these are orthogonal idempotents with sum one, acting as
identity exactly on the strings of highest ratio weight $n$.

For an original vector $x$ of weight $(u,15-u;v,15-v)$ set
$i=a-u$, $j=b-v$. Keep a pair $(a,b)$ exactly when
$0\leq i\leq n$ and $0\leq j\leq m$. Its inverse component is
\begin{equation}
 e_i\otimes f_j\otimes
 \frac{E_V^{(i)}E_W^{(j)}
       \Pi_n(B_V)\Pi_m(B_W)x}
 {\displaystyle \genfrac{[}{]}{0pt}{}{n}{i}_q
                   \genfrac{[}{]}{0pt}{}{m}{j}_q},
 \qquad \genfrac{[}{]}{0pt}{}{n}{i}_q=\frac{[n]!}{[i]![n-i]!}.
 \label{eq:gct-source-intertwiner-inverse}
\end{equation}
To check both compositions, apply this formula to
$F_V^{(r)}F_W^{(s)}z$ from each string in
\eqref{eq:gct-actual-string-intertwiner}. The projectors kill all
different highest-weight pairs. Weight equality forces $i=r,j=s$
in the surviving pair, and repeated raising gives
$E_V^{(i)}F_V^{(i)}z=\genfrac{[}{]}{0pt}{}{n}{i}_qz$ and its $W$ analogue.
Thus the formula recovers $z$ and every component, proving both
inverse identities. This is a formula on the actual source
generators, with concrete finite products; it is not a character-only
identification.

\subsection{The precise integral localization and specialization}

Let $S$ be the multiplicative set generated by $[1],\ldots,[9]$
and $\beta_n-\beta_t$ for distinct $n,t\in\mathcal N$,
and write $R=S^{-1}A$. The exact Laurent identity
\[
 [r]^2-[s]^2=[r-s][r+s]
\]
follows by expanding the four powers after multiplication by
$(q-q^{-1})^2$, including its sign when $r<s$.
For distinct $r,s\in\{1,2,3,4,5\}$ its two nonzero factors
are signed members of $[1],\ldots,[9]$. Thus this displayed
localization equals $A[([1]\cdots[9])^{-1}]$; the Casimir
denominators were retained and have now been proved to be units.
The lattice $X_A$ remains the original
lattice; this explicitly stated scalar extension is needed for
the following splitting. Every operator in
\eqref{eq:gct-actual-string-intertwiner} and
\eqref{eq:gct-source-intertwiner-inverse} is defined over $R$:
the divided generators already preserve $X_A$, the integral
$B$ does also, and the remaining denominators are units in $R$.
Put $H_{ab,R}=X_R{}_{(a,15-a;b,15-b)}\cap\ker E_V\cap\ker E_W$.
The formulas give mutually inverse $R$-linear maps
\begin{equation}
 \bigoplus_{a,b}L_{(a,15-a),R}\otimes
 L_{(b,15-b),R}\otimes H_{ab,R}\ \cong\ X_R.
 \label{eq:gct-localized-lattice-splitting}
\end{equation}
Indeed the numerator in the inverse formula lies in this highest
kernel over $K$ by the proved inverse, and its annihilation by
both $E$ holds over $R$ because $X_R$ is torsion-free. Both
composition identities also descend from $K$ by torsion-freeness.
Moreover
$H_{ab,R}=\Pi_n(B_V)\Pi_m(B_W)X_R{}_{(a,15-a;b,15-b)}$:
one inclusion follows from the inverse formula at $i=j=0$,
and a highest vector has the indicated two Casimir eigenvalues,
which proves the other inclusion. It is therefore a finite
projective $R$-module of rank $g_{ab}$.

Evaluation at $q=1$ sends $[j]$ to $j$ and the distinct
Casimir differences to differences of $1,4,9,16,25$.
These nonzero differences have prime factors only in
$\{2,3,5,7\}$, and $[1],\ldots,[9]$ already require all four
primes. The universal property of localization therefore gives
\[
 R/(q-1)R=\mathbb Z[1/210].
\]
Tensoring the two inverse maps in
\eqref{eq:gct-localized-lattice-splitting} with this ring keeps
them inverse. In particular the specialized highest kernel is
exactly the specialization of $H_{ab,R}$. One can check this
also directly on the strings: the coefficient of every
non-top raising step is one of $1,\ldots,9$, a unit in
$\mathbb Z[1/210]$, so simultaneous raising kills exactly
the two top positions. No additional highest vectors appear.
The specialized $B$ is
$FE+\sum_h((h+1)/2)^2\mathrm{pr}_h$ and retains its five
distinct eigenvalues. Thus these projectors specialize without
a pole, and the string intertwiners become the classical
divided-power intertwiners. After tensoring with $\mathbb Q$,
the source map $\sigma$ above relates them to the classical
Schur module. This does not claim that the splitting holds
over unlocalized $A$, or identify an integral Schur lattice
without its own lattice comparison.

\subsection{Exact locations of the three retained weights}

The weight $(12,3;9,6)$ has dimension two. Its two components
are $(a,b)=(12,10)$ at positions $(i,j)=(0,1)$ and
$(12,9)$ at positions $(0,0)$, each with multiplicity one.
The complete locations at the next two weights are as follows;
``copies'' is the dimension of the indicated actual $H_{ab}$.
\begin{center}\small
\begin{tabular}{c|cc|c|c}
\toprule
$(u,v)$ & $a$ & $b$ & $(i,j)$ & copies\\\midrule
$(11,9)$&12&10&$(1,1)$&1\\
&12&9&$(1,0)$&1\\
&11&11&$(0,2)$&1\\
&11&10&$(0,1)$&2\\
&11&9&$(0,0)$&3\\\midrule
$(10,9)$&12&10&$(2,1)$&1\\
&12&9&$(2,0)$&1\\
&11&11&$(1,2)$&1\\
&11&10&$(1,1)$&2\\
&11&9&$(1,0)$&3\\
&10&12&$(0,3)$&1\\
&10&11&$(0,2)$&2\\
&10&10&$(0,1)$&4\\
&10&9&$(0,0)$&4\\\bottomrule
\end{tabular}
\end{center}
The corresponding dimensions are $8$ and $19$. Every location
is obtained by $i=a-u$, $j=b-v$ from an actual source isotypic
projector and the proved string isomorphism. A specific original
honest vector such as $T,T_1,T_2,Q$ is located inside these
summands by applying \eqref{eq:gct-source-intertwiner-inverse}
to that vector. Merely knowing its weight does not assign it
to one summand. The original four-vector negative-Borel interval
is still the specified descendant quotient of this full
$1260$-dimensional module; these character and projector formulas
do not identify that interval with an invariant full quantum
group module.

For these four particular vectors the full original source matrices
now permit a stronger determination than their weights alone.
In their original order and with their original scalars they are
\begin{align*}
 T&=(123)(123)(124)(12)(12)(1)(1),\\
 T_1&=(123)(123)(124)(23)(12)(1)(1),\\
 T_2&=-[2]^{-1}(123)(123)(123)(12)(12)(4)(1),\\
 Q&=-[2]^{-1}(123)(123)(124)(13)(12)(4)(1).
\end{align*}
The source's integral height-two commutation relation moves
the columns $(23)(12)$ of $T_1$ to $(12)(23)$ and
the columns $(13)(12)$ of $Q$ to $(12)(13)$ with coefficient
$+1$. Their resulting canonical positive representatives,
including $-[2]^{-1}$ for $T_2,Q$, are basis indices
$126,146,8,144$ of the full original matrix certificate.
The quotient projection on the original ambient words proves
these four equalities, not merely a match of their weights.
The graphical raising rule pairs a $2$ with a later $1$ in
the original $V$ or $W$ word and sums over unpaired $2$'s.
For the four representatives its unpaired-word data are
\[
\begin{array}{c|cc}
 &V&W\\\hline
T&1^9&1^3\\T_1&1^7&1^3\\T_2&1^7&1^3\\Q&1^5&1^3
\end{array}
\]
Thus each raising sum is empty, proving
$E_VT=E_WT=E_VT_1=E_WT_1=E_VT_2=E_WT_2=E_VQ=E_WQ=0$
directly in the original module. Hence
\[
 T\in H_{12,9},\qquad T_1,T_2\in H_{11,9},\qquad
 Q\in H_{10,9}.
\]
The two middle vectors are independent because they are distinct
members of the original honest basis. In particular they span
two dimensions in the three-dimensional actual $H_{11,9}$.
The exact projector computation on all four original vectors
returns the vector itself for this listed pair and zero for
every other pair; its raising and lowering inverse checks use
the full original matrices over $\mathbb Q(q)$.
For example $F_V^{(2)}T$ remains in the isotypic summand
$(\lambda,\mu)=((12,3),(9,6))$, whereas $Q$ is in the different
summand $((10,5),(9,6))$. The original nonzero coefficient of
$Q$ in the honest-basis expansion of $F_V^{(2)}T$ is therefore
not an isotypic projection. Applying $\Pi_5(B_V)\Pi_3(B_W)$
to the complete expansion gives zero: the contribution of the
$Q$ basis term is canceled by the projections of the other
terms. This identity follows because the projectors commute
with $F_V$ and kill $T$ for that highest pair. It preserves
the original coefficient and explains its exact relation to
the full irreducible decomposition.

\section{The retained inverse fibre as an explicit determinantal family}
\label{sec:fc-determinantal}

All rings in the algebraic statements are over a specified field $k$ of
characteristic zero. The bit statements use $k=\mathbb Q$ and signed binary
integers. No operation changes the coefficients of the original map.
The third source coordinate is named $w$; the renaming from the source's
$t$ is the identity $(x,y,t)\mapsto(x,y,w)$. The source expressions were
read in full from the parent satellite and the independent fluid extension;
their paths and SHA--256 hashes are in \texttt{source\_hashes.json}.
Every calculation used below is reproved here and replayed by the included
exact checker.

\subsection{The original map and its localized coordinate-ring isomorphism}

Put $A=k[a,b,c]$, $B=k[x,y,w]$, and define $F^*:A\to B$ by
\begin{align*}
 F^*(a)=F_1&=(1+xy)^3w+y^2(1+xy)(4+3xy),\\
 F^*(b)=F_2&=y+3x(1+xy)^2w+3xy^2(4+3xy),\\
 F^*(c)=F_3&=2x-3x^2y-x^3w.
\end{align*}
The determinant retains the source orientation $(x,y,w)$ and target
orientation $(a,b,c)$ and equals $-2$. Here is a direct proof. On $x\ne0$
set $v=1+xy$, $h=2-3xy-x^2w$, $A_0=v^2+v-v^3h$ and
$B_0=4v+2-3v^2h$. Their inverse coordinate formulas are
$y=(v-1)/x$ and $w=(5-3v-h)/x^2$, so these are actual coordinates
on the stated open set. Substitution in the original expressions yields
$F=(A_0/x^2,B_0/x,xh)$ and
\[
 \det\frac{\partial(x,v,h)}{\partial(x,y,w)}=-x^3,
 \qquad
 \det\frac{\partial F}{\partial(x,v,h)}
 =x^{-3}\det\begin{pmatrix}
 -2A_0&2v+1-3v^2h&-v^3\\
 -B_0&4-6vh&-3v^2\\h&0&1
 \end{pmatrix}=2x^{-3}.
\]
Expansion of the last numerator gives
\[
 (4-6vh)(-2v^2-2v+3v^3h)
 +(2v+1-3v^2h)(4v+2-6v^2h)=2.
\]
Multiplication proves $\det DF=-2$ on $x\ne0$. Both sides are
polynomials, and their equality in the localization of the domain $B$
proves equality in $B$ as well.

Retain the exact polynomial and derivative
\[
 P=cr^3-2r^2+br-2a,\qquad P'=3cr^2-4r+b,\qquad \alpha=P'/2.
\]
Define the $A$-algebra
\[
 E=\bigl(A[r]/(P)\bigr)[1/P'].
\]
There is an isomorphism of $A$-algebras $\Phi:E\to B[1/x]$ given by
\[
 (a,b,c,r,1/P')\longmapsto
 (F_1,F_2,F_3,y+1/x,x/2).
\]
Its inverse $\Psi:B[1/x]\to E$ is specified on every generator by
\[
 x\longmapsto 2/P',\quad
 y\longmapsto r-P'/2,\quad
 w\longmapsto 5(P'/2)^2-3r(P'/2)-c(P'/2)^3,
 \quad x^{-1}\longmapsto P'/2.
\]
To prove these assertions before quotienting, introduce independent
coordinates $(r,\alpha,c)$ with $\alpha\ne0$ and put
\[
 H(r,\alpha,c)=
 \left(\alpha^{-1},r-\alpha,5\alpha^2-3r\alpha-c\alpha^3\right).
\]
Solving $F_3=c$ after setting $r=y+1/x$, $\alpha=1/x$ gives this
formula and its inverse exactly. Direct substitution gives
\[
 F\circ H=(r^2+r\alpha-cr^3,\ 4r+2\alpha-3cr^2,\ c).
\]
The equation $F_2=b$ says $2\alpha=P'$, and then
$2(F_1-a)=P$. Conversely those two relations give $F_1=a,F_2=b$.
They show that the displayed maps respect $P=0$, the inverse elements,
and the $A$-algebra structure. The recovery equations
$r=y+1/x$, $\alpha=1/x$ show both compositions are the identity on
all their generators. This proves the ring isomorphism, including its
excluded divisor $P'=0$.

The morphism $F^*$ is injective. Indeed $A[r]/(P)$ is isomorphic to
$k[b,c,r]$ by eliminating $a=(cr^3-2r^2+br)/2$; both substitution maps
are inverse on the generators. A nonzero
$h(a,b,c)=\sum_{j=0}^d h_j(b,c)a^j$, $h_d\ne0$, has, after this
substitution, highest $r$ coefficient $h_d(b,c)c^d/2^d$ at degree $3d$,
which is nonzero in $k[b,c]$. Thus $A$ injects into the quotient and,
since $P'\ne0$ there, into $E$ and then $B[1/x]$. Consequently it
also injects into $B$. This proves dominance without treating it as an
assumption.

\subsection{Every specialized fibre and the finite free locus}

For $(a_0,b_0,c_0)\in k^3$, set
\[
 B_{a_0,b_0,c_0}=k[x,y,w]/(F_1-a_0,F_2-b_0,F_3-c_0),
 \qquad E_{a_0,b_0,c_0}=k[r,1/P']/(P),
\]
where $P$ and $P'$ here have the specialized coefficients.
If $c_0\ne0$, then $xh=c_0$ in the fibre ring, so $x$ is a unit
with inverse $h/c_0$. The preceding isomorphism therefore proves
$B_{a_0,b_0,c_0}\simeq E_{a_0,b_0,c_0}$ on the entire fibre.
If $c_0=0$, the entire fibre ring is instead
\[
 B_{a_0,b_0,0}\simeq k\ \times\ E_{a_0,b_0,0}.
\]
Here is a proof that includes its scheme structure. Work first in
$R=k[x,y,w]/(F_1-a_0,F_2-b_0)$. The polynomials $x,h$ generate the
unit ideal since
$1=h/2+x(3y+xw)/2$. Their ideals therefore have intersection $(xh)$:
if $z=xu=hv$, multiplication by the displayed expression for $1$
shows $z\in(xh)$. The map
$R/(xh)\to R/(x)\times R/(h)$ is injective by that intersection;
given classes modulo the two ideals, their lifts combine using the
same two summands of $1$ to prove surjectivity. This proves the asserted
product decomposition. The first factor is $k$ because
$F(0,y,w)=(w+4y^2,y,0)$, with point $(0,b_0,a_0-4b_0^2)$.
In the second factor, $x(3y+xw)=2$ makes $x$ a unit and the chart
isomorphism identifies it with $E_{a_0,b_0,0}$.

Over an algebraic closure, write $P=d\prod_j(r-\rho_j)^{m_j}$,
with distinct $\rho_j$. The coprime factors yield a product of the
rings $\bar k[r]/((r-\rho_j)^{m_j})$: the elementary two-ideal
argument just given applies repeatedly. When $m_j=1$, $P'$ has
nonzero value at $\rho_j$ and is a unit in that factor. When $m_j>1$,
$P'$ belongs to its nilpotent ideal $(r-\rho_j)$, so inverting $P'$
turns that factor into the zero ring. Thus $E$ retains exactly the
simple roots, each with its reduced point. The complete finite inverse
points are
\[
 \left(2/P'(\rho_j),\rho_j-P'(\rho_j)/2,
 5(P'(\rho_j)/2)^2-3\rho_jP'(\rho_j)/2-c_0(P'(\rho_j)/2)^3\right)
\]
for $m_j=1$, together with the boundary point when $c_0=0$.
The original polynomial, its multiplicities, and its excluded roots
remain recorded by $(P,P')$ even when the localization removes a factor.

The coefficient discriminant is exactly
\[
 D=4(b^2-cb^3+18abc-16a-27a^2c^2).
\]
For direct verification, expansion of the determinant
\[
 \det\begin{pmatrix}
 c&-2&b&-2a&0\\0&c&-2&b&-2a\\
 3c&-4&b&0&0\\0&3c&-4&b&0\\0&0&3c&-4&b
 \end{pmatrix}=-cD
\]
gives the stated polynomial. This determinant is the resultant because
its matrix represents $(u,v)\mapsto uP+vP'$ on the bases of
polynomials of degrees less than $2$ and $3$, respectively, with the
indicated descending coefficient order. Over a field it vanishes
exactly when $P,P'$ share a nonconstant factor: a nonzero kernel gives
$uP=-vP'$ and the coprime case forces $P\mid v$, impossible at
$\deg v<3$; a common factor gives a kernel vector by dividing both
polynomials by it. This also proves the simple-root interpretation.

Let $A_U=A[1/(cD)]$. Then $A_U[r]/(P)$ is free with basis
$(1,r,r^2)$. Polynomial division by the unit leading coefficient $c$
gives a spanning remainder of degree at most two, and a nonzero multiple
of $P$ has degree at least three, proving independence. No replacement
of $P$ is made here. The original relation gives the multiplication matrix
\[
 C=\begin{pmatrix}
 0&0&2a/c\\1&0&-b/c\\0&1&2/c
 \end{pmatrix},\qquad
 N=3cC^2-4C+bI_3.
\]
Its columns are exactly the coefficients of multiplying the ordered
basis by $r$ and $P'$, respectively. Direct determinant expansion gives
$\det N=-D/c$, so
$N^{-1}=-(c/D)\operatorname{adj}(N)$. This explicit inverse proves that
$P'$ is a unit in the quotient, and hence
\[
 B\otimes_A A_U\simeq A_U[r]/(P).
\]
This is a finite free rank-three morphism. Its unramified behavior can
also be proved without omitting the derivative: for a square-zero ideal
$I$ of an $A_U$-algebra $T$, a root $r_0$ in $T/I$ has a lift
$\widetilde r$ in $T$, and the unique root lift is
$\widetilde r-P(\widetilde r)/P'(\widetilde r)$.
The denominator is a unit because its reduction is; Taylor expansion
has no terms of degree two in $I$. If two lifts differ by $e\in I$,
subtraction gives $P'(\widetilde r)e=0$, so $e=0$. This is the
unique infinitesimal lifting property, establishing the finite etale
claim on precisely $U$.

At $c_0=0$, the polynomial still has degree two and its quadratic
discriminant is $b_0^2-16a_0$, while $D=4(b_0^2-16a_0)$. The
boundary point remains necessary. In particular the fibre at
$(-1/4,0,0)$ consists exactly of
\[
 (1,-3/2,13/2),\quad(-1,3/2,13/2),\quad(0,0,-1/4).
\]
At $c_0\ne0$ the cubic is a triple-root polynomial exactly at
$b_0=4/(3c_0)$, $a_0=4/(27c_0^2)$, with root $2/(3c_0)$.
Indeed comparing $c_0(r-\rho)^3$ with all four coefficients gives
these three equations and conversely verifies them. Its entire finite
fibre is empty since its only root has $P'=0$ and $c_0\ne0$ excludes
$x=0$.

\subsection{Generic degree, the absence of a rational inverse, and nonfiniteness}

In $k[a,b,c,r]$, $(P)$ is a prime ideal since its quotient was explicitly
identified with the domain $k[b,c,r]$. Thus $P$ is irreducible there.
It is primitive as a polynomial over $k[a,b,c]$ in $r$, since its
coefficient $-2$ is a unit. The primitive-polynomial argument proves
irreducibility over $K=k(a,b,c)$: clearing denominators in a proposed
factorization and cancelling the common prime factors of the coefficients
would produce a factorization in the unique-factorization ring
$k[a,b,c][r]$. To justify that cancellation, the product of two primitive
polynomials is primitive, since reduction modulo any prime leaves two
nonzero polynomials over a domain and hence a nonzero product. This
proves the required primitive-polynomial argument.

The chart identities now give
\[
 k(x,y,w)=k(F_1,F_2,F_3)(y+1/x),\qquad
 [k(x,y,w):k(F_1,F_2,F_3)]=3.
\]
Both containments in the first equality follow from the displayed
inverse formulas; irreducibility of the degree-three polynomial proves
the degree. A rational inverse would make the two fields equal, so none
exists. There is also no rational section over a dense target open set:
if its $x$ coordinate were identically zero then its $F_3$ coordinate
would be zero, contrary to the generic variable $c$; otherwise its
$r=y+1/x$ would be an element of $K$ satisfying the irreducible cubic.
These proofs rule out those exact rational maps and retain the
rank-three algebraic correspondence constructed above.

The morphism $F$ is not finite. Define the following homomorphism to
$k[t,t^{-1}]$ by its values on the original coordinates:
\[
 x=t^{-1},\quad y=-3t/2,\quad w=13t^2/2.
\]
Direct substitution gives $(F_1,F_2,F_3)=(-t^2/4,0,0)$.
If $B$ were finite over $A$, then multiplication by $x$ on finitely
many $A$-module generators would give a monic relation for $x$:
write $xm_i=\sum_j a_{ij}m_j$ and multiply the matrix $xI-(a_{ij})$
by its adjugate; its monic determinant annihilates every generator and
therefore $1\in B$. Under the displayed homomorphism such a relation
would become
\[
 t^{-m}+\sum_{j=0}^{m-1}q_j(-t^2/4,0,0)t^{-j}=0.
\]
Multiplication by $t^m$ gives a polynomial with constant term $1$,
while every summand in the sum is divisible by $t$. This is impossible.
Thus nonfiniteness is proved directly from the original map and is
compatible with its finite free restriction to $U$.

\subsection{The projective binary cubic restores the entire boundary sheet}

There is a stronger global presentation retaining the $x=0$ sheet as a
projective root. Define
\[
 Z=\{cR^3-2R^2T+bRT^2-2aT^3=0\}\subset\mathbb P^1_A,
 \qquad [R:T]=[1+xy:x].
\]
The last expression defines a morphism on the entire source because
$1=(1+xy)-yx$ proves its two homogeneous coordinates generate the
unit ideal. Substitution of $F_1,F_2,F_3$ in the homogeneous equation
gives zero: on $x\ne0$ it is $x^3P(y+1/x)=0$, and equality in the
localization of the source polynomial domain proves the polynomial
identity everywhere.

Let $Z^{\rm s}$ be the open simple-root locus: on $T\ne0$ its
coordinates are $r=R/T$ and it is defined by $P=0$, $P'\ne0$; on
$R\ne0$ use the distinct root coordinate $h=T/R$ and put
\[
 Q(h)=c-2h+bh^2-2ah^3,\qquad d=Q'(h)=-2+2bh-6ah^2.
\]
Its simple-root condition is $d\ne0$. On the overlap $h=1/r$,
$Q(h)=h^3P(1/h)$ and, on that zero scheme, $d=-hP'(1/h)$.
Since $h$ is a unit on the overlap, the two simple-root conditions
agree exactly.

The morphism from the source to $Z$ is an isomorphism onto $Z^{\rm s}$.
On the first chart this is the already proved inverse-root isomorphism.
Here is the entire second chart, including $x=0$:
\[
 x=-2h/d,\qquad y=b-3ah,\qquad
 w=-ad^3/8-y^2d^2/4+3y^2d/2.
\]
To prove it, first keep $a,b,h$ independent and define $d,x,y,w$ by
these displayed formulas with $d\ne0$. Substitution and expansion give
\[
 1+xy=-2/d,\qquad
 F_1=a,\qquad F_2=b,\qquad F_3=2h-bh^2+2ah^3=c-Q(h).
\]
For example $1+xy=(d-2hb+6ah^2)/d=-2/d$ directly. For the other
formulas one may insert that identity into the original $F_i$:
with $v=-2/d$, the first is
$v^3w+y^2v(1+3v)=a$, since $4+3xy=1+3v$;
the second is
$y+3xv^2w+3xy^2(1+3v)=y+3ha=b$.
The third follows by $F_3=x(5-3v-x^2w)$, inserting the displayed
$w$ and $y=b-3ah$, and expanding to $2h-bh^2+2ah^3$.
Thus the formulas define the inverse map on $Q=0$, $d\ne0$.
Conversely on the source chart $v=1+xy\ne0$, set $h=x/v$,
$a=F_1$, $b=F_2$, $c=F_3$. The polynomial identity gives $Q=0$;
the two original equations for $F_1,F_2$ give
$y=b-3ah$ by subtracting $3hF_1$ from $F_2$.
Consequently
$1+xy=1+x(b-3ah)=1+hv(b-3ah)$ gives
$v(1-bh+3ah^2)=1$, equivalently $d=-2/v$.
Then $x=hv=-2h/d$, and solving the original $F_1=a$ for $w$
gives exactly the stated $w$. These recover every generator, proving
both compositions. The source charts $x\ne0$ and $v\ne0$ cover it;
the two inverse maps agree on their overlap because both are inverse
to the same forward coordinates. This proves the global isomorphism.
At $h=0$, $Q=0$ requires $c=0$, $d=-2$, and the formulas give
$(x,y,w)=(0,b,a-4b^2)$. Thus the boundary point is retained as
$[R:T]=[1:0]$, with no discarded sheet.

The projective scheme $Z$ is finite locally free of rank three over
$A$. An elementary affine cover proves this. Evaluate its binary cubic
at the three constant points $[j:1]$, $j=0,1,-1$, obtaining
\[
 f_0=-2a,\quad f_1=c-2+b-2a,\quad f_{-1}=-c-2-b-2a.
\]
They generate the unit ideal since $f_1+f_{-1}-2f_0=-4$.
On the base open set $f_j\ne0$, the point $[j:1]$ is absent from $Z$.
Use the explicitly invertible homogeneous change
\[
 U=T,\quad V=R-jT,\qquad R=V+jU,\quad T=U.
\]
The removed point becomes $[U:V]=[1:0]$, so all of $Z$ on this
base open lies in the affine chart $V\ne0$ with $s=U/V$.
The affine coordinate formula is $R=1+js$, $T=s$; its polynomial is
\[
 c(1+js)^3-2(1+js)^2s+b(1+js)s^2-2as^3,
\]
whose leading coefficient is exactly $f_j$. It is a unit on this
base open, so division with the original polynomial gives the free
basis $1,s,s^2$ exactly as in the earlier cubic proof. The changes
are invertible projective linear maps and retain all original
coefficients. The three base opens cover, proving finite local
freeness and its rank.

The branch locus is exactly $D=0$. When $c\ne0$ this is the
resultant calculation already proved. When $c=0$, $[1:0]$ is a simple
root because $d(0)=-2$, and the remaining affine polynomial has degree
two with repeated root exactly when $b^2-16a=0$, equivalently $D=0$.
These arguments apply in the residue field at each base point and in
its algebraic closure. Consequently over $D\ne0$ every projective
root is simple, so $Z^{\rm s}=Z$ there and the original $F$ is finite
locally free of rank three on the larger open $\{D\ne0\}$, including
$c=0$. The unique root-lifting argument on the two charts proves it is
etale there. The algebra on $cD\ne0$ used previously gives its concrete
single-chart basis; that smaller open was sufficient and was not the
entire finite etale locus.

\subsection{A size-parameterized family with complete encoding}

For each integer $n\ge1$ retain independent variables
$(x_i,y_i,w_i)$ and targets $(a_i,b_i,c_i)$ for $1\le i\le n$.
Define $F^{[n]}$ to apply the preceding three original polynomials to
each block. This is a morphism $\mathbb A_k^{3n}\to\mathbb A_k^{3n}$
with coordinate homomorphism $a_i\mapsto F_1(x_i,y_i,w_i)$,
$b_i\mapsto F_2(x_i,y_i,w_i)$, $c_i\mapsto F_3(x_i,y_i,w_i)$.
Its Jacobian is block diagonal, so its determinant is exactly $(-2)^n$.
Its complete fibre algebra is the tensor product over $k$ of the
single-block fibre algebras just proved. This assertion follows by
presenting the tensor product as a polynomial ring on the disjoint sets
of generators modulo the union of their defining relations. It includes
every boundary factor at $c_i=0$.

Over $\{\prod_i D_i\ne0\}$ the block family is finite locally free
of rank $3^n$, by the preceding projective-cover proof and tensor products.
Over its subopen $U_n=\{\prod_i c_iD_i\ne0\}$, the algebra is finite free of
rank $3^n$ with the explicit basis
$\{\prod_i r_i^{e_i}:0\le e_i\le2\}$. Spanning follows by repeated
use of each original relation
$c_ir_i^3=2r_i^2-b_ir_i+2a_i$. Independence follows by taking tensor
products of the proven single-block bases: applying the corresponding
coordinate functionals to a putative relation extracts each coefficient.
The generic algebra is a domain since eliminating the $a_i$ identifies
its polynomial presentation with $k[b_i,c_i,r_i:1\le i\le n]$ before
localizing. A finite-dimensional domain over the target rational
function field is a field, because multiplication by a nonzero element
is an injective linear map and hence surjective. Thus the generic
function-field degree is also exactly $3^n$.

At the rational target with every block $(-1/4,0,0)$, the fibre contains
exactly the $3^n$ ordered choices of the three displayed rational points.
A format requiring one separate record for each point therefore has
output length at least $3^n$ symbols. In contrast, the tuple of equations
$(P_i,P_i')$, their localizations, and the specified boundary rule has
linear length in $n$ with fixed-size coefficients at this target. The
inequality follows from the mandated output format, and is not an
inference about the difficulty of obtaining that compact representation.

Here is an exact arithmetic-circuit model. A gate is binary $+,-$, or
$\times$, with fan-out allowed, and constants are $1,2,3,4$ for $F$.
Input nodes, constant nodes, and output wires are not counted as gates.
For one block the following list is a 26-gate circuit; each assignment
is exactly one gate and no listed coefficient is suppressed:
\[
\begin{array}{llll}
 g_1=xy,&g_2=1+g_1,&g_3=g_2g_2,&g_4=g_3g_2,\\
 g_5=3g_1,&g_6=4+g_5,&g_7=yy,&g_8=g_7g_6,\\
 g_9=g_2g_8,&g_{10}=g_4w,&g_{11}=g_{10}+g_9,&g_{12}=xg_3,\\
 g_{13}=g_{12}w,&g_{14}=3g_{13},&g_{15}=xg_8,&g_{16}=3g_{15},\\
 g_{17}=y+g_{14},&g_{18}=g_{17}+g_{16},&g_{19}=xx,&g_{20}=g_{19}x,\\
 g_{21}=g_{19}y,&g_{22}=3g_{21},&g_{23}=g_{20}w,&g_{24}=2x,\\
 g_{25}=g_{24}-g_{22},&g_{26}=g_{25}-g_{23}.&&
\end{array}
\]
The outputs are $(g_{11},g_{18},g_{26})$; substitution verifies all
three original formulas. Thus $F^{[n]}$ uses $26n$ gates. Its component
degrees remain $7,6,4$ per block. The tuple is uniformly constructible:
a loop on the binary index $i$ writes the same 26 instructions with
the input and gate indices of block $i$. Using binary indices, its
bit description has length $O(n\log(n+1))$, and the loop takes that
many elementary output operations.

For bit evaluation, encode each signed integer by a sign and its binary
magnitude, and assume magnitudes of the $3n$ inputs are less than $2^L$,
$L\ge1$. Every displayed intermediate has magnitude less than
$2^{7L+10}$. For example $|g_1|<2^{2L}$,
$|g_2|<2^{2L+1}$, $|g_6|<2^{2L+2}$,
$|g_{10}|<2^{7L+3}$, and $|g_9|<2^{6L+3}$;
the remaining products in the list have smaller exponent bounds, while
at most two sums require an extra bit. These inequalities also verify
the loose common bound for all gates. Grade-school multiplication on
$O(L)$ bits takes $O(L^2)$ bit operations by summing shifted partial
products; addition and subtraction take $O(L)$ operations by one pass
with carries. Evaluation therefore takes $O(n(L+1)^2)$ bit operations,
plus the circuit description or output indexing cost. This is a uniform
Boolean-time upper bound for this specified evaluation problem.

For target integers of magnitude less than $2^L$, retain the coefficient
tuples\[
 (c_i,-2,b_i,-2a_i),\qquad (3c_i,-4,b_i).
\]
Their entries have at most $L+3$ magnitude bits. Writing all polynomial
and localization equations
\[
 P_i(r_i)=0,\qquad z_iP_i'(r_i)-1=0
\]
therefore costs $O(n(L+1))$ bits and bit operations up to the index
encoding. The ring maps above identify $z_i$ with $1/P_i'$ and hence
$x_i=2z_i$ on its stated chart. Testing $c_i=0$ supplies the already
proved boundary rule. No root selection is performed by writing these
equations. When a rational root $r=p/q$, $q\ne0$, is supplied as a
certificate, with signed numerator and denominator lengths at most $M$,
integer cross multiplication verifies $P(r)=0$ and $P'(r)\ne0$ using
a fixed number of products of $O(L+M+1)$-bit integers. The coordinate
formulas then return exact rational numerators and denominators of
$O(L+M+1)$ bits in $O((L+M+1)^2)$ operations. Reduction to coprime
numerators and denominators is unnecessary: the output represents the
same rational numbers and retains the actual expression. This proves
the stated verification and lifting complexity while keeping the
selection of an algebraic embedding distinct from its full algebra.

\subsection{Exact determinant projections that retain every cubic coefficient}

A single retained polynomial has the following affine-linear matrix:
\[
 M(a,b,c,r)=\begin{pmatrix}
 r&-1&0&0\\0&r&-1&0\\0&0&r&-1\\-2a&b&-2&c
 \end{pmatrix},\qquad \det M=P.
\]
Expansion along the last row gives the four signed contributions
$-2a,br,-2r^2,cr^3$, respectively, proving the identity including signs.
Introduce one new variable $z$ and replace each constant entry by that
constant times $z$. This gives the homogeneous linear matrix
\[
 M^h=\begin{pmatrix}
 r&-z&0&0\\0&r&-z&0\\0&0&r&-z\\-2a&b&-2z&c
 \end{pmatrix},
 \quad \det M^h=cr^3-2z^2r^2+brz^2-2az^3=:P^h.
\]
Every monomial on the right has degree four; setting $z=1$ gives the
original polynomial exactly. On $z\ne0$ its relation to the affine
coordinates is the explicitly invertible substitution
$(a,b,c,r,z)\mapsto(a/z,b/z,c/z,r/z,z)$ with inverse
$(A,B,C,R,z)\mapsto(zA,zB,zC,zR,z)$, and
$P^h=z^4P(a/z,b/z,c/z,r/z)$. The divisor $z=0$ is newly added:
there $P^h=cr^3$. It is not identified with a finite original chart.

For $n$ blocks, the block-diagonal determinant $\prod_i P_i$ has
size $4n$. Its zero condition is a union of the individual hypersurfaces,
since a product in a field is zero exactly when a factor is zero.
It is not the simultaneous system $P_1=\cdots=P_n=0$ when $n\ge2$.
For example set all $r_i=b_i=c_i=0$, $a_1=0$, $a_2=1$, and for
$i>2$ set $a_i=1$. The product vanishes but $P_2=-2$.

The following determinant instead encodes every equation by a coefficient
map. Introduce independent selectors $t_1,\ldots,t_n$ and set
\[
 Q_n=\sum_{i=1}^n t_iP_i(r_i),\quad
 T_i=\begin{pmatrix}1&-r_i&0&0\\0&1&-r_i&0\\0&0&1&-r_i\\0&0&0&1\end{pmatrix},
 \quad v_i=\begin{pmatrix}-2a_i\\b_i\\-2\\c_i\end{pmatrix}.
\]
Let $T=\operatorname{diag}(T_1,\ldots,T_n)$,
$v=(v_1^{\mathsf T},\ldots,v_n^{\mathsf T})^{\mathsf T}$, and
$u=(t_1,0,0,0,\ldots,t_n,0,0,0)$. Then
\[
 \mathcal M_n=\begin{pmatrix}T&v\\-u&0\end{pmatrix},
 \qquad \det\mathcal M_n=Q_n,
 \qquad \operatorname{size}(\mathcal M_n)=4n+1.
\]
For a complete proof, write $T_i=I-N_i$. Since $N_i^4=0$,
$T_i^{-1}=I+N_i+N_i^2+N_i^3$, whose first row is
$(1,r_i,r_i^2,r_i^3)$. Thus $e_1^{\mathsf T}T_i^{-1}v_i=P_i$.
Left multiplication of $\mathcal M_n$ by the determinant-one matrix
$\left(\begin{smallmatrix}I&0\\uT^{-1}&1\end{smallmatrix}\right)$
gives $\left(\begin{smallmatrix}T&v\\0&uT^{-1}v\end{smallmatrix}\right)$.
As $\det T=1$, its determinant is $uT^{-1}v=Q_n$.
Every entry of $\mathcal M_n$ is affine-linear in the original variables
and selectors with coefficients in $\{0,1,-1,-2\}$.

For an arbitrary coefficient ring $R$, the map
$\iota:R^n\to R[t_1,\ldots,t_n]_1$,
$(f_1,\ldots,f_n)\mapsto\sum_i t_if_i$, is an $R$-module
isomorphism onto the homogeneous selector-degree-one part, with inverse
$Q\mapsto([t_1]Q,\ldots,[t_n]Q)$. Comparing the unique monomial
coefficients proves both compositions are identities. In particular
$(P_1,\ldots,P_n)=([t_1]\det\mathcal M_n,\ldots,[t_n]\det\mathcal M_n)$
as ideals in $k[a_i,b_i,c_i,r_i]$. For a specialized point over any
extension field, all fibre equations hold if and only if the determinant
is identically zero as a polynomial in the selectors. Specializing the
selectors first can lose information: when $t_1=t_2=1$, the choices
$P_1=2$, $P_2=-2$ give zero sum without either equation holding.
The coefficient map is an explicit reversible transport of the tuple,
not an asserted isomorphism between the fibre scheme and one
hypersurface in a larger space.

Replace each constant in $\mathcal M_n$ by that constant times $z$;
call the resulting homogeneous linear matrix $\mathcal M_n^h$.
Its determinant is
\[
 \det\mathcal M_n^h
 =z^{4n-4}\sum_{i=1}^n t_i
   (c_ir_i^3-2z^2r_i^2+b_ir_iz^2-2a_iz^3).
\]
To prove this without assigning an inverse at $z=0$, first localize at
$z$. The four diagonal entries of each $T_i$ become $z$, and the first
row of its inverse is $(z^{-1},r_iz^{-2},r_i^2z^{-3},r_i^3z^{-4})$.
The last vector becomes $(-2a_i,b_i,-2z,c_i)^{\mathsf T}$, and
$\det T=z^{4n}$. Applying the same determinant-one block multiplication
gives the displayed formula in $k[\text{variables},z,z^{-1}]$.
Both sides belong to the polynomial subring, whose localization is
injective, so the identity holds also at $z=0$. The factor
$z^{4n-4}$ is explicit: for $n>1$ the entire hyperplane $z=0$ is an
additional zero component of the padded determinant. At $z=1$ all
original fibre coefficients are recovered.

The determinantal projection is an exact coordinate-ring homomorphism.
With $d=4n+1$ let $k[X_{pq}:1\le p,q\le d]$ be the coordinate ring
of all $d\times d$ matrices. The homomorphism
$\vartheta_n:k[X_{pq}]\to k[a_i,b_i,c_i,r_i,t_i,z]$ sends
$X_{pq}$ to the displayed $(p,q)$ entry of $\mathcal M_n^h$.
It sends the universal determinant to the stated padded polynomial.
It is surjective: each $a_i,b_i,c_i,r_i,t_i$ is a nonzero scalar
multiple of a specified entry, and a diagonal entry gives $z$.
Its kernel is generated by the linear relations among those entries.
To prove the last claim constructively, for each target variable choose
one entry mapping to it after its indicated scalar division. Subtract
from every other $X_{pq}$ the same linear combination of these chosen
entries as its image. These linear differences lie in the kernel.
Modulo them the source is the polynomial ring on the chosen entries;
the induced map to the target polynomial ring has the inverse sending
each target variable to its chosen entry. Therefore the kernel is
exactly the ideal generated by those differences. This proves the map,
its image, its kernel and the degree $d$ of the determinant image.

Finally, Horner evaluation of the unaltered cubic uses seven gates:
$c_ir_i$, then addition of $-2$, multiplication by $r_i$, addition of
$b_i$, multiplication by $r_i$, separately $(-2)a_i$, and addition of
that last term. Computing all $t_iP_i$ and their sum uses $9n-1$ gates.
Thus $Q_n$ has an explicitly uniform linear-size arithmetic circuit and
a linear-size affine determinantal representation. The homogeneous padded
form has a linear-size determinant representation as well. These are
constructive upper bounds for a family derived from the original fibres.
They retain their generic degree $3^n$ and their original exceptional
loci. They prove no reduction from arbitrary Boolean satisfiability to
these independent blocks, and they assert no Boolean complexity equality
or separation. The concrete complexity transport established here is
the coefficient-recoverable determinant projection and the fully specified
uniform bit algorithms above.

\section{Exact invariant-theory transports of the retained cubic and node}
\label{sec:gct-invariant-transport}

All schemes, coordinate rings and representations in this section are over
$\mathbb C$. The computations are defined over $\mathbb Q$, except for the
explicitly stated algebraic cube roots in the stabilizer. The original
inverse-fibre polynomial is retained as
\[
 P_{a,b,c}(r)=cr^3-2r^2+br-2a,
 \qquad P^h_{a,b,c}(R,T)=cR^3-2R^2T+bRT^2-2aT^3.
\]
In the source inverse chart, $r=y+1/x$, $\alpha=P'(r)/2$,
$x=\alpha^{-1}$, $y=r-\alpha$ and
$w=5\alpha^2-3r\alpha-c\alpha^3$; consequently the finite chart retains
only simple roots. The constructions below keep the algebraic completion
as a separate object. Their root at infinity when $c=0$ is a homogeneous
cubic root, not a licence to discard the source's additional point
$(0,b,a-4b^2)$.

The coordinate-ring restriction used in GCT is present explicitly in the
primary treatment of Ikenmeyer and Panova~\cite[Sections~1(a) and~2(b)]
{invariantIkenmeyerPanova}. Here we derive the particular
maps and representations directly. No determinant--permanent separation
is inferred from the following fixed-dimensional calculations.

\subsection{The binary representation, its coefficient slice and its discriminant}

Let $V$ be the four-dimensional space of binary cubics
$f=A R^3+B R^2T+C RT^2+D T^3$. Write $G_{\rm bin}=\operatorname{GL}_2(\mathbb C)$,
$z=(R,T)^{\mathsf T}$, and define the left action
\[
 (g\cdot f)(z)=f(g^{\mathsf T}z).
\]
Indeed $(g\cdot(h\cdot f))(z)=f(h^{\mathsf T}g^{\mathsf T}z)
=((gh)\cdot f)(z)$. For
$g=\left(\begin{smallmatrix}p&q\\s&t\end{smallmatrix}\right)$,
$\delta=pt-qs\ne0$, expansion gives
\begin{align*}
 A_g&=Ap^3+Bp^2q+Cpq^2+Dq^3,\\
 B_g&=3Ap^2s+B(p^2t+2psq)+C(2pqt+sq^2)+3Dq^2t,\\
 C_g&=3Aps^2+B(2pst+s^2q)+C(pt^2+2sqt)+3Dqt^2,\\
 D_g&=As^3+Bs^2t+Cst^2+Dt^3.
\end{align*}
The coefficient-slice embedding and its entire coordinate-ring map are
\[
 \iota:\mathbb A^3_{a,b,c}\longrightarrow V,
 \quad(a,b,c)\longmapsto(c,-2,b,-2a),
 \qquad
 \iota^*:\mathbb C[A,B,C,D]\longrightarrow\mathbb C[a,b,c],
\]
\[
 A\longmapsto c,\quad B\longmapsto-2,\quad C\longmapsto b,
 \quad D\longmapsto-2a.
\]
It is surjective: $a=-\iota^*D/2$, $b=\iota^*C$ and $c=\iota^*A$.
Its kernel is $(B+2)$, because evaluating a polynomial at $B=-2$
is division by the monic linear polynomial $B+2$, and the remaining
substitution $(A,C,D)=(c,b,-2a)$ is an invertible linear substitution.
The original affine slice itself is not $G_{\rm bin}$-stable.

Its saturation is the explicit equivariant morphism
\[
 \Phi:G_{\rm bin}\times\mathbb A^3\longrightarrow V,
 \qquad(g,a,b,c)\longmapsto g\cdot\iota(a,b,c),
\]
where $h$ acts on the source by $(g,a,b,c)\mapsto(hg,a,b,c)$.
Its pullback sends $A,B,C,D$ to the four displayed expressions with
$(A,B,C,D)=(c,-2,b,-2a)$, in
$\mathbb C[p,q,s,t,\delta^{-1},a,b,c]$. These formulas specify the
map without eliminating any root or selecting a branch.

\begin{proposition}\label{prop:invariant-saturation}
The image of $\Phi$ is $V\setminus\{0\}$, and $\ker\Phi^*=0$.
\end{proposition}
\begin{proof}
Every slice cubic has coefficient $-2$ and is nonzero; invertible
substitution cannot send it to zero. Conversely, take a nonzero cubic
$f$. There is a vector $z_0$ with $f(z_0)\ne0$: a polynomial vanishing
at every complex vector has all its coefficients zero, by successive
one-variable interpolation. Choose $w_0$ independent of $z_0$.
Euler's homogeneous identity, obtained by differentiating
$f(\lambda z_0)=\lambda^3f(z_0)$ at $\lambda=1$, gives
$df_{z_0}(z_0)=3f(z_0)\ne0$. Consequently
$df_{z_0}(w_0+kz_0)$ is nonzero for all but one $k\in\mathbb C$.
Choose one such $k$ and replace $w_0$ by $w=w_0+kz_0$.
The coefficient of $R^2T$ in $f(Rz_0+Tw)$ is exactly $df_{z_0}(w)$.
Replace $w$ by $-2w/df_{z_0}(w)$. The two chosen vectors remain
independent and the resulting coefficient is $-2$. Thus an explicitly
specified invertible variable substitution sends $f$ to the slice;
its inverse shows $f\in\operatorname{im}\Phi$. If a polynomial lies
in $\ker\Phi^*$, it vanishes on every nonzero vector of $V$ and hence
also at zero, by restricting to each line. Successive interpolation
again makes it the zero polynomial.
\end{proof}

Define, with no division by the leading coefficient,
\[
 \Delta(A,B,C,D)=B^2C^2-4AC^3-4B^3D-27A^2D^2+18ABCD.
\]
The substitution in $\iota$ gives exactly
\begin{equation}\label{eq:invariant-disc}
 \iota^*\Delta
 =4\bigl(b^2-cb^3+18abc-16a-27a^2c^2\bigr).
\end{equation}
For completeness, factor a binary cubic over $\mathbb C$ as
$f(z)=\prod_{i=1}^3(\alpha_iR+\beta_iT)$, allowing one factor to
carry any scalar. Direct expansion of the polynomial on the right
gives
\[
 \Delta(f)=\prod_{1\le i<j\le3}
              (\alpha_i\beta_j-\alpha_j\beta_i)^2.
\]
One may verify the identity first for
$f=A\prod_i(R-r_iT)$ by expanding the squared Vandermonde product:
it is $A^4\prod_{i<j}(r_i-r_j)^2$ and expands to the displayed
coefficient polynomial. Cubics with $A\ne0$ are dense, so equality
of these polynomials holds also at $A=0$. This proves that $\Delta=0$
is precisely a repeated projective root, including roots at infinity.
Under $g\cdot f$, the coefficient column of each linear factor is
multiplied by $g$. Each of the three two-by-two determinants is
multiplied by $\det g$, proving
\begin{equation}\label{eq:invariant-disc-weight}
 \Delta(g\cdot f)=(\det g)^6\Delta(f).
\end{equation}
The coordinate-ring representation is
$(g\cdot H)(f)=H(g^{-1}\cdot f)$, so the one-dimensional span of
$\Delta$ has character $(\det g)^{-6}$, not $(\det g)^6$.

\subsection{The actual source orbit and all of its stabilizer}

At the original target $(a,b,c)=(-1/4,0,0)$ put
\[
 q_0=-2R^2T+\tfrac12T^3,
 \qquad\Delta(q_0)=16.
\]
Its three projective roots are $[1:0]$, $[1/2:1]$ and $[-1/2:1]$.
We compute its stabilizer without identifying distinct scalar cubics.
Introduce the auxiliary polynomial $h(U,V)=UV(U-V)$ and the exact
invertible matrix
\[
 M=\begin{pmatrix}1&1/2\\0&1\end{pmatrix},\qquad
 M^{-1}=\begin{pmatrix}1&-1/2\\0&1\end{pmatrix},\qquad
 q_0(z)=-2h(Mz).
\]
This is an explicit comparison with an auxiliary polynomial; $q_0$
and its factor $-2$ remain fixed. Let
\begin{align*}
 D_1&=\begin{pmatrix}1&0\\0&1\end{pmatrix},&
 D_2&=\begin{pmatrix}-1&1\\0&1\end{pmatrix},&
 D_3&=\begin{pmatrix}0&1\\1&0\end{pmatrix},\\
 D_4&=\begin{pmatrix}1&0\\1&-1\end{pmatrix},&
 D_5&=\begin{pmatrix}0&1\\-1&1\end{pmatrix},&
 D_6&=\begin{pmatrix}1&-1\\1&0\end{pmatrix}.
\end{align*}
Substitution in the three factors of $h$ gives
\[
 h(D_jz)=\varepsilon_jh(z),\qquad
 (\varepsilon_1,\ldots,\varepsilon_6)=(1,1,-1,1,-1,-1).
\]
Their fractional-linear maps are respectively
$t$, $1-t$, $1/t$, $t/(t-1)$, $1/(1-t)$ and $(t-1)/t$.
They give all permutations of $\{0,1,\infty\}$: inspection gives
six distinct permutations, and a projective transformation fixing all
three must be the identity. To prove the latter assertion, fixing
$0$ and $\infty$ makes its matrix diagonal up to scalar; fixing $1$
makes the two diagonal entries equal.

It follows that the full stabilizer $H_0$ is the following set of
eighteen explicitly determined matrices:
\begin{equation}\label{eq:invariant-stab}
 H_0=
 \left\{\left(M^{-1}\kappa D_jM\right)^{\mathsf T}:
      1\le j\le6,\quad\kappa^3\varepsilon_j=1\right\}.
\end{equation}
Indeed a stabilizer must permute the three roots. After conjugating
its transpose by $M$, it is therefore $\kappa D_j$ for exactly one
$j$ and a nonzero scalar $\kappa$. Substitution into $-2h(Mz)$ makes
the multiplier exactly $\kappa^3\varepsilon_j$, which must be one.
Conversely every listed matrix has that multiplier and stabilizes
$q_0$. Each of the six scalar equations has three different roots,
and the six projective classes are different. The homomorphism to
root permutations is explicitly $[z]\mapsto[g^{-\mathsf T}z]$:
$(gh)^{-\mathsf T}=g^{-\mathsf T}h^{-\mathsf T}$ proves its composition
law. The transposes used to compute stabilizer membership above
give its inverse permutations. Consequently the root action gives
the exact sequence
\[
 1\longrightarrow\{\omega I:\omega^3=1\}
 \longrightarrow H_0\longrightarrow\mathfrak S_3\longrightarrow1.
\]

Every two unordered triples of distinct projective points are related
by a projective transformation. Here is a constructive proof. Choose
nonzero representatives $v_1,v_2$ for the first two points; they form
a basis. Write a representative for the third as $av_1+bv_2$;
distinctness gives $a,b\ne0$. The linear isomorphism sending
$av_1$ and $bv_2$ to $(1,0)^{\mathsf T}$ and $(0,1)^{\mathsf T}$
sends the third point to $[1:1]$. Apply this construction twice and
compose one map with the other's inverse. Consequently a squarefree
binary cubic is an invertible substitution of $q_0$ times a nonzero
scalar. Absorb that scalar by a scalar matrix with its specified
cube root. We have proved
\begin{equation}\label{eq:invariant-open-orbit}
 G_{\rm bin}q_0=\{f:\Delta(f)\ne0\},\qquad
 \overline{G_{\rm bin}q_0}=V.
\end{equation}
The closure equality also follows directly because the complement
is the zero set of a nonzero polynomial. Its coordinate ring and
the coordinate ring of the orbit are, respectively,
\[
 \mathbb C[\overline{G_{\rm bin}q_0}]=\mathbb C[A,B,C,D],
 \qquad
 \mathbb C[G_{\rm bin}q_0]=\mathbb C[A,B,C,D,\Delta^{-1}].
\]
The localization description is the defining coordinate ring of the
principal affine open set: it is represented by the hypersurface
$z\Delta-1=0$, with projection inverse $f\mapsto(f,\Delta(f)^{-1})$.

We can compute the invariant functions as well. Scalar matrices act
on $V$ by $f\mapsto\lambda^3f$. If a polynomial is $G_{\rm bin}$-invariant,
decompose it into its homogeneous terms. Comparing the distinct powers
$\lambda^{3d}$ in its value on $\lambda^3f$ makes every positive-degree
term zero. Hence $\mathbb C[V]^{G_{\rm bin}}=\mathbb C$.

In contrast,
\begin{equation}\label{eq:invariant-SL2}
 \mathbb C[V]^{\operatorname{SL}_2}=\mathbb C[\Delta].
\end{equation}
To prove this equality explicitly, first note that
$\det(H_0)=\{\zeta:\zeta^6=1\}$. Formula
\eqref{eq:invariant-disc-weight} gives containment. Scalar stabilizers
give all cube roots of unity as determinants, and the matrix with
$D_2$ and $\kappa=1$ gives determinant $-1$, so the reverse containment
holds. If $f_1$ and $f_2$ are squarefree and have equal discriminant,
choose $g$ with $g f_1=f_2$. Then $(\det g)^6=1$. The stabilizer of
$f_1$ is conjugate to $H_0$, so it has an element $h$ with
$\det h=(\det g)^{-1}$. Now $gh\in\operatorname{SL}_2$ still sends
$f_1$ to $f_2$. For an $\operatorname{SL}_2$-invariant polynomial $J$,
define the polynomial in one variable
\[
 j(t)=J\left(-2R^2T+\frac{t}{32}T^3\right).
\]
The argument of $J$ has discriminant $t$. Thus the just-proved orbit
statement gives $J(f)=j(\Delta(f))$ on $\Delta\ne0$, hence everywhere
by polynomial density. Conversely \eqref{eq:invariant-disc-weight}
makes every polynomial in $\Delta$ an $\operatorname{SL}_2$-invariant.

For later representation comparisons one need not suppress any weights.
For $g=\operatorname{diag}(\lambda,\mu)$ the full degree-$d$ character
of the coordinate ring of this orbit closure is
\begin{equation}\label{eq:invariant-cubic-character}
 \sum_{i+j+k+\ell=d}
 \lambda^{-(3i+2j+k)}\mu^{-(j+2k+3\ell)}.
\end{equation}
This follows from its monomial basis $A^iB^jC^kD^\ell$ and the four
coefficient weights. In particular every degree, every weight and
every multiplicity in this orbit closure agrees with the corresponding
ambient coordinate-ring representation, because the two rings are
equal by \eqref{eq:invariant-open-orbit}.

\subsection{The retained node as a ring and a representation}

On $b=c=0$, the source root cover is $a=-r^2$. Its completed
self-fibre product and its exact linear comparison are
\[
 B_{\rm root}=\mathbb C[r_1,r_2]/(r_1^2-r_2^2),\qquad
 R_{\rm node}=\mathbb C[u,v]/(uv),
\]
\[
 u=r_1-r_2,\quad v=r_1+r_2,\qquad
 r_1=(u+v)/2,\quad r_2=(v-u)/2.
\]
Substitution gives $uv=r_1^2-r_2^2$, and the displayed inverse proves
the ring isomorphism. In this ring,
\begin{equation}\label{eq:invariant-node-base}
 a=-r_1^2=-r_2^2=-\frac{u^2+v^2}{4}.
\end{equation}
Removing $r_1r_2=0$ removes the origin from both components, so the
crossing and the modules supported there belong to the completed object.

Let $K\subset G_{\rm node}=\operatorname{GL}_2(\mathbb C)$ be the monomial
subgroup generated by the diagonal torus
$T_{\rm node}=\{\operatorname{diag}(\lambda,\mu):\lambda\mu\ne0\}$
and $\sigma=\left(\begin{smallmatrix}0&1\\1&0\end{smallmatrix}\right)$.
It acts on node points by ordinary matrix multiplication, and on
functions by $(k\cdot f)(z)=f(k^{-1}z)$. Thus $u$ and $v$ have
torus characters $\lambda^{-1}$ and $\mu^{-1}$ and $\sigma$ exchanges
them. The relations are preserved. This $K$ does not fix the map to
$a$: a diagonal point transformation sends the expression in
\eqref{eq:invariant-node-base} to
$-(\lambda^2u^2+\mu^2v^2)/4$. All uses of $K$ below concern this
specified node action, not an unstated automorphism over $\mathbb A^1_a$.

For clarity retain $R_{\rm node}$ and define its comparison ring
\[
 \widetilde R=\mathbb C[u]\oplus\mathbb C[v],\qquad
 j(f)=(f(u,0),f(0,v)).
\]
Every element of $R_{\rm node}$ has a unique expression
$c+u p(u)+v q(v)$: the relation kills all mixed monomials, and
restriction to the two axes proves uniqueness. Thus $j$ is injective,
its image is exactly the pairs with equal constant terms, and the
inverse on that image is $(p(u),q(v))\mapsto p(u)+q(v)-p(0)$.
Every element of $\widetilde R$ is integral over the image, since
$\widetilde R$ is generated as an $R_{\rm node}$-module by $(1,1)$
and the idempotent $(1,0)$, which satisfies $e^2-e=0$.
The total ring of fractions is $\mathbb C(u)\oplus\mathbb C(v)$.
To check this assertion, the non-zero-divisor $u+v$ allows the
fractions $u/(u+v)=(1,0)$ and $v/(u+v)=(0,1)$; localizing within
each component gives its rational-function field. Each polynomial
ring in one variable is integrally closed: a rational function
$p/q$ with coprime $p,q$ satisfying a monic equation makes $q$
divide $p^{n}$, so $q$ is constant. Componentwise application proves
$\widetilde R$ integrally closed. It is therefore precisely the
algebraic normalization comparison ring, with the injection retained.

Define $\epsilon:\widetilde R\to\mathbb C$ by
$\epsilon(p,q)=p(0)-q(0)$. Its codomain is the $K$-module
$\mathbb C_{\rm sgn}$ on which the torus acts trivially and
$\sigma$ acts by $-1$. The preceding reconstruction proves the exact
$K$-equivariant sequence
\begin{equation}\label{eq:invariant-node-sequence}
 0\longrightarrow R_{\rm node}\xrightarrow{j}\widetilde R
 \xrightarrow{\epsilon}\mathbb C_{\rm sgn}\longrightarrow0.
\end{equation}
The conductor is precisely
\[
 \mathfrak c=(u,v),\qquad
 j(\mathfrak c)=u\mathbb C[u]\oplus v\mathbb C[v].
\]
Indeed a zero-constant pair remains an equal-constant pair after
multiplication by any element of $\widetilde R$. Conversely multiplying
a conductor element by $(1,0)$ forces its common constant to be zero.
The quotient in \eqref{eq:invariant-node-sequence} is annihilated by
$\mathfrak c$, and conversely a nonzero constant cannot annihilate it.

Give $u,v$ degree one. The degree-zero pieces of
\eqref{eq:invariant-node-sequence} are respectively the trivial
representation, the sum of trivial and sign representations, and the
sign representation. For every $n\ge1$ the two rings have the same
degree-$n$ basis $u^n,v^n$, with torus character
$\lambda^{-n}+\mu^{-n}$ and with $\sigma$ swapping its basis.
Therefore
\[
 R_{\rm node}^{T_{\rm node}}=\mathbb C,\quad
 \widetilde R^{T_{\rm node}}=\mathbb C\oplus\mathbb C,\quad
 R_{\rm node}^{K}=\widetilde R^K=\mathbb C.
\]
These equalities follow by comparing the independent torus characters
in the monomial bases and then taking the swap invariants. Thus the
full $K$-invariant scalar rings coincide, while the exact sign
cokernel in \eqref{eq:invariant-node-sequence} remains nonzero.

The differential module and its comparison map are
\[
 \Omega^1_{R_{\rm node}}=
 (R_{\rm node}\,du\oplus R_{\rm node}\,dv)/
       R_{\rm node}(v\,du+u\,dv),
 \qquad
 \Omega^1_{\widetilde R}=\mathbb C[u]du\oplus\mathbb C[v]dv.
\]
The presentation follows directly from the universal derivation of
the polynomial ring and $d(uv)=vdu+udv$. Its class
$\theta=vdu=-udv$ is nonzero. Indeed if $(v,0)=h(v,u)$ in the free
module, then $hu=0$, so the unique-form description gives $h\in(v)$.
The equation $hv=v$ would then imply $h(0,v)=1$, contradicting its
zero constant. Multiplying the relation by $u$ and $v$ proves
$u\theta=v\theta=0$. The non-zero-divisor $u+v$ kills $\theta$,
so this is torsion in the literal module-theoretic sense.

If $Adu+Bdv$ maps to zero in $\Omega^1_{\widetilde R}$, then
$A(u,0)=0$ and $B(0,v)=0$, giving $A=va(v)$ and $B=ub(u)$.
The relations $v^2du=0$ and $u^2dv=0$ reduce its class exactly to
$(a(0)-b(0))\theta$. Conversely $\theta$ maps to zero, and every
element of $\Omega^1_{\widetilde R}$ lifts by these two displayed
summands. We have proved the exact sequence
\begin{equation}\label{eq:invariant-differential-sequence}
 0\longrightarrow\mathbb C\theta\longrightarrow
 \Omega^1_{R_{\rm node}}\longrightarrow\Omega^1_{\widetilde R}
 \longrightarrow0.
\end{equation}
The original-coordinate identity is
$\theta=r_2dr_1-r_1dr_2$: substitute the two linear inverse formulas
and use $r_1dr_1=r_2dr_2$, the derivative of the original relation.

Differentiation commutes with the specified action on functions.
It gives
\[
 \operatorname{diag}(\lambda,\mu)\cdot\theta
 =\lambda^{-1}\mu^{-1}\theta,
 \qquad \sigma\cdot\theta=u\,dv=-\theta.
\]
Thus $\mathbb C\theta=\mathbb C_{\det^{-1}}$ as a $K$-module.
With $\deg du=\deg dv=1$, it lies in degree two, and
\[
 \operatorname{Hilb}(\Omega^1_{R_{\rm node}},t)
 =\frac{2t}{1-t}+t^2.
\]
More precisely, the degree-$n$ character for $n\ge1$ is
$\lambda^{-n}+\mu^{-n}$ plus the additional character
$\lambda^{-1}\mu^{-1}$ exactly at $n=2$; the additional swap
character is negative. No copy of $\det^{-1}$ occurs in
$R_{\rm node}$ itself, since none of its torus weights is
$\lambda^{-1}\mu^{-1}$. The exact module comparison
\eqref{eq:invariant-differential-sequence}, rather than replacement
of the module by its coefficient ring, retains this difference.

\subsection{An exact common torus and the discriminant pullback}

Define the two homomorphisms from $\mathbb G_m$ by
\[
 \beta(\eta)=\operatorname{diag}(\eta^{-1},\eta^2)\in G_{\rm bin},
 \qquad
 \nu(\eta)=\eta^3 I\in G_{\rm node}.
\]
The map from the completed node to binary cubics is
\begin{equation}\label{eq:invariant-node-to-cubic}
 \psi(u,v)=-2R^2T+\frac{u^2+v^2}{2}T^3
          =P^h_{a,0,0}(R,T),\quad a=-\frac{u^2+v^2}{4}.
\end{equation}
Its coordinate-ring map is
\[
 \psi^*: A\mapsto0,\ B\mapsto-2,\ C\mapsto0,
 \ D\mapsto(u^2+v^2)/2.
\]
Its kernel is exactly $(A,C,B+2)$. To see the unasserted part of
that equality, after division by those three generators any kernel
element would be a polynomial $H(D)$ with $H(u^2/2)=0$ on $v=0$;
its distinct powers of $u$ force every coefficient to vanish.
Its image is the polynomial subalgebra
$\mathbb C[(u^2+v^2)/2]$, not the entire node ring.
Since $\beta(\eta)$ multiplies $R^2T$ by
$\eta^{-2}\eta^2=1$ and $T^3$ by $\eta^6$, while $\nu(\eta)$
multiplies $u^2+v^2$ by $\eta^6$, this map is equivariant for the
two stated homomorphisms. Its discriminant pullback is exactly
\begin{equation}\label{eq:invariant-disc-node}
 \psi^*\Delta=16(u^2+v^2)=-64a.
\end{equation}
Both $\Delta$ restricted along $\beta$ and $\theta$ restricted along
$\nu$ have coordinate-module weight $\eta^{-6}$; their domains and
module structures remain different and have just been specified.
The discriminant retains quadratic vanishing on each branch. Its
zero scheme in the node has coordinate ring
\[
 R_{\rm node}/(u^2+v^2)
 =\mathbb C\{1,u,v,u^2\},\qquad
 v^2=-u^2,\quad u^3=v^3=0.
\]
These relations span the quotient, and independence follows by
homogeneous degree: degrees zero and one have no extra relation,
while in degree two precisely the one relation $u^2+v^2$ is imposed.
It consequently has length four. The conductor quotient
$R_{\rm node}/\mathfrak c=\mathbb C$ has length one. Their support
is the same origin because the radical of $(u^2+v^2)$ in the node
is $(u,v)$, as the displayed nilpotence proves. These are the exact
scheme and module relationships, including the multiplicity retained
by the discriminant.

\subsection{Exact induction to the connected general linear group}

We now turn the $K$-representations into explicitly defined rational
$G_{\rm node}$-representations. Put
\[
 L=\mathbb C[G_{\rm node}]
   =\mathbb C[p,q,s,t,(pt-qs)^{-1}].
\]
For any rational $K$-module $M$ define
\[
 \mathcal I(M)=(L\otimes_{\mathbb C}M)^K,
\]
where the diagonal action on a pure tensor is
$k\cdot(f\otimes m)=(g\mapsto f(gk))\otimes(k\cdot m)$.
The left $G_{\rm node}$-action is
$h\cdot(f\otimes m)=(g\mapsto f(h^{-1}g))\otimes m$.
The two actions commute by direct substitution; hence the invariant
subspace is a rational $G_{\rm node}$-module. When $M$ is a
$K$-algebra, $\mathcal I(M)$ is an algebra with componentwise
multiplication. These formulas define the invariant rings directly;
no assertion that their spectra are geometric quotients is needed.

We prove exactness rather than assume it. A rational torus module
is a sum of its integer weight spaces: expanding its coaction in the
independent Laurent monomials $\lambda^i\mu^j$ and comparing the
coaction identity proves that each coefficient is a vector of that
weight and that each vector has a finite such expansion. For a
surjective rational $T_{\rm node}$-module map and an invariant vector
in its target, choose a lift and take its weight-zero component;
the latter is an invariant lift. The normalizing swap preserves
weight zero. If the target vector is also swap invariant, average
that lift with its swap, dividing by two. This is a $K$-invariant
lift. Injections and kernels commute with invariants simply by
their definitions. Thus taking $K$-invariants is exact. Tensoring
an exact sequence of complex vector spaces with $L$ is exact, by
choosing bases or a vector-space splitting. This proves exactness
of $\mathcal I$ in full.

Apply it to the two proved sequences. With
$\mathcal A=\mathcal I(R_{\rm node})$ and
$\widetilde{\mathcal A}=\mathcal I(\widetilde R)$, we obtain
\begin{align}
 0&\longrightarrow\mathcal A\longrightarrow
 \widetilde{\mathcal A}\xrightarrow{\mathcal I(\epsilon)}
 \mathcal I(\mathbb C_{\rm sgn})\longrightarrow0,
 \label{eq:invariant-induced-ring}\\
 0&\longrightarrow\mathcal I(\mathbb C_{\det^{-1}})
 \longrightarrow\mathcal I(\Omega^1_{R_{\rm node}})
 \longrightarrow\mathcal I(\Omega^1_{\widetilde R})
 \longrightarrow0.
 \label{eq:invariant-induced-differential}
\end{align}
The maps are the original inclusion, difference of constants and
differential comparison applied coefficient by coefficient in $L$.
The first sequence is a sequence of $\mathcal A$-modules; its
quotient carries the action induced from evaluation at the node
origin. We have not identified the last two modules of
\eqref{eq:invariant-induced-differential} with absolute K\"ahler
differentials of the invariant rings. Such an identification is
unnecessary for the displayed exact transport.

The quotient and the kernel have explicit nonzero elements and
finite-dimensional $G_{\rm node}$-subrepresentations. Let
$J=\operatorname{diag}(1,-1)$ and define
\begin{equation}\label{eq:invariant-adjoint}
 Z(g)=gJg^{-1}
 =\frac1{pt-qs}
 \begin{pmatrix}pt+qs&-2pq\\2st&-(pt+qs)\end{pmatrix}.
\end{equation}
For a diagonal $k$ one has $Z(gk)=Z(g)$; for $k=\sigma$ one has
$Z(g\sigma)=-Z(g)$, since $\sigma J\sigma=-J$.
Consequently each entry of $Z(g)$ tensored with the fixed generator
of $\mathbb C_{\rm sgn}$ belongs to
$\mathcal I(\mathbb C_{\rm sgn})$. A lift through
\eqref{eq:invariant-induced-ring} is obtained by tensoring the same
entry with $(1/2,-1/2)\in\widetilde R$, whose difference of constants
is one and whose $K$-character is sign. The diagonal entry is nonzero
already at $g=I$.

The three coefficient functions $(pt+qs)/(pt-qs)$,
$-2pq/(pt-qs)$ and $2st/(pt-qs)$ are linearly independent. If a
linear combination were zero, multiplication by $pt-qs$ would give
a polynomial linear combination of $pt+qs,pq,st$; comparison of
their four distinct monomials forces its three coefficients zero.
They span a $G_{\rm node}$-stable three-dimensional module, because
$Z(h^{-1}g)=h^{-1}Z(g)h$. More exactly, with
$\mathfrak{sl}_2=\{X:\operatorname{tr}X=0\}$, the map
\[
 \mathfrak{sl}_2\longrightarrow\mathcal I(\mathbb C_{\rm sgn}),
 \qquad X\longmapsto\bigl(g\mapsto\operatorname{tr}(XZ(g))\bigr)
                 \otimes1
\]
is injective by the established coefficient independence and the
nondegenerate trace pairing on trace-zero two-by-two matrices.
It is equivariant for the adjoint action, since
\[
 \operatorname{tr}(X Z(h^{-1}g))
 =\operatorname{tr}(hXh^{-1}Z(g)).
\]
If desired its irreducibility can be checked without a classification
theorem: under diagonal conjugation the matrices $E_{12}$,
$\operatorname{diag}(1,-1)$ and $E_{21}$ have three distinct weights.
A nonzero invariant subspace therefore contains one of these weight
vectors, by coefficient projection. Differentiating conjugation by
$I+tE_{12}$ and $I+tE_{21}$ gives brackets; the identities
$[E_{12},E_{21}]=J$, $[J,E_{12}]=2E_{12}$ and
$[J,E_{21}]=-2E_{21}$ then produce all three basis vectors.

For the torsion kernel use the regular function $\det g=pt-qs$.
The tensor
\begin{equation}\label{eq:invariant-induced-torsion}
 e(g)=(\det g)\,\theta
 \in\mathcal I(\mathbb C_{\det^{-1}})
\end{equation}
is invariant because right translation multiplies its first factor
by $\det k$ and its second factor by $(\det k)^{-1}$.
It is nonzero, and the left action sends it to
$(\det h)^{-1}e$. Thus the exact connected-group transport retains
a determinant-inverse line in degree two as well as the adjoint
subrepresentation in the degree-zero conductor quotient. Degrees
refer to the node variables; $L$ has degree zero. The original
annihilator information is retained by the module maps: an element
of $\mathcal A$ whose coefficient functions vanish at the node origin
annihilates both displayed defect modules, because each original
defect is supported at that origin.

\subsection{The exact coordinate-ring obstruction interface}

For any two closed $G$-stable subsets $X,Y$ of a finite-dimensional
representation $W$, their polynomial ideals and coordinate rings
are defined by
$I(X)=\{f\in\mathbb C[W]:f|_X=0\}$ and
$\mathbb C[X]=\mathbb C[W]/I(X)$. Inclusion $X\subseteq Y$ gives
$I(Y)\subseteq I(X)$, and the entire equivariant restriction map is
\[
 \mathbb C[Y]\twoheadrightarrow\mathbb C[X],\qquad
 [f]_{I(Y)}\longmapsto[f]_{I(X)},\qquad
 \ker=I(X)/I(Y).
\]
Conversely, if a polynomial $f\in I(Y)$ has $f(x)\ne0$ at some
$x\in X$, that evaluation proves $x\notin Y$ and excludes the
inclusion. A finite-dimensional $G$-submodule containing that
polynomial may be used to record its representation content; it
is finite-dimensional because substitution in a polynomial of
bounded degree preserves that degree bound. This is the precise
polynomial-map interface used by an orbit-closure obstruction.

For the squarefree source cubic, the full ideal is zero by
\eqref{eq:invariant-open-orbit}, and its degreewise characters are
exactly \eqref{eq:invariant-cubic-character}. Thus this particular
orbit closure supplies no nonzero polynomial vanishing on it.
The saturation calculation proves exactly the same zero-kernel
statement for the whole original coefficient slice. The conductor
transport instead supplies the nonzero kernels and cokernels in
\eqref{eq:invariant-induced-ring}--\eqref{eq:invariant-induced-differential},
with explicit module elements and representations. Those sequences
are actual correspondences with connected-group representation
data; they do not assert that these modules are the coordinate ring
of a permanent or determinant orbit closure. The proved common
edge between the original node and the binary-cubic representation
is \eqref{eq:invariant-node-to-cubic}, with its exact kernel,
discriminant pullback, degrees, and torus homomorphisms. Every claim
in this fixed-dimensional transport is thus attached to a specified
map. No assertion about the complexity of an unbounded family or
about a nonstandard quantum-group positivity conjecture has been
substituted for one of these calculations.

\section{Retaining specialization layers and arithmetic operators}
\label{sec:specialization}

\subsection{The completed two-sheet algebra over its original base}
Throughout this subsection $k$ is a field of characteristic different from
$2$. Put $A=k[a]$ and retain $a=-r^2$ from the original $b=c=0$
inverse-root cover. Its self-fibre product is
\[
 B=A[r_1,r_2]/(r_1^2+a,r_2^2+a)
   \simeq k[r_1,r_2]/(r_1^2-r_2^2).
\]
The isomorphism substitutes $a=-r_1^2$; its inverse sends the displayed
$r_i$ to their residue classes. In particular it keeps the base map.
Successive division by the two monic quadratics gives the unique expression
\begin{equation}\label{eq:B-basis}
 f_0(a)+f_1(a)r_1+f_2(a)r_2+f_{12}(a)r_1r_2.
\end{equation}
Uniqueness follows first in the free $A$-module
$A[r_1]/(r_1^2+a)$ and then in its monic extension in $r_2$.
Thus $B$ is free of rank four over $A$.

The component ring is $N=k[t]\oplus k[t]$ with $a$ acting as
$(-t^2,-t^2)$. Define
\begin{equation}\label{eq:nu-r}
 \nu(r_1)=(t,t),\quad \nu(r_2)=(t,-t),\quad
 \delta(p,q)=\frac{p(0)-q(0)}2.
\end{equation}
The relations are respected. Use the invertible linear substitution
$u=r_1-r_2$, $v=r_1+r_2$, with inverse
$r_1=(u+v)/2$, $r_2=(v-u)/2$ and relation $uv=0$.
Every class has a unique expression $c+up(u)+vq(v)$.
Its restrictions to the two axes determine all three terms, so $\nu$
is injective. A pair of polynomials lies in the image precisely when
its constants agree: the constant and two positive-degree axis parts
reconstruct it. Hence
\begin{equation}\label{eq:conductor-A-exact}
 0\longrightarrow B\xrightarrow{\nu}N
 \xrightarrow{\delta}k\longrightarrow0
\end{equation}
is exact over $A$, with $ak=0$.
The conductor is $(r_1,r_2)=(u,v)$: pairs with zero constants
remain in the image after multiplication by every element of $N$;
multiplication by $(1,0)$ forces zero common constant in the converse.

In the ordered $A$-bases
\[
 \mathcal B=(1,r_1,r_2,r_1r_2),\qquad
 \mathcal N=((1,1),(t,t),(t,-t),(1,-1))
\]
of $B,N$, respectively, the matrix is
\begin{equation}\label{eq:specialization-matrix}
 [\nu]_{\mathcal N\leftarrow\mathcal B}
     =\operatorname{diag}(1,1,1,-a).
\end{equation}
The last column uses $(t^2,-t^2)=-a(1,-1)$.
The four vectors of $\mathcal N$ are a basis by even--odd decomposition
of each component and division by $2$.
Thus $\nu$ is an isomorphism over $a\ne0$, with inverse last entry
$-1/a$. No inverse is assigned at $a=0$.

\begin{proposition}[The exact specialization defect]\label{prop:spec-defect}
For $B_0=B/aB$, $N_0=N/aN$ there is an exact sequence
\[
 0\longrightarrow k\xrightarrow{\partial}B_0
 \xrightarrow{\nu_0}N_0\xrightarrow{\delta_0}k\longrightarrow0,
 \qquad \partial(1)=-r_1r_2.
\]
Here $B_0=k[r_1,r_2]/(r_1^2,r_2^2)$,
$N_0=k[t]/(t^2)\oplus k[t]/(t^2)$, and the first $k$ is
$\operatorname{Tor}^A_1(k,k)$ for the quotient in
\eqref{eq:conductor-A-exact}.
\end{proposition}
\begin{proof}
The free resolution $0\to A\xrightarrow{a}A\to k\to0$ becomes
the zero map $k\to k$ after tensoring with $k$, so its first homology
is $k$. The middle modules in \eqref{eq:conductor-A-exact} are free.
Alternatively the diagonal matrix at $a=0$ directly gives the exact
kernel and cokernel. To compute the connecting sign, lift $1\in k$
by $(1,-1)\in N$. Its product by $a$ is
$(-t^2,t^2)=\nu(-r_1r_2)$.
The class $r_1r_2$ is nonzero because \eqref{eq:B-basis} stays a basis
after tensoring. This proves every map and sign.
\end{proof}

Let $\Gamma=\mu_2\times\mu_2$ act by
$(\epsilon_1,\epsilon_2)r_i=\epsilon_i r_i$ and fix $a$.
Write $\chi_i(\epsilon_1,\epsilon_2)=\epsilon_i$.
On $N$ its first generator acts by
$(p(t),q(t))\mapsto(q(-t),p(-t))$, and the second by
$(p(t),q(t))\mapsto(q(t),p(t))$.
Checking $r_1,r_2$ proves equivariance of $\nu$.
Both generators negate the constant difference, so the quotient
character is $\chi_1\chi_2$. The Tor term has the same character
because $a$ is fixed. With $\deg r_i=\deg t=1$, $\deg a=2$, its
image has degree $2$, whereas the cokernel has degree $0$.
The ungraded character of each special fibre is
$1+\chi_1+\chi_2+\chi_1\chi_2$, but their graded characters are
\[
 \operatorname{ch}_zB_0=1+z(\chi_1+\chi_2)+z^2\chi_1\chi_2,
 \qquad
 \operatorname{ch}_zN_0=1+\chi_1\chi_2+z(\chi_1+\chi_2).
\]
The bases prove these statements completely.
Moreover $B_0$ has unique maximal ideal $\mathfrak m=(r_1,r_2)$,
$\mathfrak m^2=(r_1r_2)$, $\mathfrak m^3=0$.
Every element with nonzero constant is invertible by the finite
geometric series in its nilpotent part, which proves uniqueness.
In contrast $N_0$ contains the nontrivial idempotent $(1,0)$ and has
two residue-field quotients. Thus identical ungraded multiplicities
do not reconstruct the algebra or the specialization map.

\subsection{An explicit finite-field zeta calculation}
All preceding equations are defined over $\mathbb Z[1/2]$.
Take $k=\mathbb F_q$ with $q$ odd.
Over $\mathbb F_{q^m}$ the two lines $r_2=r_1$ and $r_2=-r_1$
intersect in one point, so
\[
 |\operatorname{Spec}B(\mathbb F_{q^m})|=2q^m-1,\qquad
 |\operatorname{Spec}N(\mathbb F_{q^m})|=2q^m.
\]
Define $Z_X(T)=\exp(\sum_{m\ge1}|X(\mathbb F_{q^m})|T^m/m)$.
The formal identity $-\log(1-vT)=\sum_{m\ge1}v^mT^m/m$ gives
\begin{equation}\label{eq:node-zeta}
 Z_{\operatorname{Spec}B}(T)=\frac{1-T}{(1-qT)^2},
 \qquad Z_{\operatorname{Spec}N}(T)=\frac1{(1-qT)^2}.
\end{equation}
The conductor point contributes the exact factor $1-T$.
The reciprocal zero is $1$ and reciprocal poles are $q,q$.
This explicitly computed singular affine curve is not assigned the
purity assertion for a smooth projective object.
The original inverse chart requires $r_1r_2\ne0$, equivalently $a\ne0$.
Its two lines each lose the origin and are disjoint, giving point count
$2(q^m-1)$ and zeta function $(1-T)^2/(1-qT)^2$.
This is the exact effect of restoring the chart exclusion.

\subsection{The arithmetic operator and its finite matrix map}
The arithmetic application uses the exact source quotient of
Connes--Consani--Moscovici~\cite[Section 3.6, Proposition 3.6]{ccmProlate2024}:
\[
 Q_{\rm ar}=\mathcal H_{\le1}/
 \overline{\operatorname{span}
 \{\mathcal F_\mu\mathcal E(\psi_\ell^\pm)\}},\qquad
 (\mathcal F_\mu\mathcal E(\psi_\ell^\pm))(s)=R_\ell^\pm(s)\Xi(s).
\]
Here $\mathcal H_{\le1}$ is their Hadamard topological ring of entire
functions of order at most one. Multiplication by the original coordinate
$s$ induces $S$ on this quotient. Their proposition identifies its
spectrum with the $s$ for which $1/2+is$ is a nontrivial zeta zero.
This arithmetic identification is a reported primary-source theorem.
The following extension identities are proved here for this exact
module and for every continuous module of the same type.
The closed subspace is retained as stated; no replacement by an
algebraic principal ideal or orthogonal complement is used.

Let $Q$ be a complex topological vector space and $S$ a continuous
linear operator. Regard $Q$ as a $\mathbb C[s]$-module by $s\cdot v=Sv$.
The base change $s\mapsto-r^2/2$ gives
\[
 \widehat Q=\mathbb C[r]\otimes_{\mathbb C[s]}Q=Q\oplus rQ,\qquad
 \mathcal R=\begin{pmatrix}0&-2S\\I&0\end{pmatrix},\qquad
 -\mathcal R^2/2=\operatorname{diag}(S,S).
\]
The unique even--odd expression
$p_0(-r^2/2)+rp_1(-r^2/2)$ proves freeness and all displayed maps.
Use the finite direct-sum topology.
The involution is $(v_0,v_1)\mapsto(v_0,-v_1)$; the trace is $2v_0$,
with kernel $rQ$. This is $\mathbb C[s]$-linear.
Projection $\pi(v_0,v_1)=v_0$ has section $\iota(v)=(v,0)$.
Here spectrum means failure of a continuous two-sided inverse.
If $T$ inverts $S-\lambda I$, then $\operatorname{diag}(T,T)$
inverts $\operatorname{diag}(S,S)-\lambda I$.
Conversely an inverse $D$ of the latter gives the continuous inverse
$\pi D\iota$ of $S-\lambda I$, by the two intertwining identities.
Thus the spectrum in the original $s$ coordinate is preserved exactly.

For each $d\ge1$ define an actual polynomial map
\begin{equation}\label{eq:matrix-Phi}
 \Phi_d:\operatorname{Mat}_d(\mathbb C)\longrightarrow
 \operatorname{Mat}_{2d}(\mathbb C),\qquad
 A\longmapsto\begin{pmatrix}0&-2A\\I_d&0\end{pmatrix}.
\end{equation}
Its image is the affine linear slice with these three fixed blocks;
its inverse there is $A=-M_{12}/2$.
The map on coordinate rings is surjective with kernel generated by
the entries of $M_{11},M_{22},M_{21}-I_d$.
The remaining block maps freely to $-2A_{ij}$, proving the kernel.
Conjugation by $g$ is intertwined with conjugation by
$\operatorname{diag}(g,g)$, by block multiplication.

\begin{proposition}[All characteristic-coefficient invariants retained]
\label{prop:matrix-invariants}
Write $\det(uI-A)=\sum_{j=0}^d(-1)^je_j(A)u^{d-j}$, $e_0=1$.
Then
\[
 \det(tI-\Phi_d(A))=\det(t^2I+2A)
   =\sum_{j=0}^d2^je_j(A)t^{2d-2j}.
\]
Pullback on conjugation-invariant rings is the surjection
\[
 \mathbb C[e_1,\ldots,e_{2d}]\longrightarrow\mathbb C[e_1,\ldots,e_d],
 \quad e_{2j}\mapsto2^je_j,\quad e_{2j-1}\mapsto0,
\]
whose kernel is exactly the ideal of odd generators in the source.
\end{proposition}
\begin{proof}
For indeterminate $t\ne0$, the Schur complement of $tI_d$ in
$\left(\begin{smallmatrix}tI&2A\\-I&tI\end{smallmatrix}\right)$ gives
$t^d\det(tI+2A/t)=\det(t^2I+2A)$. Polynomial identity extends to $t=0$.
Expansion gives every coefficient and sign.

We also prove the invariant-ring assertion. An invariant polynomial
restricted to diagonal matrices is symmetric, by permutation matrices.
Every symmetric polynomial is a polynomial in the $e_j$:
in one homogeneous component its largest lexicographic monomial has
exponents $b_1\ge\cdots\ge b_d$. Subtract its coefficient times
$e_1^{b_1-b_2}\cdots e_{d-1}^{b_{d-1}-b_d}e_d^{b_d}$.
The leading monomial disappears without a larger one being introduced.
The finite set of monomials of that degree makes this algorithm terminate.
Matrices with distinct eigenvalues form the nonempty discriminant-open
set and are diagonalizable. Independence of their eigenvectors follows
by applying the product of the other eigenvalue factors to a dependence.
The invariant and its polynomial in $e_j$ agree there and hence everywhere:
localization at a nonzero polynomial in a polynomial domain is injective.
The $e_j$ are independent since every monic complex polynomial factors
and is the characteristic polynomial of a diagonal matrix.
This proves the two invariant rings. Quotienting the source by its odd
generators leaves a polynomial ring mapped isomorphically by the nonzero
scalars $2^j$, proving the exact kernel.
\end{proof}

No unspecified finite truncation of the arithmetic quotient is used in
\eqref{eq:matrix-Phi}. It is a natural finite-dimensional version of the
same base change. The algebra $B$ also embeds into $\operatorname{End}_A(B)$
by left multiplication: an element acting as zero annihilates $1$.
In the basis $(1,r_1)\otimes(1,r_2)$ its two root operators are
\[
 R_1=\begin{pmatrix}0&-a\\1&0\end{pmatrix}\otimes I_2,\qquad
 R_2=I_2\otimes\begin{pmatrix}0&-a\\1&0\end{pmatrix}.
\]
They commute, square to $-aI_4$ and retain their nonzero nilpotent
specializations at $a=0$. Setting $a=2s$ reproduces the arithmetic
factor $-2s$. These regular representations and invariant-ring maps
connect the retained cover to concrete representation-theoretic objects.

\section{Integral coefficient lattices and the ramified conductor}
\label{sec:integral-coefficient-conductor}

The following maps retain the arithmetic of the source-displayed coefficient
\eqref{eq:gct-nrh-coefficient-original}. Put
\[
 O=\mathbb Z[1/2],\quad R=O[q,q^{-1}],\quad t=q-1,\quad
 r=q-q^{-1},\quad h=[7]_q[3]_q,\quad C=r^2h.
\]
The letters $r$ and $a$ in this section use the original $b=c=0$
root coordinate and base map $a=-r^2$, respectively. The trace coordinate
is still $p=q+q^{-1}$. None of $3,7,h$ is inverted in $R$.

\subsection{The exact lattice map over the integral base}

The two-sheet calculation of Section~\ref{sec:specialization} holds over
$O[a]$: the displayed monic bases, the component restriction maps, and
their inverses only use division by $2$. Base change by
$O[a]\to R$, $a\mapsto-r^2$, gives free modules $B_R,N_R$ with their
original ordered bases and map
\[
 V:B_R\longrightarrow N_R,\qquad [V]=\operatorname{diag}(1,1,1,r^2).
\]
This assertion is also checked directly: the last component vector
$r_1r_2$ maps to $-a(1,-1)=r^2(1,-1)$. Thus the original sign is retained.
Since $R$ is a domain and $r\ne0$, this map is injective, and projecting
onto the last coordinate modulo $r^2$ proves
\[
 0\longrightarrow B_R\xrightarrow{V}N_R
 \longrightarrow R/(r^2)\longrightarrow0.
\]

Let $E_C=R^4$ be a free module with ordered basis $e_1,e_2,e_3,e_4$.
Define all entries of the following maps:
\[
 H:E_C\longrightarrow B_R,\quad[H]=\operatorname{diag}(1,1,1,h),
 \qquad D_C:E_C\longrightarrow N_R,\quad[D_C]=\operatorname{diag}(1,1,1,C).
\]
Matrix multiplication proves $D_C=VH$; both maps are injective because
$h,C$ are nonzero. Define the two actual image lattices
$L_C=D_C(E_C)$ and $L_r=V(B_R)$ inside $N_R$.
Their first three coordinates equal $R$ and their last coordinates are
$CR$ and $r^2R$, respectively. Consequently there is an exact sequence
\begin{equation}\label{eq:integral-lattice-sequence}
 0\longrightarrow R/(h)\xrightarrow{\iota}R/(C)
 \xrightarrow{\pi}R/(r^2)\longrightarrow0,
 \quad \iota(\bar f)=\overline{r^2f},\quad\pi(\bar f)=\bar f.
\end{equation}
Here $R/(h)=L_r/L_C$, $R/(C)=N_R/L_C$ and $R/(r^2)=N_R/L_r$ under
the indicated last-coordinate maps. For well-definedness of $\iota$,
replacing $f$ by $f+hg$ changes its image by $Cg$. If $r^2f=Cg$,
cancellation of the nonzero $r^2$ gives $f=hg$, proving injectivity.
The kernel of $\pi$ consists exactly of classes $r^2f$, and $\pi$ is
surjective by lifting a representative. This proves every assertion.

The sign action $\Gamma=\mu_2\times\mu_2$ on the four original basis
vectors has characters $1,\chi_1,\chi_2,\chi_1\chi_2$.
Let it act trivially on $R$ and by these same characters on $E_C$.
All diagonal maps above commute with it, so every last-coordinate
quotient in \eqref{eq:integral-lattice-sequence} has character
$\chi_1\chi_2$. The semilinear involution $q\mapsto q^{-1}$, fixing
the four chosen basis vectors, also commutes with the maps: it sends
$r$ to $-r$ and fixes $r^2,h,C$. These are explicit symmetries of this
lattice diagram. No unspecified Chevalley action on $E_C$ is inferred.

\subsection{The full specialization sequence and its factor 21}

Specialization means tensoring over $R$ with $O=R/(t)$.
Write $u=(q+1)/q$, so $r=tu$ and $u(1)=2$.
We have $h(1)=21$. Both $V$ and $D_C$ specialize to
$\operatorname{diag}(1,1,1,0)$. Their kernels and cokernels are therefore
free copies of $O$ with basis given by the fourth coordinate.
The induced kernel map from $H$ is multiplication by $21$; the induced
cokernel map from the identity on $N_R$ is the identity of $O$.
In particular this diagram preserves more than the ranks of the maps.

Here is the complete derived specialization of
\eqref{eq:integral-lattice-sequence}, including the connecting map.
For any $R$-module $M$, the two-term free resolution
$0\to R\xrightarrow{t}R\to O\to0$ proves
\[
 \operatorname{Tor}_1^R(M,O)=\{m\in M:tm=0\},\qquad
 \operatorname{Tor}_j^R(M,O)=0\quad(j\ge2).
\]
If $d=td_1\ne0$, cancellation in $R$ shows
\[
 \operatorname{Tor}_1^R(R/(d),O)=O\,\overline{d_1}.
\]
Indeed $tf=dg=td_1g$ implies $f=d_1g$, and $d_1g\in(d)$ holds
exactly when $g\in(t)$. The ideal $(t)$ is prime because $R/(t)=O$
is a domain. It does not contain $h$, whose value is $21\ne0$ in $O$.
Thus $tf=hg$ forces $g=tg_1$, then $f=hg_1$, and
$\operatorname{Tor}_1^R(R/(h),O)=0$. Applying these facts to $d=C,r^2$ proves the
exact specialized sequence
\begin{equation}\label{eq:integral-specialization-21}
 0\longrightarrow O\xrightarrow{\,21\,}O
 \xrightarrow{\mathrm{mod}\ 21}O/(21)
 \xrightarrow{\,0\,}O\xrightarrow{\,1\,}O\longrightarrow0.
\end{equation}
We verify the maps, so that no exactness assertion depends on a hidden
choice of generators. The two Tor generators are $C/t=tu^2h$ and
$r^2/t=tu^2$, respectively. Projection sends the first to $h$ times the
second, and hence gives $21$ after reduction by $t$.
To compute the connecting map, lift $r^2/t$ from $R/(r^2)$ to
$R/(C)$. Its product by $t$ is $r^2=\iota(1)$; the connecting class
is therefore $1$ in $R/(h,t)=O/(21)$, with positive sign. The next map
is induced by multiplication by $r^2$, which vanishes modulo $t$.
The final map is the identity on constant classes. Multiplication by
$21$ is injective on $O$, with cokernel exactly $O/(21)$, verifying
exactness of the displayed sequence directly as well.

Over $D=\mathbb Q[q,q^{-1}]_{(q-1)}$ the element $h$ is a unit: its value
$21$ is nonzero in the residue field $\mathbb Q$. Thus $H$ has the explicit
inverse $\operatorname{diag}(1,1,1,h^{-1})$ there, and it identifies
the two injections over $D$. Over $R$ the additional quotient $R/(h)$
and the residual $O/(21)$ remain. These are the additional comparison
modules killed by passage to $D$. This is not an exhaustive inventory
of localization effects: the same localization also removes the $q=-1$
branch of $r^2$.

For completeness the entire $t$-saturation can be computed while retaining
the integral constants. Work in
$R_+=O[q,q^{-1},(q+1)^{-1}]$, where $u$ is a unit and $(t)$ is still
prime. The ideal generated by $C$ is $(t^2h)$ and the ideal generated
by $r^2$ is $(t^2)$. Their $t$-saturations are $(h)$ and $R_+$.
For the first assertion, if $t^ef=t^2hg$, then for $e\ge2$ cancellation
gives $t^{e-2}f=hg$. Primality of $t$ and $t\nmid h$ imply successively
that $g$ is divisible by $t^{e-2}$, so $f\in(h)$. For $e<2$ the same
conclusion follows directly by cancellation. Conversely $t^2hR_+$ lies
in $(C)$, proving the asserted saturation; the second assertion follows
from $t^2R_+=(r^2)$. Thus the saturated image lattices are
$R_+^3\oplus hR_+$ and $R_+^4$, respectively.

\subsection{The retained integral jets and the two exceptional primes}

The exact equality $C=t^2u^2h$ gives
\[
 C\equiv84t^2\pmod{t^3R_+}.
\]
This follows by evaluating $u^2h$ at $q=1$: $2^2\cdot7\cdot3=84$.
Over $\mathbb Q$ the coefficient has exact order two, as previously proved.
The integral residue $84$ also determines why that order changes in
characteristics $3$ and $7$; these characteristics are allowed over $O$.

For an odd prime $\ell$, in $\mathbb F_\ell[q,q^{-1}]$ one has the polynomial
identity
\begin{equation}\label{eq:integral-quantum-prime}
 [\ell]_q=(q-q^{-1})^{\ell-1}.
\end{equation}
To prove it, multiply the finite sum defining $[\ell]_q$ by $q^{\ell-1}$.
The result is $1+q^2+\cdots+q^{2(\ell-1)}$.
Its product with $q^2-1$ is $q^{2\ell}-1=(q^2-1)^\ell$ in
characteristic $\ell$. Cancellation in the polynomial domain gives
$(q^2-1)^{\ell-1}$, and multiplication by $q^{-(\ell-1)}$ proves
\eqref{eq:integral-quantum-prime}. In particular the original coefficient
reduces exactly to
\[
 C=r^4[7]_q\quad\text{over }\mathbb F_3,\qquad
 C=r^8[3]_q\quad\text{over }\mathbb F_7.
\]
Because $r=tu$, $u(1)=2$, $[7]_1=1$ in $\mathbb F_3$ and $[3]_1=3$ in
$\mathbb F_7$, the exact orders of $C$ at $q=1$ are four and eight,
respectively. Their first nonzero coefficients are
$2^4\cdot1=1$ in $\mathbb F_3$ and $2^8\cdot3=5$ in $\mathbb F_7$.
For every other odd characteristic $\ell$, both $21$ and $4$ are
nonzero, so the exact order is two and its first coefficient is
$84\bmod\ell$. The node map still has last entry $r^2$, of exact
order two in every odd characteristic. The larger orders therefore
belong to the retained residual factor of the GCT source-displayed polynomial,
and the injection $H$ records their relation to the node map.

These results construct an integral coefficient-lattice correspondence.
They specify every map and every residual module in that correspondence.
Section~\ref{sec:gct-generator-interval} supplies the actual Cartan-twisted
divided-square endpoint from the original GCT-IV basis subquotient.
Adding its three specified identity summands gives exactly $D_C$ here,
so $D_C=VH$ is also an explicit comparison of that generator endpoint
with the original node. The trace obstruction proved there prevents
extending the four-dimensional interval to the full raising action.
Section~\ref{sec:full-source-generators} now constructs that full original
source module, all its raising and lowering operators, and the exact
integral quotient map to the retained interval. Its raw-word projection
has a specified section and free kernel; the square comparing its full
domain with the earlier restricted four-word lattice retains the latter's
nonzero cokernel. Thus the larger source action and the conductor
comparison here are related by proved integral maps.

\section{The original Laurent coefficient and its geometric objects}
\label{sec:cg-original}
Throughout this construction the quantum parameter is the indeterminate
$q$, and $Q$ denotes a prime power. They are different objects. The retained
source-displayed Laurent polynomial, with its original four terms, is
\begin{equation}\label{eq:cg-C}
 C(q)=q^{10}-q^4-q^{-4}+q^{-10}.
\end{equation}
Define, as a Laurent polynomial over $\mathbb Z$,
\[
 [n]_q=q^{n-1}+q^{n-3}+\cdots+q^{-(n-1)},\qquad
 h(q)=[7]_q[3]_q,
\]
and set
\[
 B(T)=(1+T+\cdots+T^6)(1+T+T^2).
\]
We construct schemes over $\mathbb Z$:
\[
 U_n=\mathbb A^n_{\mathbb Z}\setminus V(x_0,\ldots,x_{n-1}),
 \qquad X=U_7\times_{\mathbb Z}U_3,\qquad
 Y=\mathbb P^6_{\mathbb Z}\times_{\mathbb Z}\mathbb P^2_{\mathbb Z}.
\]
Here $V(x_0,\ldots,x_{n-1})$ is the entire zero section
$\operatorname{Spec}\mathbb Z$, not a chosen point over one field.
Products and complements in what follows commute with base change to any
field. These definitions build a geometry for the exact scalar
\eqref{eq:cg-C}; they do not specify a quantum-group representation or a
canonical-basis convolution space. The source audit in the accompanying
research record questions the assignment of this displayed polynomial to
the printed target tableau. The construction here takes exactly the
displayed polynomial as its input and asserts no resolution of that
matrix-entry assignment.

\begin{proposition}\label{prop:cg-factor}
The following identities hold in $\mathbb Z[q,q^{-1}]$:
\begin{align}
 C(q)&=q^{-10}(q^{14}-1)(q^6-1),\label{eq:cg-factor}\\
 h(q)&=q^{-8}B(q^2),\label{eq:cg-h}\\
 C(q)&=(q-q^{-1})^2h(q).\label{eq:cg-rh}
\end{align}
Furthermore
\begin{equation}\label{eq:cg-B}
 B(T)=1+2T+3T^2+3T^3+3T^4+3T^5+3T^6+2T^7+T^8.
\end{equation}
\end{proposition}
\begin{proof}
The right side of \eqref{eq:cg-factor} expands to
$q^{-10}(q^{20}-q^{14}-q^6+1)$, preserving all four terms of
\eqref{eq:cg-C}. Directly from the definition,
$[7]_q=q^{-6}\sum_{i=0}^6q^{2i}$ and
$[3]_q=q^{-2}\sum_{j=0}^2q^{2j}$, which proves
\eqref{eq:cg-h}. The finite polynomial identities
$(q^2-1)\sum_{i=0}^6q^{2i}=q^{14}-1$ and
$(q^2-1)\sum_{j=0}^2q^{2j}=q^6-1$, together with
$(q-q^{-1})^2=q^{-2}(q^2-1)^2$, prove \eqref{eq:cg-rh}
without inverting $q^2-1$. The coefficient of $T^d$ in $B$ is the
number of ordered pairs $(i,j)$ with $0\leq i\leq6$, $0\leq j\leq2$
and $i+j=d$. Listing these nine counts gives \eqref{eq:cg-B}.
\end{proof}

\section{The quotient morphism, its inverses on charts, and its fibres}
\label{sec:cg-torsor}
\begin{proposition}\label{prop:cg-torsor}
There is a morphism
\[
 \pi_n:U_n\longrightarrow\mathbb P^{n-1}_{\mathbb Z},\qquad
 (x_0,\ldots,x_{n-1})\longmapsto[x_0:\cdots:x_{n-1}],
\]
which is a Zariski locally trivial $\mathbb G_m$-torsor for the action
$\lambda\cdot(x_i)_i=(\lambda x_i)_i$. Its product
\[
 \pi=\pi_7\times\pi_3:X\longrightarrow Y
\]
is a Zariski locally trivial $(\mathbb G_m)^2$-torsor. For any field $k$
and any point of $Y(k)$, its scheme fibre is isomorphic over $k$ to
$(\mathbb G_{m,k})^2$.
\end{proposition}
\begin{proof}
Let $D_+(z_i)$ be the standard chart on $\mathbb P^{n-1}$, with
coordinates $u_j=z_j/z_i$ for $j\ne i$, and let $D(x_i)\subset U_n$.
Define on these charts the ring homomorphism
\[
 \mathbb Z[u_j:j\ne i]\longrightarrow
 \mathbb Z[x_0,\ldots,x_{n-1},x_i^{-1}],\qquad u_j\longmapsto x_j/x_i.
\]
On $D(x_i)\cap D(x_k)$ the two projective coordinates obey
$u^{(k)}_j=u^{(i)}_j/u^{(i)}_k$, with the $i$th coordinate equal to
$1/u^{(i)}_k$. Thus the homomorphisms glue to $\pi_n$.
An explicit isomorphism and inverse are
\begin{align*}
 D(x_i)&\longrightarrow D_+(z_i)\times\mathbb G_m,
 &(x_j)_j&\longmapsto ((x_j/x_i)_{j\ne i},x_i),\\
 D_+(z_i)\times\mathbb G_m&\longrightarrow D(x_i),
 &((u_j)_{j\ne i},t)&\longmapsto x_i=t,\quad x_j=tu_j\ (j\ne i).
\end{align*}
Every expression is regular on the indicated chart. Substitution in each
direction returns every original coordinate, and the action becomes
$t\mapsto\lambda t$. The charts cover $U_n$ and $\mathbb P^{n-1}$.
On a double overlap the fibre coordinate changes by
$t_k=t_i u^{(i)}_k$, with an invertible multiplier, so the local torsor
identifications agree.

For completeness, the action criterion holds as an isomorphism of schemes:
\[
 \mathbb G_m\times U_n\longrightarrow U_n\times_{\mathbb P^{n-1}}U_n,
 \qquad(\lambda,x)\longmapsto(\lambda x,x).
\]
On the chart where $x_i$ is invertible, its inverse sends $(x',x)$
to $(x'_i/x_i,x)$. Equality of the projective coordinates gives
$x'_j=(x'_i/x_i)x_j$ for every $j$, and $x'_i/x_i$ is invertible.
These ratios agree on overlaps, proving the global inverse. The local
sections with $t_i=1$ prove that $\pi_n$ is locally surjective. Equally,
an invariant morphism out of $U_n$ factors on each chart through
$(u,t)\mapsto u$, because its value at $(u,t)$ equals its value at
$(u,1)$; these factorizations agree on overlaps. This proves the
geometric quotient property for invariant morphisms to schemes.

Taking products gives all statements for $\pi$. A $k$-rational point
has some nonzero projective coordinate in each factor, so it lies in
one of these trivializing charts defined over $k$. Pulling that chart
isomorphism back to the point gives the asserted scheme fibre and
preserves the two separate scalar actions.
\end{proof}

\begin{remark}\label{rem:cg-rpoints}
The chart proof is a scheme proof and covers arbitrary base rings. It does
not identify arbitrary projective $R$-points with a single globally chosen
coordinate vector. A point of $U_n(R)$ is represented by coordinates
generating the unit ideal, and the chart trivializations hold after the
corresponding localizations. For finite fields, ordinary nonzero vectors
provide exactly the required representatives.
\end{remark}

\section{Classes, point counts and the quantum specialization map}
\label{sec:cg-counts}
Let $k$ be a perfect field. The ring $K_0(\operatorname{Var}_k)$ used here is
generated by isomorphism classes $[V]$ of finite-type reduced separated
$k$-schemes, with relations $[V]=[Z]+[V\setminus Z]$ for closed reduced
$Z$, and product $[V][W]=[V\times_k W]$. All schemes used in the
calculation and their products are reduced. Write
$\mathbb L=[\mathbb A^1_k]$. We use the localization in which
$\mathbb L$ is invertible only where explicitly displayed.

\begin{proposition}\label{prop:cg-classes}
In this ring,
\begin{align}
 [Y_k]&=B(\mathbb L),\label{eq:cg-classY}\\
 [X_k]&=(\mathbb L^7-1)(\mathbb L^3-1)
       =(\mathbb L-1)^2B(\mathbb L).\label{eq:cg-classX}
\end{align}
The polynomial lift of these identities under the injective ring map
\[
 \iota:\mathbb Z[T,T^{-1}]\longrightarrow\mathbb Z[q,q^{-1}],
 \qquad T\longmapsto q^2,
\]
is precisely
\[
 \iota\bigl(T^{-5}(T^7-1)(T^3-1)\bigr)=C(q),\qquad
 \iota\bigl(T^{-4}B(T)\bigr)=h(q).
\]
\end{proposition}
\begin{proof}
In $\mathbb P^{n-1}$ stratify by the last nonzero coordinate. The stratum
with last nonzero coordinate $z_d$ has the explicit isomorphism
\[
 \mathbb A^d_k\longrightarrow
 \{z_{d+1}=\cdots=z_{n-1}=0,\ z_d\ne0\},\quad
 (a_0,\ldots,a_{d-1})\longmapsto[a_0:\cdots:a_{d-1}:1:0:\cdots:0].
\]
The inverse consists of the regular ratios $z_i/z_d$. The successive
closed subspaces $\mathbb P^0\subset\cdots\subset\mathbb P^{n-1}$
show by the defining relation that
$[\mathbb P^{n-1}]=\sum_{d=0}^{n-1}\mathbb L^d$. Taking the two
products proves \eqref{eq:cg-classY}. Removing the zero section over
$k$ gives $[U_n]=\mathbb L^n-1$, proving the first equality for $X$.
The torsor is trivial over each displayed stratum, so its inverse image
has class $(\mathbb L-1)\mathbb L^d$. Applying this separately in
the two factors proves the second equality directly from the geometric
map. The same equality also follows from finite polynomial multiplication.
The map $\iota$ is injective because distinct powers $T^d$ map to the
distinct Laurent monomials $q^{2d}$, which are linearly independent over
$\mathbb Z$. Substitution proves the two asserted lifts using
Proposition~\ref{prop:cg-factor}. No injectivity of
$\mathbb Z[T,T^{-1}]\to K_0(\operatorname{Var}_k)[\mathbb L^{-1}]$
is asserted or needed.
\end{proof}

\begin{theorem}\label{thm:cg-counts}
For every prime power $Q$ and every integer $m\geq1$, set $N=Q^m$.
Then
\begin{align}
 \#Y(\mathbb F_N)&=B(N),\label{eq:cg-NY}\\
 \#X(\mathbb F_N)&=(N^7-1)(N^3-1)
                =(N-1)^2B(N),\label{eq:cg-NX}\\
 C(Q^{m/2})&=Q^{-5m}\#X(\mathbb F_{Q^m}),\label{eq:cg-countC}\\
 h(Q^{m/2})&=Q^{-4m}\#Y(\mathbb F_{Q^m}).\label{eq:cg-counth}
\end{align}
Here $Q^{m/2}$ is the positive real square root of $Q^m$ used to
evaluate a Laurent polynomial. The values in the last two equations
are rational and independent of the choice of sign of that square root.
\end{theorem}
\begin{proof}
There are exactly $N^n$ ordered $n$-tuples over $\mathbb F_N$, one of
which is zero; this proves the first expression for $\#X$. In the
projective stratum with last nonzero coordinate indexed by $d$, the
displayed isomorphism gives $N^d$ rational points. Summing over $d$
and multiplying the two factors proves \eqref{eq:cg-NY}. Every rational
base point of $Y$ has a fibre isomorphic to $(\mathbb G_m)^2$ over that
same field by Proposition~\ref{prop:cg-torsor}; the fibre has
$(N-1)^2$ rational points. The rational-point map is surjective, since
each point lies in a trivializing chart with the section $(1,1)$.
Partitioning $X(\mathbb F_N)$ by those fibres proves the second
expression in \eqref{eq:cg-NX}. Substituting $q^2=N$ in
\eqref{eq:cg-factor} and \eqref{eq:cg-h} gives the last two formulas
with exactly the powers $N^{-5}$ and $N^{-4}$. Every exponent of $q$
in both expressions is even. Thus these substitutions require no square
root in a finite field and their values do not change under $q\mapsto-q$.

The corresponding counting map is the ring homomorphism
\[
 \nu_m:K_0(\operatorname{Var}_{\mathbb F_Q})[\mathbb L^{-1}]
 \longrightarrow\mathbb Q,\qquad
 [V]\longmapsto\#V(\mathbb F_{Q^m}),\quad\mathbb L^{-1}\longmapsto Q^{-m}.
\]
Indeed, a closed subvariety and its complement partition the finite set
of rational points, and the rational points of a product are the Cartesian
product of the two finite sets. These facts prove that the defining
relations and multiplication are preserved. Since $\nu_m(\mathbb L)=Q^m$
is invertible in $\mathbb Q$, the displayed localization map exists.
\end{proof}

\section{Zeta functions with every sign retained}
\label{sec:cg-zeta}
For a finite-type scheme $V/\mathbb F_Q$, its point-count zeta series is
defined in $1+t\mathbb Q[[t]]$ by
\[
 Z(V,t)=\exp\left(\sum_{m\geq1}\#V(\mathbb F_{Q^m})\frac{t^m}{m}\right).
\]
All exponentials and logarithms in this section are formal power series;
there is no analytic convergence assumption.

\begin{theorem}\label{thm:cg-zeta}
The schemes above have zeta functions
\begin{align}
 Z(X,t)&=\frac{(1-Q^7t)(1-Q^3t)}{(1-Q^{10}t)(1-t)},\label{eq:cg-ZX}\\
 Z(Y,t)&=\prod_{d=0}^8(1-Q^dt)^{-a_d},\label{eq:cg-ZY}\\
 (a_0,\ldots,a_8)&=(1,2,3,3,3,3,3,2,1).
\end{align}
The series of the original Laurent coefficient is
\begin{equation}\label{eq:cg-ZC}
 \exp\left(\sum_{m\geq1}C(Q^{m/2})\frac{t^m}{m}\right)
 =Z(X,Q^{-5}t)
 =\frac{(1-Q^2t)(1-Q^{-2}t)}{(1-Q^5t)(1-Q^{-5}t)}.
\end{equation}
Similarly the series of $h$ equals
$Z(Y,Q^{-4}t)=\prod_{d=0}^8(1-Q^{d-4}t)^{-a_d}$.
\end{theorem}
\begin{proof}
For any $a\in\mathbb Q$, termwise formal differentiation gives
$\frac{d}{dt}\sum_{m\geq1}a^mt^m/m=a/(1-at)$. Both
$\sum_{m\geq1}a^mt^m/m$ and $-\log(1-at)$ have that derivative
and constant term zero; equality follows coefficient by coefficient
in characteristic zero. The expression for $\#X$ is
$Q^{10m}-Q^{7m}-Q^{3m}+1$, so applying this identity to its four
separate terms gives \eqref{eq:cg-ZX} with the two negative terms
in the numerator. Formula \eqref{eq:cg-B} gives \eqref{eq:cg-ZY}
by the same argument for all nine terms. Equation \eqref{eq:cg-countC}
replaces $t^m$ by $(Q^{-5}t)^m$ in the defining series. Substitution
in \eqref{eq:cg-ZX} gives every exponent in \eqref{eq:cg-ZC}.
Equation \eqref{eq:cg-counth} proves the last statement identically.
\end{proof}

\section{Actual compact-support Frobenius groups}
\label{sec:cg-cohomology}
Fix $k=\mathbb F_Q$ of characteristic $p$, an algebraic closure
$\bar k$, and a prime $\ell\ne p$. In this section $H^i_c(V)$ means
the $\ell$-adic group $H^i_c(V_{\bar k},\mathbb Q_\ell)$.
We use geometric Frobenius $\Phi_Q$, the inverse of the field
automorphism $a\mapsto a^Q$, on the Galois representations.
Thus $\Phi_Q$ acts on $\mathbb Q_\ell(-d)$ by $Q^d$.
The four foundational cohomological results used here are affine-space
acyclicity, open--closed localization, K\"unneth, and Poincar\'e duality.
Precise author-hosted references are Milne, \emph{Lectures on \'Etale
Cohomology}, v2.21, Example~16.3, Proposition~18.3 together with the
open--closed sheaf sequence of Proposition~8.15, Theorem~22.4,
and Theorem~24.1,
\url{https://www.jmilne.org/math/CourseNotes/LEC.pdf}.
These are established foundations of the cohomology theory, used with
their stated hypotheses; none is a proposed GCT hypothesis.

We specify the coefficient passage in those references. Put
$\Lambda_r=\mathbb Z/\ell^r\mathbb Z$ for $r\geq1$.
The affine acyclicity, localization, and duality statements first apply
with $\Lambda_r$ coefficients because $\bar k$ is algebraically closed
and $\ell$ is invertible on every scheme here. Compact-support
cohomology of constructible finite sheaves on these finite-type schemes
is finite by Milne, Proposition~18.3(c). At each cohomological degree,
the images in any fixed finite group along the inverse system form a
descending sequence of subgroups and therefore stabilize. This is the
Mittag--Leffler condition, so inverse limits preserve these long exact
sequences. Taking the inverse limit over $r$ and then tensoring over
$\mathbb Z_\ell$ with its fraction field $\mathbb Q_\ell$ is precisely
the cohomology convention of Milne, Section~19, pages~123--126;
the latter tensor product is exact because localization is exact.
The affine top groups obtained from the compatible duality trace maps
are $\Lambda_r(-n)$, whose limit is $\mathbb Z_\ell(-n)$.
Consequently the identifications used below hold as Galois
representations with their full integral Tate indices before tensoring.
The trace maps commute with Galois: Galois permutes the closed-point
classes specifying the trace, and uniqueness in Theorem~24.1(a)
preserves that trace. Finally the arithmetic automorphism acts on
$\mu_{\ell^r}$ by raising to the $Q$th power; its inverse acts on
$\mathbb Q_\ell(1)$ by $Q^{-1}$, proving the stated action on
$\mathbb Q_\ell(-d)$. This Frobenius convention agrees with
Milne, Proposition~29.10 and Section~30.

\begin{lemma}\label{lem:cg-Un}
For $n\in\{3,7\}$ the only nonzero compact-support groups of $U_n$ are
\[
 H^1_c(U_n)=\mathbb Q_\ell,\qquad
 H^{2n}_c(U_n)=\mathbb Q_\ell(-n).
\]
The first identification is the localization connecting map from the
origin; the second is the map induced by the open immersion into
$\mathbb A^n$. Both are Galois-equivariant.
\end{lemma}
\begin{proof}
Affine-space acyclicity says that $H^0(\mathbb A^n)=\mathbb Q_\ell$
and that every positive-degree group vanishes. The scheme is smooth,
geometrically connected, of dimension $n$. Its duality trace therefore
gives $H^{2n}_c(\mathbb A^n)=\mathbb Q_\ell(-n)$ and
$H^i_c(\mathbb A^n)=0$ for $i\ne2n$. The origin is a rational
proper point; its only cohomology group is $H^0_c(\{0\})=\mathbb Q_\ell$
with trivial Galois action. For the open immersion $j:U_n\to\mathbb A^n$
and closed immersion $i:\{0\}\to\mathbb A^n$, the sheaf sequence
\[
 0\longrightarrow j_!\mathbb Q_\ell
 \longrightarrow\mathbb Q_\ell
 \longrightarrow i_*\mathbb Q_\ell\longrightarrow0
\]
is exact: a geometric stalk in $U_n$ gives $0\to\mathbb Q_\ell
\xrightarrow{1}\mathbb Q_\ell\to0$, and a stalk at the origin gives
$0\to0\to\mathbb Q_\ell\xrightarrow{1}\mathbb Q_\ell\to0$.
Its compact-support long exact sequence contains
\[
 0=H^0_c(\mathbb A^n)\longrightarrow\mathbb Q_\ell
 \xrightarrow{\delta}H^1_c(U_n)
 \longrightarrow H^1_c(\mathbb A^n)=0
\]
and
\[
 0=H^{2n-1}_c(\{0\})\longrightarrow H^{2n}_c(U_n)
 \longrightarrow\mathbb Q_\ell(-n)
 \longrightarrow H^{2n}_c(\{0\})=0.
\]
At every other index all adjacent groups are zero, forcing the remaining
groups of $U_n$ to vanish. All schemes and morphisms are defined over
$k$, so the maps, including $\delta$, commute with the Galois action.
The assumptions $n=3,7$ in particular give $2n>1$, keeping the two
displayed degrees distinct.
\end{proof}

\begin{theorem}\label{thm:cg-fullFrob}
The entire compact-support graded Galois representation of $X$ is
\[
\begin{array}{c|c|c|c|c}
 i&\text{K\"unneth summand}&H^i_c(X)&\Phi_Q&(-1)^i\\\hline
 2&H^1_c(U_7)\otimes H^1_c(U_3)&\mathbb Q_\ell&1&+1\\
 7&H^1_c(U_7)\otimes H^6_c(U_3)&\mathbb Q_\ell(-3)&Q^3&-1\\
 15&H^{14}_c(U_7)\otimes H^1_c(U_3)&\mathbb Q_\ell(-7)&Q^7&-1\\
 20&H^{14}_c(U_7)\otimes H^6_c(U_3)&\mathbb Q_\ell(-10)&Q^{10}&+1
\end{array}
\]
Every group in an unlisted degree vanishes. After tensoring every
group with $\mathbb Q_\ell(5)$, its alternating trace is exactly
\[
 \sum_i(-1)^i\operatorname{Tr}
 \left(\Phi_Q^m\mid H^i_c(X)(5)\right)
 =Q^{5m}-Q^{2m}-Q^{-2m}+Q^{-5m}=C(Q^{m/2}).
\]
The alternating determinant of the same twisted representation is
the right side of \eqref{eq:cg-ZC}.
\end{theorem}
\begin{proof}
Over the field $\mathbb Q_\ell$, K\"unneth identifies $H^i_c(X)$
with the direct sum of the tensor products $H^a_c(U_7)\otimes H^b_c(U_3)$
for $a+b=i$; there is no torsion term. Lemma~\ref{lem:cg-Un} leaves
exactly the four pairs $(1,1),(1,6),(14,1),(14,6)$. Tensoring Tate
representations adds their twist indices. This gives exactly the four
degrees, representations and eigenvalues in the table. K\"unneth is
natural over $k$, so these identifications preserve Frobenius, not only
dimensions. The twist $(5)$ multiplies every eigenvalue by $Q^{-5}$.
Taking the indicated signed $m$th powers gives the original four terms
in their original order. Finally each nonzero group is one dimensional,
so $\det(1-t\Phi_Q)=1-t\lambda$ for its listed eigenvalue $\lambda$.
In $\prod_i\det(1-t\Phi_Q\mid H^i_c(X)(5))^{(-1)^{i+1}}$,
the even-degree factors occur in the denominator and the odd-degree
factors in the numerator. This is \eqref{eq:cg-ZC} directly; no
general trace-formula assertion is needed to match the independently
proved point-count zeta function in this case.
\end{proof}

\begin{proposition}\label{prop:cg-YFrob}
The only nonzero groups of $Y$ are
\[
 H^{2d}(Y)=\mathbb Q_\ell(-d)^{\oplus a_d}\quad(0\leq d\leq8),
\]
with $a_d$ as in \eqref{eq:cg-B}. Consequently $h(Q^{m/2})$ is the
trace of $\Phi_Q^m$ on $\bigoplus_dH^{2d}(Y)(4)$.
\end{proposition}
\begin{proof}
We spell out the projective-space calculation from affine cohomology.
The pair $\mathbb P^{n-1}\subset\mathbb P^n$, with open complement
$\mathbb A^n$, gives the localization sequence. Since $\mathbb P^n$
and $\mathbb P^{n-1}$ are proper, compact-support cohomology is ordinary
cohomology for them. The affine compact-support calculation in the proof
of Lemma~\ref{lem:cg-Un} applies to every $n\geq1$. Hence restriction
$H^i(\mathbb P^n)\to H^i(\mathbb P^{n-1})$ is an isomorphism for
$i<2n-1$. At degree $2n-1$, induction from $\mathbb P^0$ gives zero
for the target and therefore for the source. The next part is
\[
 0\longrightarrow\mathbb Q_\ell(-n)
 \longrightarrow H^{2n}(\mathbb P^n)\longrightarrow0.
\]
At degrees greater than $2n$, the same sequence forces vanishing.
Induction gives $H^{2j}(\mathbb P^n)=\mathbb Q_\ell(-j)$ for
$0\leq j\leq n$, with odd groups zero; every map is defined over $k$.
K\"unneth now gives one summand $\mathbb Q_\ell(-(i+j))$ for each
pair $0\leq i\leq6$, $0\leq j\leq2$. Counting the pairs gives
$a_d$ and all claimed Frobenius eigenvalues. The twist $(4)$ changes
$Q^{dm}$ to $Q^{(d-4)m}$. Their sum is $Q^{-4m}B(Q^m)$, already
identified with $h(Q^{m/2})$ in \eqref{eq:cg-counth}.
\end{proof}

\begin{remark}\label{rem:cg-weight}
These explicit eigenvalues determine the weights without a general
Riemann-hypothesis argument. On $H^i_c(X)$ the weights at
$i=2,7,15,20$ are respectively $0,6,14,20$, because an eigenvalue
$Q^d$ has complex absolute value $Q^{(2d)/2}$. On $H^{2d}(Y)$ the
weight is $2d$. Thus the negative terms of $C$ come from the two
odd compact-support degrees of a nonproper variety. They do not
describe a negative dimension or a failure of a Frobenius absolute-value
law. The actual groups of $X$ were computed by localization and
K\"unneth on $U_7\times U_3$; the torsor identity in the class ring
does not assert that $H^*_c(X)$ is $H^*(Y)\otimes H^*_c((\mathbb G_m)^2)$.
\end{remark}

\section{Exact specialization order and the limits of the construction}
\label{sec:cg-specialization}
\begin{proposition}\label{prop:cg-jets}
The original Laurent coefficient has an exact order-two zero at $q=1$:
\[
 C(1)=0,\qquad C'(1)=0,\qquad C''(1)=168,
 \qquad \left.\frac{C(q)}{(q-1)^2}\right|_{q=1}=84.
\]
The residual factor has $h(1)=21$. The full pre-specialization relation is
\begin{equation}\label{eq:cg-jetfactor}
 C(q)=(q-1)^2\,q^{-10}(q+1)^2B(q^2),
\end{equation}
so these numbers arise from the same projective base and the two
separately retained torus factors of $\pi$.
\end{proposition}
\begin{proof}
Equations \eqref{eq:cg-factor}--\eqref{eq:cg-h} give
$C(q)=q^{-10}(q^2-1)^2B(q^2)$, which is exactly
\eqref{eq:cg-jetfactor}. The factor after $(q-1)^2$ is a Laurent
polynomial regular at $q=1$ with value
$1\cdot4\cdot(7\cdot3)=84$, which is nonzero in $\mathbb Z$.
This proves the exact order, and differentiating the product twice
gives $C''(1)=2\cdot84=168$. Directly differentiating the original
four terms checks the preserved signs:
\[
 C'(1)=10-4+4-10=0,\qquad
 C''(1)=90-12-20+110=168.
\]
The formula for $h(1)$ follows from $[7]_1=7$ and $[3]_1=3$,
or from the 21 affine strata of $Y$. The two factors $q^2-1$
are the polynomial images under $T\mapsto q^2$ of the two
$\mathbb G_m$ classes $T-1$ in Proposition~\ref{prop:cg-classes}.
\end{proof}

The exact realization obtained here consists of the morphism
$X\to Y$, its two scalar actions, the class identities, the counting
homomorphisms, the even-Laurent substitution $T\mapsto q^2$, and the
specified graded Frobenius representations and Tate twists. In particular,
the original coefficient is realized as an alternating trace of an actual
scheme's compact-support cohomology, while its residual $[7]_q[3]_q$
is the centered trace of the displayed projective product. All their
maps and constants have been specified above. No half-integral Tate
twist is used: the formal ring $\mathbb Z[q^2,q^{-2}]$ maps to the
ring generated by integral Tate representations by
$q^{2d}\mapsto[\mathbb Q_\ell(-d)]$.
More precisely, the target is the Grothendieck ring of finite-dimensional
continuous $\mathbb Q_\ell$-representations of
$\operatorname{Gal}(\bar k/k)$, restricted to the subring generated by
the Tate classes. Tensor products satisfy
$\mathbb Q_\ell(-d)\otimes\mathbb Q_\ell(-e)
=\mathbb Q_\ell(-(d+e))$, and the class of
$\mathbb Q_\ell(d)$ is the multiplicative inverse of the class of
$\mathbb Q_\ell(-d)$. These identities establish the stated ring
homomorphism. On any invariant exact sequence, the matrix of
$\Phi_Q^m$ is block triangular and its trace is additive; on a tensor
product its trace is the product of the two traces. Thus taking this
trace defines a ring homomorphism to $\mathbb Q_\ell$, and its
composition on $q^{2d}$ is exactly $Q^{dm}$. It agrees with the
even-Laurent specialization already used in the point-count formulas.

The scheme and torsor realize this retained scalar with all stated
cohomological degrees. Section~\ref{sec:fgb} constructs actual
additive generator functors, their integral decategorification map
to the source interval, and a separately specified cross-path relation
cone. Its explicit Frobenius-equivariant endpoint map uses precisely
these cohomology groups after two-periodization; the integer grading
and the chosen extra relation data are recorded there. Compatibility
with canonical-basis convolution in the full original source module
remains unproved. The source endpoint itself is resolved by the
original action paths in Section~\ref{sec:gct-generator-interval},
with the full source module in Section~\ref{sec:full-source-generators}.

\section{Graded generator functors and an explicit Frobenius relation cone}
\label{sec:fgb}

We retain the four actual source classes in their fixed order
\(e_0=T,e_1=T_1,e_2=T_2,e_3=Q\), with precisely the source rescalings
and column words proved in the generator-interval calculation. Put
\[
 \begin{gathered}
 A=\mathbb Z[q,q^{-1}],\quad
 \alpha=[7]-[3],\quad\beta=[8],\quad\gamma=[6],\quad d=[3],\\
 \delta=-d,\quad p=[2],\quad C=q^{10}-q^4-q^{-4}+q^{-10}.
 \end{gathered}
\]
The source quotient has lowering matrix
\[
 N=\begin{pmatrix}0&0&0&0\\\alpha&0&0&0\\
 \beta&0&0&0\\0&\gamma&-d&0\end{pmatrix},\qquad
 D=CE_{30},\qquad N^2=pD.
\]
The construction below gives actual additive functors defined by tensor
products, an integral map
between their two source-labelled length-two branches, and a deformation
retraction after two-periodization to the previously computed Frobenius
object. Its relation map is specified explicitly; it is not inferred from
the equality of Laurent polynomials.

\subsection{The state category and its exact integral decategorification}

Let \(\mathcal G\) be the category of finite free internally
\(\mathbb Z\)-graded abelian groups, each furnished with a homogeneous
basis. Write \(\mathbb Z\langle w\rangle\) for a copy of \(\mathbb Z\)
whose basis vector has internal degree \(w\). Morphisms preserve internal
degree. Let \(\mathcal K=K^b(\mathcal G)\) be its bounded homotopy category,
and let \(\mathcal S=\mathcal K^4\). The tensor differential convention is
\[
 d(x\otimes y)=dx\otimes y+(-1)^i x\otimes dy
 \quad(x\text{ of cohomological degree }i).
\]
Internal degrees add under tensor product. The internal shift
\(\langle w\rangle\) increases every internal degree by \(w\).
The cochain translation convention is \(V[1]^i=V^{i+1}\), with
translated differential \(-d\); in particular a degree-zero group
translated by \([1]\) lies in degree \(-1\).

Define the decategorification group using isomorphism classes in
\(\mathcal K\), direct-sum relations, and
\([\operatorname{Cone}(f)]=[Y]-[X]\) for every cochain map
\(f:X\to Y\). Its internal shift gives an \(A\)-module structure.
There is an isomorphism
\[
 \chi_q:K_0(\mathcal K)\longrightarrow A,
 \qquad [V]\longmapsto
 \sum_{i,w}(-1)^i\operatorname{rank}_{\mathbb Z}(V^i_w)q^w.
 \tag{*}
\]
Here and below sums are finite. To prove the assertion, a direct sum adds
ranks, and \(\operatorname{Cone}(f)^i=Y^i\oplus X^{i+1}\) gives exactly
the required subtraction. A contractible complex has Euler rank zero
in each internal degree: its contracting homotopy splits its successive
boundaries, so the adjacent boundary contributions cancel. Alternatively,
the cone of a homotopy equivalence is contractible, and its finite
alternating chain ranks vanish by the same splitting. Thus the formula
is unchanged under isomorphism in the homotopy category. Filtering a
bounded complex by its cohomological degrees and using the cone relation
expresses its class as the alternating sum of the classes of its chain
groups. Each such group is a direct sum of the specified
\(\mathbb Z\langle w\rangle\). This proves that the map
\(q^w\mapsto[\mathbb Z\langle w\rangle]\) is inverse to
\((\ast)\). The tensor formula and the displayed tensor differential
show that \(\chi_q\) also preserves products. The same argument in the
four factors gives a fixed isomorphism
\[
 \Psi:K_0(\mathcal S)\longrightarrow M_A=\bigoplus_{i=0}^3Ae_i,
 \qquad[(V_0,V_1,V_2,V_3)]\longmapsto
 \sum_{i=0}^3\chi_q(V_i)e_i.
\]
In particular the four labels have not been replaced or permuted.

\subsection{Objects on the four original arrows}

For each indicated finite set, retain the displayed label as part of
the basis. Define degree-zero complexes
\[
 \begin{aligned}
 \mathsf A&=\bigoplus_{a\in\{-6,-4,4,6\}}\mathbb Z a,
       &\deg_q(a)&=a,\\
 \mathsf B&=\bigoplus_{b\in\{-7,-5,-3,-1,1,3,5,7\}}\mathbb Z b,
       &\deg_q(b)&=b,\\
 \mathsf G&=\bigoplus_{g\in\{-5,-3,-1,1,3,5\}}\mathbb Z g,
       &\deg_q(g)&=g,\\
 \mathsf J&=\bigoplus_{j\in\{-2,0,2\}}\mathbb Z j,
       &\deg_q(j)&=j.
 \end{aligned}
\]
The equality \(\chi_q(\mathsf A)=[7]-[3]\) follows by retaining
the seven terms of \([7]\) and subtracting its three terms with
exponents \(2,0,-2\); the four surviving terms are exactly the four
displayed summands. This equality does not rescale any source class.
The other three positive characters are \([8],[6],[3]\), by their
complete finite sums. The negative source sign is implemented by the
cochain translation \(\mathsf J[1]\), which has character \(-[3]\).

Define an exact additive functor \(\mathsf F:\mathcal S\to\mathcal S\)
using tensor products by
\[
 \mathsf F(V)_0=0,\qquad
 \mathsf F(V)_1=\mathsf A\otimes V_0,\qquad
 \mathsf F(V)_2=\mathsf B\otimes V_0,\qquad
 \mathsf F(V)_3=(\mathsf G\otimes V_1)
                       \oplus(\mathsf J[1]\otimes V_2).
\]
Its action on a tuple of cochain maps is the corresponding identity
tensor those maps. Tensoring by these finite free complexes preserves
mapping cones, with the usual tensor-differential signs, so it induces
an endomorphism of the stated decategorification group. Direct use of
the preceding formulas proves the intertwining identity
\[
                 \Psi\,[\mathsf F]=N\,\Psi.
\]
This is an integral identity over the original ring \(A\).

Let \(\mathsf K_V\) shift the four components internally by
\((9,7,7,5)\), respectively, and let \(\mathsf K_W\) shift each
component by \(3\). Their inverses are the opposite shifts. These are
actual invertible functors. For every nonzero edge of \(\mathsf F\),
the target \(V\)-shift is the source shift minus two. Identifying equal
internal shifts on each tensor factor therefore gives a natural
isomorphism
\[
 \mathsf K_V\mathsf F\mathsf K_V^{-1}
                  \simeq\mathsf F\langle-2\rangle.
\]
The \(W\)-shifts are equal on both ends of each edge and commute
with \(\mathsf F\). Set \(\mathsf F_W=0\). Its commutation and
Cartan relations hold as equal zero functors. The two Cartan functors
commute componentwise. Each of the two original central total-weight
generators is represented by the common internal shift \(\langle15\rangle\)
on all four components; these commute with every displayed functor.
Finally \(\mathsf F^3=0\): the only length-two
paths end in component three, from which there is no outgoing arrow.
Thus the negative-Borel relations involving the first lowering
generators and Cartan operators hold already as specified natural
isomorphisms, not only as scalar equalities.

\subsection{The retained geometric scalar object}

Let \(\mathsf C_X\) be the zero-differential integral bigraded complex
with exactly four basis vectors
\[
\begin{array}{c|rrrr}
\text{label}&x_{00}&x_{03}&x_{70}&x_{73}\\\hline
\text{cohomological degree}&2&7&15&20\\
\text{internal degree}&-10&-4&4&10
\end{array}
\]
These labels retain, in order, the K\"unneth factors
\((H^1_c(U_7),H^1_c(U_3))\),
\((H^1_c(U_7),H^6_c(U_3))\),
\((H^{14}_c(U_7),H^1_c(U_3))\), and
\((H^{14}_c(U_7),H^6_c(U_3))\).
The actual groups and their localization and trace identifications
were computed for
\(X=(\mathbb A^7\setminus0)\times(\mathbb A^3\setminus0)\).
Their Tate twists after the common twist \((5)\) are respectively
\((5),(2),(-2),(-5)\). Thus the internal degrees retain twice the
Frobenius exponents and \(\chi_q(\mathsf C_X)=C\), with every
cohomological degree and sign retained.

Define \(\mathsf D(V)\) to have only component three nonzero, equal
to \(\mathsf C_X\otimes V_0\). Then
\(\Psi[\mathsf D]=D\Psi\). Its Cartan relation uses the original
weight difference \(-4\). The support also gives
\(\mathsf F\mathsf D=\mathsf D\mathsf F=\mathsf D^2=0\).
Put \(\mathsf H=\mathbb Z\langle1\rangle\oplus
\mathbb Z\langle-1\rangle\) in cohomological degree zero, so that
\(\chi_q(\mathsf H)=p\).

The following obstruction is exact for these specified objects.
On the state \((\mathbb Z,0,0,0)\), the composite
\(\mathsf F^2\) has component-three complex
\[
  (\mathsf G\otimes\mathsf A)
     \oplus(\mathsf J\otimes\mathsf B)[1],
\]
whose zero differential leaves 24 even and 24 odd free generators.
The complex \(\mathsf H\otimes\mathsf C_X\) has four even and four
odd generators. Tensoring with \(\mathbb Q\), an isomorphism in either
the bounded homotopy category or the two-periodic derived category
would induce an isomorphism of its cohomology. The indicated dimensions
are unequal. Therefore these specified functors have no natural
isomorphism \(\mathsf F^2\simeq\mathsf H\otimes\mathsf D\).
This obstruction proves the failure of that lift; the next map supplies
an exact stronger correspondence between its two branches and the
geometric endpoint.

\subsection{A complete integral map between the two length-two branches}

Write
\[
 \mathsf P=\mathsf G\otimes\mathsf A,\qquad
 \mathsf R=\mathsf J\otimes\mathsf B.
\]
For basis naming use pairs \((a,g)\) in \(\mathsf P\) and \((b,j)\)
in \(\mathsf R\), preserving their source intermediate states
\(T_1\) and \(T_2\), respectively. At each internal weight \(w\),
list the pairs of either type in ascending lexicographic order and
let their lengths be \(p_w,r_w\). Define
\[
 \theta:\mathsf R\longrightarrow\mathsf P,
 \qquad
 \theta(R_{w,k})=
 \begin{cases}P_{w,k},&1\leq k\leq\min(p_w,r_w),\\
 0,&k>p_w.
 \end{cases}
\]
Here \(R_{w,k}\) denotes the \(k\)-th source pair and \(P_{w,k}\)
the \(k\)-th target pair of the given weight. This defines an actual
matrix over \(\mathbb Z\) consisting of twenty distinct entries one
and all remaining entries zero. It preserves internal and cohomological
degree. Its exact ranks and unmatched vectors are
\[
\begin{array}{c|r|l|l}
w&\operatorname{rank}\theta_w&\text{unmatched }(a,g)
                                      &\text{unmatched }(b,j)\\\hline
11&0&(6,5)&\\
9&1&(6,3)&\\
7&2&&\\
5&2&&(7,-2)\\
3&2&&(5,-2)\\
1&3&&\\
-1&3&&\\
-3&2&&(-1,-2)\\
-5&2&&(-3,-2)\\
-7&2&&\\
-9&1&(-4,-5)&\\
-11&0&(-6,-5)&
\end{array}
\]
For verification, the \(\mathsf P\)-basis is the Cartesian product
of the four explicitly given values of \(a\) and the six values of
\(g\); the \(\mathsf R\)-basis is the product of the eight values
of \(b\) and the three values of \(j\). Adding the two entries gives
exactly the twelve rows in the table. Listing and matching within one
row cannot affect any other row. The rank column sums to twenty,
and each side has 24 basis vectors, so each side has exactly the
four listed residual generators. Because all matched entries are
one on distinct rows and columns, the image is a direct summand and
both the kernel and cokernel are free, with exactly those residual
bases. There is no unrecorded integral torsion.

Let \(\mathsf L=\operatorname{Cone}(\theta)\). It has
\(\mathsf R\) in cohomological degree \(-1\), \(\mathsf P\) in
degree zero, and differential exactly \(\theta\). In this convention
the sign of its character is
\[
 \chi_q(\mathsf L)=\gamma\alpha-d\beta
                 =pC.
\]
More strongly, the complex splits integrally as twenty two-term
complexes \(\mathbb Z\xrightarrow{1}\mathbb Z\) and eight isolated
residual generators. The split is the displayed basis permutation,
not a rational rank argument. Map each matched target basis vector
back to its matched source basis vector and all other basis vectors
to zero; call this cochain homotopy \(h\). If \(\rho\) is projection
to the eight residual vectors and \(\iota\) is their inclusion, then
on a matched pair \(dh+hd\) is the identity on both vectors and on
a residual vector it is zero. Consequently
\[
                    1-\iota\rho=dh+hd,
                    \qquad \rho\iota=1.
\]
This proves an explicit integral deformation retraction and all its
signs.

Regard \(\mathsf P,\mathsf R,\mathsf L\) also as endpoint functors
on \(\mathcal S\), tensoring their arguments from component zero
into component three. The map \(\theta\otimes1\) is natural in
every state and morphism. Its mapping-cone triangle is therefore
a triangle of explicitly specified functors
\[
             \mathsf R\xrightarrow{\theta}\mathsf P
                    \longrightarrow\mathsf L\longrightarrow\mathsf R[1].
\]
The original composite \(\mathsf F^2\) is the split sum
\(\mathsf P\oplus\mathsf R[1]\), not this cone with its nonzero
cross-branch differential. Their equality in the decategorification
group follows from this actual triangle. Their nonisomorphism follows
from the already computed cohomology dimensions. In particular,
the cross-branch relation has not silently been installed as the
tensor differential of \(\mathsf F^2\).

\subsection{The endpoint map with all cohomological degrees recorded}

Let \(\operatorname{Per}\) send a bounded complex to its two-periodic
complex by summing its even and odd cohomological degrees separately,
retaining the differential, internal degree, and individual summand
labels. This operation explicitly forgets the distinction between
cohomological degrees of the same parity. The eight residual
generators of \(\mathsf L\) identify with the eight indicated tensor
generators of \(\mathsf H\otimes\mathsf C_X\) as follows. The symbol
\(+\) means the \(\mathsf H\)-basis of internal degree one and
\(-\) means its basis of internal degree minus one:
\[
\begin{array}{c|c|r|r}
\text{residual path}&\mathsf H\otimes\mathsf C_X
 &\text{internal degree}&\text{original endpoint cohomological degree}\\\hline
(6,5)&+\otimes x_{73}&11&20\\
(6,3)&-\otimes x_{73}&9&20\\
(-4,-5)&+\otimes x_{00}&-9&2\\
(-6,-5)&-\otimes x_{00}&-11&2\\
(7,-2)&+\otimes x_{70}&5&15\\
(5,-2)&-\otimes x_{70}&3&15\\
(-1,-2)&+\otimes x_{03}&-3&7\\
(-3,-2)&-\otimes x_{03}&-5&7
\end{array}
\]
The first four residual vectors are in cohomological degree zero
in \(\mathsf L\), while the last four are in degree \(-1\).
Each row preserves parity and internal degree. Extend this table
linearly to an isomorphism of the residual two-periodic complex
with \(\operatorname{Per}(\mathsf H\otimes\mathsf C_X)\).
Composition with \(\rho\), and in the other direction with
\(\iota\), gives inverse homotopy equivalences
\[
 \operatorname{Per}(\mathsf L)
       \simeq\operatorname{Per}(\mathsf H\otimes\mathsf C_X),
\]
with precisely the contraction \(h\) proved above. This supplies
the actual morphism relating all surviving branch summands to the
scalar geometry. The cohomological columns also prove that this
map is not a degree-preserving map to the unperiodized scalar
cohomology: for example its weight-eleven class has degree zero
on the cone side and degree twenty on the geometric side.
Moreover the differences for the weights eleven and minus eleven
are twenty and two, respectively, so no single common cohomological
shift repairs this failure.

\subsection{Frobenius realization and the integral-Tate parity obstruction}

Fix a prime power \(Q\), a prime \(\ell\) not dividing \(Q\), and
choose a root \(z\) of \(z^2=Q\) in
\(E=\mathbb Q_\ell(z)\). This is a specified extra coefficient
choice for the odd internal degrees; it was unnecessary for the
even-degree scalar object. The element \(z\) is a unit in the
ring of integers of \(E\). For each power of its maximal ideal,
its reduction has finite multiplicative order, hence
\(n\mapsto z^n\) extends from \(\mathbb Z\) to a continuous
homomorphism from \(\widehat{\mathbb Z}\) to that finite unit
group. The extensions commute with reduction because they agree
on the dense integer subgroup. Their inverse limit is a continuous
character \(\widehat{\mathbb Z}\to E^\times\).
Identifying the Galois group of \(\overline{\mathbb F}_Q/\mathbb F_Q\)
with \(\widehat{\mathbb Z}\), with geometric Frobenius sent to one,
therefore defines an actual one-dimensional continuous representation
\(\mathcal L\) on which geometric Frobenius acts as \(z\).
Its square is identified, on their specified basis vectors, with
\(E(-1)\); both characters send geometric Frobenius to \(Q\).

Define the tensor realization on every graded basis summand by
\[
 \mathbb Z\langle w\rangle\longmapsto\mathcal L^{\otimes w},
\]
using dual powers for negative \(w\), extending scalars in the
integer matrices to \(E\), and retaining every cohomological degree.
Maps preserving internal degree become Frobenius-equivariant maps,
because every nonzero matrix entry joins the same character on
source and target. In particular the twenty entries of \(\theta\),
the projection, inclusion and contraction, and every entry of the
residual table are Frobenius equivariant. The actual geometric
object \(H^*_c(X,E)(5)\) is the realization of \(\mathsf C_X\)
by its previously proved localization and K\"unneth identifications.
Here \(H^*_c\) explicitly denotes the direct sum of the actual
cohomology groups placed in their original cohomological degrees with
zero differential. It does not denote a chosen cochain model for
\(R\Gamma_c(X,E)\). No canonical splitting of that derived object
into its cohomology groups is asserted or used.
We obtain a specified equivalence of two-periodic complexes of
continuous Frobenius representations
\[
 \operatorname{Per}(\mathsf L_E)
 \simeq\operatorname{Per}\bigl((\mathcal L\oplus\mathcal L^{-1})
                                  \otimes H^*_c(X,E)(5)\bigr).
\]
Every entry of the map and its inverse has been given above. Its
alternating trace of the \(m\)-th Frobenius power is
\[
 (z^m+z^{-m})(Q^{5m}-Q^{2m}-Q^{-2m}+Q^{-5m}).
\]
If the chosen complex embedding of \(z\) is the positive square
root, this is precisely \([2]_{Q^{m/2}}C(Q^{m/2})\). The two choices
of \(z\) are retained choices of odd-degree realization; the four
even-degree scalar summands are independent of that choice.

The need for this extension has a precise parity obstruction. Suppose
one tries to realize all four arrow coefficients and the divided-square
coefficient by only integral Tate degrees after replacing the four
state generators by \(q^{s_i}e_i\), for integers \(s_i\). A coefficient
from state \(i\) to state \(j\) becomes multiplied by
\(q^{s_i-s_j}\). Integral Tate characters have only even powers of
\(q\). The nonzero exponent parities of \(\alpha,\beta,\gamma,-d\)
are respectively \(0,1,1,0\). Thus all four arrows being even requires
\[
 s_0-s_1=0,\quad 1+s_0-s_2=0,\quad
 1+s_1-s_3=0,\quad s_2-s_3=0\pmod2.
\]
The first and third imply \(s_3-s_0=1\pmod2\). But every term of
\(C\) has even exponent, so the divided-square arrow from zero to
three being even would require \(s_3-s_0=0\pmod2\), a contradiction.
This proves that diagonal state regradings cannot provide that
all-integral-Tate lift. The character \(\mathcal L\) just constructed
is an explicit alternative, not an unmentioned half-integral Tate
twist in the original scalar calculation.

The construction therefore gives an integral four-state generator
functor, all its first-generator negative-Borel relations, a precise
decategorification morphism to the actual source subquotient, a complete
rank-twenty relation map, and its two-periodic Frobenius endpoint
equivalence. It also records the exact obstruction to the naive divided
power functor isomorphism and to restoring the full cohomological
grading by a common shift. It does not identify the added relation map
with a canonical-basis convolution map in the full source module.
Constructing that source geometric compatibility remains a further
mathematical problem; no nonstandard Riemann-hypothesis or complexity
conclusion is supplied by the finite interval alone.

\section{The source GCT coefficient on the retained arithmetic cover}
\label{sec:conic-arithmetic}
The GCT IV calculation quoted and checked in the preceding source
section uses the Laurent parameter $q$ and the original coefficient
\[
 C(q)=q^{10}-q^4-q^{-4}+q^{-10}
 =(q-q^{-1})^2[7]_q[3]_q,\qquad
 [j]_q=\frac{q^j-q^{-j}}{q-q^{-1}}.
\]
Here $[j]_q$ denotes its Laurent-polynomial expression at $q=\pm1$
as well; division in the displayed expression is an identity away
from these points, not deletion of those points from the polynomial.
We now construct an exact map to the original retained cover.

\begin{proposition}[A base change with all coordinates retained]
\label{prop:conic-base-change}
For a field $k$ of characteristic different from $2$, set
\[
 D=k[a,r]/(r^2+a),\qquad
 E=k[p]\otimes_{k[a]}D,\quad a\longmapsto4-p^2.
\]
There is an isomorphism
\[
 E=k[p,r]/(p^2-r^2-4)\ \simeq\ k[q,q^{-1}]
\]
with maps
\[
 p\longmapsto q+q^{-1},\quad r\longmapsto q-q^{-1},
 \qquad q\longmapsto(p+r)/2,\quad q^{-1}\longmapsto(p-r)/2.
\]
Its original base coordinate is exactly
$a=4-p^2=-(q-q^{-1})^2$.
\end{proposition}
\begin{proof}
Tensoring the displayed presentation of $D$ replaces $a$ by $4-p^2$,
giving the stated relation. The Laurent expressions satisfy that
relation. Conversely $(p+r)(p-r)/4=1$, so the proposed inverse images
of $q,q^{-1}$ have product one and define a Laurent ring map.
The sums and differences of these expressions recover $p,r$.
Starting with $q$, the same compositions recover $q,q^{-1}$.
Thus the maps are inverse on generators, proving the isomorphism
and every stated base-coordinate identity.
\end{proof}

The variable $p$ is new; it does not rename or replace the original
third source coordinate $w$. On the original simple-root chart with
$b=c=0$, the retained formulas become
\[
 \alpha=-2r,\quad x=-\frac1{2r},\quad y=3r,\quad w=26r^2,
 \qquad a=-r^2,\quad r\ne0.
\]
Indeed substitution into $\alpha=P'/2$ and the full inverse formulas
gives each equality with its original coefficient.
Substitution into the original $F$ then gives $(a,b,c)=(-r^2,0,0)$.
After Proposition~\ref{prop:conic-base-change}, these are rational
functions of $q$ defined precisely on $q\ne0,\pm1$:
$q=0$ was already absent from the Laurent parameter space, and
$q=\pm1$ is equivalent to $r=0$.
For each such $q$, the finite boundary point $(0,0,a)$ in the original
fibre is retained separately. It is not a point of this $x\ne0$ chart.

The exact positive Laurent factor has the polynomial expression
\[
 [7]_q=p^6-5p^4+6p^2-1,\qquad [3]_q=p^2-1.
\]
For a direct verification, set $U_0=1,U_1=p$ and
$U_j=pU_{j-1}-U_{j-2}$. Substituting $p=q+q^{-1}$ makes
$U_j=q^j+q^{j-2}+\cdots+q^{-j}$, by cancellation of the interior
terms in the recurrence. Applying this through $j=6$ proves the
two displayed polynomials.
Consequently the original coefficient, without discarding its
vanishing factor, is
\begin{align}
 C(q)&=-a\bigl((4-a)^3-5(4-a)^2+6(4-a)-1\bigr)(3-a)\notag\\
     &=-a(7-14a+7a^2-a^3)(3-a).
 \label{eq:C-a}
\end{align}
The second expression follows by expanding the first, which remains
displayed to record the substitution.
The coefficient $-a$ is part of the original object in both expressions.
At $q=1$, $a=-(q-1)^2(q+1)^2/q^2$.
The other factors take values $7$ and $3$, so
\[
 \left.\frac{C(q)}{(q-1)^2}\right|_{q=1}=4\cdot7\cdot3=84.
\]
In characteristic zero this proves a zero of order exactly two.
It proves a zero of order one in the original base coordinate $a$:
the coefficient of $a$ in \eqref{eq:C-a} is $-21$.
The order change is precisely the displayed ramification, not a
lost specialization layer.

\subsection{The resulting rank-four arithmetic extension}
Now take $k=\mathbb C$ and retain $a=2s$.
The same ring is finite free over $\mathbb C[s]$:
\begin{equation}\label{eq:quartic-q}
 \mathbb C[q,q^{-1}]
 \simeq\mathbb C[s,q]/\bigl(q^4+(2s-2)q^2+1\bigr).
\end{equation}
To prove this, the relation gives
$q^{-1}=-q^3+(2-2s)q$, so $q$ is already invertible in the quotient.
Conversely the Laurent map $s=-(q-q^{-1})^2/2$ gives the quartic.
The two substitutions recover $s$ and $q$, proving inverse ring maps.
Division by the monic quartic gives the free basis $(1,q,q^2,q^3)$.
In the conic presentation the alternative basis $(1,r,p,rp)$ is
free over $\mathbb C[s]$ by the two monic relations
$r^2=-2s$, $p^2=4-2s$.
The conversion identities are
\[
 r=q^3+(2s-1)q,\quad
 p=-q^3+(3-2s)q,\quad rp=2q^2+2s-2,\quad q=(p+r)/2.
\]
They follow from the expression for $q^{-1}$ and prove the basis
conversion in both directions (also $q^2=(rp+2-2s)/2$).

For the actual topological $\mathbb C[s]$-module $(Q,S)$ used in
the retained arithmetic construction, form
$Q^{(4)}=\mathbb C[q,q^{-1}]\otimes_{\mathbb C[s]}Q$.
Give it the fourfold direct-sum topology in the monic basis.
Multiplication by $q$ has the continuous block matrix
\begin{equation}\label{eq:quartic-operator}
 K=
 \begin{pmatrix}
 0&0&0&-I\\
 I&0&0&0\\
 0&I&0&2I-2S\\
 0&0&I&0
 \end{pmatrix},\qquad
 \widehat S=\operatorname{diag}(S,S,S,S).
\end{equation}
The last column is exactly the quartic relation. Thus
$K^4+(2\widehat S-2I)K^2+I=0$ and
$K^{-1}=-K^3+(2I-2\widehat S)K$.
Both are continuous polynomial expressions in commuting block
operators. The two original root and trace coordinates are
\[
 R=K-K^{-1},\quad P=K+K^{-1},\qquad
 R^2=-2\widehat S,\quad P^2=4I-2\widehat S.
\]
These identities follow either by the proved ring maps or by direct
substitution of the inverse polynomial.
In the alternative ordering $(1,r,p,rp)$, multiplication by $r$
is the direct sum of two copies of
$\left(\begin{smallmatrix}0&-2S\\I&0\end{smallmatrix}\right)$.
Thus this is the exact two-stage extension of the parent's two-sheet
module, with its factors and topology retained.
Projection to the constant coefficient and its section show, by the
same inverse construction used earlier, that
$\operatorname{Spec}\widehat S=\operatorname{Spec}S$.

Equation~\eqref{eq:C-a} also supplies a completely specified operator:
\begin{equation}\label{eq:coefficient-operator}
 C(K)=
 -2\widehat S(7I-28\widehat S+28\widehat S^2-8\widehat S^3)
 (3I-2\widehat S).
\end{equation}
Laurent evaluation is legitimate because $K$ has the explicit inverse.
This proves a concrete transport of a source GCT coefficient into the
retained arithmetic module. It does not assert that all nonstandard
quantum-group generators act on $Q^{(4)}$: only this scalar Laurent
parameter algebra and its proved action have been constructed.

\subsection{The original coefficient heat flow on this cover}
Use material time $\tau$, which generates the forward coefficient heat
flow of the retained polynomial, with $a'=2$, $b'=6c$, $c'=0$.
The arithmetic heat parameter $t$ acts instead by $-\partial_s^2/4$;
Section~\ref{sec:full-material-interface} proves its exact pullback and
retains the original Euclidean diffusion coefficients.
For the entire original cubic,
$\partial_\tau(cr^3-2r^2+br-2a)=6cr-4=\partial_r^2P$.
On $b=c=0$ this is $P=-2(r^2+a)$ with $a'=2$.
Implicit differentiation at a simple root gives $r'=-1/r$.
On the conic the second relation $p^2=4-a$ gives $p'=-1/p$
when $p\ne0$. Therefore the exact lifted coefficient derivation is
\[
 D_\tau=-\frac{q}{rp}\frac{d}{dq},\qquad
 D_\tau a=2,\quad D_\tau s=1,\quad
 D_\tau r=-\frac1r,\quad D_\tau p=-\frac1p.
\]
Indeed $dr/dq=1+q^{-2}=p/q$ and $dp/dq=1-q^{-2}=r/q$;
multiplication by the displayed coefficient proves all four identities.
Its domain is the Laurent curve with $rp\ne0$, equivalently
$q\notin\{0,1,-1,i,-i\}$ over $\mathbb C$.
The points $q=\pm1$ have the original excluded root $r=0$.
At $q=\pm i$ one has $p=0$ and $r=\pm2i$; these are the new
ramification points of the additional $p$-cover over $s=2$,
while the original $r$-chart remains finite.
Thus both exceptional loci and their different coordinate causes
are recorded by explicit equations.

For the original source coordinates this heat derivative is
\[
 D_\tau x=-\frac1{2r^3},\qquad
 D_\tau y=-\frac3r,\qquad
 D_\tau w=-52.
\]
Differentiating the proved identity gives
\[
 F(-1/(2r),3r,26r^2)=(-r^2,0,0),\qquad
 DF\,(D_\tau x,D_\tau y,D_\tau w)^{\mathsf T}=(2,0,0)^{\mathsf T}.
\]
Since $\det DF=-2$, this is exactly the original inverse-Jacobian
heat lift on the stated slice. The boundary point has derivative
$(0,0,2)$ separately.
For the GCT source-displayed coefficient polynomial the chain rule, with the full
polynomial $C_a(a)$ in \eqref{eq:C-a}, gives
$D_\tau C(q)=2C_a'(a)$. No analytic unit or vanishing factor is
deleted in this derivation. This is an action on the localized
coefficient ring. A differentiation action on the topological
arithmetic quotient is not inferred from its multiplication action;
the latter is exactly the operator action already proved in
\eqref{eq:quartic-operator}--\eqref{eq:coefficient-operator}.

The source identifies the classical spectral coordinate through
$\Xi(s)=\xi(1/2+is)=8H_0(2s)$.
Nothing in the finite-free extension changes that original coordinate.
A certified off-real spectral point of $S$ would remain off-real in
each direct summand; a nonreal eigenvalue of $K$ alone has no such
implication. The exact identity $\widehat S=-(K-K^{-1})^2/2$ gives
$S=\pi[-(K-K^{-1})^2/2]\iota$, with constant-coefficient projection
$\pi:Q^{(4)}\to Q$ and section $\iota:Q\to Q^{(4)}$.
This is the required test for the original coordinate.
None of these identities supplies
an off-real $s$, a positivity proof for all GCT coefficients, or a
complexity lower bound. They supply the exact common algebra and the
specific transported coefficient on which further calculations can act.

\subsection{The exact positivity test under two real structures}
On $\mathbb C[q,q^{-1}]$ define two conjugate-linear involutions by
$q^{*_{\rm real}}=q$ and $q^{*_{\rm circle}}=q^{-1}$, conjugating
complex coefficients in each case. Multiplicativity on Laurent
monomials and the fact that applying either map twice fixes $q$
prove both involution identities on the entire commutative algebra.
The conic isomorphism gives
\[
 p^{*_{\rm real}}=p,\quad r^{*_{\rm real}}=r,\qquad
 p^{*_{\rm circle}}=p,\quad r^{*_{\rm circle}}=-r.
\]
Both fix $a=-r^2$ and $s=a/2$. Thus they extend the same real structure
on the arithmetic base with different root-coordinate adjoints.

Evaluation at $q_0\ne0$ preserves $*_{\rm real}$ exactly when
$\overline{q_0}=q_0$, and preserves $*_{\rm circle}$ exactly when
$\overline{q_0}=q_0^{-1}$. These say, respectively, that $q_0$ is
nonzero real and that $|q_0|=1$.
For positive real $q_0$, all monomials of $[7]_q[3]_q$ are positive
and $(q-q^{-1})^2\ge0$, so $C(q_0)\ge0$.
For a precise circle character instead take
\[
 q_0=\frac{\sqrt{14}+i\sqrt2}{4},\quad
 q_0^{-1}=\frac{\sqrt{14}-i\sqrt2}{4},\quad
 p_0=\frac{\sqrt{14}}2,\quad r_0=\frac{i}{\sqrt2}.
\]
Products and sums verify $|q_0|=1$, $a_0=1/2$, $s_0=1/4$.
The full coefficient in \eqref{eq:C-a} then gives
\[
 C(q_0)=-\frac12
 \left(7-7+\frac74-\frac18\right)\frac52=-\frac{65}{32}.
\]
This exactly evaluates the source's mixed-sign coefficient, retaining
its positive residual Laurent factor and signed vanishing factor.
It rules out the particular claim of nonnegativity at every circle
character. It proves neither a failure of the source's stated
canonical-basis positivity nor a negative value on an established
zeta spectral point: $s_0=1/4$ has not been asserted to belong to that
spectrum. The maps and involutions above specify the entire comparison.

\section{The full material generator and the extended coefficient connection}
\label{sec:full-material-interface}
The following calculation incorporates the parent's complete material
generator at its original coefficients. Its source path, reading scope
and hash are recorded in \texttt{sources.json}. We rederive every
algebraic and differential identity used here from the displayed $F$.
The material coordinate is denoted $\mathbf z=(x,y,w)$ in this section;
the letter $q$ continues to denote the GCT Laurent parameter.

\subsection{The drift, its exact pullback, and the Euclidean operator}
Put $s(\mathbf z)=F_1(\mathbf z)/2$, $b(\mathbf z)=F_2(\mathbf z)$,
$c(\mathbf z)=F_3(\mathbf z)$. Since $\det DF=-2$, the vector fields
$W_j=(DF)^{-1}e_j$ are defined everywhere and satisfy
$W_jF_k=\delta_{jk}$. Set $U=2W_1+6F_3W_2$.
For a $C^2$ function $\phi$ on an open real target set, or a holomorphic
function on an open complex target set,
the ordinary chain rule gives the complete intertwining identity
\begin{equation}\label{eq:gct-full-material-drift}
 U\phi(s,b,c)=(\partial_s\phi+6c\partial_b\phi)(s,b,c),
 \qquad Us=1,\quad Ub=6c,\quad Uc=0.
\end{equation}
Every coefficient is retained. In particular $U(b-6cs)=6c-6c=0$.
Applying the same product rule twice gives
\[
 U^2\phi(s,b,c)=
 (\partial_s^2\phi+12c\partial_s\partial_b\phi
                +36c^2\partial_b^2\phi)(s,b,c).
\]
There are no missing derivatives of the coefficients: the derivative
of $c$ along $U$ is zero by \eqref{eq:gct-full-material-drift}.

The algebra map $\mathcal I f=f\circ s$ on entire functions is injective:
the polynomial section $\sigma(s)=(0,0,2s)$ satisfies
$s\circ\sigma=s$, so $\sigma^*\mathcal I=I$. Both maps are continuous
for the topology of uniform convergence on compact subsets, because
the continuous image of a compact set is compact. The preceding
chain rule proves inductively
\[
 U^k\mathcal I f=\mathcal I f^{(k)},\qquad
 -\tfrac14U^2\mathcal I f=-\tfrac14\mathcal I f''.
\]
Thus the original arithmetic differential operator has this exact
material pullback on its specified image. This proves an operator
identity, without asserting that differentiation descends to a
particular quotient by a closed subspace. The material time translates
$s$ at unit speed; the arithmetic heat operator is its negative
second derivative divided by four.

The full inverse-root polynomial in these same coordinates is
$P=cr^3-2r^2+br-4s$, with $P_r=3cr^2-4r+b$.
On the localization at $P_r$, the derivation
\[
 \mathcal D=\partial_s+6c\partial_b+
                 \frac{4-6cr}{P_r}\partial_r,\qquad \mathcal D c=0
\]
annihilates $P$: its three contributions are $-4$, $6cr$ and
$4-6cr$. It therefore defines a derivation of the localized quotient
ring. It is the unique lift with the stated derivatives of $s,b,c$,
because differentiating $P=0$ determines $P_r\mathcal D r=4-6cr$.
Through the previously proved inverse-chart isomorphism it equals
$U$, since both have the same derivatives of all target coordinates
and $DF$ is invertible. The $b=c=0$ derivation used earlier is
exactly its restriction, for which $P_r=-4r$ and
$\mathcal D r=-1/r$.

For the original real Euclidean metric put $v=1+xy$. The original
polynomial is $s=(v^3w+y^2(v+3v^2))/2$. Differentiation gives
\begin{align*}
 s_x&=\tfrac12\bigl(3yv^2w+y^3(1+6v)\bigr),\\
 s_y&=\tfrac12\bigl(3xv^2w+2y(v+3v^2)+xy^2(1+6v)\bigr),\\
 s_w&=v^3/2,\\
 s_{xx}&=3y^2vw+3y^4,\\
 s_{yy}&=3x^2vw+v+3v^2+2xy(1+6v)+3x^2y^2,
 \qquad s_{ww}=0.
\end{align*}
These identities follow by the product rule with $v_x=y$, $v_y=x$;
in particular the two displayed second derivatives sum to
\[
 B(\mathbf z):=\Delta_Es
 =3wv(x^2+y^2)+18x^2y^2+21xy+3y^4+4.
\]
Define $A(\mathbf z)=s_x^2+s_y^2+s_w^2$ using the three full
expressions above. Applying the scalar chain rule in each of the
three source directions and summing proves
\begin{equation}\label{eq:gct-Euclidean-coefficients}
 \Delta_E\mathcal I f=A\,\mathcal I f''+B\,\mathcal I f'.
\end{equation}
For real coordinates $A$ is the squared Euclidean norm of $ds$.
Its holomorphic extension is the indicated sum of squares, with no
complex conjugation inserted.

The original common-fibre points
$(1,-3/2,13/2)$ and $(0,0,-1/4)$ both map under $F$ to
$(-1/4,0,0)$, as already proved. Substitution in the last formulas
gives, respectively,
\[
 (s_x,s_y,s_w)=(-9/32,-3/16,-1/16),\quad
 (A,B)=(121/1024,-7/2),
\]
and
\[
 (s_x,s_y,s_w)=(0,0,1/2),\quad (A,B)=(1/4,4).
\]
Consequently no operator on functions of $s$ can satisfy
$\Delta_E\mathcal I=\mathcal I L$ on all twice differentiable
functions: applying such an identity to $f(s)=s$ would require
$-7/2=4$. The full relation that continues the transport is
\eqref{eq:gct-Euclidean-coefficients}, retaining its two
source-dependent coefficients. It is compatible with the separate
exact identity for $-U^2/4$ above and does not change that identity.

\subsection{The conic extension with the original parameter \texorpdfstring{$b$}{b} retained}
\label{sec:retained-b-conic}
All the constructions in this section are over $\mathbb C$. The
polynomial variables $b,s$ remain independent until a scalar value of
$b$ is explicitly fixed for the topological module construction.
The original coefficient is $a=2s$, and the original cubic on $c=0$ is
$-2r^2+br-4s$. Set
\[
 A=\mathbb C[b,s],\qquad
 B=A[r]/(r^2-br/2+2s),\qquad
 \beta=b/4,\quad \eta=r-\beta,
\]
and retain the inverse substitution $r=\eta+\beta$. Thus
\begin{equation}\label{eq:rb-original-base}
 s=br/4-r^2/2=b^2/32-\eta^2/2,\qquad
 \delta=b^2-32s=16\eta^2,\qquad
 L=4+b^2/16-2s=4+\eta^2.
\end{equation}
Here $\eta,\beta,\delta,L$ record explicit expressions in the original
objects; none changes the meanings of $b,r,s$ or $a$.

\subsubsection{The finite cover and its inverse maps}
Form the rank-two extension of $B$
\[
 E=B[p]/(p^2-L).
\]
Successive division by the two monic quadratic polynomials gives
the free $A$-basis $(1,r,p,rp)$: division first gives coefficients
of $1,p$ in $B$, and division in $B$ gives coefficients of $1,r$;
uniqueness at each step proves linear independence. There are inverse
$\mathbb C[b]$-algebra maps $E\longleftrightarrow\mathbb C[b,q,q^{-1}]$,
given on generators by
\begin{equation}\label{eq:rb-conic-maps}
 \begin{aligned}
 r&\longmapsto\beta+q-q^{-1},& p&\longmapsto q+q^{-1},\\
 s&\longmapsto b^2/32-(q-q^{-1})^2/2,& b&\longmapsto b,\\
 q&\longmapsto(p+r-\beta)/2,&q^{-1}&\longmapsto(p-r+\beta)/2.
 \end{aligned}
\end{equation}
To prove that these are well-defined, the two Laurent expressions
have $p^2-\eta^2=4$, and the proposed images of $q,q^{-1}$ have
product $(p^2-\eta^2)/4=1$. The image of $s$ then satisfies
\eqref{eq:rb-original-base}, hence both defining relations of $E$.
Conversely the sums and differences of the images of $q,q^{-1}$
recover $p,\eta$, and hence $r$ and $s$; the reverse compositions
recover $q,q^{-1}$ and $b$. This proves every map and its inverse.

For a temporary displayed coefficient put $h=2s-b^2/16$.
Then there are inverse $A$-algebra maps
\begin{equation}\label{eq:rb-quartic}
 E\simeq A[q]/\bigl(q^4+(2s-b^2/16-2)q^2+1\bigr),\qquad
 q^{-1}=-q^3+(2-2s+b^2/16)q.
\end{equation}
Indeed multiplying the proposed inverse of $q$ by $q$ gives one
by the quartic relation. In this quotient
$(q-q^{-1})^2=-h$, which follows by multiplying the quartic by
$q^{-2}$. Substituting in \eqref{eq:rb-conic-maps} therefore gives
the inverse map from $E$, and both compositions fix their generators.
In particular the quartic quotient is already a Laurent algebra.
Monic division gives its free basis $(1,q,q^2,q^3)$.

The complete conversion between these two ordered bases is
\begin{align}
 r&=\beta+q^3+(h-1)q,&
 p&=-q^3+(3-h)q,\notag\\
 rp&=h-2+\beta(3-h)q+2q^2-\beta q^3,&
 q&=(p+r-\beta)/2,\notag\\
 q^2&=(rp-\beta p+2-h)/2,&
 q^3&=r-\beta-(h-1)(p+r-\beta)/2.
 \label{eq:rb-basis-conversion}
\end{align}
The first two identities follow from the inverse of $q$ in
\eqref{eq:rb-quartic}. The identity
$\eta p=q^2-q^{-2}=2q^2+h-2$ proves the third and fifth.
The first identity proves the sixth after the fourth is inserted.
Thus no inverse basis map is omitted. The matrix whose columns are
$(1,r,p,rp)$ in the quartic basis is
\[
 T_b=\begin{pmatrix}
 1&\beta&0&h-2\\
 0&h-1&3-h&\beta(3-h)\\
 0&0&0&2\\
 0&1&-1&-\beta
 \end{pmatrix},\qquad \det T_b=4.
\]
Expanding along its first column and then the row containing $2$
gives $-2(-(h-1)-(3-h))=4$.

\subsubsection{Every original source coordinate on the cover}
Retain the original map, with its third source coordinate called $w$:
\begin{align*}
 F_1&=(1+xy)^3w+y^2(1+xy)(4+3xy),\\
 F_2&=y+3x(1+xy)^2w+3xy^2(4+3xy),\\
 F_3&=2x-3x^2y-x^3w.
\end{align*}
The simple-root inverse on $c=0$ is
\begin{equation}\label{eq:rb-source-H}
 \alpha=(b-4r)/2=-2\eta,\qquad
 H_b(\eta)=
 \left(-\frac1{2\eta},\frac b4+3\eta,
                 26\eta^2+\frac{3b\eta}{2}\right),\qquad\eta\ne0.
\end{equation}
These are the original formulas $x=1/\alpha$, $y=r-\alpha$,
$w=5\alpha^2-3r\alpha$ after the displayed invertible substitution.
Here $p$ is an additional coordinate, and never replaces $w$.
For a direct check of the full map set
$v=1+xy=-(\eta+\beta)/(2\eta)$ and note
$4+3xy=-(\eta+3\beta)/(2\eta)$. The polynomial identity
\[
 (\beta+3\eta)^2(\eta+3\beta)
 -(\eta+\beta)^2(13\eta+3\beta)=4\eta^2(\beta-\eta)
\]
follows by expanding both degree-three polynomials. Since
$w=2\eta(13\eta+3\beta)$, substitution in $F_1$ gives
$(\eta+\beta)(\beta-\eta)=\beta^2-\eta^2=2s$.
Substitution in $F_2$ gives
$\beta+3\eta+3(\beta-\eta)=4\beta=b$. Finally
\[
 F_3=-\frac1\eta-\frac{3(\beta+3\eta)}{4\eta^2}
                       +\frac{26\eta^2+6\beta\eta}{8\eta^3}=0.
\]
This proves $F(H_b(\eta))=(2s,b,0)$ with all coefficients retained.
Conversely the original inverse-chart coordinate is
$r=y+1/x$, so $\eta=y+1/x-b/4$; adjoining $p$ then recovers
$q=(p+y+1/x-b/4)/2$. The finite point
\[
 (x,y,w)=(0,b,2s-4b^2)
\]
lies on the same original fibre, because direct substitution gives
$F=(w+4y^2,y,0)=(2s,b,0)$. It remains separate from the $x\ne0$
chart, including at $s=b^2/32$.

The original inverse chart requires $\eta\ne0$, equivalently
$\delta\ne0$, or $q\ne\pm1$. The Laurent coordinate already
requires $q\ne0$. The extra quadratic extension has its additional
branch at $p=0$, equivalently $L=0$, or $q=\pm i$.
For completeness the relative differential module is exactly
\[
 \Omega_{E/A}=
 (E/(\eta))\,d\eta\ \oplus\ (E/(p))\,dp.
\]
Indeed $E=A[\eta,p]/(\eta^2-b^2/16+2s,p^2-L)$; the presentation
of relative differentials gives precisely the relations
$2\eta\,d\eta=0$ and $2p\,dp=0$, and $2$ is invertible.
Thus these are all the ramification loci, not just exhibited ones.
Their base values are
\[
 s=b^2/32\quad\hbox{and}\quad s=b^2/32+2,
\]
respectively. They are disjoint since their difference is $2$.
At the latter points $\eta=\pm2i$, so the original inverse chart
is still defined. At the former points $p=\pm2$; it is exactly the
original chart which excludes the point. These identities also
specify the domains of the derivative calculation below.

\subsubsection{Multiplication on the same arithmetic module}
Fix any scalar $b\in\mathbb C$, retaining its chosen value.
Let $Q$ be the original topological $\mathbb C[s]$-module with
continuous multiplication $S$ by its original coordinate $s$.
Give $Q$ the $A$-module structure $b\mapsto bI$, $s\mapsto S$ and
form $Q_b^{(4)}=E\otimes_AQ$. Since $E$ is free in the quartic basis,
the coefficient map gives an algebraic isomorphism $Q_b^{(4)}\simeq Q^4$;
we equip it with the fourfold direct-sum topology through this map.
With $\widehat S=\operatorname{diag}(S,S,S,S)$, multiplication by $q$ is
\begin{equation}\label{eq:rb-K}
 K_b=\begin{pmatrix}
 0&0&0&-I\\
 I&0&0&0\\
 0&I&0&(2+b^2/16)I-2S\\
 0&0&I&0
 \end{pmatrix}.
\end{equation}
The last column is the original quartic relation; the other three
columns shift the monic basis. Therefore
\[
 K_b^4+(2\widehat S-(b^2/16+2)I)K_b^2+I=0,\qquad
 K_b^{-1}=-K_b^3+((2+b^2/16)I-2\widehat S)K_b.
\]
Multiplication by this proposed inverse on either side gives $I$
by the quartic, because all its coefficients commute with $K_b$.
Both operators are continuous, being finite matrices of polynomials
in $S$. Define the original-root and added-coordinate operators by
\[
 \mathcal R_b=\frac b4I+K_b-K_b^{-1},\qquad
 \mathcal P_b=K_b+K_b^{-1}.
\]
The proved ring maps give
\begin{align}
 \mathcal R_b^2-\frac b2\mathcal R_b+2\widehat S&=0,&
 \mathcal P_b^2&=(4+b^2/16)I-2\widehat S,\notag\\
 \widehat S&=\frac b4\mathcal R_b-\frac12\mathcal R_b^2
 =\frac{b^2}{32}I-\frac12(K_b-K_b^{-1})^2.
 \label{eq:rb-original-S}
\end{align}
In the basis $(1,r,p,rp)$ multiplication by $r$ is
\[
 R_b\oplus R_b,\qquad
 R_b=\begin{pmatrix}0&-2S\\I&bI/2\end{pmatrix}.
\]
Indeed $r\cdot1=r$ and $r\cdot r=-2s+(b/2)r$; multiplying both
equations by $p$ gives the second block. The operator obtained by
substituting $S$ in $T_b$ is a continuous isomorphism: the inverse
formulas \eqref{eq:rb-basis-conversion} are finite polynomial
expressions in $S$ with scalar denominators $2$. Thus this is an
exact extension of the original two-sheet arithmetic module.

The continuous-invertibility spectrum of $\widehat S$ is precisely
that of $S$. If $S-\lambda I$ has a continuous two-sided inverse,
four copies give an inverse for $\widehat S-\lambda I$.
Conversely let $V$ be a continuous two-sided inverse of the latter.
For the first-coordinate inclusion $\iota:Q\to Q^4$ and projection
$\pi:Q^4\to Q$, the identities
$\widehat S\iota=\iota S$ and $\pi\widehat S=S\pi$ show that
$\pi V\iota$ is a continuous left and right inverse of
$S-\lambda I$. Also \eqref{eq:rb-original-S} gives the typed identity
\[
 S=\pi\left[\frac{b^2}{32}I
                 -\frac12(K_b-K_b^{-1})^2\right]\iota.
\]
Thus the recovered spectral coordinate is the original $s$.
An eigenvalue or spectral value of $K_b$ alone does not assert that
its recovered $s$ is off the original real axis.

Retain the GCT source-displayed Laurent coefficient polynomial
\[
 C(q)=q^{10}-q^4-q^{-4}+q^{-10}
      =(q-q^{-1})^2[7]_q[3]_q.
\]
The recurrence $U_0=1,U_1=p,U_j=pU_{j-1}-U_{j-2}$, after
$p=q+q^{-1}$, gives $U_j=q^j+q^{j-2}+\cdots+q^{-j}$ by cancellation
of the interior terms. Hence $[7]_q=p^6-5p^4+6p^2-1$ and
$[3]_q=p^2-1$. The factorization follows directly by multiplying
$(q-q^{-1})[7]_q=q^7-q^{-7}$ and
$(q-q^{-1})[3]_q=q^3-q^{-3}$; their product is exactly the
four-term Laurent expression above. Every factor of the coefficient
therefore descends by the explicit formula
\begin{align}
 C(q)&=(b^2/16-2s)(L^3-5L^2+6L-1)(L-1)\notag\\
 &=-(2s-b^2/16)\bigl(3-2s+b^2/16\bigr)\notag\\
 &\quad\cdot
 \bigl(7-14(2s-b^2/16)+7(2s-b^2/16)^2-(2s-b^2/16)^3\bigr).
 \label{eq:rb-C}
\end{align}
The second equality substitutes $L=4-(2s-b^2/16)$ and expands
only that polynomial; the first formula preserves its direct
conic substitution. In particular the coefficient has not become
the previous $b=0$ polynomial at the same $s$ unless $b=0$ is imposed.
For $\widehat L=(4+b^2/16)I-2\widehat S$, the actual transported
operator is exactly
\[
 C(K_b)=(b^2I/16-2\widehat S)
 (\widehat L^3-5\widehat L^2+6\widehat L-I)(\widehat L-I).
\]
Evaluation is legitimate by the explicit continuous inverse of $K_b$;
the identity follows from the proved Laurent algebra map. No factor
vanishing at a branch point has been divided out.

\subsubsection{The complete localized coefficient connection and its square}
Localize the base at both branch factors: $A^\circ=A[1/(\delta L)]$
and $E^\circ=E\otimes_AA^\circ$. The original material derivation on
the requested $c=0$ slice has $Db=0$, $Ds=1$, and consequently
$Da=2$. There is exactly one derivation of $E^\circ$ extending
these values, and it satisfies
\begin{equation}\label{eq:rb-derivation}
 D\eta=-1/\eta,\quad
 Dr=\frac4{b-4r}=\frac{4b-16r}{\delta},\quad
 Dp=-1/p,\quad
 Dq=-q/(\eta p).
\end{equation}
In fact differentiating the two monic quadratic relations forces
the first and third values because $\eta$ and $p$ are invertible
on this localization. Defining them this way annihilates both
relations, proving existence as well as uniqueness. Since
$Dr=D\eta$ and $b-4r=-4\eta$, the second formula follows, including
its expression in the free basis. In Laurent coordinates
$d\eta/dq=p/q$ and $dp/dq=\eta/q$; thus
\[
 D=-\frac{q}{\eta p}\frac{\partial}{\partial q}\quad(b\text{ fixed})
\]
has precisely these values and also gives $Ds=1$ in
\eqref{eq:rb-original-base}. Its Laurent domain is
$q\notin\{0,1,-1,i,-i\}$. This records separately the absence of
$q=0$, the original inverse-chart branch, and the added $p$ branch.

For a column $f=(f_0,f_1,f_2,f_3)^{\mathsf T}$ of coefficients in
$A^\circ$, interpret it as $f_0+rf_1+pf_2+rpf_3$.
The product rule and \eqref{eq:rb-derivation} give exactly
\begin{equation}\label{eq:rb-full-connection}
 Df=(\partial_sI_4+\Gamma_b)f,\qquad
 \Gamma_b=\begin{pmatrix}A_b&0\\0&A_b-I_2/L\end{pmatrix},\qquad
 A_b=\begin{pmatrix}0&4b/\delta\\0&-16/\delta\end{pmatrix}.
\end{equation}
Here the left side denotes the coefficient column of the derived
element, rather than componentwise differentiation alone.
The first two columns follow from $D1=0$ and
$Dr=(4b-16r)/\delta$. For the lower columns, $Dp=-p/L$ because
$p^2=L$, and
$D(rp)=p(4b-16r)/\delta-rp/L$. These verify all four columns.
Thus the $p$ contribution $-I_2/L$ is present in the full connection.

All zeroth-order terms of its square are
\begin{equation}\label{eq:rb-full-square}
 \begin{aligned}
 D^2&=\partial_s^2I_4+2\Gamma_b\partial_s+Z_b,\\
 Z_b&=\begin{pmatrix}
 (16/\delta)A_b&0\\
 0&(16/\delta-2/L)A_b-I_2/L^2
 \end{pmatrix}\\
 &=\begin{pmatrix}
 0&64b/\delta^2&0&0\\
 0&-256/\delta^2&0&0\\
 0&0&-1/L^2&64b/\delta^2-8b/(\delta L)\\
 0&0&0&-256/\delta^2+32/(\delta L)-1/L^2
 \end{pmatrix}.
 \end{aligned}
\end{equation}
For a full proof, $\partial_s\delta=-32$ and $\partial_sL=-2$ give
$A_b'=(32/\delta)A_b$ and $(-I_2/L)'=-2I_2/L^2$.
Direct matrix multiplication gives $A_b^2=-(16/\delta)A_b$.
It follows that the upper block of $\Gamma_b'+\Gamma_b^2$ is
$(16/\delta)A_b$. The lower block is
\[
 (32/\delta)A_b-2I_2/L^2
 -(16/\delta)A_b-2A_b/L+I_2/L^2
 =(16/\delta-2/L)A_b-I_2/L^2.
\]
Applying the product rule twice gives the displayed differential
operator and multiplying the scalars into $A_b$ gives every entry
of $Z_b$. In particular neither $-2A_b/L$ nor $-I_2/L^2$ disappears.

As a further compatibility check, the original root multiplication
obeys $R_b^2-(b/2)R_b+2sI_2=0$ and
\[
 [\partial_sI_2+A_b,R_b]
 =R_b'+A_bR_b-R_bA_b=4(bI_2-4R_b)/\delta.
\]
The last equality follows by substituting its four matrix entries:
the result is
$\delta^{-1}\left(\begin{smallmatrix}4b&32s\\-16&-4b\end{smallmatrix}\right)$.
It is exactly multiplication by $Dr$, establishing compatibility
with the original quadratic relation rather than only with its shift.

Restriction to $b=0$ is a quotient of this entire coefficient
calculation. Then $\eta=r$, $\delta=-32s$, $L=4-2s$, and the
connection becomes
\[
 \Gamma_0=\operatorname{diag}
 \left(0,\frac1{2s},-\frac1{4-2s},
                  \frac1{2s}-\frac1{4-2s}\right).
\]
Its full zeroth-order square term becomes
\[
 Z_0=\operatorname{diag}
 \left(0,-\frac1{4s^2},-\frac1{(4-2s)^2},
 -\frac1{4s^2}-\frac1{s(4-2s)}-\frac1{(4-2s)^2}\right).
\]
These follow by direct substitution in \eqref{eq:rb-full-square}.
The restricted domain is $s\ne0,2$, and the quartic, inverse source
map, original-root operator and coefficient in the preceding formulas
all specialize by setting $b=0$. Thus the earlier two-sheet connection
is exactly the first block of this four-sheet connection.

Finally the original coordinates in \eqref{eq:rb-source-H} have
the retained derivatives
\[
 Dx=-\frac1{2\eta^3},\qquad
 Dy=-\frac3\eta,\qquad
 Dw=-52-\frac{3b}{2\eta}.
\]
For example $Dw=(52\eta+3b/2)(-1/\eta)$; differentiation of the
other two displayed source functions proves their values.
Differentiating the already proved complete identity
$F(H_b(\eta))=(2s,b,0)$ now gives
\[
 DF(H_b(\eta))\,(Dx,Dy,Dw)^{\mathsf T}=(2,0,0)^{\mathsf T}.
\]
The separate finite point has derivative $(0,0,2)$.
For the original GCT coefficient, \eqref{eq:rb-C} gives the exact
identity $DC(q)=\partial_s C_b(s)$ with $b$ fixed, including all
its factors and the $b$-dependence. This derivation is on
$E^\circ$. The topological construction above proves multiplication
on $Q_b^{(4)}$ and its spectral-coordinate recovery; it does not
by itself define $\partial_s$ on $Q$, prove the invertibility there
of $\delta(S)$ or $L(S)$, or descend this localized connection to
an arithmetic quotient. No such derivative action is used in any
operator identity proved here.


\section{An explicit Boolean-to-algebraic edge with its encoding}
\label{sec:boolean}
The inclusion $\mathrm P\subseteq\mathrm{NP}$ follows directly:
a deterministic machine is a nondeterministic machine with one
available move at each step and the same running time.
The reverse inclusion is the equality question.
We retain a concrete edge by arithmetizing a three-literal conjunctive
formula's verifier.

For $n,m\ge1$, use $6mn$ commuting variables $E_{j,k,i,\sigma}$,
where $1\le j\le m$, $1\le k\le3$, $1\le i\le n$,
$\sigma\in\{+,-\}$, and auxiliary variables $b_1,\ldots,b_n$.
Over $\mathbb Z$ define
\begin{align*}
 L_{j,k}(E,b)&=\sum_{i=1}^n
   (E_{j,k,i,+}b_i+E_{j,k,i,-}(1-b_i)),\\
 V_{n,m}(E,b)&=\prod_{j=1}^m
   \left(1-\prod_{k=1}^3(1-L_{j,k}(E,b))\right),\\
 Q_{n,m}(E)&=\sum_{b\in\{0,1\}^n}V_{n,m}(E,b).
\end{align*}
A formula $\varphi$ with $n$ explicitly named variables and $m$
ordered three-literal clauses gives the zero--one array $E(\varphi)$
with one $1$ in each slot $(j,k)$, at its literal's index and sign.
Repeated literals are allowed. Input consists of explicitly named
variables and ordered signed binary indices with delimiters.
Writing $E(\varphi)$ uses $O(mn)$ bits and operations.
This convention does not allow exponentially many unused variables
to be specified solely by a binary upper bound: alternatively relabel
the $r$ occurring variables and record the unused count separately.
The projection from all $n$ Boolean variables to those $r$ coordinates
has exactly $2^{n-r}$ elements in each fibre and preserves the verifier.
Thus the original count is $2^{n-r}$ times the relabelled count.
A binary specification of an enormous unused count would require this
factored output format; writing all bits of the expanded count need
not be polynomial in that shorter input encoding.

\begin{proposition}[Exact counting reduction]\label{prop:counting}
The integer $Q_{n,m}(E(\varphi))$ is exactly the number of satisfying
assignments of $\varphi$. Deciding satisfiability of these formulas
therefore reduces by this polynomial array map and a nonzero test
to evaluation of this family on the specified arrays.
\end{proposition}
\begin{proof}
For Boolean $b$, each selected $L_{j,k}$ is exactly the literal's
truth value, $b_i$ or $1-b_i$.
The product of their three complements is $1$ exactly when the
clause is false. Its complement is its truth value.
The outer product is $1$ exactly when all clauses hold, and $0$
otherwise. Summation proves the integer equality and the bound
$0\le Q_{n,m}(E(\varphi))\le2^n$.
The stated array construction and zero test give the reduction.
\end{proof}

The displayed construction gives $V_{n,m}$ a uniform division-free
arithmetic circuit of $O(mn)$ gates, constants $0,1$, total degree
at most $6m$ in $(E,b)$ and degree at most $3m$ in $E$.
Each literal has degree at most two, each clause at most six, and
the products give these bounds. The loops over the displayed indices
construct the circuit in polynomial time. To obtain the usual one-index
family, assign
$t=(n+m-2)(n+m-1)/2+n$. The consecutive diagonals of pairs $(n,m)$
show this is a bijection from positive pairs to positive integers.
Trying consecutive diagonals through $n+m\le t+1$ decodes it in
polynomial time, and $n,m\le t$, so all variable, degree and gate bounds
remain polynomial in $t$. If a definition requires exactly $t$ witness
bits, multiply the verifier by $\prod_{i=n+1}^{t}(1-b_i)$ before
summing; only the all-zero added bits contribute, leaving $Q_{n,m}$
unchanged. By the definition of $\mathrm{VNP}$
as Boolean exponential summation of a polynomial-size,
polynomially bounded-degree arithmetic circuit family, this proves
$Q_{n,m}\in\mathrm{VNP}_{\mathbb Q}$ by construction.

There is also an explicit bit interface for a supplied, topologically
ordered, fan-in-two division-free circuit $C$ over $+,-,\times$, with
signed binary integer constants. On input $E(\varphi)$, evaluate its gates
modulo $M=2^{n+1}$. For $L$ gates and total constant-bit length $B_C$,
this takes $O(B_C+L(n+1)^2)$ arithmetic bit operations, plus graph and
input reading and the chosen machine's memory-access overhead,
by binary addition and school multiplication. Reduction to low
$n+1$ bits handles products, long constants and signs.
Induction over gates proves exact agreement with the integer value
modulo $M$. If the supplied circuit computes $Q_{n,m}$, its residue
is its exact value on this input since $0\le Q\le2^n<M$.
This gives the precise evaluation reduction without a bound on
intermediate integer magnitudes.

No polynomial-size circuits for $Q_{n,m}$ have been constructed here.
Polynomial-time generation of their full binary descriptions would
make the displayed evaluation uniform and bound both gate count and
constant-bit length. Descriptions supplied as size-dependent advice
give a polynomial-advice interface only when their total binary
description length is polynomial.
Arbitrary complex constants do not admit this binary modular
operation without additional input data.
The retained inverse-fibre selector polynomial already has the
explicit small determinant proved earlier; the present counting
polynomial has the different, fully specified joint verifier product
and Boolean summation. The next reduction must transport those
joint constraints with their parameter growth. None is erased by
calling both objects polynomial families.

\section{Shared Boolean constraints on the full original retained fibre}
\label{sec:sv-shared-verifier}

This section constructs the joint constraints and their counting map for
arbitrarily many witnesses. The base field $k$ has characteristic zero.
The integral and bit statements use $\mathbb Z$ and $\mathbb Q$ with the
encodings specified below. A witness variable is denoted $B_i$; it is
different from the original target coordinate $b_i$.

\subsection{The full three-point source fibre and its Boolean corner}

Retain the original source map, including all coefficients,
\begin{align*}
 F_1&=(1+xy)^3w+y^2(1+xy)(4+3xy),\\
 F_2&=y+3x(1+xy)^2w+3xy^2(4+3xy),\\
 F_3&=2x-3x^2y-x^3w,
\end{align*}
with source orientation $(x,y,w)$ and target $(a,b,c)$.
To retain the Jacobian claim with a direct proof, on $x\ne0$ put
$v=1+xy$, $h=2-3xy-x^2w$,
$A_0=v^2+v-v^3h$, $B_0=4v+2-3v^2h$.
The inverse coordinates are $y=(v-1)/x$, $w=(5-3v-h)/x^2$,
and substitution gives $F=(A_0/x^2,B_0/x,xh)$.
The coordinate-change determinant is $-x^3$, while the determinant
of the latter map is $x^{-3}$ times
\[
 (4-6vh)(-2v^2-2v+3v^3h)
 +(2v+1-3v^2h)(4v+2-6v^2h)=2.
\]
Their product is $-2$. The identity holds everywhere because the
polynomial coordinate ring embeds in its localization at $x$.
Thus $\det DF=-2$ in the original orientations. Its retained polynomial is
$P=cr^3-2r^2+br-2a$ and $P'=3cr^2-4r+b$.
We use the explicitly specified target $(-1/4,0,0)$ and set
\[
 A_1=k[x,y,w]/(F_1+1/4,F_2,F_3),\qquad
 T_1=k[X]/(X^3-X).
\]

\begin{proposition}\label{prop:sv-full-fibre}
There is an isomorphism $A_1\longrightarrow T_1$ given by
\[
 x\longmapsto X,\qquad y\longmapsto-3X/2,\qquad
 w\longmapsto(27X^2-1)/4.
\]
Evaluation at $X=-1,0,1$ identifies $T_1$ with $k^3$ in that order.
The three idempotents, in the same order, are
\[
 e_-=(X^2-X)/2,\qquad e_0=1-X^2,\qquad e_+=(X^2+X)/2.
\]
The factor $e_0$ is the original projective boundary point
$(x,y,w)=(0,0,-1/4)$.
\end{proposition}
\begin{proof}
Put $h=2-3xy-x^2w$ and
$R=k[x,y,w]/(F_1+1/4,F_2)$. The identity
$1=h/2+x(3y+xw)/2$ proves $(x)+(h)=R$.
Consequently $(x)\cap(h)=(xh)$: for $z=xu=hv$ multiply the
displayed expression for $1$ by $z$. As $F_3=xh$, the map
$R/(F_3)\to R/(x)\times R/(h)$ is an isomorphism. Its injectivity
is the intersection equality; multiplying lifts of the two components
by the two displayed summands of $1$ proves surjectivity.
The first factor is $k$, at $(0,0,-1/4)$.

In the second factor $x(3y+xw)=2$ makes $x$ invertible. Put
$r=y+1/x$ and $\alpha=1/x$. Solving $F_3=0$ gives
$x=\alpha^{-1}$, $y=r-\alpha$,
$w=5\alpha^2-3r\alpha$. Direct insertion into the original map gives
\[
 (F_1,F_2,F_3)=(r^2+r\alpha,4r+2\alpha,0).
\]
Thus $F_2=0$ gives $\alpha=-2r$, and $F_1=-1/4$ gives
$r^2=1/4$. Conversely these equations, with
$\alpha=-2r$, give the original three target equations and recover
$r=y+1/x$, $\alpha=1/x$. The factor is therefore exactly
$k[r]/(r^2-1/4)$. In it $x=-2r$, $y=3r=-3x/2$,
$x^2=1$, and $w=13/2=(27x^2-1)/4$.
The two factors together prove the stated map, since
$k[X]/(X^3-X)\simeq k[X]/(X)\times k[X]/(X^2-1)$ by the
same coprime-ideal argument. The displayed idempotents have values
$(1,0,0),(0,1,0),(0,0,1)$; their values show that they sum to one,
are orthogonal, and recover every triple. This also proves evaluation
is an isomorphism, without discarding the boundary factor.
\end{proof}

For $n\ge1$ let $A_n=A_1^{\otimes n}$ and write $X_i$ for its
corresponding source coordinate. This is the full original fibre of
$F^{[n]}$ and has the basis
$\{X_1^{d_1}\cdots X_n^{d_n}:0\le d_i\le2\}$.
Repeated division by $X_i^3-X_i$ proves spanning; evaluations at
$\{-1,0,1\}^n$ and products of the three idempotents prove independence.
In particular its dimension is exactly $3^n$. Put
\[
 e_i=X_i^2,\quad \beta_i=(X_i^2+X_i)/2,
 \quad e=\prod_{i=1}^n e_i.
\]
The elements $e_i,\beta_i,e$ are idempotents. The element $\beta_i$
has values $0,0,1$ at $X_i=-1,0,1$, respectively, while $e_i$
has values $1,0,1$.
The quotient and corner maps are
\[
 A_n/(1-e)\xrightarrow{\ \sim\ }eA_n,\quad [f]\longmapsto ef,
 \qquad
 eA_n\xrightarrow{\ \sim\ }k[B_1,\ldots,B_n]/(B_i^2-B_i),
 \quad eX_i\longmapsto2B_i-1.
\]
The corner has identity $e$. For the first map $ef=0$ implies
$f=(1-e)f$, proving exactly its kernel and injectivity; its image
is its definition. For the second map the inverse sends
$B_i$ to $e\beta_i=e(1+X_i)/2$. The relations $eX_i^2=e$
prove both compositions on the generators, including the identities.
Also $(1-e)=(1-e_1,\ldots,1-e_n)$ as ideals:
$1-e=\sum_i(1-e_i)\prod_{h<i}e_h$, while
$1-e_i=(1-e_i)(1-e)$. Thus this corner keeps exactly the
$2^n$ choices with $X_i\ne0$; the complementary boundary factors
remain recorded by $1-e$ in the full algebra.

There is a useful exact integral distinction here. The evaluation map
\[
 \mathbb Z[X]/(X^3-X)\longrightarrow\mathbb Z^3
\]
has image $\{(u,v,w):u\equiv w\pmod2\}$. Indeed
$A+BX+CX^2$ has values $(A-B+C,A,A+B+C)$; conversely use
$A=v$, $B=(w-u)/2$, $C=(u+w-2v)/2$.
The inclusion has index two, and its conductor in $\mathbb Z^3$ is
$2\mathbb Z\times\mathbb Z\times2\mathbb Z$: multiplication
by $(1,0,0)$ and $(0,0,1)$ proves necessity of the two evenness
conditions, and the image criterion proves sufficiency. The split
lattice $\mathbb Ze_-\oplus\mathbb Ze_0\oplus\mathbb Ze_+$
is consequently an explicit integral lattice inside $A_1$;
tensor products give a lattice $\mathbb Z^{3^n}$ in $A_n$.
This retains the factor two and the original rational source coordinates.
It does not reduce their denominators modulo an even modulus.

\subsection{A trace that carries actual shared verifier constraints}

Let $n,m\ge1$ and retain the $6mn$ independent variables
$E_{j,k,i,\sigma}$, where $1\le j\le m$, $1\le k\le3$,
$1\le i\le n$, and $\sigma\in\{+,-\}$. In
$\mathbb Z[E,B]$ put
\begin{align*}
 L_{j,k}(E,B)&=\sum_{i=1}^n
   \bigl(E_{j,k,i,+}B_i+E_{j,k,i,-}(1-B_i)\bigr),\\
 H_j(E,B)&=\prod_{k=1}^3(1-L_{j,k}(E,B)),\\
 V_{n,m}(E,B)&=\prod_{j=1}^m(1-H_j(E,B)),\\
 Q_{n,m}(E)&=\sum_{B\in\{0,1\}^n}V_{n,m}(E,B).
\end{align*}
Every occurrence of a witness index uses the same $B_i$, across all
clauses and positions. Over $k[E]$, define
\[
 u(E)=e\,V_{n,m}(E,\beta)\in A_n\otimes_k k[E],
 \qquad m_u:f\longmapsto uf.
\]

\begin{theorem}[Full-fibre counting trace]\label{thm:sv-trace}
The following identity is an identity of polynomials in all the
independent $E$ coordinates:
\[
 \operatorname{Tr}_{A_n\otimes k[E]/k[E]}(m_u)=Q_{n,m}(E).
\]
For a three-literal conjunctive formula $\varphi$, let $E(\varphi)$
have a single $1$ in each $(j,k)$ slot at its selected signed literal,
and zero in every other position. Then $u(E(\varphi))$ is an
idempotent, its multiplication map is a projection of rank exactly
$\#\mathrm{SAT}(\varphi)$, and its trace is that same integer.
\end{theorem}
\begin{proof}
The product idempotents from Proposition~\ref{prop:sv-full-fibre}
are a basis after extending scalars to $k[E]$. Multiplication by $u$
is diagonal in this basis, with entry
\[
 \left(\prod_i X_i^2\right)
 V_{n,m}\left(E,\left((X_i^2+X_i)/2\right)_i\right)
 \quad\text{at }(X_1,\ldots,X_n)\in\{-1,0,1\}^n.
\]
The entry is zero if any $X_i=0$. On the remaining points
$B_i=(X_i+1)/2$ runs bijectively over $\{0,1\}^n$.
Adding the diagonal entries proves the polynomial identity.
On the specified array each $L_{j,k}$ has its literal's truth value.
Thus $H_j=1$ exactly when all three literals are false, and
$\prod_j(1-H_j)$ has value one exactly on the satisfying assignments.
All eigenvalues of $m_u$ are therefore zero or one, and its
idempotence, rank and trace assertions follow in the displayed basis.
The equality with the integer count uses the canonical embedding
$\mathbb Z\to k$, which is injective in characteristic zero.
\end{proof}

This also provides the integral matrix needed for reduction modulo
an integer. In the split lattice's product-idempotent basis, $m_u$
is diagonal with entries in $\mathbb Z[E]$, by the preceding
display: each $\beta_i$ has the integer value zero or one.
Its trace can thus be reduced modulo $2^{n+1}$ entry by entry
without interpreting $1/2$ in that quotient ring. The lattice still
has rank $3^n$, and its Boolean corner has rank $2^n$.

The construction respects explicit witness symmetries. Permuting the
$X_i$ and simultaneously permuting the index $i$ of $E$ permutes
the product-idempotent basis and leaves the trace unchanged.
On the corner, $X_i\mapsto-X_i$ sends $B_i$ to $1-B_i$;
exchanging $E_{j,k,i,+}$ and $E_{j,k,i,-}$ then leaves every literal
sum invariant. These maps commute for distinct $i$, square to the
identity, and permutations conjugate the $i$th map to the permuted
one. They give the exact action of
$(\mathbb Z/2\mathbb Z)^n\rtimes\mathfrak S_n$ on the corner
and the parameter ring, under which the trace polynomial is invariant.
On the full algebra the bit flip sends $\beta_i$ to $e_i-\beta_i$.
The difference from $1-\beta_i$ is $e_i-1$, annihilated by $e$.
Thus the same trace invariance holds on the full algebra with its
boundary factors retained. This proves the action and the map it
preserves, rather than identifying two presentations by notation.

\subsection{Universal incidence rings retaining the original cubics}

Use the polynomial ring
\[
 R_{n,m}=k[E,a_i,b_i,c_i,r_i,B_i:1\le i\le n]
\]
and, before imposing any target equations, retain
\[
 P_i=c_ir_i^3-2r_i^2+b_ir_i-2a_i,\qquad
 P_i'=3c_ir_i^2-4r_i+b_i.
\]
Let $I_{n,m}$ be the ideal generated, for each $i$, by the six
specified expressions
\[
 P_i,\quad a_i+1/4,\quad b_i,\quad c_i,\quad
 B_i^2-B_i,\quad 2r_i+2B_i-1,
\]
and additionally by all $H_j(E,B)$, $1\le j\le m$.
All target specializations are therefore equations in the original
coordinate ring; no original coefficient is silently replaced.

\begin{proposition}\label{prop:sv-incidence}
There are explicit isomorphisms of $k[E]$-algebras
\begin{align*}
 R_{n,m}/I_{n,m}
 &\simeq k[E,B]/(B_i^2-B_i,H_j(E,B))\\
 &\simeq (A_n\otimes k[E])/(1-e,H_j(E,\beta)).
\end{align*}
The inverse-fibre localization by $\prod_i P_i'$ changes neither
ring. Specializing $E=E(\varphi)$ gives a reduced algebra with one
copy of $k$ for each satisfying assignment, and no other factors.
\end{proposition}
\begin{proof}
The first map sends
\[
 a_i\mapsto-1/4,\quad b_i,c_i\mapsto0,\quad
 r_i\mapsto(1-2B_i)/2,
\]
and fixes $E,B$. Under it the original polynomial becomes
$-2(B_i^2-B_i)$, so every defining relation maps to zero.
The inverse fixes $E,B$ and is inverse on $a,b,c,r$ by their
listed relations, proving the first isomorphism and its entire kernel.
The second is the previously proved Boolean-corner map, with
$B_i\leftrightarrow e\beta_i$; it sends the clause relations to
their displayed counterparts.
In the first quotient $P_i'=4B_i-2$ and
\[
 (4B_i-2)(B_i-1/2)=1+4(B_i^2-B_i)=1.
\]
Thus its inverse is exactly $B_i-1/2$, proving the localization
claim. The original source recovery formulas give
\begin{align*}
 x_i&=2/P_i'=2B_i-1,&
 y_i&=r_i-P_i'/2=3/2-3B_i,\\
 w_i&=5(P_i'/2)^2-3r_iP_i'/2-c_i(P_i'/2)^3=13/2.
\end{align*}
They are the selected source sheets of the original fibre. Their
inverse on this quotient is $B_i=(x_i+1)/2$, as required.

For completeness the Boolean algebra over $k[E]$ splits by the
orthogonal idempotents
$\prod_i B_i^{\epsilon_i}(1-B_i)^{1-\epsilon_i}$ into
$\prod_{\epsilon\in\{0,1\}^n}k[E]$. The clause ideal is sent
componentwise to $(H_j(E,\epsilon):1\le j\le m)$.
After specializing the one-hot array, each generator has value zero
or one. A component survives as $k$ exactly when all of these values
are zero; otherwise it becomes the zero ring. These are exactly
the satisfying assignments. This proves reducedness and the full
factor description, including inconsistent and repeated clauses.
\end{proof}

\subsection{A single determinant with reversible coefficient recovery}

Introduce independent selectors
$t_i,\lambda_i,\mu_i,\nu_i,\rho_i,\tau_i$ for the six
expressions in their displayed order, and $s_j$ for the clauses. Put
\begin{align*}
 S_{n,m}={}&\sum_{i=1}^n\bigl(t_iP_i+\lambda_i(a_i+1/4)
   +\mu_i b_i+\nu_i c_i+\rho_i(B_i^2-B_i)
   +\tau_i(2r_i+2B_i-1)\bigr)\\
 &+\sum_{j=1}^m s_jH_j(E,B).
\end{align*}
The coefficient map from the homogeneous selector-degree-one part
of $R_{n,m}[t,\lambda,\mu,\nu,\rho,\tau,s]$ to
$R_{n,m}^{6n+m}$ is an isomorphism of $R_{n,m}$-modules: its
inverse forms the sum of each tuple entry times its own selector.
The uniqueness of polynomial coefficients proves both compositions.
In particular the coefficient ideal of $S_{n,m}$ is exactly
$I_{n,m}$. At any point over an extension field, the incidence
relations hold if and only if $S_{n,m}$ is identically zero in all
selectors. This retains the shared variables, not merely a product
of independent hypersurface equations.

We give a complete determinant construction with a polynomial size
bound. Write the displayed expression as a binary arithmetic formula
with operations addition and multiplication and signed rational
constants; retain every occurrence. For the cubic use
$c_i r_i r_i r_i+(-2)r_i r_i+b_i r_i+(-2)a_i$.
Write $B_i^2-B_i=B_iB_i+(-1)B_i$, and in each literal write
$1-B_i=1+(-1)B_i$. Each clause complement is
$1+(-1)L_{j,k}$. Parenthesize sums and products in their displayed
index order. This formula has exactly
\[
 \ell=30n+(18n+7)m
\]
leaves. For each $i$ the six selected terms have respectively
$12,3,2,2,5,6$ leaves. Each $L_{j,k}$ has $6n$ leaves, its
complement has $6n+2$, and a selected three-factor product has
$18n+7$. This proves the count without hiding repeated witnesses.
Let $M$ be this formula's number of multiplication nodes. The six
selected terms for each $i$ contribute respectively $8,1,1,1,3,3$
multiplications, while each literal sum contributes $3n$, each
complement contributes one more, and combining the selector and its
three complements contributes three. Thus
\[
 M=17n+(9n+6)m\le\ell-1.
\]

Construct a directed graph with a source $v_s$ and sink $v_t$ as
follows. To implement a sum between given endpoints, implement both
summands with the same endpoints and fresh internal vertices. To
implement a product, create one fresh middle vertex and implement
its first factor from the source to that vertex and its second factor
from there to the sink. To implement a leaf with label $z_0$, create
a fresh private vertex, with edges labelled $z_0$ from source to
private vertex and $1$ from private vertex to sink. Here $z_0$ is
one original indeterminate or one signed rational constant, not the
homogenizing variable introduced below. Induction proves the graph
is acyclic and that the sum of products of labels over all source-to-sink
paths is exactly the formula: sums partition paths between the two
subgraphs, and products uniquely concatenate paths at their middle
vertex. Its number of vertices is
$v=2+\ell+M$ and its number of edges is $2\ell$.

Order vertices topologically, let $N$ be the strictly upper triangular
matrix of edge labels, and let $e_s,e_t$ be the corresponding coordinate
vectors. No two leaf-labelled edges are merged: each has its own
private target vertex. With
\[
 d=v+1=3+\ell+M=3+47n+(27n+13)m\le2\ell+2,\qquad
 \mathcal D_{n,m}=
 \begin{pmatrix}I_v-N&e_t\\-e_s^{\mathsf T}&0\end{pmatrix},
\]
every entry is affine-linear over $\mathbb Q$ and
\[
 \det\mathcal D_{n,m}=S_{n,m}.
\]
Indeed $N^v=0$, so $(I_v-N)^{-1}=\sum_{j=0}^{v-1}N^j$,
whose $(s,t)$ entry is the path sum. Multiplication on the left by
$\left(\begin{smallmatrix}I&0\\e_s^{\mathsf T}(I-N)^{-1}&1
\end{smallmatrix}\right)$ has determinant one and produces the
upper block triangular matrix with last diagonal entry
$e_s^{\mathsf T}(I-N)^{-1}e_t$. Its other diagonal block has
determinant one. This proves the exact determinant identity and sign.

Let $Z$ be a new homogenizing variable and replace each constant
matrix entry by that constant times $Z$, leaving all variable entries
linear. Call the resulting matrix $\mathcal D^h_{n,m}$.
There are $q=11n+6mn+m$ affine variables in $S_{n,m}$, including
all selectors. Let $w$ denote their ordered tuple, with order fixed
by the above index loops, and put
\[
 H_{n,m}(w,Z)=\det\mathcal D^h_{n,m}(w,Z)
             =Z^dS_{n,m}(w/Z).
\]
The last identity follows on $Z\ne0$ by factoring $Z$ from each
matrix row. Both sides are polynomials, since $\deg S_{n,m}\le7<d$,
so injectivity of localization proves equality at $Z=0$ also.
In fact $\deg S_{n,m}=7$: for $n,m\ge1$ the monomial using
$s_1E_{1,1,1,+}E_{1,2,1,+}E_{1,3,1,+}B_1^3$ has nonzero
coefficient. Consequently the explicit padding factor is
$Z^{d-7}$ times the degree-seven homogenization of $S_{n,m}$.
The hyperplane $Z=0$ is an added zero component; at $Z=1$ every
original incidence coefficient is recovered.

The coordinate-ring map is fully specified:
\[
 \Theta:\mathbb Q[Y_{ab}:1\le a,b\le d]\longrightarrow
 \mathbb Q[w,Z],\qquad Y_{ab}\longmapsto
    (\mathcal D^h_{n,m})_{ab}.
\]
Every indeterminate in $w$ occurs as a leaf, hence as an individual
matrix entry with coefficient $-1$; a diagonal entry is $Z$.
Thus $\Theta$ is surjective. Choose one such entry for each variable,
divide by its nonzero scalar, and subtract from every remaining
$Y_{ab}$ the exact linear combination of chosen entries given by its
image. The differences generate the kernel: modulo them the source
is the polynomial ring on the chosen entries, and the map has the
inverse sending each target variable to its chosen entry. This proves
the image, kernel, polynomial degree and all variable identifications.

\subsection{The orbit and representation consequence of these joint constraints}

Work now over $\mathbb C$. Set
$V_d=\operatorname{Mat}_{d\times d}(\mathbb C)$,
$W_d=\operatorname{Sym}^d(V_d^*)$, and
$G_d=\operatorname{GL}(V_d)$, with
$(g\cdot f)(Y)=f(g^{-1}Y)$.
Let $L:\mathbb C^{q+1}\to V_d$ be the linear map with value
$\mathcal D^h_{n,m}$. The coordinate calculation proves $L$ is
injective. Also $q+1\le\ell+1<d^2$, since each different variable
occurs among the $\ell$ leaves and $d\ge\ell+3$.
Let $p:V_d\to\mathbb C^{q+1}$ select the first $q+1$ entries
in row order, let $i$ fill those entries and set the rest to zero,
and put $A=Lp$. Then $pi=1$,
\[
 \widetilde H=H_{n,m}\circ p=\det_d\circ A,
 \qquad i^*\widetilde H=H_{n,m}.
\]
The map $p^*$ is injective and $i^*p^*=1$, while $L^*=\Theta$.

Write $\Omega_d=\overline{G_d\cdot\det_d}$. Then
$\widetilde H\in\Omega_d$. Indeed $A+\epsilon I$ is invertible
except at the finitely many zeros of its monic determinant polynomial,
and $\det_d\circ(A+\epsilon I)$ is an orbit point there.
Every polynomial equation vanishing on the orbit evaluates to a
polynomial in $\epsilon$ vanishing away from that finite set, hence
vanishes identically and at zero. This is exactly closure membership.
It is a strict boundary point. The universal determinant has zero
directional-annihilator subspace, because its cofactors are linearly
independent: each cofactor has a distinct pair of row and column
multidegrees. The chain rule preserves the dimension of this subspace
under $G_d$, whereas $\ker p$ is a nonzero subspace of the
annihilator of $\widetilde H$. The polynomial is nonzero because
the recovered selector coefficient of $\lambda_1$ is $a_1+1/4$.

Let $Z_{n,m}=\overline{G_d\cdot\widetilde H}$. Both this and
$\Omega_d$ are cones: scalar coordinate substitutions realize every
nonzero scalar multiple of a degree-$d$ form. The containment
$Z_{n,m}\subseteq\Omega_d$ gives, in every coordinate-ring degree
$h\ge0$, a surjective $G_d$-equivariant map
\[
 \mathbb C[\Omega_d]_h\longrightarrow\mathbb C[Z_{n,m}]_h.
\]
Here the map is restriction, and its kernel is exactly
$I(Z_{n,m})_h/I(\Omega_d)_h$. Homogeneity follows by evaluating
each equation on scalar multiples and comparing coefficients in the
scalar. This proves the exact representation quotient. In particular,
for every finite-dimensional irreducible representation $U$, precomposing
with this surjection is an injection
\[
 \operatorname{Hom}_{G_d}(\mathbb C[Z_{n,m}]_h,U)
 \hookrightarrow
 \operatorname{Hom}_{G_d}(\mathbb C[\Omega_d]_h,U),
\]
because a linear map vanishing on the surjection's image vanishes
everywhere. This inequality uses the exact modules without importing
an unproved multiplicity calculation or a semisimplicity assertion.

\subsection{Unbounded family, bit evaluation, and the size of the counting map}

Index all positive pairs by
$t=(n+m-2)(n+m-1)/2+n$. Consecutive diagonals $n+m=2,3,\ldots$
give consecutive disjoint intervals of integers of lengths
$1,2,\ldots$, and $n$ labels positions within each interval.
This proves a bijection with the positive integers. Sequentially
testing the diagonals up to $t+1$ decodes it in polynomial time,
and $n,m\le t$. The construction of the formula, graph, topological
order, matrix, variable list and selector coefficient map is uniform:
the displayed loops write $O(mn+n)$ formula symbols and vertices,
and their binary labels have $O(\log(mn+n+1))$ bits. Constants
belong to $\{0,1,-1,2,-2,1/4,-1/4\}$; the sign in $I_v-N$
accounts for the entry $-1/4$. The sparse matrix description
has $O((mn+n)\log(mn+n+1))$ bits; a dense description has
$O((mn+n)^2\log(mn+n+1))$ bits. Both are polynomial in $t$.

The verifier $V_{n,m}$ has a uniformly generated division-free
arithmetic circuit with $O(mn)$ gates, degree at most $6m$ in
$(E,B)$ and degree at most $3m$ in $E$. Each literal sum has total
degree at most two, a clause at most six, and the product has the
stated bounds. Therefore the one-index family $Q_t=Q_{n,m}$ is
in $\mathrm{VNP}_{\mathbb Q}$ by the definition as a Boolean sum
of such a polynomial-size, bounded-degree circuit family. If exactly
$t$ witness bits are required, multiply the verifier by
$\prod_{i=n+1}^t(1-B_i)$; only the all-zero additional bits survive,
so its sum stays exactly $Q_t$ and its degree and size stay polynomial.

Represent a three-literal formula by its explicitly listed variables
and ordered triples of signed binary variable indices, with delimiters.
Its one-hot array can be written in $O(mn)$ bit positions plus index
processing; this is polynomial in that representation. The equality
$Q_{n,m}(E(\varphi))=\#\mathrm{SAT}(\varphi)$ and the test for
a nonzero integer give the exact reduction to evaluating this
one-index family. This convention does not encode exponentially many
unused variables by a binary upper bound. If only $r$ of the explicitly
listed $n$ variables occur, restricting an assignment has exactly
$2^{n-r}$ lifts, so the original count is precisely $2^{n-r}$ times
the relabelled count, with the factor retained.

For a supplied topologically ordered division-free circuit with
fan-in at most two and integer constants that computes $Q_{n,m}$,
let its gate count be
$G$ and total constant-bit length be $C$. Evaluate its gates modulo
$M_0=2^{n+1}$. Induction over $+,-,\times$ proves the result is the
integer value modulo $M_0$. Since the value on a formula array lies
in $[0,2^n]$, its least nonnegative residue is its exact value.
Binary addition and school multiplication give
$O(C+G(n+1)^2)$ bit operations, plus reading the circuit graph and
the chosen machine's indexing costs. Integer constants can be reduced
by reading their low $n+1$ bits and their sign. This bound does not
assume bounds on unreduced intermediate integers. A rational or complex
circuit requires its actual constant representation; arbitrary complex
constants cannot be interpreted by this binary operation.

The determinant $S_{n,m}$ and its padding have the proved polynomial
size, and coefficient extraction from them recovers the actual coupled
incidence ideal. To obtain the count, the exact operation proved here
forms $\prod_j(1-H_j)$, maps it to the Boolean corner of the
original fibre, and takes the multiplication trace. That full algebra
has dimension $3^n$ and its corner has dimension $2^n$. A direct
diagonal computation therefore visits exponentially many basis entries;
the small determinant for the constraints does not supply a smaller
matrix for this trace. The construction proves an explicit bridge from
the original retained fibres to the unbounded shared-verifier VNP
family and its incidence orbit modules. It proves neither polynomial
circuits for $Q_t$ nor a Boolean or algebraic complexity separation.
The next size question is the exact complexity of this now specified
trace family, not the loss of the joint constraints.

\section{A direct permanent projection of the shared-verifier trace}
\label{sec:permanent-trace-projection}

The polynomial $Q_{N,M}(E)$ and its full-fibre trace now have a concrete
relation to the permanent, with all witness extensions and sizes specified.
For an integer $n\ge1$ introduce the $n^2$ independent variables
$Z_{ij}$, $1\le i,j\le n$, and define
\[
 \operatorname{per}_n(Z)=\sum_{\sigma\in\mathfrak S_n}
                  \prod_{i=1}^n Z_{i,\sigma(i)}
 \quad\text{in }\mathbb Z[Z_{ij}].
\]
We will use exactly
\begin{equation}\label{eq:per-trace-size}
 N=3n^2,\qquad M_0=n^3+5n^2,\qquad M=n^3+6n^2
\end{equation}
witness bits and clauses. The original source map in
Proposition~\ref{prop:sv-full-fibre} and its target $(-1/4,0,0)$
remain unchanged in every one of these $N$ factors.

\subsection{The uniquely extended Boolean witness}

Order the $N$ bits by the following three groups, each in lexicographic
index order:
\[
 B_{ij}\ (1\le i,j\le n),\qquad
 R_{ik}\ (1\le i,k\le n),\qquad
 C_{jk}\ (1\le j,k\le n).
\]
The bit $B_{ij}$ specifies whether the permutation uses entry $(i,j)$.
For each row $i$, impose the equivalences
\[
 R_{i1}\leftrightarrow B_{i1},\qquad
 R_{ik}\leftrightarrow(R_{i,k-1}\vee B_{ik})\ (2\le k\le n),
 \qquad R_{in}=1.
\]
Also impose $\neg B_{ij}\vee\neg B_{il}$ for every $j<l$.
For each column $j$, impose the equivalences
\[
 C_{j1}\leftrightarrow B_{1j},\qquad
 C_{jk}\leftrightarrow(C_{j,k-1}\vee B_{kj})\ (2\le k\le n),
 \qquad C_{jn}=1,
\]
and $\neg B_{ij}\vee\neg B_{lj}$ for every $i<l$.

We specify every clause in these equivalences. Encode $w\leftrightarrow v$
as $(\neg w\vee v)$ and $(w\vee\neg v)$. Encode
$w\leftrightarrow(u\vee v)$ as
\[
 (\neg u\vee w),\qquad(\neg v\vee w),\qquad
 (u\vee v\vee\neg w).
\]
The first two clauses say that either true input forces $w=1$;
the third says that $w=1$ forces at least one true input.
These two implications prove the equivalence for every Boolean triple.
The two clauses for equality prove both implications in the same way.
Encode the final truth requirement by the unary clause $w$.
Pad a binary clause $(a\vee b)$ to $(a\vee b\vee b)$ and a unary
clause $a$ to $(a\vee a\vee a)$. Since $0,1$ satisfy $v^2=v$, this
padding preserves each clause's truth value under its polynomial
arithmetization $1-\prod_{k=1}^3(1-L_k)$.

Order clauses first by rows, then by columns, and within each row or
column by the equivalences in increasing prefix index, the final truth
requirement, and the exclusion pairs in lexicographic order. This fixes
an ordered list of three-literal clauses without an unspecified choice.
Each row or column contributes
\[
 2+3(n-1)+1+\binom n2=3n+\binom n2
\]
clauses, including when $n=1$. There are $2n$ such lists, yielding
$M_0=2n(3n+\binom n2)=n^3+5n^2$.

For each assignment to the $B_{ij}$, the equivalences force the unique
extension
\[
 R_{ik}=\bigvee_{j=1}^k B_{ij},\qquad
 C_{jk}=\bigvee_{i=1}^k B_{ij}.
\]
The formulas follow by induction on $k$ starting with the equality clause;
they also satisfy every equivalence, proving existence and uniqueness.
The final truth requirements and pair exclusions then say that each
row and each column has exactly one true entry. Reading the true entry
of row $i$ as $\sigma(i)$ gives a map to $\mathfrak S_n$: no two rows
choose the same column, and every column is chosen. Conversely every
permutation yields exactly those entries and exactly the displayed
prefix flags. These maps are inverse. In particular auxiliary witnesses
do not multiply the number of contributions.

\subsection{The exact parameter-ring map}

Use the $6MN$ independent parameters $E_{a,k,h,\pm}$ of
Section~\ref{sec:sv-shared-verifier}, with clause index $1\le a\le M$,
slot $1\le k\le3$, and witness index $1\le h\le N$ in the order above.
Define a ring homomorphism
\begin{equation}\label{eq:per-trace-parameter-map}
 \Psi_n:\mathbb Z[E_{a,k,h,\pm}]
       \longrightarrow\mathbb Z[Z_{ij}]
\end{equation}
by giving the image of every parameter. In the first $M_0$ clauses,
put $1$ at the unique parameter corresponding to the signed literal
in that slot and $0$ at all other parameters of that slot. For each
$(i,j)$, append one clause, in lexicographic order. In its first slot,
put $Z_{ij}$ at the positive parameter for the witness $B_{ij}$,
put $1$ at the negative parameter for that same witness, and put $0$
at all its other parameters. Put $0$ at every parameter in its
second and third slots. These $n^2$ appended clauses give exactly
the value of $M$ in \eqref{eq:per-trace-size}. The map uses only
variables and the integer constants $0,1$.

This map is surjective: $Z_{ij}$ is the image of its chosen appended
positive parameter, denoted $E^*_{ij}$. Its kernel is explicit.
For each of the other parameters whose image is $0$ or $1$, take
the relation $E-\Psi_n(E)$. These relations generate the kernel:
the quotient is the polynomial ring on the $E^*_{ij}$, and its
inverse to the induced map sends $Z_{ij}$ to $[E^*_{ij}]$.
Thus \eqref{eq:per-trace-parameter-map} is the coordinate map of a
closed affine embedding, with the indicated linear ideal.

Write $\mathcal A(B,R,C)$ for the product of the first $M_0$
arithmetized clauses. The appended clause has first literal
$Z_{ij}B_{ij}+(1-B_{ij})$ and other two literals zero. Therefore
its arithmetized value is exactly
\[
 1-\bigl(1-Z_{ij}B_{ij}-(1-B_{ij})\bigr)(1-0)(1-0)
       =1-B_{ij}+Z_{ij}B_{ij}.
\]
This calculation is a polynomial identity before Boolean substitution.
It proves
\[
 \Psi_n(V_{N,M})=
 \mathcal A(B,R,C)\prod_{i,j}(1-B_{ij}+Z_{ij}B_{ij}).
\]
When evaluated at Boolean witnesses, $\mathcal A$ is one precisely
at the uniquely extended permutation assignments, and zero at all
others. At the permutation $\sigma$ the remaining product is
$\prod_i Z_{i,\sigma(i)}$: the selected entries contribute $Z_{ij}$
and every unselected entry contributes $1$. Hence
\begin{equation}\label{eq:per-trace-projection}
 \Psi_n\bigl(Q_{N,M}(E)\bigr)=\operatorname{per}_n(Z).
\end{equation}
All these equalities hold over $\mathbb Z$. Indeed both sides are
the same finite sum of the explicitly computed integer-coefficient
polynomials, term by term. Base change gives the identity over any
commutative ring; the interpretation of an integer count as a rank
uses the characteristic-zero fibre in the earlier trace theorem.

\subsection{Transport through the original full fibre and its trace}

Work over a characteristic-zero field $k$ and retain $A_N=A_1^{\otimes N}$,
its original source coordinates, its product idempotent $e$, and the
Boolean idempotents $\beta_h$. Apply $\Psi_n$ to the coefficients of
$u(E)=eV_{N,M}(E,\beta)$, and call its image $u_n(Z)$.
There is a commutative square of actual linear maps
\[
\begin{array}{ccc}
 A_N\otimes_k k[E]&\xrightarrow{\ m_{u(E)}\ }&A_N\otimes_k k[E]\\
 \big\downarrow{1\otimes\Psi_n}&&\big\downarrow{1\otimes\Psi_n}\\
 A_N\otimes_k k[Z]&\xrightarrow{\ m_{u_n(Z)}\ }&A_N\otimes_k k[Z].
\end{array}
\]
For a pure tensor the equality says
$\Psi_n(u(E)f)=\Psi_n(u(E))\Psi_n(f)$, which is the multiplicativity
of the coefficient map. Linearity proves it on all inputs.
In the retained product-idempotent basis the multiplication matrices
are diagonal, and coefficient substitution acts on each diagonal entry.
Thus taking the finite sum of those entries commutes with $\Psi_n$,
giving the fully specified equality
\begin{equation}\label{eq:per-original-trace}
 \operatorname{Tr}_{A_N\otimes k[Z]/k[Z]}(m_{u_n(Z)})
        =\operatorname{per}_n(Z).
\end{equation}
The full algebra has dimension $3^{3n^2}$; its Boolean corner has
dimension $2^{3n^2}$. The $x=0$ factors of the original source fibre
are still present and receive weight zero through $e$.
For each nonzero contribution, the $B$ entries encode a permutation
and the $R,C$ entries are the forced flags. This describes every
contributing original fibre point, not merely the final trace value.

The same formulas give a complete finite-support generating series.
In the split idempotent basis, $m_{u_n}$ has one diagonal entry
$\prod_i Z_{i,\sigma(i)}$ for each $\sigma\in\mathfrak S_n$ and
zero at all other basis positions. Consequently
\[
 \det(I-Tm_{u_n})=
 \prod_{\sigma\in\mathfrak S_n}
       \left(1-T\prod_i Z_{i,\sigma(i)}\right),
\]
and, for each integer $d\ge1$,
\[
 \operatorname{Tr}(m_{u_n}^{d})
       =\operatorname{per}_n\bigl((Z_{ij}^{d})_{ij}\bigr).
\]
These follow by multiplying or summing diagonal entries, retaining
their multiplicities even when two monomials acquire the same value
after a later specialization of $Z$. The first coefficient of the
displayed characteristic determinant is $-\operatorname{per}_n(Z)$,
with that sign fixed by choosing the single $-T$ factor in the product.

\subsection{The corresponding class-variety morphism}

The affine parameter projection also gives a precise homogeneous
orbit-closure map. Retain $N,M$ from \eqref{eq:per-trace-size}, let
$a=6MN$, and set $D=3M$. The total degree of $Q_{N,M}$ in its
$a$ parameters is exactly $D$ over a characteristic-zero field.
The upper bound follows because each of the $M$ clauses has parameter
degree at most three. For the lower bound, take the monomial
\[
 \mathfrak m=\prod_{b=1}^{M}\prod_{k=1}^3 E_{b,k,1,-}.
\]
At a fixed Boolean witness its coefficient in $V_{N,M}$ is
$(1-B_1)^{3M}$: the degree-three part of each clause is
$L_{b,1}L_{b,2}L_{b,3}$, with positive sign. Summing gives coefficient
$2^{N-1}$ in $Q_{N,M}$, since exactly that many witnesses have $B_1=0$.
This coefficient is nonzero in characteristic zero, proving the exact
degree without any possible cancellation being suppressed.

Introduce a new coordinate $Z_0$ and define
\[
 f(E,Z_0)=Z_0^D Q_{N,M}(E/Z_0).
\]
Every term has nonnegative $Z_0$ exponent because the degree is at
most $D$, so this defines a degree-$D$ polynomial on
$V=\mathbb C^{a+1}$. Let $W=\mathbb C^{n^2+1}$ have coordinates
$(Z_{ij},Z_0)$. Homogenize every substitution in
\eqref{eq:per-trace-parameter-map}: replace its image $0$ by $0$,
$1$ by $Z_0$, and $Z_{ij}$ by itself. Together with the identity on
$Z_0$, these formulas define a linear map $L:W\to V$.
The projection $p:V\to W$ reads the chosen coordinates
$(E^*_{ij},Z_0)$, so $pL=1_W$. In particular $L$ is injective and
$A=Lp$ is an explicitly specified idempotent endomorphism of $V$.

On $Z_0\ne0$ the proven permanent projection yields
\[
 f(L(Z,Z_0))
 =Z_0^D\operatorname{per}_n(Z/Z_0)
 =Z_0^{D-n}\operatorname{per}_n(Z).
\]
Since $D\ge n$, both sides are polynomials; their difference vanishes
after localization at the nonzero polynomial $Z_0$ in a polynomial
domain, and hence is zero everywhere. Put
\[
 g=f\circ A=p^*\bigl(Z_0^{D-n}\operatorname{per}_n(Z)\bigr)
   \in\operatorname{Sym}^D(V^*).
\]
The section $L^*$ gives $L^*p^*=1$, so no coefficients of the
padded permanent are lost when it is written on this larger space.

Let $G=\operatorname{GL}(V)$ act by $(h\cdot u)(v)=u(h^{-1}v)$.
Then
\[
 \overline{G\cdot g}\ \subseteq\ \overline{G\cdot f}.
\]
To prove the inclusion, write
$v=Lp(v)+(v-Lp(v))$. The two summands belong to $\operatorname{im}L$
and $\ker p$; their intersection is zero because $pL=1$.
Thus $V=\operatorname{im}L\oplus\ker p$, and $A$ acts by $1$ and
$0$ on those summands, respectively. In particular
\[
 \det(A+\epsilon I_V)
   =(1+\epsilon)^{n^2+1}\epsilon^{a-n^2}.
\]
This proves invertibility for $\epsilon\ne0,-1$, retaining both
exceptional values and their full multiplicities.
At every other scalar value,
$f\circ(A+\epsilon I_V)$ belongs to $G\cdot f$, by taking
$h=(A+\epsilon I_V)^{-1}$. A polynomial equation of that orbit
evaluated on this polynomial family vanishes at infinitely many
$\epsilon$, so vanishes identically and at $\epsilon=0$.
This proves $g\in\overline{G\cdot f}$. The latter is closed and
$G$-stable: multiplication by a fixed element of $G$ is an
invertible linear map of the ambient space and sends the orbit to
itself, hence its closure to itself. It therefore contains the
orbit closure of $g$, proving the full inclusion.

Both orbit closures are affine cones. Nonzero scalar multiples of
a degree-$D$ form are obtained from scalar coordinate substitutions,
because every nonzero complex number has a $D$th root. The zero
form is the limit of the polynomial family $\epsilon^D u$, so also
belongs to each closure. Comparing powers of a scalar in an
equation on either cone shows that each homogeneous component of
that equation vanishes there. Thus restriction gives, in every
coordinate-ring degree $j\ge0$, the exact $G$-equivariant sequence
\[
 0\longrightarrow
 \frac{I(\overline{G\cdot g})_j}{I(\overline{G\cdot f})_j}
 \longrightarrow\mathbb C[\overline{G\cdot f}]_j
 \longrightarrow\mathbb C[\overline{G\cdot g}]_j
 \longrightarrow0.
\]
The maps are induced by the identity on the ambient polynomial
coordinate ring; inclusion of the two ideals proves their kernels
and surjectivity. Equivariance follows because both ideals are
$G$-stable. For every finite-dimensional $G$-module $U$,
precomposition with the last surjection is an injection
\[
 \operatorname{Hom}_G(\mathbb C[\overline{G\cdot g}]_j,U)
 \hookrightarrow
 \operatorname{Hom}_G(\mathbb C[\overline{G\cdot f}]_j,U).
\]
A map in the kernel vanishes on the entire image of a surjection
and hence on its whole domain, proving injectivity directly.
This is an exact representation consequence for the stated padded
permanent and verifier class varieties. It does not assert an
inclusion of either class variety in a determinant class variety.

\subsection{Uniform sizes and the exact circuit implication}

The one-index parameter for the preceding family is exactly
\[
 t_n=(N+M-2)(N+M-1)/2+N,
 \qquad N+M=n^3+9n^2.
\]
In particular $t_n=O(n^6)$. The clause construction writes $M=O(n^3)$
clauses with $O(\log(n+1))$-bit witness indices; the full dense parameter
map lists $6MN=O(n^5)$ entries, each either $0$, $1$, or a named
$Z_{ij}$. All lists are generated by the explicit bounded loops above
in polynomial time. A sparse list uses only the three selected literal
positions per ordinary clause and two nonzero positions per appended
clause. This proves uniformity as well as the numerical bounds.

Substitution of \eqref{eq:per-trace-parameter-map} in any division-free
arithmetic circuit computing $Q_{N,M}$ replaces input leaves by
$0,1,Z_{ij}$ and leaves every gate intact. Induction on the gate order
proves the resulting circuit computes the permanent, with no increase
in the number of arithmetic-operation gates. Existing constants remain
unchanged; at most two shared constant vertices for $0,1$ add $O(1)$
constant storage. In a description that repeats constants at every
input leaf, the additional total constant-bit length is at most a
constant times the number of replaced leaves. In either convention
the maximum bit length of a constant becomes at most the maximum of
the previous bound and $1$. If the supplied circuit family is
polynomial in its one-index parameter $t$, its restriction at $t=t_n$
has polynomial size in $n$ because the displayed $t_n$ is polynomial.
An algorithm producing those descriptions remains polynomial after this
explicit substitution. These are exact maps of supplied circuit
descriptions; no small circuit for $Q$ or for the permanent has been
assumed or constructed by them.

This permanent projection makes the trace complexity question concrete.
For a binary input matrix $Z$, every surviving diagonal entry is $0$
or $1$, so the integer permanent lies in $[0,n!]$. Also
$n!\le n^n\le2^{n^2}<2^{N+1}$: the inequality $n\le2^n$ follows
by induction from $1\le2$, since $n+1\le2n$ for $n\ge1$.
Consequently evaluating a supplied integer-constant circuit for the
projected trace modulo $2^{N+1}$ recovers the exact count as its
least nonnegative residue. The preceding gate-substitution map
and the earlier school-multiplication bound apply with this
unchanged modulus, including its exact $N+1=3n^2+1$ bits.
The polynomial-size incidence determinant from the preceding section
still encodes the shared equations. The trace and characteristic
determinant above retain their explicitly exponential ambient sizes.
Obtaining smaller representations or proving the required lower bounds
is additional mathematics, and no Boolean or algebraic complexity
equality or separation follows solely from these exact morphisms.

\section{Exact obstructions for two specified trace representations}
\label{sec:per-trace-obstructions}

Retain the parameter projection, source fibre, and multiplication operator
of the preceding section. Write $\mathcal A_N$ for its full split source
algebra and $u_Z$ for the projected multiplier, where
$N=3n^2$, $M=n^3+6n^2$. In its original point-idempotent basis the
diagonal entries of multiplication by $u_Z$ are zero except at the
uniquely flagged permutation witnesses. Their entries are
$\prod_i Z_{i,\sigma(i)}$. Thus the already constructed polynomial is
\[
 P_n(Z)=\operatorname{Tr}_{\mathcal A_N}(m_{u_Z})
       =\sum_{\sigma\in S_n}\prod_{i=1}^n Z_{i,\sigma(i)}.
\]
The following calculations concern this same trace. No change of source
coordinates or removal of its zero eigenvalues is needed. The original
rational-coordinate fibre retains its characteristic-zero base here.
The integral split idempotent lattice constructed above carries the
same diagonal multiplication matrix with entries in $\mathbb Z[Z]$;
its trace polynomial therefore has integer coefficients. It is that
integral matrix and polynomial, and hence their coefficient matrices,
which can be base changed to any field, including characteristic two.
This does not reduce the source coordinates containing $1/2$ modulo two.

\subsection{Coefficient contraction and its exact rank}

Fix a field $k$, $n\ge1$, and $0\le k_0\le n$. To avoid using the
field symbol as an integer index, denote the cut by $k_0$ throughout.
Let $H_{k_0}$ be the vector space with basis
\[
 h_{\boldsymbol j}=\prod_{i=1}^{k_0}Z_{i,j_i},
 \qquad \boldsymbol j\in\{1,\ldots,n\}^{k_0},
\]
and let $T_{k_0}$ have basis
$t_{\boldsymbol l}=\prod_{i=k_0+1}^{n}Z_{i,l_i}$ indexed by the
remaining $n-k_0$ choices. Empty products give the one-dimensional
spaces at the two endpoints. Since the variable sets are disjoint,
multiplication identifies $H_{k_0}\otimes_k T_{k_0}$ with the space
of polynomials homogeneous of degree one in each row of $Z$.
Indeed, multiplication maps the displayed tensor basis bijectively
onto its distinct monomial basis. Define the coefficient contraction
\[
 \mathcal F_{k_0}:H_{k_0}^*\longrightarrow T_{k_0},\qquad
 \mathcal F_{k_0}(h_{\boldsymbol j}^*)
  =\sum_{\boldsymbol l}
   [h_{\boldsymbol j}t_{\boldsymbol l}]P_n\,t_{\boldsymbol l}.
\]
Here $[m]$ denotes the coefficient of the exact monomial $m$.

\begin{theorem}\label{thm:per-exact-row-rank}
For every field and every displayed cut,
$\operatorname{rank}\mathcal F_{k_0}=\binom n{k_0}$.
There is an integral factorization through the free module on the
$k_0$-element subsets of $\{1,\ldots,n\}$, and an integral identity
minor of that size.
\end{theorem}
\begin{proof}
Let $U_{k_0}$ have basis $e_S$ for those subsets. Define
$a:H_{k_0}^*\to U_{k_0}$ by
$a(h_{\boldsymbol j}^*)=e_{\{j_1,\ldots,j_{k_0}\}}$ if the
indices $j_i$ are distinct, and by zero otherwise. Define
\[
 b(e_S)=\sum_{\{l_{k_0+1},\ldots,l_n\}=S^c
                   \atop\text{all indices distinct}}t_{\boldsymbol l}.
\]
A monomial in $P_n$ uses each column exactly once. Conversely every
choice using each row and column exactly once is the monomial of a
unique permutation and has coefficient one. These two facts give
$\mathcal F_{k_0}=ba$ on every displayed basis vector, including
those with repeated columns. The factorization proves the upper bound.

For each $S$, take $\boldsymbol j(S)$ to be the elements of $S$ in
increasing order assigned to the first $k_0$ rows, and take
$\boldsymbol l(S)$ to be the elements of $S^c$ in increasing order
assigned to the remaining rows. The coefficient of
$h_{\boldsymbol j(S)}t_{\boldsymbol l(S')}$ is one exactly when
$S=S'$: disjointness and cardinalities force $S=(S'^c)^c$.
It is zero otherwise. Thus these rows and columns give the identity
matrix over $\mathbb Z$, proving the lower bound over every field.
No division by a factorial or by a characteristic-dependent integer
has occurred.
\end{proof}

This is also an exact map from the original trace construction.
For a matrix $B(Z)$ over a polynomial ring, finite coefficient extraction
commutes with its trace:
$[m]\operatorname{Tr}B=\sum_v[m]B_{vv}
=\operatorname{Tr}([m]B)$. This follows directly from the definition
of the diagonal sum, for every coefficient and over every base ring.
Apply it to the full $3^N\times3^N$ matrix $m_{u_Z}$.
The projection to its uniquely flagged permutation diagonal entries
therefore gives precisely the maps $a,b$ above, while every retained
zero diagonal entry contributes zero. This proves the coefficient
contraction correspondence; the matrix size $3^N$ is unchanged.

For every integer $d\ge1$ the same argument applied to the previously
proved power trace $\operatorname{Tr}(m_{u_Z}^d)
=\operatorname{per}_n((Z_{ij}^d))$ gives the same rank.
Use the head and tail bases with each exponent $1$ replaced by $d$;
these are still distinct monomials in every characteristic. Their
coefficient matrix is the same integral zero-one matrix.

\subsection{Optimal widths of row-ordered branching programs}

A row-ordered algebraic branching program here means a finite directed
layered graph with layers $0,1,\ldots,n$, one source at layer zero,
one sink at layer $n$, and edges only from layer $i-1$ to layer $i$.
Every such edge is labelled by a homogeneous linear form in
$Z_{i1},\ldots,Z_{in}$ over $k$. The output is the sum over all
source-to-sink paths of the product of their labels. The width $w_i$
is the number of vertices at layer $i$; distinct parallel edges may
equivalently be combined by adding their labels.

\begin{theorem}\label{thm:per-row-abp}
Every such program for $P_n$ has $w_i\ge\binom ni$ at every layer.
These bounds hold simultaneously with equality. The minimum total
number of vertices, including both endpoints, is $2^n$.
An attaining program has exactly $n2^{n-1}$ edges.
\end{theorem}
\begin{proof}
For a vertex $v$ in layer $i$, let $p_v$ be the sum of path products
from the source to $v$, and $q_v$ the corresponding sum from $v$
to the sink. Each full path has exactly one such $v$ and decomposes
uniquely into its prefix and suffix. Distributivity therefore gives
$P_n=\sum_{v\text{ in layer }i}p_vq_v$, with $p_v\in H_i$ and
$q_v\in T_i$. Its coefficient contraction factors as
\[
 H_i^*\longrightarrow k^{w_i}\longrightarrow T_i,
 \qquad\lambda\longmapsto(\lambda(p_v))_v,
 \qquad(c_v)_v\longmapsto\sum_v c_vq_v.
\]
Its rank is at most $w_i$. The preceding theorem proves the lower bound.

For attainment, put at layer $i$ all subsets $S\subseteq\{1,\ldots,n\}$
of size $i$. For every $j\notin S$ put an edge from $S$ to
$S\cup\{j\}$ labelled $Z_{i+1,j}$. A full path chooses the columns
$j_1,\ldots,j_n$ without repetition, so determines one permutation.
Every permutation determines exactly this path. Multiplying its
labels and then summing gives $P_n$. The number of vertices is
$\sum_i\binom ni=2^n$, because each subset has one cardinality and
there are two membership choices for each of the $n$ elements.
The number of edges is
$\sum_{S\subseteq[n]}(n-|S|)=n2^{n-1}$: count pairs $(S,j)$ with
$j\notin S$, choosing $j$ in $n$ ways and an arbitrary subset of
its $n-1$ other elements. Summing the width lower bounds proves
minimality of the stated total vertex number.
\end{proof}

The row condition is part of this precise representation class.
The factorization just proved does not apply to a branching program
that may reuse a row at several layers or intermix row variables
inside a layer. Thus the result gives no lower bound for that larger
class from an unproved identification of its states with $H_i$.
The proved exact map is the prefix--suffix contraction above.

\subsection{An exponential obstruction for nonnegative arithmetic}

Now work over $\mathbb R$. A monotone arithmetic circuit means a
finite acyclic circuit with leaves that are variables $Z_{ij}$ or
nonnegative real constants, and binary addition and multiplication
gates. Its size $s$ is the number of arithmetic-operation gates.
Fan-out is allowed. No subtraction or division is allowed. A parse
tree at a gate is defined recursively: at an addition choose one
of its two children; at a multiplication take both; copy occurrences
when a gate is used more than once. A nonzero parse tree is one
whose constant-leaf product is positive. It contributes that positive
coefficient times the product of its variable leaves.
Induction in a topological ordering proves that a gate polynomial
is the sum of its parse-tree contributions. There are only finitely
many parse trees, since the circuit is finite and acyclic.

\begin{lemma}\label{lem:per-live-gate-profile}
Suppose such a circuit computes $P_n$. A gate occurrence is called
live if it occurs in a nonzero output parse tree, and a gate is live
if it has such an occurrence. For every live gate $g$ there are
uniquely determined subsets $I_g,J_g\subseteq[n]$ of the same size
such that every monomial at $g$ with positive coefficient is a
matching from $I_g$ onto $J_g$. In particular its degree is the
fixed integer $d_g=|I_g|=|J_g|$.
\end{lemma}
\begin{proof}
Fix one nonzero output parse tree and one occurrence of $g$ in it.
Remove the subtree rooted at that occurrence and retain the rest,
including every other occurrence of $g$ if there is one. Let $c m$
be the product of the retained constant and variable leaves, so
$c>0$. Every monomial $t$ with positive coefficient at $g$ has a
nonzero parse tree there, by the parse expansion. Substituting that
tree in the retained occurrence produces a nonzero output parse
tree contributing a positive multiple of $mt$. Nonnegative
coefficients cannot cancel, so $mt$ occurs in $P_n$.

Every monomial of $P_n$ has total exponent one in each row and
column, and exponent at most one on each variable. Consequently
$m$ is a partial matching; all such $t$ use exactly the rows and
columns missing from $m$, once each. These two complementary
sets are $I_g$ and $J_g$. They have the same size because any
monomial $t$ uses one variable per used row and column. At least
one $t$ exists, since the fixed occurrence was nonzero. That
monomial already determines its sets of rows and columns,
showing uniqueness and independence of the chosen live occurrence.
\end{proof}

At a live addition, every child chosen in a nonzero parse tree has
the same profile and degree as the parent: each of its monomials
also occurs at the parent. At a live multiplication with nonzero
children, both children are live by extending the fixed context,
their row sets and their column sets are disjoint, and their unions
are those of the parent. This follows by multiplying any positive
monomial of the first child by one of the second and using the
lemma. In particular their degrees add. These conclusions establish
the needed homogeneity and disjointness from positivity; they have
not been added as circuit hypotheses.

\begin{theorem}\label{thm:per-monotone-lower-bound}
For every $n\ge4$, every monotone arithmetic circuit computing
the original projected trace $P_n$ satisfies
\[
 s\ \ge\ \binom n{\lceil n/3\rceil}
       \ \ge\ 2^{\lceil n/3\rceil}.
\]
\end{theorem}
\begin{proof}
Each of the $n!$ permutation monomials has positive coefficient,
so choose one of its nonzero output parse trees, using any fixed
ordering of the finite tree set to make the choice definite.
Start at its output occurrence, of degree $n$. At an addition
follow its chosen child. At a multiplication follow a child of
larger degree, with a fixed rule for ties. Continue until reaching
the first occurrence of degree at most $2n/3$.

This procedure terminates: it follows a descending path in a finite
tree, and leaf degrees are zero or one. An addition cannot be the
first degree decrease across the threshold, because its chosen
child has the same degree. Hence that first crossing follows a
multiplication of degree greater than $2n/3$. Its selected child
has at least half the parent degree. The selected gate $g$ thus has
\[
 n/3<d_g\le2n/3,
 \quad\text{and in particular}\quad
 \lceil n/3\rceil\le d_g\le\lfloor2n/3\rfloor.
\]
For $n\ge4$ this degree is at least two, so $g$ is an
arithmetic-operation gate, not an input leaf.

Assign the original permutation to this gate. Its chosen parse tree
contains a monomial from $g$, which matches precisely $I_g$ to $J_g$.
Therefore every permutation assigned to $g$ sends $I_g$ onto $J_g$.
For $d_g=d$ the number of all permutations with this property is
exactly $d!(n-d)!$: choose independently a bijection $I_g\to J_g$
and one $I_g^c\to J_g^c$. Each has $d!$ and $(n-d)!$ choices,
respectively, by successive choices of images. The assignments
partition the $n!$ permutations among at most $s$ gates. Thus
\[
 n!\le s\max_{\lceil n/3\rceil\le d\le\lfloor2n/3\rfloor}
                         d!(n-d)!.
\]
The ratio $\binom n{d+1}/\binom nd=(n-d)/(d+1)$ proves that
the binomial coefficients increase up to their middle and then
decrease symmetrically, since $\binom nd=\binom n{n-d}$.
The two endpoints of the displayed range are complementary
integers, so its minimum binomial coefficient is
$\binom n{\lceil n/3\rceil}$. This proves the first bound.
Finally set $r_0=\lceil n/3\rceil$. For $n\ge4$ we have
$r_0\le n/2$: write $n=3a+b$ with $b=0,1,2$ and check the
three inequalities directly, with $a\ge1$ and the case $b=0$
having $a\ge2$ if necessary. For $0\le i<r_0$,
$(n-i)/(r_0-i)\ge n/r_0\ge2$, since
$r_0(n-i)-n(r_0-i)=i(n-r_0)\ge0$. Multiplication gives
\[
 \binom n{r_0}
 =\prod_{i=0}^{r_0-1}\frac{n-i}{r_0-i}
 \ge2^{r_0},
\]
as claimed.
\end{proof}

This theorem constrains positive arithmetic representations of the
specific original-fibre trace after the proved permanent projection.
The universal verifier polynomial before that projection contains
signed coefficients, and the support lemma was proved only for the
nonnegative circuit just defined. The projection is nevertheless the
exact gate-preserving input-substitution map of the preceding section.
If a nonnegative circuit description survives that map as a circuit
of this class, the displayed lower bound applies to its arithmetic
gate count. For signed circuits the cancellation step in the proof
is unavailable; the present theorem makes no assertion that it holds.
The row-ordered and nonnegative bounds provide two explicit
asymptotic obstructions with their maps and domains. The broader
unrestricted obstruction in the research programme remains to be
constructed.

\section{A smaller sign-algebra trace and its exact source correspondence}
\label{sec:compressed-permanent-trace}

The permanent trace of Section~\ref{sec:permanent-trace-projection}
admits a second, smaller finite-algebra representation. We retain the
same $n^2$ coordinates $Z_{ij}$ and the same permanent polynomial;
we construct the correspondence between the two representations explicitly.
The sign sum used below is the Ryser--Glynn expression recorded in
\cite[Equation (8)]{LandsbergRessayre2017SignTrace}; the original
Glynn article is \cite{Glynn2010SignTrace}. All identities and maps
needed here are proved below, independently of those references.

\subsection{The sign algebra, its integral basis, and its regular trace}

Fix an integer $n\ge1$. Put $O=\mathbb Z[1/2]$, $R=O[Z_{ij}]$,
$r_n=2^{n-1}$, and
\[
 B_n=O[d_2,\ldots,d_n]/(d_j^2-1:2\le j\le n),\qquad d_1=1.
\]
For $n=1$ this means $B_1=O$. For each subset
$S\subseteq\{2,\ldots,n\}$ write $d_S=\prod_{j\in S}d_j$,
including $d_\varnothing=1$. The elements $d_S$ form an $O$-basis.
Indeed, successive division by the monic polynomials $d_j^2-1$
gives a unique remainder with exponents zero or one: viewing the
ring as a polynomial ring in the last variable over the preceding
quotient, monic division gives a free module with basis $1,d_j$.
Induction proves the assertion over $O$ and equally over $\mathbb Z$.
Multiplication is
\[
 d_Sd_T=d_{S\mathbin{\triangle}T},
\]
where $S\mathbin{\triangle}T=(S\setminus T)\cup(T\setminus S)$.
There are exactly $r_n$ basis elements, one for each independent
membership choice in a set of size $n-1$.

Write $\tau_B(b)=\operatorname{Tr}_{B_n/O}(m_b)$ for the regular
trace, where $m_b(x)=bx$. Multiplication by $d_S$ permutes the basis
by $T\mapsto S\mathbin{\triangle}T$. A basis vector is fixed if
and only if $S=\varnothing$, since cancelling $T$ under symmetric
difference gives that equivalence. Thus
\begin{equation}\label{eq:compressed-regular-trace}
 \tau_B\Bigl(\sum_S b_Sd_S\Bigr)=r_n b_\varnothing.
\end{equation}
This identity holds first over $\mathbb Z$ and then over every
coefficient ring by taking the same finite diagonal sum.

There is also an explicit split basis. Let
$\Delta_n=\{\delta\in\{-1,1\}^n:\delta_1=1\}$ and set
\[
 \epsilon_\delta=\prod_{j=2}^n\frac{1+\delta_jd_j}{2}.
\]
Each factor is idempotent, since
$(1+\delta_jd_j)^2/4=(1+\delta_jd_j)/2$.
Opposite factors multiply to $(1-d_j^2)/4=0$, and summing the
two factors for one coordinate gives one. Multiplying these identities
shows that the $\epsilon_\delta$ are orthogonal idempotents summing
to one. Evaluation at $d_j=\eta_j$ gives
$\epsilon_\delta(\eta)=1$ when $\eta=\delta$ and zero otherwise.
The identity
$d_S=\sum_\delta(\prod_{j\in S}\delta_j)\epsilon_\delta$
follows by multiplying $d_j\epsilon_\delta=\delta_j\epsilon_\delta$
and summing the idempotents. It proves the mutually inverse algebra maps
\[
 B_n\longrightarrow O^{\Delta_n},\quad b\longmapsto(b(\delta))_\delta,
 \qquad(c_\delta)_\delta\longmapsto\sum_\delta c_\delta\epsilon_\delta.
\]
In particular $\tau_B(b)=\sum_\delta b(\delta)$.

For completeness the factor inverted here has exact integral content.
Before localizing, the evaluation matrix from the $d_S$ basis to
the integral point basis is
$H_{\delta,S}=\prod_{j\in S}\delta_j$. For two subsets $S,T$,
the sum $\sum_\delta H_{\delta,S}H_{\delta,T}$ is $r_n$ if
$S=T$ and zero otherwise: in the latter case choose a coordinate
in $S\mathbin{\triangle}T$ and pair each $\delta$ with the sign
vector obtained by flipping that coordinate. Therefore
\[
 H^{\mathsf T}H=r_n I,\qquad H^{-1}=r_n^{-1}H^{\mathsf T}
 \quad\text{over }\mathbb Q.
\]
The image of $\mathbb Z[d_2,\ldots,d_n]/(d_j^2-1)$ in
$\mathbb Z^{\Delta_n}$ consists exactly of integer vectors $v$
such that every component of $H^{\mathsf T}v$ is divisible by $r_n$:
necessity follows from the first identity, and sufficiency follows
by applying the displayed inverse. Thus the integral evaluation map
is injective and becomes an isomorphism after inverting two. For
$n\ge2$ it is not surjective: the coordinate vector supported at
a single $\delta$ has $H^{\mathsf T}v$ with every component $1$
or $-1$, which is not divisible by $r_n\ge2$.
No integral identification of these two lattices has been made.

\subsection{The exact parity computation, with both orientations}

Put $\chi_n=\prod_{j=1}^n d_j=d_{\{2,\ldots,n\}}$ and define
row-linear elements and their multiplier by
\begin{equation}\label{eq:compressed-row-multiplier}
 L_i(Z)=\sum_{j=1}^n d_j Z_{ij},\qquad
 v_n(Z)=r_n^{-1}\chi_n\prod_{i=1}^n L_i(Z)
 \quad\text{in }B_n\otimes_O R.
\end{equation}

\begin{theorem}\label{thm:compressed-permanent-trace}
The regular multiplication trace on the rank-$r_n$ algebra satisfies
\[
 \operatorname{Tr}_{B_n\otimes R/R}(m_{v_n(Z)})
 =\sum_{\sigma\in\mathfrak S_n}\prod_{i=1}^n Z_{i,\sigma(i)}
 =\operatorname{per}_n(Z).
\]
The identity holds over $O$ and after any homomorphism from $O$
to a commutative ring, in particular over $\mathbb Q$ and every
field of odd characteristic.
\end{theorem}
\begin{proof}
Expand the product in \eqref{eq:compressed-row-multiplier} without
changing its row order:
\[
 v_n(Z)=r_n^{-1}
 \sum_{(j_1,\ldots,j_n)\in[n]^n}
      \chi_n d_{j_1}\cdots d_{j_n}\prod_{i=1}^n Z_{ij_i}.
\]
For a tuple let $m_j$ be the number of its entries equal to $j$.
The group-basis monomial $\chi_n\prod_i d_{j_i}$ is one exactly
when $m_j+1$ is even for each $j\ge2$, or equivalently when each
of $m_2,\ldots,m_n$ is odd. In that case write $m_j=1+2a_j$
for $j\ge2$, with $a_j\ge0$. The equality $\sum_jm_j=n$ becomes
\[
 m_1+2\sum_{j=2}^n a_j=1.
\]
It forces $m_1=1$ and every $a_j=0$. Conversely, if every column
is chosen once, all the exponents are even as required. Thus exactly
the permutation tuples contribute to the coefficient of $d_\varnothing$,
each with coefficient $r_n^{-1}$ times its original row monomial.
Equation~\eqref{eq:compressed-regular-trace} multiplies that coefficient
by the exact factor $r_n$, proving the formula. For $n=1$ the
same equation says $m_1=1$; directly $r_1=\chi_1=d_1=1$ and
$v_1(Z)=Z_{11}$ in the rank-one algebra. Finally all the calculations
are identities in the displayed free algebra over $O$, so applying
a coefficient homomorphism preserves them, including their traces.
\end{proof}

The split-basis form of this equality is the fully signed expression
\[
 \operatorname{per}_n(Z)=2^{1-n}
 \sum_{\delta\in\Delta_n}
 \Bigl(\prod_{j=1}^n\delta_j\Bigr)
 \prod_{i=1}^n\Bigl(\sum_{j=1}^n\delta_j Z_{ij}\Bigr).
\]
The prefactor and product of signs are retained. To give the other
orientation explicitly, use a second copy of the sign algebra with
generators $e_2,\ldots,e_n$, $e_1=1$, and put
\[
 \widehat v_n(Z)=2^{1-n}
 \Bigl(\prod_{i=1}^n e_i\Bigr)
 \prod_{j=1}^n\Bigl(\sum_{i=1}^n e_i Z_{ij}\Bigr).
\]
Let the coefficient-ring involution send $Z_{ij}$ to $Z_{ji}$,
and send $d_j$ in the first sign algebra to $e_j$ in the second.
This is an isomorphism of the underlying $O$-algebras with polynomial
coefficients; its square, with the reverse generator renaming, is
the identity. It sends $v_n(Z)$ to $\widehat v_n(Z)$ by direct
substitution. The trace of the latter is
$\operatorname{per}_n(Z^{\mathsf T})=\operatorname{per}_n(Z)$:
the term $\prod_i Z_{\sigma(i),i}$ becomes
$\prod_j Z_{j,\sigma^{-1}(j)}$, and inversion permutes
$\mathfrak S_n$ bijectively. Thus transposition is a specified
coordinate map, not an implicit interchange of the original rows
and columns in \eqref{eq:compressed-row-multiplier}.

Over $\mathbb Z$ the unscaled numerator
$w_n(Z)=\chi_n\prod_iL_i(Z)$ is an element of the integral group
algebra and satisfies $\operatorname{Tr}(m_{w_n})=r_n\operatorname{per}_n$.
For $n\ge2$ the scaled element $v_n$ is not in that integral
group algebra with polynomial coefficients. The monomial
$Z_{11}\cdots Z_{n1}$ occurs exactly once in the expansion and has
coefficient $\chi_n/r_n$, which is not integral in the $d_S$ basis.
The resulting trace polynomial nevertheless belongs to $\mathbb Z[Z]$
by the proved identity. It can be reduced to characteristic two
through that integer polynomial or the original split integral trace
matrix. There is no homomorphism $O\to k$ for a field $k$ of
characteristic two, because the unit two would map to zero.
The scaled sign-algebra operator is therefore not specialized by
such a nonexistent coefficient map. Its integer numerator has trace
zero there for $n\ge2$, which cannot by itself recover the permanent;
at $Z=I_n$ the permanent is one.

\subsection{The original source fibre over the same coefficient ring}

Retain the original map with its source orientation $(x,y,w)$ and
target orientation $(a,b,c)$:
\begin{align*}
 F_1&=(1+xy)^3w+y^2(1+xy)(4+3xy),\\
 F_2&=y+3x(1+xy)^2w+3xy^2(4+3xy),\\
 F_3&=2x-3x^2y-x^3w.
\end{align*}
It still has $\det DF=-2$, and its retained polynomial and derivative
are $P=cr^3-2r^2+br-2a$ and $P'=3cr^2-4r+b$. Neither expression
is altered by the sign algebra. The original fibre at $(-1/4,0,0)$
is defined over $O$, and the isomorphism
\[
 O[x,y,w]/(F_1+1/4,F_2,F_3)
 \longrightarrow O[X]/(X^3-X),\qquad
 (x,y,w)\longmapsto\bigl(X,-3X/2,(27X^2-1)/4\bigr)
\]
continues to hold there. To check the base ring explicitly, put
$h=2-3xy-x^2w$. The identity
$1=h/2+x(3y+xw)/2$ makes $(x)$ and $(h)$ comaximal after imposing
$F_1+1/4,F_2$. As $F_3=xh$, the quotient splits into its $x=0$
and $h=0$ factors by the same two-ideal intersection argument.
The first is $O$, with $(x,y,w)=(0,0,-1/4)$. In the second,
$x(3y+xw)=2$ makes $x$ a unit; putting $r=y+1/x$, $\alpha=1/x$
gives $F_1=r^2+r\alpha$, $F_2=4r+2\alpha$, hence
$\alpha=-2r$, $r^2=1/4$. Conversely these equations make $r$
a unit and recover the original source coordinates and target
equations. This factor is $O[r]/(r^2-1/4)$, with
$x=-2r$, $y=-3x/2$, $w=13/2$. All inverses used are powers of
two or the stated source unit. The two factors give exactly
$O[X]/(X^3-X)$, with its original $x=0$ point retained.

The idempotents $e_-=(X^2-X)/2$, $e_0=1-X^2$,
$e_+=(X^2+X)/2$ give evaluation at $X=-1,0,1$ over $O$ as
in Proposition~\ref{prop:sv-full-fibre}. Take $N=3n^2$ tensor
factors and write $\mathcal A_N$ for the resulting full source
algebra over $O$. Its point set is $\Omega_N=\{-1,0,1\}^N$;
write $E_\omega$ for its product idempotents. It has rank $3^N$,
and $\Lambda_N=\bigoplus_{\omega\in\Omega_N}\mathbb Z E_\omega$
is the original split integral lattice, with its coordinatewise
multiplication. In particular $\mathcal A_N=\Lambda_N\otimes O$.
This lattice differs from the unsplit integral presentation with
source-coordinate denominators, as already proved for each factor.

For a permutation $\sigma$, let $\omega(\sigma)$ be the original
source point with $B_{ij}=1$ exactly when $j=\sigma(i)$ and with
the unique row and column prefix flags of the preceding construction.
In every witness coordinate $h$ this point has
\[
 X_h=2B_h-1,\qquad y_h=-3X_h/2,\qquad
 w_h=(27X_h^2-1)/4.
\]
Here $B_h$ denotes the appropriate original witness bit, including
the two groups of flags. Distinct permutations have distinct
$B$-arrays and therefore distinct points. Put
\[
 f_n=\sum_{\sigma\in\mathfrak S_n}E_{\omega(\sigma)},
 \qquad A_{\mathrm{perm}}=f_n\mathcal A_N.
\]
Orthogonality proves $f_n^2=f_n$ and gives the algebra decomposition
$\mathcal A_N=A_{\mathrm{perm}}\oplus(1-f_n)\mathcal A_N$.
The first ideal has identity $f_n$ and rank $n!$. The map
$a\mapsto f_na$ is a surjective algebra map to this ideal with
kernel $(1-f_n)\mathcal A_N$. It retains the complementary ideal
as its explicitly specified kernel, including every boundary point
with some $X_h=0$.

The projected original multiplier is exactly
\begin{equation}\label{eq:compressed-original-u}
 u_n(Z)=\sum_{\sigma\in\mathfrak S_n}
           E_{\omega(\sigma)}\prod_{i=1}^n Z_{i,\sigma(i)}.
\end{equation}
Indeed the unweighted constraint-clause product is one precisely on those uniquely
flagged assignments, each appended weighted clause supplies the
specified entry, and the original product idempotent kills every
boundary factor. These are precisely the evaluations of the
projected multiplier established above; the split evaluation
isomorphism proves equality as an element, not only as a trace.
Since $u_n=f_nu_n$, its multiplication is zero on the full
complementary ideal. Consequently
\[
 \operatorname{Tr}_{\mathcal A_N}(m_{u_n})
 =\operatorname{Tr}_{A_{\mathrm{perm}}}(m_{u_n})
 =\operatorname{per}_n(Z).
\]
The full original rank $3^{3n^2}$ and all its zero eigenvalues
remain present in the first trace.

\subsection{A common tensor module and an exact sequence of traces}

Let $T_n=(O^n)^{\otimes n}$, with factor $i$ assigned to row $i$.
For $\boldsymbol j=(j_1,\ldots,j_n)\in[n]^n$ denote its tensor
basis vector by $t_{\boldsymbol j}$. Define three $O$-linear maps
on every basis vector by
\[
 \begin{split}
 S_n:T_n&\longrightarrow A_{\mathrm{perm}},\qquad
 S_n(t_{\boldsymbol j})=
 \begin{cases}E_{\omega(\sigma)},&j_i=\sigma(i)\text{ for all }i,\\
 0,&\text{otherwise},\end{cases}\\
 K_n:T_n&\longrightarrow B_n,\qquad
 K_n(t_{\boldsymbol j})=r_n^{-1}\chi_n\prod_i d_{j_i},\\
 \varepsilon_n:T_n&\longrightarrow O,\qquad
 \varepsilon_n(t_{\boldsymbol j})=
 \begin{cases}1,&\boldsymbol j\text{ is a permutation tuple},\\
 0,&\text{otherwise}.\end{cases}
 \end{split}
\]
Writing $\tau_A$ for regular trace on $A_{\mathrm{perm}}$, we have
the equalities of actual linear maps
\begin{equation}\label{eq:compressed-trace-square}
 \tau_AS_n=\varepsilon_n=\tau_BK_n.
\end{equation}
The first follows because a coordinate idempotent has trace one;
the second is exactly the parity computation for one tuple.
The universal row-choice tensor
\[
 t(Z)=\bigotimes_{i=1}^n\Bigl(\sum_{j=1}^n Z_{ij}t_j\Bigr)
     =\sum_{\boldsymbol j}\Bigl(\prod_i Z_{ij_i}\Bigr)t_{\boldsymbol j}
 \quad\text{in }T_n\otimes_O R
\]
is sent by $S_n$ to \eqref{eq:compressed-original-u}, by $K_n$
to \eqref{eq:compressed-row-multiplier}, and by $\varepsilon_n$
to the permanent. Thus \eqref{eq:compressed-trace-square} relates
every coefficient before taking the trace of the polynomial.

The kernels and surjections in this correspondence are explicit.
The map $S_n$ is surjective, and its kernel has as basis the
nonpermutation tuples, because the remaining images are the
distinct basis idempotents. For a tuple put
\[
 p(\boldsymbol j)=\{j\in\{2,\ldots,n\}:
            j\text{ occurs an odd number of times in }\boldsymbol j\}.
\]
For each subset $U\subseteq\{2,\ldots,n\}$ choose
$\boldsymbol a(U)$ to consist first of the elements of $U$ in
increasing order and then of $n-|U|$ entries equal to one.
Its parity set is $U$. Thus
$K_n(t_{\boldsymbol a(U)})=r_n^{-1}\chi_nd_U$,
and these elements form a basis of $B_n$ because multiplication
by the unit $\chi_n/r_n$ permutes and rescales its basis.
This proves surjectivity with a specified section. The differences
\[
 t_{\boldsymbol j}-t_{\boldsymbol a(p(\boldsymbol j))}
 \quad\bigl(\boldsymbol j\ne\boldsymbol a(p(\boldsymbol j))\bigr)
\]
form a basis of $\ker K_n$: subtracting these differences reduces
an arbitrary vector to the selected representatives, whose images
are independent; each difference has coefficient one on its own
unselected basis vector and zero on the other unselected vectors,
which proves independence.

The parity proof also shows that $p(\boldsymbol j)=\{2,\ldots,n\}$
holds exactly for permutation tuples. Every other selected representative
$\boldsymbol a(U)$ is therefore nonpermutational. The regular trace
kernel in $B_n$ is the free span of $d_V$ with $V\ne\varnothing$,
by \eqref{eq:compressed-regular-trace} and the invertibility of $r_n$.
The preceding representatives consequently prove
\begin{equation}\label{eq:compressed-kernel-trace}
 K_n(\ker S_n)=\ker\tau_B.
\end{equation}

\begin{proposition}\label{prop:compressed-exact-sequence}
There is an exact sequence of $O$-modules
\[
 T_n\xrightarrow{\ (S_n,-K_n)\ }
 A_{\mathrm{perm}}\oplus B_n
 \xrightarrow{\ (a,b)\mapsto\tau_A(a)+\tau_B(b)\ }O
 \longrightarrow0.
\]
The first image has rank $n!+r_n-1$, and its kernel
$\ker S_n\cap\ker K_n$ is free of rank $n^n-n!-r_n+1$.
\end{proposition}
\begin{proof}
Equation~\eqref{eq:compressed-trace-square} says the composite is
zero. The second map is surjective because it sends
$(E_{\omega(\mathrm{id})},0)$ to one. Suppose
$\tau_A(a)+\tau_B(b)=0$. Choose $t_0\in T_n$ with $S_n(t_0)=a$.
Then $b+K_n(t_0)$ has trace
$\tau_B(b)+\varepsilon_n(t_0)=\tau_B(b)+\tau_A(a)=0$.
Equation~\eqref{eq:compressed-kernel-trace} supplies
$t_1\in\ker S_n$ such that $K_n(t_1)=-b-K_n(t_0)$.
It follows that $(S_n,-K_n)(t_0+t_1)=(a,b)$, proving exactness.

There is an explicit basis for the kernel of the second arrow:
take $(E_{\omega(\sigma)}-E_{\omega(\mathrm{id})},0)$ for
$\sigma\ne\mathrm{id}$, together with
$(0,d_V)$ for $V\ne\varnothing$ and the single vector
$(-r_nE_{\omega(\mathrm{id})},1)$. Their traces vanish.
Subtracting the nonidentity $A_{\mathrm{perm}}$ coefficients and
the nonempty $B_n$ coefficients from any trace-zero pair reduces
it to $(a_0E_{\omega(\mathrm{id})},b_0)$ with
$a_0+r_nb_0=0$, which is the last vector multiplied by $b_0$.
The same coordinates prove independence. There are $n!+r_n-1$
such vectors. Choose lifts of these basis vectors under the first
surjection just proved; their linear extension gives a section.
Then $T_n$ is the direct sum of their free span and the first
kernel. For freeness of that kernel without an unstated theorem,
decompose $T_n$ into permutation and nonpermutation basis vectors.
An element of $\ker S_n$ uses only the latter. Group those vectors
by their parity sets $U\ne\{2,\ldots,n\}$. On each group $K_n$
sends every vector to the same nonzero distinct basis multiple.
Its kernel therefore has basis the differences against the selected
representative in that group, by the argument above. The total
number is $(n^n-n!)-(r_n-1)$, as asserted.
\end{proof}

For $n\ge2$ this common-tensor construction cannot be replaced by
a linear map making either of its two surjections factor through
the other. The tuple $(1,\ldots,1)$ is nonpermutational, so its
$S_n$ image is zero whereas its $K_n$ image is the nonzero unit
$\chi_n/r_n$. This disproves a factorization $K_n=L S_n$.
For two distinct permutations, the difference of their tensor
vectors has $K_n$ image zero, since each image is $1/r_n$,
whereas its $S_n$ image is a nonzero difference of basis idempotents.
This disproves a factorization $S_n=M K_n$. The exact sequence
and coefficient trace maps above supply the proved relation despite
these specific obstructions. For $n=1$ both coefficient maps are
isomorphisms between rank-one modules, and the original full source
still has rank $3^3=27$ with its stated complement.
More explicitly, at $n=1$ the unique flagged point is
$(X_{B_{11}},X_{R_{11}},X_{C_{11}})=(1,1,1)$. Evaluation there
is a surjective $O$-algebra map $\mathcal A_3\to O=B_1$;
its kernel is the span of the other 26 point idempotents, and
it sends $u_1(Z)=Z_{11}E_{(1,1,1)}$ to $v_1(Z)=Z_{11}$.

\subsection{A multiplication obstruction and the unequal power traces}

Equality of the first traces does not give equality of their spectra.
Specialize the original coordinates to the identity matrix $I_n$.
In \eqref{eq:compressed-original-u} only $\sigma=\mathrm{id}$
survives, so $u_n(I_n)=E_{\omega(\mathrm{id})}$ is an idempotent
in the full original algebra and its permutation ideal. On the sign
side $L_i(I_n)=d_i$, and hence
\[
 v_n(I_n)=r_n^{-1}\chi_n\prod_i d_i=r_n^{-1}\chi_n^2=r_n^{-1}1.
\]
For $n\ge2$ over $\mathbb Q$ this element is not idempotent,
since $r_n^{-2}-r_n^{-1}=(1-r_n)/r_n^2\ne0$.
Every multiplicative map, whether unital or not, sends an idempotent
to an idempotent. Thus no $\mathbb Q[Z]$-linear multiplicative map
from the original full source algebra or its permutation ideal
to $B_n\otimes\mathbb Q[Z]$ sends $u_n(Z)$ to $v_n(Z)$:
such a map would specialize at $Z=I_n$ to the impossible map just
described. This is an obstruction for the stated compatible base
map. No assertion about arbitrary changes of the $Z$ coordinates
is needed. In odd characteristic the specific test may disappear
when $r_n=1$ in the field; the rational assertion retains its domain.

There is a stronger generic obstruction that covers every field $k$
of characteristic different from two, even after passing to the
rational function field $K=k(Z_{ij})$. Both source algebras and the
sign algebra split as products of copies of $K$, by their stated
point idempotents. Compose a $K$-linear multiplicative map from
either source algebra to the sign algebra with any sign-coordinate
evaluation. Each primitive source idempotent must map to an
idempotent of the field $K$, hence to zero or one: the equation
$x(x-1)=0$ in a field forces these two possibilities. Orthogonality
allows at most one of them to map to one. Therefore the image
of $u_n$ in this coordinate is either zero or one of its permutation
monomials. At a fixed sign vector $\delta$, however, the coefficient
of the nonpermutation monomial $\prod_i Z_{i1}$ in $v_n(\delta)$
is $r_n^{-1}\prod_j\delta_j\ne0$, and the coefficient of every
permutation monomial is $r_n^{-1}\ne0$. These statements follow
by selecting the indicated unique term in each row sum; all row
monomials are distinct. For $n\ge2$ this polynomial is consequently
neither zero nor a single permutation monomial. The injection
$k[Z]\hookrightarrow K$ preserves these inequalities. This proves
that no such $K$-linear multiplicative map sends $u_n$ to $v_n$,
including nonunital maps. It also proves the compatible polynomial
base-map obstruction over every such field, since a map over
$k[Z]$ would extend to $K$ by scalar extension.

For every integer $d\ge1$ the original power trace and the new one
are, respectively,
\[
 \operatorname{Tr}(m_{u_n}^d)=
       \sum_{\sigma}\prod_i Z_{i,\sigma(i)}^d,
 \qquad
 \operatorname{Tr}(m_{v_n}^d)=
       r_n^{-d}\sum_{\delta\in\Delta_n}
       \Bigl(\prod_j\delta_j\Bigr)^d
       \prod_i\Bigl(\sum_j\delta_j Z_{ij}\Bigr)^d.
\]
These follow by taking powers of the diagonal entries in their
specified point bases. At $Z=I_n$ they equal $1$ and $r_n^{1-d}$.
In particular $d=2$ distinguishes them over $\mathbb Q$ for $n\ge2$.
Their characteristic determinants at this same specialization are
$1-T$ and $(1-T/r_n)^{r_n}$, respectively: the full original
operator has one eigenvalue one and all other eigenvalues zero,
whereas the sign operator has $r_n$ copies of $1/r_n$.
The common first trace, the exact coefficient maps, and these
different higher traces are all retained simultaneously.

\subsection{The exact dimension obstruction for a fixed algebra}

Let $k$ now be any field and $D$ a finite-dimensional associative
unital $k$-algebra of dimension $q_D$. Fix $c\in D$ and elements
$a_{ij}\in D$, all independent of the polynomial coordinates $Z$.
No commutativity of $D$ is assumed. Define
\[
 A_i(Z)=\sum_{j=1}^n a_{ij}Z_{ij},\qquad
 Q_D(Z)=\operatorname{Tr}_{D\otimes k[Z]/k[Z]}
                 (m_{c A_1(Z)\cdots A_n(Z)}),
\]
where the product has the displayed order, the $Z$ coordinates are
central, and $m_b(x)=bx$ is left multiplication. Associativity
implies $m_bm_{b'}=m_{bb'}$ directly on every $x\in D$.

\begin{theorem}\label{thm:compressed-fixed-algebra-bound}
For every cut $0\le k_0\le n$, the row coefficient contraction
of $Q_D$ has rank at most $q_D$. Consequently, if
$Q_D(Z)=\operatorname{per}_n(Z)$ as polynomials, then
\[
 q_D\ge\binom n{k_0}\quad(0\le k_0\le n),\qquad
 q_D\ge\binom n{\lfloor n/2\rfloor}.
\]
\end{theorem}
\begin{proof}
Use the head and tail monomial spaces $H_{k_0},T_{k_0}$ of
Theorem~\ref{thm:per-exact-row-rank}. Define a linear map on
the dual head basis by
\[
 a:H_{k_0}^*\longrightarrow D,\qquad
 h_{\boldsymbol j}^*\longmapsto
 a_{1j_1}\cdots a_{k_0j_{k_0}},
\]
with the empty product equal to $1_D$. Define
\[
 b:D\longrightarrow T_{k_0},\qquad
 z\longmapsto\sum_{\boldsymbol l}
 \operatorname{Tr}_D\bigl(m_{c z
      a_{k_0+1,l_{k_0+1}}\cdots a_{n,l_n}}\bigr)
 t_{\boldsymbol l}.
\]
Both are $k$-linear, since multiplication by a fixed factor is
linear and the matrix trace is linear. Expanding the ordered product
in $Q_D$ shows that its coefficient at
$h_{\boldsymbol j}t_{\boldsymbol l}$ is exactly the displayed
trace with $z=a_{1j_1}\cdots a_{k_0j_{k_0}}$. Its coefficient
contraction is therefore $ba$ and has rank at most $\dim_kD=q_D$.
If it is the permanent contraction, the integral identity minor
of Theorem~\ref{thm:per-exact-row-rank} gives rank $\binom n{k_0}$
over this same field. The middle choice gives the last inequality.
\end{proof}

For any field of odd characteristic, or of characteristic zero,
\eqref{eq:compressed-row-multiplier} is a representation of this
precise form with $D=B_n\otimes k$, $c=\chi_n/r_n$,
$a_{ij}=d_j$, and $q_D=2^{n-1}$. Thus the construction gives an
upper bound $2^{n-1}$ for this algebra dimension, while the theorem
gives its universal lower bound $\binom n{\lfloor n/2\rfloor}$.
This lower bound is exponential: the binomial coefficients count
all subsets by their cardinalities and sum to $2^n$; the ratio
$\binom n{k+1}/\binom nk=(n-k)/(k+1)$ shows the middle coefficient
is maximal, so it is at least $2^n/(n+1)$.
The factorization assumes a fixed algebra with coordinate-independent
structure and coefficients. It places no restriction on representing
an arbitrary polynomial as a single unconstrained element of a
rank-one algebra; in that larger description the element itself
could already contain the permanent.

The sign algebra is smaller than the original source algebra for
every $n\ge1$, since $2^{n-1}\le3^{n-1}<3^{3n^2}$, and smaller
than its Boolean corner $2^{3n^2}$. Its stated dimension remains
exponential. Enumerating its $2^{n-1}$ signs and directly evaluating
the displayed $n$ row sums and product uses $O(n^2 2^{n-1})$
arithmetic operations: each sign requires $n^2$ signed entries,
$n(n-1)$ additions for the row sums, $n-1$ multiplications for
their product, and a bounded number of sign and scale operations;
the resulting values are summed over that many signs. These loops
give a uniform explicit algorithm. For integer $L\ge0$ and entries with
$|Z_{ij}|<2^L$, each row sum has absolute value at most $n2^L$,
and the integer signed numerator is bounded in absolute value by
$2^{n-1}n^n2^{nL}$. Its exact division by $2^{n-1}$ yields the
integer permanent by the proof above; a sufficient signed bit bound
for that numerator is
$n-1+n\lceil\log_2 n\rceil+nL+2$.
No polynomial-size representation or unrestricted circuit separation
is inferred from this compressed trace or its fixed-algebra bound.

\section{The determinant orbit interface, with its exact obstruction scope}
\label{sec:gct-obstruction-interface}

All varieties and representations in this section are over $\mathbb C$.
Put $V_D=\operatorname{Mat}_{D\times D}(\mathbb C)$,
$W_D=\operatorname{Sym}^D(V_D^*)$, and $G_D=\operatorname{GL}(V_D)$.
The action, including its inverse, is
\[
 (g\cdot f)(X)=f(g^{-1}X),\qquad
 \Omega_D=\overline{G_D\cdot\det_D}\subset W_D.
\]
Closures below are Zariski closures. The affine coordinate ring of $W_D$
is $B_D=\operatorname{Sym}(W_D^*)$, with degree-$e$ part
$\operatorname{Sym}^e(\operatorname{Sym}^D V_D)$.
Here $D$ is the matrix size; $e$ is a coordinate-ring degree and is
not the polynomial degree $D$.

\subsection{Original padding and stabilizer conventions}

GCT I, Section 4, Proposition 4.1, constructs the determinant orbit-closure
map from a formula through a possibly singular linear substitution;
GCT II, Section 3.1, uses an independent padding variable
\cite{MS2001Interface,MS2006Interface}. Explicitly, for $m<D$, let
$X$ be the bottom-right $m\times m$ submatrix of $Y\in V_D$ and
$y=Y_{11}$, which is outside that submatrix. The linear map on forms is
\[
 j_{D,m}:\operatorname{Sym}^m(V_m^*)\longrightarrow W_D,
 \qquad h\longmapsto y^{D-m}h(X).
\]
It is injective: the ring substitution retaining $X$ and putting $y=1$
and all other entries of $Y$ equal to zero is a left inverse on its
image. There is also a homogeneous coefficient left inverse, obtained
by taking the coefficient of $y^{D-m}$ and then setting the other
outside variables to zero. Thus this map has no kernel and explicitly
retains every coefficient of $h$. In the notation of GCT II, its $n$
is the present $m$, and its $m$ is the present $D$.

Those papers use projective $\operatorname{SL}(V_D)$ orbit closures.
The precise comparison with the affine convention above is as follows.
For $g\in G_D$, choose $\mu\in\mathbb C^*$ with
$\mu^{D^2}=\det(g)$ and set $s=\mu^{-1}g$. Then
$s\in\operatorname{SL}(V_D)$ and
$g\cdot f=\mu^{-D}(s\cdot f)$. Hence the projective orbits agree.
Moreover every nonzero scalar multiple of $f$ is in $G_D\cdot f$,
because $\lambda I$ acts by $\lambda^{-D}$. The affine closure is
therefore a cone, and its homogeneous vanishing ideal is precisely the
homogeneous ideal of the projective orbit closure: a homogeneous equation
vanishes on the projective orbit exactly when it vanishes on every
nonzero vector above it. This proves the comparison of coordinate rings,
with their degrees retained.

The original stabilizer discussion distinguishes a form from its line.
For example the linear substitution $T_{A,B}(Y)=AYB^{-1}$, with
$A,B\in\operatorname{GL}_D$, satisfies
\[
 \det(T_{A,B}Y)=\frac{\det A}{\det B}\det Y,
 \qquad
 T_{A,B}\cdot\det_D=\frac{\det B}{\det A}\det_D.
\]
The first identity follows by multiplication of determinants and the
second by inserting $T_{A,B}^{-1}$. Thus every such substitution fixes
the determinant line, and it fixes the form exactly when
$\det A=\det B$. Transposition fixes the form as well. To intersect
these substitutions with the original ambient $\operatorname{SL}(V_D)$,
retain the additional equation
$(\det A/\det B)^D=1$; with transposition the left side is multiplied
by $(-1)^{D(D-1)/2}$. Indeed left multiplication acts independently on
$D$ columns, right multiplication on $D$ rows, and transposition exchanges
the $D(D-1)/2$ off-diagonal coordinate pairs. These equations prevent
identifying the line stabilizer with the form stabilizer.
For the permanent, row and column permutations and transposition preserve
the form, and row scalars $\alpha_i$ and column scalars $\beta_j$
multiply it by $\prod_i\alpha_i\prod_j\beta_j$: each defining monomial
uses each row and column once. These verified subgroups supply the
stabilizer context used here; no classification of a new fibre stabilizer
is inferred from them.

\subsection{An explicit endomorphism and a strict boundary point}

For the retained-fibre family take $n\ge1$, $D=4n+1$, and
$E_n=\mathbb C^{5n+1}$ with ordered coordinates
\[
 (a_1,b_1,c_1,r_1,t_1,\ldots,a_n,b_n,c_n,r_n,t_n,z).
\]
Write $L_n:E_n\to V_D$ for the linear map whose matrix value is
the previously constructed $\mathcal M_n^h$, retaining all entries:
\[
 T_i^h=\begin{pmatrix}z&-r_i&0&0\\0&z&-r_i&0\\0&0&z&-r_i\\0&0&0&z\end{pmatrix},
 \quad v_i^h=\begin{pmatrix}-2a_i\\b_i\\-2z\\c_i\end{pmatrix},
 \quad
 \mathcal M_n^h=\begin{pmatrix}\operatorname{diag}(T_i^h)&v^h\\-u&0\end{pmatrix},
\]
where $u=(t_1,0,0,0,\ldots,t_n,0,0,0)$ and $v^h$ stacks the
$v_i^h$. Its determinant, already proved by block multiplication, is
\[
 H_n=z^{4n-4}\sum_{i=1}^n t_i
 (c_ir_i^3-2z^2r_i^2+b_ir_iz^2-2a_iz^3).
\]
It has degree $D$. Every coordinate of $E_n$ is a nonzero scalar
multiple of an entry of $L_n$, so $L_n$ is injective.

Order the $D^2$ entries of $Y\in V_D$ row by row. Let
$p_n:V_D\to E_n$ take the first $5n+1$ entries in exactly the displayed
order, and let $i_n:E_n\to V_D$ fill those entries and put zero in
the remaining positions. The inequality $5n+1<D^2$ holds for $n\ge1$,
and $p_ni_n=I$. These are specified linear maps, rather than an
identification of two variable spaces. Set
\[
 A_n=L_np_n\in\operatorname{End}(V_D),\qquad
 \widetilde H_n=H_n\circ p_n=\det_D\circ A_n\in W_D.
\]
On coordinate rings, $p_n^*$ is injective, $i_n^*p_n^*=I$, and
$L_n^*(X_{pq})=(\mathcal M_n^h)_{pq}$. The latter is the previously
proved surjection $\vartheta_n$, with its stated ideal of linear
relations as kernel. Thus $A_n^*=p_n^*L_n^*$ retains that kernel,
and $i_n^*(\widetilde H_n)=H_n$ recovers the full determinant polynomial.

\begin{proposition}\label{prop:retained-determinant-boundary}
For every $n\ge1$, $\widetilde H_n\in\Omega_{4n+1}$, and
$\widetilde H_n\notin G_{4n+1}\cdot\det_{4n+1}$.
\end{proposition}
\begin{proof}
Let $A_n(\epsilon)=A_n+\epsilon I_{V_D}$. Its determinant as a
$D^2\times D^2$ endomorphism is a monic polynomial of degree $D^2$
in $\epsilon$. Away from its finite root set it is invertible, and
\[
 f_\epsilon(Y)=\det_D(A_n(\epsilon)Y)
             =A_n(\epsilon)^{-1}\cdot\det_D(Y)
\]
belongs to the determinant orbit. The coefficients of $f_\epsilon$
are polynomials in $\epsilon$. Every equation vanishing on that orbit,
when evaluated on $f_\epsilon$, is consequently a polynomial vanishing
for all but finitely many complex values of $\epsilon$. It is zero
identically and hence at $\epsilon=0$. This proves the claimed closure
membership, without assigning an inverse to $A_n(0)$.

For strictness, define the linear subspace
$\operatorname{Ann}(f)=\{v:\partial_v f=0\}$, using formal directional
derivatives. The chain rule gives
$\operatorname{Ann}(g\cdot f)=g\operatorname{Ann}(f)$.
For the universal determinant,
$\partial_v\det_D=\sum_{i,j}v_{ij}\operatorname{cof}_{ij}(Y)$.
The cofactors are linearly independent: each is nonzero and has row
multidegree one except in row $i$, and column multidegree one except
in column $j$. Distinct pairs $(i,j)$ have distinct combined
multidegrees, so their monomials cannot cancel. Therefore
$\operatorname{Ann}(\det_D)=0$. In contrast,
$\ker p_n\subseteq\operatorname{Ann}(\widetilde H_n)$ and
$\dim\ker p_n=D^2-(5n+1)>0$. The orbit membership is impossible.
Finally $H_n$ is nonzero: set $z=t_1=c_1=r_1=1$, $a_1=b_1=0$
and $t_i=0$ for $i>1$, obtaining $H_n=-1$.
\end{proof}

\subsection{The coordinate quotient and the obstruction direction}

Put $Z_n=\overline{G_D\cdot\widetilde H_n}$. The proposition and
$G_D$-invariance of $\Omega_D$ give $Z_n\subseteq\Omega_D$.
Both are cones by the scalar calculation above. For any such closed
cone $Z$, its ideal $I(Z)\subset B_D$ is homogeneous: if
$F=\sum_eF_e$ vanishes on $Z$, then for $z\in Z$ the polynomial
$\sum_e t^eF_e(z)$ vanishes for every $t\in\mathbb C$, so every
$F_e(z)=0$. Consequently restriction is the explicit graded map
\[
 \rho_n:\mathbb C[\Omega_D]=B_D/I(\Omega_D)
 \longrightarrow B_D/I(Z_n)=\mathbb C[Z_n],\qquad
 F+I(\Omega_D)\longmapsto F+I(Z_n).
\]
It is well defined since $I(\Omega_D)\subseteq I(Z_n)$, is surjective
in each degree by taking a polynomial representative, and has kernel
$I(Z_n)/I(\Omega_D)$. For the action on functions
$(g\cdot F)(f)=F(g^{-1}\cdot f)$, stability of both sets proves
$\rho_n(g\cdot F)=g\cdot\rho_n(F)$ by evaluation at every point.

In each degree the quotient also gives the exact multiplicity inequality
\[
 \dim\operatorname{Hom}_{G_D}(S,\mathbb C[Z_n]_e)
 \le \dim\operatorname{Hom}_{G_D}(S,\mathbb C[\Omega_D]_e)
\]
for every irreducible finite-dimensional rational $G_D$-module $S$.
Here is the splitting argument. Choose a linear section of $(\rho_n)_e$
and average $u^{-1}s u$ over the compact unitary group $U(D^2)$ with
its invariant probability measure. The averaged map is a section since
$\rho_n$ is equivariant, and is unitary-equivariant by change of integration
variable. It is $G_D$-equivariant: differentiating its intertwining
identity on skew-Hermitian matrices and taking complex linear spans
gives the identity on all matrices, and exponentiating gives it on
$\operatorname{GL}_{D^2}(\mathbb C)$. Every invertible complex matrix has a logarithm,
as follows block by block from its Jordan form and the finite logarithm
series on the nilpotent part. Composing maps $S\to\mathbb C[Z_n]_e$
with this section is an injection into the other Hom space, proving
the inequality. This proof also shows exactly why a multiplicity in the
smaller orbit closure exceeding the corresponding determinant multiplicity
would obstruct inclusion in a different instance.

The stabilizer restriction has its own exact map. For $f\in W_D$ let
$H_f=\{h\in G_D:h\cdot f=f\}$. The orbit map
$\alpha_f(g)=g\cdot f$ induces
\[
 \alpha_f^*:\mathbb C[\overline{G_D\cdot f}]
 \lhook\joinrel\longrightarrow \mathbb C[G_D]^{H_f,\mathrm{right}},
 \qquad F\longmapsto(g\mapsto F(g\cdot f)).
\]
Regularity follows from the entries of $g^{-1}$ being regular on $G_D$;
right invariance follows from $gh\cdot f=g\cdot f$; injectivity follows
because a function vanishing on the orbit vanishes on its closure.
For $e$-homogeneous $F$ and a line stabilizer with
$h\cdot f=\chi_f(h)f$, the precise identity instead is
$\alpha_f^*F(gh)=\chi_f(h)^e\alpha_f^*F(g)$.
Thus form stabilizers give invariants and line stabilizers give this
specified character. No surjectivity onto all right-invariant functions
is asserted; that would require extension across the orbit boundary.

\subsection{The later no-occurrence result at its stated scope}

For a separate pair of matrix sizes $N>m$, B\"urgisser--Ikenmeyer--Panova
use the top-left projection $\pi:V_N\to V_m$ and the form
\[
 p_{N,m}(Y)=Y_{11}^{N-m}\operatorname{per}_m(\pi Y),\qquad
 Z^{\mathrm{BIP}}_{N,m}=\overline{\operatorname{GL}(V_N)\cdot p_{N,m}}.
\]
Here $Y_{11}$ is one of the permanent variables. Their source theorem,
\cite[Theorem 1.4]{BIP2018Interface}, states that for positive integers
$N,e,m$ with $N\ge m^{25}$ and $\lambda\vdash Ne$, occurrence of
$\lambda$ in $\mathbb C[Z^{\mathrm{BIP}}_{N,m}]$ implies occurrence
in $\mathbb C[\Omega_N]$. This is their proved theorem, not a result
derived from the present fibres. We retain their same-variable padding
and exponent $25$; the independent padding above is not silently
substituted into that statement. The central scalar $tI$ acts on
$\mathbb C[W_N]_e$ by $t^{Ne}$, explaining the partition size $Ne$.
The theorem compares occurrence, and does not assert the multiplicity
inequality needed to force a quotient.

A later concrete multiplicity example has a different, specified ambient
space. In $\mathbb C[x_1,x_2,x_3]_6$ put
$\mathrm{Ch}_3^6=\{\ell_1\cdots\ell_6\}$ and
$\mathrm{Pow}_{3,k}^6=\overline{\{\ell_1^6+\cdots+\ell_k^6\}}$,
where all $\ell_i$ are homogeneous linear forms. D\"orfler--Ikenmeyer--Panova
prove, in their partition labeling, at coordinate degree $7$,
\[
 \operatorname{mult}_{(34,6,2)}\mathbb C[\mathrm{Ch}_3^6]_7
 =7<8=
 \operatorname{mult}_{(34,6,2)}\mathbb C[\mathrm{Pow}_{3,4}^6]_7.
\]
Their Main Theorem (2a) and Section 3 additionally prove that every
partition occurring in the ambient coordinate ring occurs on
$\mathrm{Ch}_3^6$, in every coordinate degree. Thus no occurrence
obstruction separates $\mathrm{Pow}_{3,k}^6$ from this Chow variety,
for any $k$, although the displayed multiplicities separate them for
$k=4$ \cite{DIP2019Interface}. For the quantifiers we use the proof's
equation labeled \texttt{eq:posplethimpliesposchow}, which retains
the same partition and degree on both sides; the theorem's final
sentence has inconsistent indices. This is a source theorem at fixed
$(m,n)=(3,6)$, not a statement in all dimensions or about permanent
orbit closures. Our finite plethysm checker does not certify the
paper's Chow ranks or its all-degree no-occurrence proof.

Our concrete family has the stronger quotient $\rho_n$ already proved,
so neither an occurrence obstruction nor a multiplicity obstruction
in the direction $Z_n\not\subseteq\Omega_{4n+1}$ can arise.
This conclusion concerns precisely the coefficient-recoverable family
$\widetilde H_n$: setting $z=1$ and taking selector coefficients still
recovers every original cubic, while forgetting the selectors or
replacing the tuple by one product has the information losses proved
earlier. The displayed maps provide an actual endpoint in GCT orbit
geometry for the retained-fibre construction. They do not identify this
family with padded permanents or supply a Boolean satisfiability reduction.

\section{Exact tensor-square transport on the full quantum source}
\label{sec:quantum-tensor-symmetry}

Throughout this section $A=\mathbb Z[q,q^{-1}]$, $K=\mathbb Q(q)$,
$D=\mathbb Q[q,q^{-1}]_{(q-1)}$, and $X_A$ is the original
$1260$-dimensional honest lattice of
Section~\ref{sec:full-source-generators}. Put $X_D=D\otimes_A X_A$.
The rational specialization is the already proved map
\[
 \sigma:X_D/(q-1)X_D\xrightarrow{\sim}
 S_{(7,5,3)}(V_{\mathbb Q}\otimes W_{\mathbb Q}).
\]
All $15$ boxes and both degree-fifteen central characters remain.
The passage to $D$ here is explicit: integral defects at the primes
$2,3,7$ are retained in the earlier integral chapters and are not
asserted to survive this localization.

The coproduct used in the original GCT-IV source is
\begin{equation}\label{eq:qt-source-coproduct}
 \Delta(K)=K\otimes K,\qquad
 \Delta(E)=E\otimes K^{-1}+1\otimes E,\qquad
 \Delta(F)=F\otimes1+K\otimes F.
\end{equation}
These are the source's equation \texttt{e U\_q coproduct}
in \cite{blasiakmulmuleysohoni2013gct4},
TeX lines 1766--1770 in the archived April 2013 edition.
They apply separately to the two commuting colours $V,W$.
The formulas also directly show coassociativity: for example both
iterated coproducts of $F$ equal
$F\otimes1\otimes1+K\otimes F\otimes1+K\otimes K\otimes F$.
Every tensor action below uses these formulas, with this orientation.

\subsection{A failure on actual source vectors}

Let $\tau(x\otimes y)=y\otimes x$ be the ordinary flip. Use the
original vector $T$ of Section~\ref{sec:gct-generator-interval}, put
$z=T$, and set $v=F_VT$. The full-source calculation proves
$E_Vz=0$, $K_Vz=q^9z$, and $E_Vv=[9]z$. Thus $v\ne0$.
The vectors $z,v$ have different weights and are independent over
$K$, and at $q=1$ they remain independent over $\mathbb Q$ since
$E_Vv=9z$ there. The coproduct gives exactly
\begin{align}
 \Delta(F_V)(z\otimes z)&=v\otimes z+q^9z\otimes v,\label{eq:qt-flip-failure}\\
 (\tau\Delta(F_V)-\Delta(F_V)\tau)(z\otimes z)
   &=(q^9-1)(v\otimes z-z\otimes v),\nonumber\\
 \Delta(E_V)(z\otimes v-v\otimes z)
   &=[9](1-q^{-9})z\otimes z.\label{eq:qt-symmetric-ideal-failure}
\end{align}
All displayed nonconstant coefficients are nonzero in $K$.
Consequently the usual symmetric subspace of $X_K^{\otimes2}$
is not a submodule: its vector $z\otimes z$ has the nonsymmetric
image in \eqref{eq:qt-flip-failure}. The usual antisymmetric
subspace is not a submodule either, by
\eqref{eq:qt-symmetric-ideal-failure}. Over characteristic zero
the symmetric-square quotient kills the antisymmetric subspace,
so that quotient does not inherit the original tensor action.
Likewise the exterior-square quotient does not inherit it.
These are statements about the specified tensor action, not about
the possible existence of other quantum symmetric or exterior powers.

At $q=1$ the two coproducts become
$E\otimes1+1\otimes E$ and $F\otimes1+1\otimes F$.
They commute with every ordinary permutation of tensor factors;
the diagonal Cartan actions do also. For the classical polynomial
group action the same assertion follows directly from
$g^{\otimes d}\tau_i=\tau_i g^{\otimes d}$.
Thus for every partition $\lambda$ of $d$, applying its row
symmetrizations and column antisymmetrizations to
$\sigma^{\otimes d}$ gives an actual classical Schur-functor
isomorphism. The inverse is induced by $(\sigma^{-1})^{\otimes d}$:
both tensor maps commute with the same permutation operators,
so their two compositions on the corresponding images are the identity.
This specifies the classical morphism and why the same construction
does not directly supply its generic quantum version.

\subsection{Explicit Clebsch--Gordan maps, with their specialization}

Write $L_m$ for the ratio-weight module with divided-power basis
$e_0,\ldots,e_m$, actions
\[
 Ke_i=q^{m-2i}e_i,\quad Fe_i=[i+1]e_{i+1},\quad
 Ee_i=[m-i+1]e_{i-1}.
\]
Terms outside the basis range are zero. Its central degree is supplied
separately by its original $\mathfrak{gl}_2$ label; no central
character is changed by this ratio-weight notation. For $L_n$ use
$f_0,\ldots,f_n$. For $0\le r\le\min(m,n)$ define
\begin{align}
 a_i(m,n,r)&=(-1)^i q^{i(n-2r+i+1)}
       \prod_{j=0}^{i-1}\frac{[n-r+j+1]}{[m-j]},\label{eq:qt-highest-coefficients}\\
 z_{mnr}&=\sum_{i=0}^r a_i(m,n,r)e_i\otimes f_{r-i},\qquad
 h_r=m+n-2r,\nonumber\\
 w_{mnr,s}&=\Delta(F)^{(s)}z_{mnr},\qquad 0\le s\le h_r.
 \label{eq:qt-clebsch-strings}
\end{align}
An empty product equals one. All denominators here are units of $D$,
since $[j](1)=j\ne0$ for $j>0$.

The coefficient of $e_i\otimes f_{r-1-i}$ in $\Delta(E)z_{mnr}$ is
\[
 a_{i+1}[m-i]q^{-n+2(r-1-i)}+a_i[n-r+i+1]=0.
\]
The equality follows by dividing the explicit consecutive coefficients
in \eqref{eq:qt-highest-coefficients}. Thus $z_{mnr}$ is a highest
vector of weight $h_r$; it is nonzero since $a_0=1$.
The string proof of Lemma~\ref{lem:gct-explicit-string-splitting}
now gives exactly
\[
 \Delta(F)w_{mnr,s}=[s+1]w_{mnr,s+1},\quad
 \Delta(E)w_{mnr,s}=[h_r-s+1]w_{mnr,s-1}.
\]
The boundary statements include $\Delta(F)w_{mnr,h_r}=0$.
The resulting strings are independent: the central Casimir
$(q-q^{-1})^2FE+qK+q^{-1}K^{-1}$ has distinct eigenvalues
$q^{h_r+1}+q^{-h_r-1}$ for the distinct nonnegative integers $h_r$,
and each individual string has distinct weights. Finally
\[
 \sum_{r=0}^{\min(m,n)}(m+n-2r+1)=(m+1)(n+1).
\]
For example, when $m\le n$ the sum is
$(m+1)(m+n+1)-m(m+1)=(m+1)(n+1)$; the other case interchanges
$m,n$. These vectors therefore give a basis over $K$.

They also give a basis over $D$, not only over its fraction field.
Specializing the coefficient recurrence at $q=1$ gives the same
highest-vector proof for the ordinary $\mathfrak{sl}_2$ action.
The ordinary Casimir $FE+H(H+2)/4$ is central (substitution of
$[H,E]=2E$, $[H,F]=-2F$, $[E,F]=H$ verifies its commutators)
and has eigenvalue $h(h+2)/4$ on a highest vector of weight $h$.
These numbers are distinct for $h\ge0$, so the specialized strings
are independent by the same argument. The defining commutator gives
$EF^sz=s(h-s+1)F^{s-1}z$, hence their lengths and nonzero
coefficients are exactly as stated. Their dimension sum is unchanged.
The square matrix of the $w_{mnr,s}$ thus has determinant with
nonzero value at $q=1$. Its determinant is a unit of the local ring
$D$, and its inverse is explicitly its adjugate divided by that
determinant. This proves both inverse compositions over $D$.

\subsection{An equivariant involution on all $1260^2$ tensor coordinates}

Set the nonzero rational constant
\[
 \gamma_{mnr}=(-1)^r\frac{n!\,(m-r)!}{(n-r)!\,m!}.
\]
Define the $D$-linear map on the preceding basis by
\begin{equation}\label{eq:qt-commutor}
 J_{m,n}:L_m\otimes L_n\longrightarrow L_n\otimes L_m,
 \qquad w_{mnr,s}\longmapsto\gamma_{mnr}w_{nmr,s}.
\end{equation}
This is a completely specified finite matrix: take the matrix of
the target strings, multiply each $r$ block by $\gamma_{mnr}$,
and multiply by the adjugate inverse of the source string matrix.
The string action formulas prove that it intertwines $E,F,K$;
the retained total central degrees agree on source and target.
Moreover $\gamma_{mnr}\gamma_{nmr}=1$, which proves
$J_{n,m}J_{m,n}=1$ on every basis vector.

At $q=1$ the ordinary flip sends $z_{mnr}$ to
$\gamma_{mnr}z_{nmr}$. Indeed it is a highest vector of the same
weight; that highest space is one-dimensional in this tensor
product, as follows from the string basis and the distinct $h_r$.
Its coefficient of $f_0\otimes e_r$ is $a_r(m,n,r)|_{q=1}$,
which is exactly $\gamma_{mnr}$. Commutation with the classical
lowering operator proves the same assertion for every descendant.
Hence $J_{m,n}|_{q=1}$ is precisely the flip.

Now use the actual source isomorphism $\Phi$ of
\eqref{eq:gct-actual-string-intertwiner}. Its explicit projectors and
inverse have denominators nonzero at $q=1$, so they extend over $D$.
Writing $H_{ab,D}=X_D\cap\ker E_V\cap\ker E_W$ in weight
$(a,15-a;b,15-b)$ gives
\[
 X_D=\bigoplus_{a,b=8}^{12}
 L_{2a-15}\otimes L_{2b-15}\otimes H_{ab,D}
\]
via this same $\Phi_D$; the original determinant twists on each
$L$ factor remain those in its full highest-weight label.
To justify base change of the highest spaces, the displayed
projector and raising inverse extract their components over $D$.
The two inverse identities proved on strings thus hold over $D$.
In particular the $H_{ab,D}$ are direct summands of finite free
$D$-modules; a submodule of a finite free module over this DVR
is free, so the same multiplicities $g_{ab}$ apply.

For a pair of summands $(a,b),(c,d)$ in $X_D\otimes X_D$, first
regroup their $V,W$, and multiplicity factors in their indicated
order. Apply
\[
 J_{2a-15,2c-15}\otimes J_{2b-15,2d-15}
 \otimes\bigl(z\otimes z'\longmapsto z'\otimes z\bigr),
\]
then undo the regrouping into the swapped pair of summands.
The regrouping commutes with the action because the two colours
act only on their own factors and the multiplicity factors have
trivial action. Conjugate the resulting direct-sum map by
$\Phi_D\otimes\Phi_D$. Denote the resulting actual full-source
operator by $J_X$.

All its entries are fixed by the formulas above and the original
generator/projector matrices. It preserves the full tensor action,
including both degree-thirty central characters, and
\[
 J_X^2=1,\qquad J_X|_{q=1}=\tau_{X_1,X_1}.
\]
The inverse identities for $\Phi_D$ and $J_{m,n}$ prove the first
statement on every direct summand; their specialization statements
prove the second. Thus
\begin{equation}\label{eq:qt-square-projectors}
 P_+=(1+J_X)/2,\qquad P_-=(1-J_X)/2
\end{equation}
are orthogonal equivariant idempotents with sum one. Their images
are finite free $D$-modules, and specialization commutes with the
direct-sum decomposition. Counting the symmetric basis
$b_i\otimes b_i$, $b_i\otimes b_j+b_j\otimes b_i$ for $i<j$,
and the antisymmetric basis for $i<j$, proves
\[
 \operatorname{rank}_D\operatorname{im}P_+=794430,
 \qquad\operatorname{rank}_D\operatorname{im}P_-=793170.
\]
Under $\sigma^{\otimes2}$ their specializations are exactly the
symmetric and exterior squares of
$S_{(7,5,3)}(V_{\mathbb Q}\otimes W_{\mathbb Q})$.
This constructs actual quantum-equivariant square modules on the
full source, not just modules with the same character.

\subsection{The degree-three coherence calculation and a braided repair}

The involution just constructed must not be assigned ordinary Young
symmetrizers in higher degree without checking the permutation
relations. Here is an exact failure inside the actual source.
The entry $g_{8,8}=1$ of
\eqref{eq:gct-full-multiplicity-matrix} supplies a summand with
$V$ and $W$ ratio weights both one. Choose a $D$-basis of its
rank-one multiplicity space. Fix the $W$ highest vector in each
of three tensor positions. On this subspace the two $W$ commutors
act as identity, so the following $V$ calculation embeds in
$X_D^{\otimes3}$ by the injective map $\Phi_D^{\otimes3}$.

In the ordered $L_1\otimes L_1$ basis
$(00,01,10,11)$, the construction gives
\begin{equation}\label{eq:qt-small-commutor}
 J=\begin{pmatrix}
 1&0&0&0\\
 0&(q^2-1)/(q^2+1)&2q/(q^2+1)&0\\
 0&2q/(q^2+1)&(1-q^2)/(q^2+1)&0\\
 0&0&0&1
 \end{pmatrix}.
\end{equation}
Indeed the weight-zero highest vector of the trivial string is
$01-q10$, with eigenvalue $-1$, and the weight-zero vector of
the length-two string is $q01+10$, with eigenvalue $+1$.
In the order $r=0,r=1$, these columns have determinant $-(1+q^2)$ and yield exactly
the displayed middle block. On the triple tensor product put
$J_{12}=J\otimes1$, $J_{23}=1\otimes J$. Direct multiplication gives
\begin{equation}\label{eq:qt-braid-defect}
 (J_{12}J_{23}J_{12}-J_{23}J_{12}J_{23})(001)
 =\frac{2(q^2-1)^2}{(q^2+1)^3}(-001+q010)\ne0.
\end{equation}
The two output basis tensors are independent and the rational
coefficient is nonzero. Therefore this specific full-source
involution does not give a symmetric-group action in degree three.
Its exact defect has order two at $q=1$, with leading term
$(q-1)^2(-001+010)$. This follows from
$(q^2-1)^2=(q-1)^2(q+1)^2$ and evaluation of the residual factors.

There is also an explicit braided repair of this test sector. Put
\[
 B=\frac{q-q^{-1}}2 I+\frac{q+q^{-1}}2 J
 =\begin{pmatrix}
 q&0&0&0\\0&q-q^{-1}&1&0\\0&1&0&0\\0&0&0&q
 \end{pmatrix}.
\]
It is an intertwiner for the $V$-coloured $L_1\otimes L_1$ action
because $J$ and $I$ are. The fixed $W$ highest slice used to embed
the test is not asserted to be stable under $W$ lowering.
Its two eigenvalues
are $q$ and $-q^{-1}$, so
$(B-qI)(B+q^{-1}I)=0$, and
$B^{-1}=B-(q-q^{-1})I$. To verify the braid relation without an
unproved matrix assertion, put $a=q-q^{-1}$. On
$(001,010,100)$ the common value of
$B_{12}B_{23}B_{12}$ and $B_{23}B_{12}B_{23}$ is respectively
\[
 q^2a\,001+qa\,010+q\,100,\qquad
 qa\,001+q\,010,\qquad q\,001.
\]
The equality in the first column uses $a^2q+a=q^2a$, obtained
by substituting $a=q-q^{-1}$. On $(011,101,110)$ the same three
formulas hold in that ordered basis. On $000$ and $111$ both
products act by $q^3$. This proves equality on all eight basis
vectors. Also $B|_{q=1}=\tau$.

Thus the actual source admits the stated equivariant square
decomposition, its proposed ordinary permutation coherence fails
by \eqref{eq:qt-braid-defect}, and a Hecke braid is constructed
on the precise sector detecting that failure. No extension of this
last $4$-by-$4$ braid to the full nonstandard Schur/plethysm
construction has been assumed. These exact maps specify the next
representation-level input needed to transport the conductor
filtration to higher quantum plethysm. No GCT obstruction or
Boolean complexity separation follows merely from the nonzero
coherence defect.

\section{A diagonal-gauge obstruction and the full signed source cover}
\label{sec:positive-source-cover}

This section concerns the original honest source basis of
$\check X_{(7,5,3)}$, not an unspecified prospective canonical basis.
It gives a pointwise obstruction for every real $q>0$, including $q=1$,
and an integral positive cover with an exact quotient to all four source
generators. The cover's nonzero quantum-relation defect is also computed.
The ordinary sign-readjustment obstruction is already stated in
GCT~IV~\cite{blasiakmulmuleysohoni2013gct4}, immediately before its
Example \texttt{ex not positive}. The additional statements here are
the arbitrary diagonal-gauge invariant, its precise effectiveness
consequence, and the full source cover and defect maps.

\subsection{Original coefficients and a pointwise invariant}

Retain $A=\mathbb Z[q,q^{-1}]$, $K=\mathbb Q(q)$ and
$[m]=\sum_{j=0}^{m-1}q^{m-1-2j}$ for positive integers $m$.
Put $[0]=0$ and $[-m]=-[m]$. The fixed basis of the honest lattice
$X_A$ is $b_0,\ldots,b_{1259}$. In the original source indexing,
\[
 T=b_{126},\qquad T_1=b_{146},\qquad T_2=b_8,\qquad Q=b_{144}.
\]
Their actual column expressions, with the original scaling, are
\begin{align*}
 T&=(123)(123)(124)(12)(12)(1)(1),\\
 T_1&=(123)(123)(124)(23)(12)(1)(1),\\
 T_2&=-[2]^{-1}(123)(123)(123)(12)(12)(4)(1),\\
 Q&=-[2]^{-1}(123)(123)(124)(13)(12)(4)(1).
\end{align*}
The quotient calculation of the original words supplies these identities
inside $X_A$; the displayed $[2]^{-1}$ is part of the representative,
not an inversion newly made in $A$. The $W$-weight of each vector is
$(9,6)$. Their $V$-weights are, in order, $(12,3),(11,4),(11,4),(10,5)$.
The printed endpoint with a $(23)$ in place of $(13)$ has $W$-weight
$(8,7)$ and is a different vector; it is not used as $Q$.

The following are entries in the full $1260$ by $1260$ matrix $F_V$,
before any four-state quotient:
\[
\begin{array}{c|c|c}
\text{source}&\text{target}&\text{coefficient}\\\hline
126&146&\alpha=[7]-[3]=q^6+q^4+q^{-4}+q^{-6}\\
126&8&\beta=[8]\\
146&144&\gamma=[6]\\
8&144&-d=-[3]
\end{array}
\]
Other full-matrix entries are retained. In particular the argument
below does not require these four vectors to form a full-group
submodule. Define
\[
 p=[2],\qquad C=q^{10}-q^4-q^{-4}+q^{-10}.
\]
Expansion of the finite Laurent sums gives
\begin{equation}\label{eq:positive-source-path-identity}
 \alpha\gamma-\beta d=pC,
 \qquad C=(q-q^{-1})^2[7][3].
\end{equation}
For example $[6]=[3](q^3+q^{-3})$ and direct multiplication gives
$(q^3+q^{-3})([7]-[3])-[8]=[2](q-q^{-1})^2[7]$;
these two identities prove the first formula with no specialization.
The second formula expands to exactly the four terms defining $C$.

Fix a real $t>0$ and evaluate $q=t$. Every term of $[m](t)$ is
positive for $m>0$, and the four displayed terms of $\alpha(t)$ are
positive. Hence the loop invariant of the two original length-two paths is
\begin{equation}\label{eq:positive-source-holonomy}
 \mathcal H(t)=\frac{\alpha(t)\gamma(t)}{-\beta(t)d(t)}
 =-1-\frac{p(t)C(t)}{\beta(t)d(t)}\leq-1.
\end{equation}
Equality holds exactly at $t=1$, since $[7](t)[3](t)>0$ and
$(t-t^{-1})^2=0$ exactly there. Thus the invariant is strictly negative
on the entire positive real axis. At $t=1$ the four coefficients are
$4,8,6,-3$, respectively, and the two products are $24,-24$.
The vanishing of $C(1)$ therefore retains the opposite signs of these
two nonzero paths.

\begin{theorem}\label{thm:positive-source-arbitrary-gauge}
For every $t>0$ there is no diagonal matrix
$D=\operatorname{diag}(d_0,\ldots,d_{1259})$ with
$d_i\in\mathbb C^\times$ such that all entries of
$D^{-1}F_V(t)D$ are nonnegative real numbers. In particular there is
no such simultaneous diagonal gauge for the four original Chevalley
generators. The assertion still holds after a permutation of the source
basis. No diagonal gauge over $\mathbb C(q)$ can make all these entries
nonnegative Laurent polynomials, or rational functions that are
nonnegative wherever defined on a nonempty positive real interval.
\end{theorem}
\begin{proof}
The entry from $i$ to $j$ is multiplied by $d_i/d_j$.
Each of the two path products from $126$ to $144$ is consequently
multiplied by the same nonzero factor $d_{126}/d_{144}$.
Their quotient remains $\mathcal H(t)<0$. If all four transformed
entries were nonnegative real numbers, they would be strictly positive
because none vanishes, so the quotient would be positive. This is a
contradiction even when the $d_i$ are complex. A permutation merely
relocates the same four entries. For a rational diagonal gauge, choose
$t$ in the proposed interval away from the finite set of poles and
zeros of its diagonal entries. The preceding pointwise contradiction
then applies. A nonzero Laurent polynomial with nonnegative
coefficients is strictly positive at every positive real $t$, which
proves the polynomial assertion as a special case.
\end{proof}

\subsection{The exact sign class, and categorical effectiveness}

For a family of real matrices with fixed nonzero support, form an
undirected multigraph: each nonzero matrix entry gives a labelled edge
between its source and target, retaining its matrix colour and sign.
Loops and parallel edges are retained. Let the edge sign be
$\sigma_e\in\{1,-1\}$. A traversal in either direction uses the same
sign. A closed walk has sign equal to the product of these edge signs.
In additive notation $\sigma_e=(-1)^{\epsilon_e}$, the assignment
$\epsilon$ is a $\mathbb Z/2$-valued graph one-cochain. Since a graph
has no two-cells, its class modulo vertex coboundaries is an element
of $H^1(\Gamma;\mathbb Z/2)$, whose definition here is precisely that
quotient of edge assignments.

\begin{lemma}\label{lem:positive-source-switching}
All entries of a family can be made nonnegative by a common diagonal
complex gauge if and only if every closed walk of its sign graph has
positive sign. When this holds, a diagonal gauge with entries in
$\{1,-1\}$ already suffices.
\end{lemma}
\begin{proof}
Write a putative complex diagonal entry as $d_i=r_i z_i$, where
$r_i>0$ and $|z_i|=1$. A real nonzero coefficient of sign $\sigma_e$
from $i$ to $j$ becomes positive only if
$\sigma_e z_i/z_j=1$, or $z_j=\sigma_e z_i$.
Multiplying these identities around a closed walk gives product sign
one. Conversely, in each connected component choose a root and set
its sign to one. Define the sign of a vertex to be the product of edge
signs along any root-to-vertex walk. Two choices have quotient the
sign of a closed walk, so the value is independent of the walk. The
endpoint signs then satisfy $s_j=\sigma_e s_i$ on every edge.
The diagonal entries $s_i$ make each coefficient positive. This also
proves that the obstruction is exactly $[\epsilon]$, not merely a
necessary condition extracted from one failed gauge.
\end{proof}

All $8162$ nonzero entries of the four original matrices are
coefficientwise sign coherent, in agreement with GCT~IV's theorem
\texttt{t positivity}. Thus their signs do not vary with $t>0$.
The exact census in the original basis is
\[
\begin{array}{c|r|r|r}
 &\text{positive entries}&\text{negative entries}&\text{total}\\\hline
 F_V&1648&350&1998\\
 E_V&1837&246&2083\\
 F_W&1848&150&1998\\
 E_W&1963&120&2083
\end{array}
\]
The graph is connected. Its $8162$ labelled edges give cycle rank
$8162-1260+1=6903$; it has $6146$ unordered vertex pairs after
parallel edges are identified. These finite statements are certified
by a spanning traversal of every original matrix entry, retained in
the accompanying independent graph certificate.

The negative cycle displayed in
\eqref{eq:positive-source-holonomy} has minimum possible length.
Indeed every edge changes exactly one of the first weight coordinates
by one. It therefore reverses the parity of their sum, so there are no
loops or odd cycles. Direct comparison of all parallel entries gives
the same sign for every edge on a fixed unordered pair, so a two-edge
closed walk has positive sign. Our four-edge walk is negative, proving
minimality. An additional two-colour witness uses the actual entries
\[
 (F_V)_{20,0}=[7],\quad (F_W)_{20,10}=[5],\quad
 (F_W)_{136,10}=[3],\quad (F_V)_{136,0}=-[3].
\]
The oriented loop $0\to20\leftarrow10\to136\leftarrow0$ has
coefficient quotient $-[7]/[5]<0$ for $t>0$, and value $-7/5$ at
$t=1$. Cancellation in a divided-power endpoint is not needed for
this second witness.

For an exact formulation of effectiveness, call a labelled graded
category effective when its object classes are finite sums
$\sum_{i,w}n_{i,w}q^w[B_i]$ with $n_{i,w}\in\mathbb Z_{\geq0}$,
and these labels and shifts freely generate its split Grothendieck
group. This includes finite direct sums of labelled graded free
groups, and a graded semisimple category with the stated simples.
An additive endofunctor compatible with the internal grading shifts
and taking objects to objects then has an $A$-linear matrix
over $\mathbb Z_{\geq0}[q,q^{-1}]$: its $i$-th column is the
decomposition of the image of $B_i$.

\begin{corollary}\label{cor:positive-source-effective}
There is no such effective category and grading-shift-compatible
additive source-generator
functor whose induced $F_V$ is identified with the original $F_V$
by $[B_i]\mapsto d_i(q)b_i$ for invertible diagonal rational factors.
This remains impossible after arbitrary internal shifts, parity
changes of individual source labels, or permutations of those labels,
provided the images of objects are required to remain effective in
the resulting labels. At a fixed $t>0$ the coordinate simplicial cone
generated by any individually rescaled source basis is not preserved
by all the generator images.
\end{corollary}
\begin{proof}
An effective matrix evaluates to a nonnegative matrix at each positive
real $t$. Apply Theorem~\ref{thm:positive-source-arbitrary-gauge} at
a nonsingular value. Internal shift and individual parity change
give diagonal factors $q^{w_i}$ and $(-1)^{a_i}$, covered by the
same theorem. The cone assertion is the real matrix statement:
a linear map preserves the coordinate nonnegative cone exactly when
every one of its columns lies in that cone.
\end{proof}

This is an effectiveness obstruction with an explicitly stated
basis constraint. It is not a statement that a Frobenius trace must
be positive: an alternating trace retains its cohomological signs.
Nor does it contradict the nonstandard positivity formulations that
explicitly retain signs and powers of $q-q^{-1}$.

\subsection{An integral positive cover of every original generator}

For each of the four full source matrices $M$, split its entries
coefficientwise as $M=M^+-M^-$ with
$M^+,M^-\in\operatorname{Mat}_{1260}(\mathbb Z_{\geq0}[q,q^{-1}])$
and disjoint coefficient supports. Define
\begin{equation}\label{eq:positive-source-cover-matrices}
 \widetilde X_A=X_A^+\oplus X_A^-,\qquad
 \widetilde M=\begin{pmatrix}M^+&M^-\\M^-&M^+\end{pmatrix},
 \qquad |M|=M^++M^-.
\end{equation}
Here $b_i^+,b_i^-$ are new labelled copies of every original vector;
the signs in their names specify the projection below. No change is
made to the underlying original $b_i$. The formula uses all the
original entries and all their exponents. Since they are sign coherent,
an original edge of positive sign stays on its sheet and an edge of
negative sign switches sheets, always with its original absolute
Laurent coefficient.

Define $A$-linear maps
\[
 \iota:X_A\longrightarrow\widetilde X_A,\quad
       x\longmapsto(x,x),\qquad
 \pi:\widetilde X_A\longrightarrow X_A,\quad
       (x,y)\longmapsto x-y.
\]
The following exact diagram holds for each of the four generators:
\begin{equation}\label{eq:positive-source-cover-exact}
 0\longrightarrow (A^{1260},|M|)
 \xrightarrow{\ \iota\ }(A^{2520},\widetilde M)
 \xrightarrow{\ \pi\ }(X_A,M)\longrightarrow0.
\end{equation}
To prove this, $\iota$ is injective, $\pi(x,0)=x$, and
$\pi(x,y)=0$ exactly when $x=y$. Direct block multiplication gives
\[
 \pi\widetilde M=M\pi,\qquad
 \widetilde M\iota=\iota|M|.
\]
The underlying sequence of free $A$-modules is split by
$s(x)=(x,0)$, with no division by two; the section is not asserted
to commute with the generators. Thus there is neither an unmentioned
$2$-inversion nor an integral torsion kernel.

Give both copies of $b_i$ its original four weight coordinates.
Every original Cartan diagonal matrix has Laurent monomial entries,
so it lifts to the same diagonal matrix on both sheets. Its inverse
lifts in the same way. The two original total degree-fifteen central
characters remain $q^{15}$. The displayed inclusion and projection
commute with these diagonal operators. The Cartan conjugation rules
hold on the cover because each nonzero lifted edge has exactly the
original source and target weights. In particular, the sign cover
does not replace a raising edge by a lowering edge.

To describe its exact algebraic scope, let $\mathcal T$ be the free
associative $A$-algebra on the four Chevalley symbols, the original
Cartan symbols and their named inverses, and the integral Cartan
operators $H_V,H_W$ defined below. On the two sheets and on the
kernel these last operators act by the same original diagonal
matrices. Acting by the specified
matrices makes all three terms of
\eqref{eq:positive-source-cover-exact} actual $\mathcal T$-modules,
and the displayed maps are $\mathcal T$-linear. For any polynomial
relation $r\in\mathcal T$ whose original source matrix is zero,
\[
 \pi r(\widetilde M)=0.
\]
Consequently there is a unique, explicitly defined $A$-linear defect
map $\Delta_r:\widetilde X_A\to A^{1260}$ with
\begin{equation}\label{eq:positive-source-relation-defect}
 r(\widetilde M)=\iota\Delta_r,
 \qquad \Delta_r=\operatorname{pr}_+r(\widetilde M).
\end{equation}
Existence follows from the equality of the plus and minus output
coordinates, and uniqueness from injectivity of $\iota$. If
$r(\widetilde M)\iota$ is evaluated instead, the same identities give
$\Delta_r\iota=r(M_{\ker})$, where $M_{\ker}$ denotes the four
absolute Chevalley matrices and the original Cartan diagonals.
Thus the complete relation loss is an actual
map into the displayed kernel. The quotient by that kernel recovers
all original source relations.

This doubled graph is the covering determined by the sign class.
It is connected: paths from a root lift to one sheet at any target;
a negative closed walk switches the root sheet, and then the same
paths reach the other sheet at every target. A vertex sign on the
cover, positive on the plus sheet and negative on the minus sheet,
trivializes the pullback of the sign class. Conversely a one-sheet
cover cannot trivialize its nonzero class. A finite graph covering
of a connected graph has a constant number of vertices in each
fibre: lift any path and its inverse to obtain bijections of fibres.
Therefore this degree-two, $2520$-vertex cover has the least possible
degree among nonempty graph covers that trivialize this sign class. This
minimality statement is about such graph covers, not arbitrary
changes of representation basis or unrestricted categorical models.

\subsection{A nonzero full-source quantum-relation defect}

Let $H_V$ be the original diagonal operator with entry
$[w_{V,1}(b_i)-w_{V,2}(b_i)]$ on $b_i$; equivalently over $K$ it is
$(K_V-K_V^{-1})/(q-q^{-1})$. Its integral diagonal formula, with its
signs, defines the operator already over $A$. Let
$\widetilde H_V=\operatorname{diag}(H_V,H_V)$, the actual same
Cartan expression on both sheets.

The exact source matrices give
\begin{equation}\label{eq:positive-source-nonzero-defect}
 (\widetilde E_V\widetilde F_V-
   \widetilde F_V\widetilde E_V-\widetilde H_V)b_{132}^+
       =[2][3](b_{148}^++b_{148}^-).
\end{equation}
Here the original source labels are
\[
 b_{132}=(123)(123)(124)(12)(12)(3)(3),\qquad
 b_{148}=(123)(123)(124)(12)(23)(1)(3).
\]
Both have $V$-weight $(10,5)$ and $W$-weight $(9,6)$.
For completeness, all paths from $132$ to $148$ contributing to
the two compositions are the following; coefficients are recorded
in traversal order:
\[
\begin{array}{c|r|c|c}
\text{composition}&\text{intermediate index}&\text{first coefficient}
                                         &\text{second coefficient}\\\hline
 E_VF_V&152&-[3]&[2]\\
 E_VF_V&158&[7]&[2]\\
 F_VE_V&128&[2]&[7]-[3]
\end{array}
\]
On the original source their sum is
$-[3][2]+[7][2]-[2]([7]-[3])=0$.
On the cover the first path ends on the minus sheet and the second
on the plus sheet. The subtracted third path ends on the plus sheet.
Each output sheet therefore has coefficient $[2][3]$. The full
column calculation verifies that all other target entries in the
commutator defect vanish, proving
\eqref{eq:positive-source-nonzero-defect}. The diagonal term cannot
affect target $148\ne132$. The original full matrices, including
the same diagonal, satisfy the zero commutator relation in every
column; thus this is precisely a failure upstairs, with zero image
under $\pi$. Its coefficient expands to
\[
 [2][3]=q^3+2q+2q^{-1}+q^{-3},
\]
which is nonzero over $A$ and strictly positive at every $t>0$.
The $W$-coloured calculation gives
\[
 (\widetilde E_W\widetilde F_W-
   \widetilde F_W\widetilde E_W-\widetilde H_W)b_{256}^+
       =[2][3](b_{263}^++b_{263}^-).
\]
Each colour has exactly $52$ even-sheet input columns with nonzero
commutator defect. These claims use all target coordinates, not a
principal submatrix. The independent Laurent checker verifies them
column by column. It follows that the cover in
\eqref{eq:positive-source-cover-matrices} is not a full quantum-group
representation, although its exact quotient is the original one.

The failure already includes the commuting lowering generators:
\begin{equation}\label{eq:positive-source-mixed-defect}
 [\widetilde F_V,\widetilde F_W]b_0^+
       =-[3]^2(b_{146}^++b_{146}^-),
 \qquad b_0=(123)(123)(123)(12)(12)(1)(1).
\end{equation}
All contributing original paths from $0$ to $146$ are
\[
\begin{array}{c|r|c|c}
\text{composition}&\text{intermediate index}&\text{first coefficient}
                                         &\text{second coefficient}\\\hline
 F_VF_W&126&[3]&[7]-[3]\\
 F_WF_V&20&[7]&[3]\\
 F_WF_V&136&-[3]&[3]
\end{array}
\]
The original difference is $[3]([7]-[3])-[7][3]+[3]^2=0$.
On the cover the third path ends on the minus sheet and is
subtracted with its positive absolute coefficient. The plus-sheet
difference is also $-[3]^2$. Every other target entry cancels,
as the full original-path enumeration verifies. In total the
numbers of nonzero even-sheet input columns for the six
commutator defects are
\[
\begin{array}{c|rrrrrr}
\text{commutator}&[E_V,F_V]&[E_W,F_W]&[E_V,E_W]&[E_V,F_W]
                  &[F_V,E_W]&[F_V,F_W]\\\hline
\text{columns}&52&52&68&52&52&50
\end{array}
\]
For the first two columns the stated Cartan diagonal is subtracted;
for the four mixed columns the original commutator is zero.
The independent checker verifies all six original identities and
all their exact diagonal-kernel defects over $A$.

\subsection{An actual categorical map and the retained cone}

Let $\mathcal G$ be the category of finite free internally graded
abelian groups with finite support. A Laurent polynomial
$a(q)=\sum_w a_wq^w$ with nonnegative integer coefficients defines
the actual group $L_a=\bigoplus_w\mathbb Z^{a_w}\langle w\rangle$.
Use the ordered basis labels $(w,1),\ldots,(w,a_w)$ and set $L_0=0$.
Define an additive functor on $\mathcal G^{2520}$ by
\[
 (\widetilde{\mathsf M}V)_i
       =\bigoplus_j L_{\widetilde M_{ij}}\otimes_{\mathbb Z}V_j.
\]
On morphisms use the identity on each first tensor factor. Its
Grothendieck matrix is exactly $\widetilde M$, since tensor products
add degrees and multiply ranks. This formula realizes every one of
the positive-cover coefficients as actual objects, over $\mathbb Z$.

Let $\mathcal S$ be internally graded finite free super abelian
groups, with even morphisms, and let $\Pi$ exchange parity.
It satisfies $\Pi^2=1$ on objects and maps. Its super Grothendieck
group is defined by direct-sum relations and $[\Pi V]=-[V]$;
the signed rank character identifies that group with $A$.
The parity-summing functor
\[
 \mathsf Q:\mathcal G^{2520}\to\mathcal S^{1260},\qquad
 (V_i^+,V_i^-)_i\longmapsto(V_i^+\oplus\Pi V_i^-)_i
\]
has induced map $\pi$, with every plus component placed in even
parity before taking the direct sum. Define
\[
 (\mathsf M W)_i=\bigoplus_j
   \left(L_{M^+_{ij}}\otimes W_j\right)
   \oplus\left(L_{M^-_{ij}}\otimes\Pi W_j\right).
\]
There is a natural parity-preserving isomorphism
\begin{equation}\label{eq:positive-source-super-map}
 \mathsf Q\widetilde{\mathsf M}\simeq\mathsf M\mathsf Q.
\end{equation}
To verify it, expand the two sides. The even output summands on
each side are $L_{M^+_{ij}}\otimes V_j^+$ and
$L_{M^-_{ij}}\otimes V_j^-$; the odd output summands are
$L_{M^+_{ij}}\otimes V_j^-$ and
$L_{M^-_{ij}}\otimes V_j^+$. Match identical labelled tensor
factors by the identity. This specifies every component, its inverse,
and naturality. Thus the signed source action has an actual integral
super-object correspondence with the full effective cover, not only
an equality of matrix dimensions.

The quantum relations have not been imposed as functor isomorphisms
by this definition. The nonzero defect just computed precludes that
interpretation for the positive cover. The full generators on
super Grothendieck groups do satisfy the source relations, through
\eqref{eq:positive-source-super-map} and the proven matrix action.
No assertion of a canonical-basis convolution model follows merely
from the existence of these additive functors.

There is also a precise relation to the previously constructed
four-state Frobenius cone. Let $\mathsf A,\mathsf B,\mathsf G,\mathsf J$
be the graded groups with respective exponent sets
\[
\begin{split}
 &\{-6,-4,4,6\},\qquad\{-7,-5,-3,-1,1,3,5,7\},\\
 &\{-5,-3,-1,1,3,5\},\qquad\{-2,0,2\}.
\end{split}
\]
The two lifted paths from $T^+$ end in $Q^+$ with
$\mathsf P=\mathsf G\otimes\mathsf A$ and in $Q^-$ with
$\mathsf R=\mathsf J\otimes\mathsf B$.
Applying $\mathsf Q$ gives $\mathsf P\oplus\Pi\mathsf R$,
with character $\alpha\gamma-d\beta=pC$ and with the intermediate
labels $T_1,T_2$ still attached. The actual integral map
$\theta:\mathsf R\to\mathsf P$ used in the Frobenius cone matches,
in each weight, equal-position basis pairs in lexicographic order.
Its twenty unit entries cancel twenty matched pairs; the cone has
four residual even and four residual odd basis vectors. This
cross-path differential is additional map data, not the differential
of the split path object. The cover and the parity map therefore
recover exactly its two branch objects before that differential is
installed. More explicitly, place an arbitrary input $V$ at $T^+$
and project the square functor to the two labelled output sheets
at $Q$. The extracted endpoint functors are
$V\mapsto\mathsf P\otimes V$ and
$V\mapsto\mathsf R\otimes V$. The map
$\theta\otimes1_V:\mathsf R\otimes V\to\mathsf P\otimes V$
is natural, and its cone is exactly the specified endpoint cone
functor. Its underlying graded parity object is the image under
$\mathsf Q$ of the two extracted paths. The identity of those
graded objects is not claimed to be a chain map from the split
zero-differential path complex to the cone.

\subsection{What changes when the basis is allowed to mix}

For each fixed $t>0$ an effective cone for the four full generators
does exist after an explicit type of nondiagonal source change.
Here the exact source commutator is read as the integral formula
$[E,F]|_{X_h}=[h](t)$, including its value at $t=1$.
This assertion can be proved without assuming a positive basis in
the original labels.

Consider one colour on any of the present finite-dimensional
specializations. If $z$ is highest of weight $n$, the commutator
gives by induction
\[
 E F^jz=[j](t)[n-j+1](t)F^{j-1}z.
\]
The induction uses the finite-sum identity
$[j+1][n-j]=[j][n-j+1]+[n-2j]$, verified by multiplying the
Laurent sums, so it also holds at $t=1$.
Let $k$ be the largest integer with $F^kz\ne0$. Applying the
formula to $F^{k+1}z=0$ gives
$[k+1](t)[n-k](t)F^kz=0$. For integral $h$, $[h](t)=0$ exactly
when $h=0$. Thus $k=n\geq0$. It follows that the divided string
$z_i=F^iz/[i]!(t)$, $0\leq i\leq n$, has all its stated vectors
nonzero and satisfies
\[
 Fz_i=[i+1](t)z_{i+1},\qquad
 Ez_i=[n-i+1](t)z_{i-1},
\]
where the out-of-range terms vanish. All nonzero coefficients are
strictly positive.

To prove these strings span the entire representation, take its
largest occurring weight $n$ and a basis of its highest space.
Their strings are independent: different levels have different
weights, and applying the corresponding power of $E$ to a relation
at one level reduces it to a relation in that highest basis with a
nonzero common coefficient. Their sum $N$ is a submodule. Induct
on dimension in the quotient, whose largest weight is smaller.
For a quotient highest vector of weight $m$, lift it to $z$ of
weight $m$. Then $Ez\in N_{m+2}$. On every top string,
$E:N_m\to N_{m+2}$ is surjective when the target occurs, because
the coefficient just displayed is nonzero; the quotient string
has $m\geq0$, so there is no endpoint exception. Subtract a
preimage to make the lift highest. Its length-$m$ string maps
isomorphically to the selected quotient string. These lifted
strings are independent modulo $N$ and split off the quotient.
This completes the induction.

Apply the construction first to the $V$ generators. The commuting
$W$ generators preserve each $V$ highest subspace, where the same
argument applies. If $H_{n,m}(t)$ is the simultaneous highest
subspace of weight ratios $(n,m)$ and a real basis of it is chosen,
define $L_n(t)$ to have the divided-string basis $e_i$, $0\leq i\leq n$,
the displayed $E,F$ action, and original $\mathfrak{gl}_2$ weights
\[
 \left(\frac{15+n}{2}-i,\frac{15-n}{2}+i\right).
\]
The occurring $n$ are odd, since all original total weights are
fifteen, so these coordinates are integers. Define $L_m(t)$ in
the same way for $W$. Then the map
\[
 \Phi_t:\bigoplus_{n,m} L_n(t)\otimes L_m(t)\otimes H_{n,m}(t)
       \longrightarrow X_A\otimes_{q\mapsto t}\mathbb R,
 \quad e_i\otimes f_j\otimes z\longmapsto
              F_V^{(i)}F_W^{(j)}z
\]
is therefore an isomorphism. The multiplicity factor has trivial
quantum-group action in the domain. The formulas above prove that
the map intertwines all raising and lowering generators; the
original weight coordinates prove the Cartan and degree-fifteen
central characters as well. This is the specialization of the
actual string construction, now justified at every positive $t$.

The nonnegative span of these string vectors is a simplicial cone
preserved by all four generators and their Cartan diagonals: their
matrices have nonnegative entries in these strings, and each Cartan
eigenvalue is a positive power of $t$. Its matrix $\Phi_t$ is
obtained from actual simultaneous highest kernels and the original
divided powers, so it is an exact correspondence with the source.
Theorem~\ref{thm:positive-source-arbitrary-gauge} proves that
$\Phi_t$ cannot be a diagonal matrix times a permutation. Thus the
fixed-source obstruction and this effective mixed-basis cone are
related by an explicit proved intertwiner. The construction depends
on choices in the highest multiplicity spaces and is not asserted
to be the source-compatible canonical geometric realization sought
in the nonstandard Riemann-hypothesis programme.

\paragraph{Reproduction and scope.}
The complete original matrices and basis, with their original
coefficient and word data, are the frozen inputs in
\texttt{agents/full\_source\_generators}. The accompanying
\texttt{check\_positive\_cover.py} verifies their hashes, every edge
sign, all four retained coefficients, both Laurent identities, every
column of the inclusion and projection intertwiners, the two full
original commutator identities, and every cover commutator defect.
The independent sign audit checks graph connectivity, the cycle
class, the shortest-cycle assertion and the same retained entries.
The primary-source loci are GCT~IV TeX lines $7240$--$7323$
for sign coherence, $7324$--$7326$ for the already stated
sign-readjustment obstruction, and $7328$--$7350$ for the printed
example and its two paths; its frozen source hash is recorded in
the machine receipt. The new obstruction is to a stated effective
realization preserving the original basis up to diagonal gauge.
The exact cover and mixed-string maps do not prove a nonstandard
Riemann hypothesis, an asymptotic class-variety obstruction, or a
Boolean or algebraic complexity-class separation.

\section{An integral tableau basis for the exact plethystic maps}
\label{sec:plethysm-integral-tableau-basis}

All constructions in this section are over an arbitrary commutative ring
$R$ with identity. In particular, neither an integer nor a determinant
is inverted. A partition is a finite sequence
$\lambda=(\lambda_1\geq\cdots\geq\lambda_\ell>0)$.
Write $d=|\lambda|$, $t=\lambda_1$, and
$h_j=\#\{r:\lambda_r\geq j\}$ for its column lengths.
The empty partition is allowed, with $d=t=\ell=0$.

\subsection{The intrinsic image and its polynomial realization}

For a finite free $R$-module $E$ define
\begin{align}
 C_\lambda(E)&=\bigotimes_{j=1}^{t}\bigwedge^{h_j}E,
 &P_\lambda(E)&=\bigotimes_{r=1}^{\ell}\operatorname{Sym}^{\lambda_r}E.
 \label{eq:plethysm-column-row-modules}
\end{align}
An empty tensor product in this section is $R$.
For each column, send an exterior product to the alternating tensor
\[
 v_1\wedge\cdots\wedge v_h\longmapsto
 \sum_{\sigma\in\mathfrak S_h}\operatorname{sgn}(\sigma)
 v_{\sigma(1)}\otimes\cdots\otimes v_{\sigma(h)}.
\]
This formula is multilinear and alternating over $R$: when two inputs
are equal, transposing their labels pairs equal tensors with opposite
coefficients. It therefore defines a map out of the exterior power,
also in characteristic two. Tensor these maps in the column order
$1,\ldots,t$, relabel their tensor positions by the boxes $(r,j)$,
reorder the positions into row order, and take the symmetric quotient
in each row. Denote the resulting map by
\begin{equation}
 \delta_{\lambda,E}:C_\lambda(E)\longrightarrow P_\lambda(E),
 \qquad
 \mathsf S_\lambda(E)=\operatorname{im}\delta_{\lambda,E}.
 \label{eq:plethysm-schur-image}
\end{equation}
Equation~\eqref{eq:plethysm-schur-image} is the Schur module convention
used throughout this section. There is no additional row
symmetrizer, factorial, or Young-idempotent normalization.

Let $e_1,\ldots,e_n$ be an ordered basis of $E$.
The module $\operatorname{Sym}^a(E)$ has the basis
$e_1^{\alpha_1}\cdots e_n^{\alpha_n}$ with
$\sum_i\alpha_i=a$. Indeed, the tensor-word basis is partitioned into
permutation orbits, and the defining symmetry relations identify the
words in each orbit; the free module on these orbits gives an inverse
to this quotient description. The module $\bigwedge^a E$ has the basis
$e_{i_1}\wedge\cdots\wedge e_{i_a}$ for
$i_1<\cdots<i_a$. To see this without an assumption on $R$, sort each
word with its permutation sign and send a word with repeated indices
to zero. This alternating multilinear rule gives an inverse to the
map from the free module on the displayed increasing words.

Introduce independent variables $X_{r,i}$ for
$1\leq r\leq\ell$ and $1\leq i\leq n$.
The row monomial bases identify $P_\lambda(E)$ with the component of
$R[X_{r,i}]$ having degree $\lambda_r$ in the variables with first
index $r$. For an ordered list $I=(i_1,\ldots,i_a)$, where
$a\leq\ell$, put
\[
 [I]_a=\det(X_{r,i_s})_{1\leq r,s\leq a},
 \qquad [\varnothing]_0=1.
\]
Repeated entries give zero, and permuting the list multiplies this
polynomial by the permutation sign.
For a filling $T$ of the diagram of $\lambda$, write $I_j(T)$ for its
$j$th column, read from top to bottom, and set
\begin{equation}
 b_T=\prod_{j=1}^{t}[I_j(T)]_{h_j}.
 \label{eq:plethysm-flag-polynomial}
\end{equation}
Expanding the alternating tensors in the definition of $\delta$
shows term by term that $b_T$ is the image of the column exterior
tensor indexed by $T$. Consequently the $b_T$ with strictly increasing
columns generate $\mathsf S_\lambda(E)$. A filling is called
\emph{semistandard} when its columns strictly increase and its rows
weakly increase, using the alphabet $1,\ldots,n$.

\subsection{An alternating identity and an integral straightening proof}

The following identity is proved in the polynomial ring, so it
remains valid after every specialization of its coefficients.

\begin{lemma}[The tail--head shuffle identity]
\label{lem:plethysm-shuffle}
Let $\ell\geq a\geq b\geq k\geq1$ and let $m=a-k+1$.
Let $A$ be a list of column indices of length $k-1$,
let $U=(u_1,\ldots,u_{a+1})$,
and let $B$ be a list of length $b-k$. For a subset
$S=\{s_1<\cdots<s_m\}\subseteq\{1,\ldots,a+1\}$, let $U_S$
be the sublist in the original position order; the complementary
sublist has the same convention. Set
\[
 \epsilon(S)=(-1)^{\sum_{v=1}^{m}(s_v-v)}.
\]
Then
\begin{equation}
 \sum_{\substack{S\subseteq\{1,\ldots,a+1\}\\|S|=m}}
 \epsilon(S)[A,U_S]_a[U_{S^c},B]_b=0.
 \label{eq:plethysm-shuffle}
\end{equation}
Here the minors use initial row segments of the same matrix.
\end{lemma}

\begin{proof}
First let every entry of the lists denote a vector of $R^a$, and
let the second determinant use its projection onto the first $b$
coordinates. Keep $A$ and $B$ fixed. The left side is a multilinear
function of the $a+1$ vectors of $U$. It is alternating, as follows.
Suppose two of these vectors have equal values. A term putting both
in the first determinant, or both in the second determinant, is
zero. Pair each remaining subset with the subset obtained by
exchanging those two positions between the subset and its complement.
The two determinant products in such a pair agree up to the signs
needed to restore the position orders inside the two sublists.
The shuffle sign is the sign of the permutation that lists $S$
first and $S^c$ second. Including the two internal reordering signs,
the exchange is a single transposition of the full list. Thus the
two paired terms have opposite coefficients and cancel. This proves
vanishing on equal inputs by an explicit pairing, without deducing
it from skew symmetry or dividing by two.

Expand each of the $a+1$ vectors in the $a$ standard basis vectors
of $R^a$. Multilinearity expresses the value as a sum of values on
standard basis vectors. Each of these tuples repeats a vector, by
the pigeonhole principle, so each value is zero by the preceding
argument. The multilinear function therefore vanishes identically.
Apply this result over the polynomial ring to the column vectors
$(X_{1,i},\ldots,X_{a,i})^{\mathsf t}$. Its projection onto the first
$b$ coordinates is precisely the column occurring in the second
flag minor. This gives~\eqref{eq:plethysm-shuffle} with its stated
signs and row sets.
\end{proof}

\begin{theorem}[Integral semistandard basis]
\label{thm:plethysm-integral-basis}
For every finite free $R$-module $E$ with its specified ordered basis,
the elements $b_T$ indexed by semistandard tableaux of shape $\lambda$
form an $R$-basis of $\mathsf S_\lambda(E)$.
The image in~\eqref{eq:plethysm-schur-image} is a direct summand as an
$R$-module of $P_\lambda(E)$; the splitting need not be natural.
For every ring homomorphism $R\to R'$, the natural base-change map
\begin{equation}
 R'\otimes_R\mathsf S_\lambda(E)
 \longrightarrow
 \mathsf S_\lambda(R'\otimes_R E)
 \label{eq:plethysm-basechange}
\end{equation}
is an isomorphism. No flatness hypothesis is needed.
\end{theorem}

\begin{proof}
If $\lambda$ is empty, all three claims concern the identity map on
$R$ and follow directly. If $h_1>n$, the first exterior factor is
zero, and there are no strictly increasing first columns. Both the
image and the claimed basis are empty, giving the zero module.
We may therefore suppose that $\lambda$ is nonempty and $h_1\leq n$.

First prove spanning. Order all strictly column-increasing fillings
by lexicographic order of the word obtained by reading their columns
from left to right, and each column from top to bottom. This is a
finite totally ordered set. A filling $T$ that is not semistandard
has adjacent columns
\[
 I=(i_1<\cdots<i_a),\qquad J=(j_1<\cdots<j_b),
 \qquad a\geq b,
\]
and a row $k\leq b$ at which $i_k>j_k$.
Apply Lemma~\ref{lem:plethysm-shuffle} with
\[
 A=(i_1,\ldots,i_{k-1}),\quad
 U=(i_k,\ldots,i_a,j_1,\ldots,j_k),\quad
 B=(j_{k+1},\ldots,j_b).
\]
The subset $S_0=\{1,\ldots,a-k+1\}$ has sign $+1$ and gives
exactly $[I]_a[J]_b$. Move all the other summands to the other side
of~\eqref{eq:plethysm-shuffle}, and multiply by the unchanged minors
in the remaining columns of $T$.

Every member of $(i_k,\ldots,i_a)$ is larger than every member of
$(j_1,\ldots,j_k)$. Thus the entries of $U$ are distinct. A subset
$S\ne S_0$ replaces at least one entry of the left-column tail by
a strictly smaller entry from the right-column head. If a resulting
column has a repeated entry, its determinant is zero and the term
is discarded. In any other term, sort the two columns numerically
and multiply by the two sorting signs. The fixed left-column prefix
is unchanged as a multiset; all inserted values are smaller than
all removed values. At the first place where the sorted old and
new left columns differ, the new column therefore has the smaller
entry. All preceding columns of the entire filling are unchanged.
The resulting filling is strictly smaller than $T$ in the specified
total order. The coefficient of the original term was exactly one,
so this replacement uses only integer addition and subtraction.
Induction on the finite total order proves that every generator
$b_T$ is an integral linear combination of the semistandard $b_T$.
This proof works after passage to $R$ as well.

For independence, order the variables by
\[
 X_{1,1}>X_{1,2}>\cdots>X_{1,n}>
 X_{2,1}>\cdots>X_{2,n}>\cdots>X_{\ell,n},
\]
and use lexicographic monomial order. If
$i_1<\cdots<i_a$, the leading monomial of $[I]_a$ is
$X_{1,i_1}\cdots X_{a,i_a}$, with coefficient $+1$.
Indeed, among determinant permutations the largest monomial must
choose the smallest available column index for its first row, then
the smallest remaining index for its second row, and so on.
Monomial order is preserved by multiplication, so
\begin{equation}
 \operatorname{in}(b_T)=
 \prod_{(r,j)\in\lambda}X_{r,T(r,j)},
 \qquad
 \text{with coefficient }+1.
 \label{eq:plethysm-leading-monomial}
\end{equation}
For a semistandard filling, this monomial determines each row's
multiset of entries. A weakly increasing row is uniquely recovered
from that multiset. Hence distinct semistandard tableaux have
distinct leading monomials. In a nonzero linear combination over
$R$, choose the largest leading monomial among the tableaux whose
coefficients are nonzero. No other summand has a monomial that
large. Its coefficient in the sum is the chosen nonzero coefficient
times $1$, and cannot vanish. This proves independence even when
$R$ has zero divisors.

For completeness, this monicity also produces a splitting. Work
first over $\mathbb Z$ and list the semistandard polynomials in
decreasing order of their distinct leading monomials. Given any
polynomial in the finite row-degree component $P_\lambda(\mathbb Z^n)$,
process this list in order. At each stage, extract the coefficient
of the current leading monomial, record it, and subtract that
coefficient times the corresponding $b_T$. Later subtractions
cannot change a leading monomial already processed, since all
their monomials are smaller. Return the sum of the recorded
multiples of $b_T$. This is a $\mathbb Z$-linear map into their
span, and it is the identity on each $b_T$. It is therefore a
retraction of the inclusion of that span. All its matrix entries
are integers, so base change gives the stated retraction over
every $R$.

Finally the map~\eqref{eq:plethysm-basechange} sends
$1\otimes b_T$ to the polynomial $b_T$ with coefficients in $R'$.
The basis just proved on both sides shows that it is an
isomorphism. The splitting explains in particular why taking this
image commutes with arbitrary, possibly nonflat, base change.
\end{proof}

\subsection{Functoriality, all matrix entries, and the exact diagonal}

\begin{proposition}[Functorial maps without denominators]
\label{prop:plethysm-functorial}
Every $R$-linear map $f:E\to F$ between finite free modules induces
a map $\mathsf S_\lambda(f):\mathsf S_\lambda(E)\to
\mathsf S_\lambda(F)$. These maps preserve identities and
composition and commute with arbitrary base change. In the tableau
bases, every matrix entry is an integer polynomial, homogeneous of
degree $|\lambda|$ in the entries of $f$.
\end{proposition}

\begin{proof}
The alternating tensor map commutes with $f^{\otimes h}$ term by
term, tensor reordering commutes with applying $f$ in every slot,
and symmetric quotient maps are natural. Thus the square
\[
 \begin{array}{ccc}
 C_\lambda(E)&\xrightarrow{\delta_{\lambda,E}}&P_\lambda(E)\\
 \big\downarrow C_\lambda(f)&&\big\downarrow P_\lambda(f)\\
 C_\lambda(F)&\xrightarrow{\delta_{\lambda,F}}&P_\lambda(F)
 \end{array}
\]
commutes. Restricting its right arrow to the images defines the
claimed map. The identities and composition laws hold on the
ambient row-symmetric modules and hence on their submodules.
All constructions use tensor, exterior, and symmetric operations
given by the displayed integral formulas. Together with
Theorem~\ref{thm:plethysm-integral-basis}, this proves the base-change
statement.

Here is a formula that also proves the claim about every matrix
entry. Choose bases $e_i$ and $f_j$, and write a matrix $M$ for the
map, so that $f(e_i)=\sum_jM_{j,i}f_j$. For a strictly increasing
column $I$ of length $a$, direct multilinear expansion and signed
sorting in the exterior power give
\begin{equation}
 (\bigwedge^a f)(e_{i_1}\wedge\cdots\wedge e_{i_a})
 =\sum_{\substack{J\text{ increasing}\\|J|=a}}
       \det(M_{J,I})\,f_{j_1}\wedge\cdots\wedge f_{j_a}.
 \label{eq:plethysm-exterior-matrix}
\end{equation}
Terms with repeated row indices vanish. For each distinct row
set the terms with the possible row orders sum, with their sorting
signs, to the determinant in~\eqref{eq:plethysm-exterior-matrix}.
For the columns $I_c(T)$ of a tableau $T$, take the product of
these expansions and then apply $\delta_{\lambda,F}$. This gives
\begin{equation}
 \mathsf S_\lambda(f)b_T
 =\sum_{J_1,\ldots,J_t}
       \left(\prod_{c=1}^t\det(M_{J_c,I_c(T)})\right)
       b_{(J_1,\ldots,J_t)}.
 \label{eq:plethysm-general-matrix}
\end{equation}
Apply the terminating integer straightening procedure already
proved to each filling on the right. Every product of minors in
its coefficient has total degree $\sum_c h_c=|\lambda|$.
Straightening multiplies by integers and adds these products, so
each final tableau coefficient is an integer homogeneous
polynomial of that degree. This is a finite formula: there are
$\prod_c\binom{\operatorname{rank}F}{h_c}$ choices in the sum,
and straightening strictly decreases an order on the finite set
of column-increasing fillings. No complexity bound beyond this
finite procedure is asserted here.
\end{proof}

\begin{theorem}[The exact signed diagonal map]
\label{thm:plethysm-diagonal}
Let $E$ and $F$ have ordered bases of the same cardinality $n$,
and let $D:E\to F$ have the specified matrix
$\operatorname{diag}(d_1,\ldots,d_n)$, with arbitrary $d_i\in R$.
For a semistandard tableau $T$, let $c_i(T)$ count its entries
equal to $i$, and put
\[
 w_T=\prod_{i=1}^n d_i^{c_i(T)}.
\]
With the same tableau order on the two bases, the exact matrix of
$\mathsf S_\lambda(D)$ is $\operatorname{diag}(w_T)_T$.
In particular, there are the following specified maps and modules:
\begin{align}
 \ker\mathsf S_\lambda(D)
   &=\bigoplus_T\operatorname{Ann}_R(w_T)\,b_T(E),\\
 \operatorname{im}\mathsf S_\lambda(D)
   &=\bigoplus_T(w_T)\,b_T(F),\\
 \operatorname{coker}\mathsf S_\lambda(D)
   &\cong\bigoplus_T R/(w_T).
 \label{eq:plethysm-diagonal-modules}
\end{align}
The last isomorphism sends the class of
$\sum_T a_Tb_T(F)$ to $(a_T\bmod (w_T))_T$.
\end{theorem}

\begin{proof}
Apply $D$ in the column exterior tensor defining $b_T$.
Each occurrence of $e_i$ contributes exactly the scalar $d_i$,
and the product of these scalars is $w_T$. The commutative
square in Proposition~\ref{prop:plethysm-functorial} therefore
gives $\mathsf S_\lambda(D)b_T(E)=w_Tb_T(F)$ before any
straightening. The tableau basis theorem proves the asserted
matrix identity. A vector is in its kernel precisely when every
coordinate $a_T$ satisfies $w_Ta_T=0$, which gives the first
formula. Each target coordinate of an image is independently an
element of $(w_T)$, giving the second. The displayed residue map
is surjective and has exactly this image as its kernel, proving
the third. This also treats $d_i=0$ and nonzero zero divisors.
Signed elements $d_i$ have not been replaced by associates or
absolute values.
\end{proof}

\subsection{Iteration and the plethystic exponent matrix}

Let $\mathcal T_\lambda(n)$ be the tableau set in
Theorem~\ref{thm:plethysm-integral-basis}, ordered by its
column-reading words. Put $N=|\mathcal T_\lambda(n)|$.
Regard this ordered set as the alphabet for semistandard tableaux
$U$ of a second partition $\mu$. Write $c_T(U)$ for the number
of entries of $U$ equal to the alphabet letter $T$, and define
\begin{equation}
 \Gamma_i(U)=\sum_{T\in\mathcal T_\lambda(n)}c_T(U)c_i(T).
 \label{eq:plethysm-nested-exponents}
\end{equation}
Here the plethysm is the composed functor
$\mathsf S_\mu\circ\mathsf S_\lambda$, with this integral
Schur-image convention at both stages.

\begin{corollary}[Nested Schur maps over the original ring]
\label{cor:plethysm-nested-diagonal}
The outer tableau elements $b_U$ form an $R$-basis of
$\mathsf S_\mu(\mathsf S_\lambda(E))$.
For the specified diagonal map $D$ of
Theorem~\ref{thm:plethysm-diagonal}, the exact nested map is
\begin{equation}
 \mathsf S_\mu(\mathsf S_\lambda(D))b_U(E)
   =\left(\prod_{i=1}^n d_i^{\Gamma_i(U)}\right)b_U(F).
 \label{eq:plethysm-nested-diagonal}
\end{equation}
Every exponent is a nonnegative integer, and
$\sum_i\Gamma_i(U)=|\mu|\,|\lambda|$.
The kernel, image, and cokernel are given by
\eqref{eq:plethysm-diagonal-modules} with the index $U$ and
scalar $\prod_i d_i^{\Gamma_i(U)}$.
For a general matrix $M$, every matrix entry of the composed
functor is an integer polynomial homogeneous of degree
$|\mu|\,|\lambda|$ in $M$.
\end{corollary}

\begin{proof}
The inner module is finite free with the explicitly ordered basis
$(b_T)_{T\in\mathcal T_\lambda(n)}$ by the basis theorem.
Apply the same basis theorem to that free module and $\mu$.
The inner map is diagonal, with entries
$w_T=\prod_i d_i^{c_i(T)}$. Apply the diagonal theorem to this
map and the outer shape $\mu$. Its scalar on $b_U$ is exactly
\[
 \prod_T w_T^{c_T(U)}
   =\prod_T\prod_i d_i^{c_T(U)c_i(T)}
   =\prod_i d_i^{\Gamma_i(U)}.
\]
The number of letters of an inner tableau is $|\lambda|$, so
$\sum_i c_i(T)=|\lambda|$; the number of letters of an outer
tableau is $|\mu|$, so $\sum_T c_T(U)=|\mu|$. Summing
\eqref{eq:plethysm-nested-exponents} proves the degree identity.
The coordinate kernel and quotient formulas follow from the
diagonal theorem, with no additional ring hypotheses.
For a general $M$, apply the integer homogeneous matrix formula
first to $\lambda$, and then substitute those degree-$|\lambda|$
entries into the degree-$|\mu|$ formula for the outer functor.
This proves the last assertion. Empty shapes and zero inner
rank follow directly from the conventions for empty tensor
products and the absence of nonempty tableaux on an empty
alphabet.
\end{proof}

The construction above retains the entire composed functor and
its tableau matrix. It does not require a decomposition into
irreducible Schur summands. The proof gives ordinary integral
Schur operations on the specified linear maps; identifying them
with a particular nonstandard quantum-group construction requires
the corresponding specified representation map.
For a primary comparison with the flag-minor standard monomial
model, see D.~C.~Lax, \emph{Accessible Proof of Standard Monomial
Basis for Coordinatization of Schubert Sets of Flags},
\href{https://arxiv.org/abs/1508.02744}{arXiv:1508.02744},
Sections~3--4. All basis, straightening, functoriality, and
diagonal assertions used here were proved above over $R$.

\section{Exact Schur transport of the integral generator endpoint}
\label{sec:plethysm-endpoint-transport}

Retain the coefficient ring, original coefficient and original endpoint
factorization of Section~\ref{sec:integral-coefficient-conductor}:
\begin{gather*}
 O=\mathbb Z[1/2],\quad R=O[q,q^{-1}],\quad t=q-1,\quad
 r=q-q^{-1},\quad h=[7]_q[3]_q,\\
 C=q^{10}-q^4-q^{-4}+q^{-10}=r^2h.
\end{gather*}
The ordered rank-four modules and maps are
\[
 E_C\xrightarrow{H}B_R\xrightarrow{V}N_R,\qquad
 H=\operatorname{diag}(1,1,1,h),\quad
 V=\operatorname{diag}(1,1,1,r^2),\quad D_C=VH.
\]
The fourth coordinate of $E_C$ is the twisted source top line,
and that of $N_R$ is the bottom line, with the three specified identity
summands retained. Thus $D_C$ is the stabilized actual divided-square
endpoint, not a new full quantum-group representation. The displayed
bases identify each of the three modules with $R^4$ as an $R$-module;
we keep their names to distinguish domains and codomains.

All Schur functors in this section are the integral, column-exterior to
row-symmetric Schur functors described in
Theorem~\ref{thm:plethysm-integral-basis}. No factorial or hook
length is inverted. We apply these functors to ordinary module maps.
They retain the stated $\mu_2^2$ characters and the semilinear bar
involution, by functoriality. They also retain the explicitly diagonal
Cartan action and zero lowering actions on the stabilized endpoint
modules. This assertion does not apply an ordinary symmetric-group
action to a generic full quantum tensor power.

\subsection{The exact diagonal map for every partition}

For a partition $\lambda$, let $\operatorname{Tab}_d(\lambda)$ be its
tableaux with entries in $\{1,\ldots,d\}$ weakly increasing along rows
and strictly increasing down columns. The empty shape has one empty
tableau. A shape with more than $d$ rows has no such tableaux and its
Schur module is zero. For $T\in\operatorname{Tab}_4(\lambda)$ write
$w(T)$ for the number of entries equal to $4$. Put
\[
 d_\lambda=|\operatorname{Tab}_4(\lambda)|,\qquad
 m_\lambda(k)=|\{T:w(T)=k\}|,
 \qquad P_\lambda(z)=\sum_{k\ge0}m_\lambda(k)z^k.
\]
These definitions are finite and effective; they do not presume a
plethysm decomposition or a positivity conjecture.

\begin{theorem}\label{thm:plethysm-endpoint-diagonal}
In the corresponding tableau bases of the three modules, the maps
$S_\lambda(H)$, $S_\lambda(V)$ and $S_\lambda(D_C)$ have diagonal entries
$h^{w(T)}$, $r^{2w(T)}$ and $C^{w(T)}$, respectively. Their composition is exactly
$S_\lambda(D_C)=S_\lambda(V)S_\lambda(H)$. Each map is injective over
$R$, and
\[
 \operatorname{coker}S_\lambda(D_C)
 =\bigoplus_{k\ge1}\big(R/(C^k)\big)^{m_\lambda(k)}.
\]
Writing $\mathcal L_C,\mathcal L_r$ for their actual image lattices
inside $S_\lambda(N_R)$ gives the exact sequence
\begin{equation}\label{eq:plethysm-conductor-sequence}
 0\longrightarrow\bigoplus_{k\ge1}(R/(h^k))^{m_\lambda(k)}
 \xrightarrow{\ i\ }
 \bigoplus_{k\ge1}(R/(C^k))^{m_\lambda(k)}
 \xrightarrow{\ \pi\ }
 \bigoplus_{k\ge1}(R/(r^{2k}))^{m_\lambda(k)}
 \longrightarrow0,
\end{equation}
where $i(\bar f_T)=\overline{r^{2w(T)}f_T}$ and
$\pi(\bar f_T)=\bar f_T$. The three terms are, in that order,
$\mathcal L_r/\mathcal L_C$, $S_\lambda(N_R)/\mathcal L_C$,
and $S_\lambda(N_R)/\mathcal L_r$.
\end{theorem}

\begin{proof}
Each term in the column alternation followed by row multiplication
has exactly the content of $T$. A diagonal map therefore multiplies
the tableau vector by the product of its diagonal entries, giving
the three stated scalars. The integral tableau basis ensures these
are matrices over $R$ itself. Naturality, or multiplication of their
individual diagonal entries, gives the composition. Nonzero diagonal
entries are non-zero-divisors in the domain $R$, proving injectivity.
The cokernel of multiplication by a scalar $a$ on $R$ is $R/(a)$;
the coordinates with $w=0$ are identity maps and contribute zero.
For each $k>0$, multiplication by $r^{2k}$ sends $R/(h^k)$ into
$R/(C^k)$ since $r^{2k}h^k=C^k$. If its image is zero, the equality
$r^{2k}f=C^kg$ cancels to $f=h^kg$. Projection onto $R/(r^{2k})$
is onto and its kernel consists exactly of these multiples of
$r^{2k}$. Taking the finite direct sum proves every map and exactness
assertion, including the image-lattice interpretations.
\end{proof}

Pad $\lambda$ with zero rows to length four. The exact smallest and
largest exponents, for a nonzero Schur module, are
\begin{equation}\label{eq:plethysm-single-floors}
 \min_T w(T)=\lambda_4,\qquad \max_Tw(T)=\lambda_1.
\end{equation}
Every column of height four must contain $4$, so the lower bound is
$\lambda_4$. Fill row $i$ constantly with $i$ to attain it. At most
one $4$ occurs per column, giving the upper bound $\lambda_1$.
In a column of height $a$, place $5-a,\ldots,4$ in order. As column
heights weakly decrease from left to right, these entries weakly
increase along each row. This is a tableau attaining the upper bound.

The complete multiplicities also have an explicit smaller formula.
Deleting the entries $4$ leaves a shape $\mu$ with at most three rows
and
\[
 \lambda_1\ge\mu_1\ge\lambda_2\ge\mu_2\ge
 \lambda_3\ge\mu_3\ge\lambda_4.
\]
Conversely each tableau of that shape on $1,2,3$ extends uniquely by
filling $\lambda/\mu$ with $4$. Indeed the displayed inequalities say
precisely that the added boxes lie at row ends and no two share a
column; the row and column conditions then follow. Thus
\begin{equation}\label{eq:plethysm-branching}
 m_\lambda(k)=
 \sum_{\substack{\lambda_i\ge\mu_i\ge\lambda_{i+1}\ (1\le i\le3)\\
                  |\lambda|-|\mu|=k}}
 \frac{(\mu_1-\mu_2+1)(\mu_2-\mu_3+1)(\mu_1-\mu_3+2)}2.
\end{equation}
The displayed summand is an integer count, not a division in $R$.
To prove it directly, delete the $3$ entries next. The remaining
two-row shape $(a,b)$ satisfies
$\mu_2\le a\le\mu_1$ and $\mu_3\le b\le\mu_2$.
Deleting its entries $2$ leaves a one-row shape of any length from
$b$ to $a$, giving $a-b+1$ possibilities. Hence its total is
\[
 \sum_{a=\mu_2}^{\mu_1}\sum_{b=\mu_3}^{\mu_2}(a-b+1)
 =\frac{(\mu_1-\mu_2+1)(\mu_2-\mu_3+1)(\mu_1-\mu_3+2)}2,
\]
by the finite arithmetic-progression sums. This proves
\eqref{eq:plethysm-branching} without a dimension formula assumption.

\subsection{Nested Schur maps and the precise plethysm character}

Fix another partition $\mu$ with $d_\mu>0$, enumerate its integral
tableau basis $T_1,\ldots,T_{d_\mu}$ so that the weights
$w_i=w(T_i)$ weakly increase, and choose any fixed order within a
tie. A tableau $U\in\operatorname{Tab}_{d_\mu}(\lambda)$ indexes an
integral basis vector of $S_\lambda(S_\mu(E_C))$. Its exponent is
\[
 W(U)=\sum_{\text{boxes }b\text{ of }\lambda}w_{U(b)},\qquad
 M_{\lambda,\mu}(k)=|\{U:W(U)=k\}|.
\]
Apply Theorem~\ref{thm:plethysm-endpoint-diagonal} to a diagonal map
whose entries now are $C^{w_1},\ldots,C^{w_{d_\mu}}$. It proves, over
the unchanged ring $R$, the exact matrices
\begin{equation}\label{eq:plethysm-endpoint-nested-diagonal}
 S_\lambda(S_\mu(D_C))=\operatorname{diag}_{U}(C^{W(U)}),
 \quad S_\lambda(S_\mu(H))=\operatorname{diag}_{U}(h^{W(U)}),
 \quad S_\lambda(S_\mu(V))=\operatorname{diag}_{U}(r^{2W(U)}).
\end{equation}
All injectivity, cokernel and conductor formulas above hold with
$m_\lambda(k)$ replaced by $M_{\lambda,\mu}(k)$. This is a direct
application of the proved basis and scalar calculation; it assumes
neither exactness of Schur functors on arbitrary short exact sequences
nor an integral direct-sum decomposition of a plethysm.

For completeness define the symmetric polynomial
$s_\mu(x_1,\ldots,x_4)=\sum_T\prod_{b}x_{T(b)}$ by its tableau
character. The multivariable diagonal action proved above shows that
the character of the nested module is obtained by evaluating
$s_\lambda$ on the list of monomials $\prod_bx_{T_i(b)}$, with each
tableau contributing its own entry, including repeated monomials.
This substitution is precisely $s_\lambda[s_\mu]$ in four variables.
Consequently the exact one-variable identity is
\begin{equation}\label{eq:plethysm-character-specialization}
 \sum_kM_{\lambda,\mu}(k)z^k
 =s_\lambda(z^{w_1},\ldots,z^{w_{d_\mu}})
 =s_\lambda[s_\mu](1,1,1,z).
\end{equation}
Each coefficient on this restriction is a proved nonnegative integer
count. This equation does not determine the decomposition into
irreducible $GL_4$ characters: it is the restriction of a character
along one specified one-parameter subgroup. In particular this
computation supplies a transport map and its invariant factors, not
a plethysm-multiplicity obstruction for a class variety.

If $\lambda$ has at most $d_\mu$ rows, the exact extreme exponents are
\begin{equation}\label{eq:plethysm-nested-floors}
 a_{\lambda,\mu}=\sum_i\lambda_iw_i,
 \qquad b_{\lambda,\mu}=\sum_i\lambda_iw_{d_\mu+1-i}.
\end{equation}
In a column of height $a$, the entry in its $i$th box is at least
$i$ and at most $d_\mu-a+i$. Summing the corresponding ordered
weights proves the lower and upper bounds. Filling row $i$ by $i$
attains the lower bound. Filling each height-$a$ column with
$d_\mu-a+1,\ldots,d_\mu$ attains the upper bound and is row weakly
increasing, for the same reason as in
\eqref{eq:plethysm-single-floors}. Grouping the lower and upper
column sums by the row lengths gives exactly
\eqref{eq:plethysm-nested-floors}.

Let $d_0=m_\mu(0)=|\operatorname{Tab}_3(\mu)|$. Weight zero in the
nested module occurs precisely when every entry labels an inner
tableau of weight zero. Therefore
\begin{equation}\label{eq:plethysm-zero-rank}
 M_{\lambda,\mu}(0)=|\operatorname{Tab}_{d_0}(\lambda)|.
\end{equation}
This also handles $d_0=0$: only the empty outer shape has a
nonzero rank. At $C=0$ the matrices are diagonal with an identity
block of this size and zero elsewhere, giving their exact rank,
kernel and cokernel over every base-change ring, not only over fields.

For example $S_{(2)}(R^4)$ has inner weights $0^6,1^3,2^1$.
An outer two-box row chooses an unordered pair with repetitions.
The five weight counts for $S_{(2)}(S_{(2)}(D_C))$ are
\[
 (21,18,12,3,1):\quad
 \binom72,\quad6\cdot3,\quad6\cdot1+\binom42,\quad3\cdot1,\quad1.
\]
Its diagonal therefore has $21$ identity entries, $18$ entries $C$,
$12$ entries $C^2$, $3$ entries $C^3$, and one entry $C^4$.
No irreducible decomposition was needed for this exact matrix.

There is also a uniform determinant formula. Let $F$ be either
$S_\lambda$ or $S_\lambda\circ S_\mu$, let $n$ be its polynomial
degree and $d$ its rank on $R^4$. Then
\begin{equation}\label{eq:plethysm-determinant}
 \det F(D_C)=C^{nd/4},\qquad
 \sum_k k\,\operatorname{rank}F(R^4)_k=nd/4\in\mathbb Z.
\end{equation}
To prove the divisibility and exponent, grade the integral tableau
basis by the four original contents. Every vector has total content
$n$. A permutation matrix exchanging any two original basis vectors
induces, by functoriality, an integral isomorphism exchanging the
corresponding weight spaces. The sums of the four content coordinates
over all basis vectors are thus equal integers, whose total is $nd$.
The fourth sum is the sum of the exponents in the diagonal matrix.
Taking its determinant proves the formula, including its positive sign
in the matched ordered bases.

\subsection{Integral specialization, residual primes and exact orders}

For a weight $k>0$, put $Q_{h,k}=R/(h^k)$,
$Q_{C,k}=R/(C^k)$ and $Q_{r,k}=R/(r^{2k})$.
The resolution $0\to R\xrightarrow{t}R\to O\to0$ computes every
derived specialization used here. The ideal $(t)$ is prime, and
$h(1)^k=21^k\ne0$ in the domain $O$. Hence
$\operatorname{Tor}_1^R(Q_{h,k},O)=0$: if $tf=h^kg$, reduction modulo
$t$ forces $g(1)=0$, and cancellation gives $f\in(h^k)$.
The two other first Tor groups are free of rank one over $O$,
with their exact generators $C^k/t$ and $r^{2k}/t$. Indeed
$tf=C^kg$ forces $f=(C^k/t)g$, and this representative vanishes
modulo $C^k$ exactly when $g\in(t)$; the same proof applies to
$r^{2k}$. All higher Tor groups vanish by the length-one resolution.
The specialized sequence for this weight is therefore
\begin{equation}\label{eq:plethysm-derived-21}
 0\longrightarrow O\xrightarrow{21^k}O
 \xrightarrow{\mathrm{mod}\ 21^k}O/(21^k)
 \xrightarrow{0}O\xrightarrow{1}O\longrightarrow0.
\end{equation}
Projection sends $C^k/t$ to $h^k$ times $r^{2k}/t$, giving the
first displayed scalar. Lift $r^{2k}/t$ into $Q_{C,k}$ and multiply
by $t$: the result $r^{2k}$ is the image of $1\in Q_{h,k}$.
Thus the connecting map has the stated positive sign and sends the
chosen generator to $1$. The next map multiplies by $r^{2k}$ and
vanishes after specialization; the last is the identity.
The entire Schur or nested Schur sequence is the direct sum of
these sequences with the exact weight multiplicities already proved.

The specialized map $F(D_C)_0$ has kernel and cokernel
$O^{d-m(0)}$ in the positive-weight coordinates. The induced map
between its kernel and that of $F(V)_0$ is
$\operatorname{diag}_{k>0}(21^k)$ with repetitions $m(k)$.
It is injective and its cokernel is
\[
 \bigoplus_{k>0}(O/(21^k))^{m(k)}
 \simeq\bigoplus_{k>0}
       \big((O/(3^k))\oplus(O/(7^k))\big)^{m(k)}.
\]
For the final isomorphism use the map taking a class to its two
residues. Since $3^k$ and $7^k$ are coprime, choose integers $a,b$
with $a3^k+b7^k=1$; an inverse sends $(x,y)$ to
$b7^kx+a3^ky$ modulo $21^k$. Its two compositions are the identity.
This proves the two residual prime modules without discarding the
original factor $21^k$.

It is essential to specify which valuation is being used.
At the height-one vertical prime $(\ell)$ of $R$, for $\ell=3$ or
$7$, the polynomial $C$ is a unit of $R_{(\ell)}$, since its Laurent
coefficients include $1$. Every transported map is consequently an
isomorphism there; its vertical valuations are zero. The ideal
$(\ell,t)$ gives a two-dimensional local ring and does not itself
specify a discrete valuation. The exceptional orders below instead
refer to the discrete valuation rings
\[
 D_0=\mathbb Q[q,q^{-1}]_{(t)},\qquad
 D_\ell=\mathbb F_\ell[q,q^{-1}]_{(t)}\quad(\ell\text{ odd}).
\]
Write $u=(q+1)/q$, so $r=tu$ and $u(1)=2$. The exact orders and
leading coefficients of the original $C$ are
\[
\begin{array}{c|c|c}
\text{coefficient field}&\operatorname{ord}_t C&
          t^{-\operatorname{ord}_tC}C\bmod t\\\hline
\mathbb Q&2&84\\
\mathbb F_3&4&1\\
\mathbb F_7&8&5\\
\mathbb F_\ell, \ell\ne2,3,7&2&84\bmod\ell.
\end{array}
\]
Here $C=t^2u^2h$ gives the characteristic-zero and ordinary-prime
rows. For an odd prime $\ell$, multiply the finite sum for
$[\ell]_q$ by $q^{\ell-1}$ and then by $q^2-1$. In characteristic
$\ell$ the result is $q^{2\ell}-1=(q^2-1)^\ell$.
Cancellation in the polynomial domain gives
$[\ell]_q=(q-q^{-1})^{\ell-1}$. Hence $C=r^4[7]_q$ in
characteristic $3$ and $C=r^8[3]_q$ in characteristic $7$.
Their leading residues are $2^4\cdot7=1$ in $\mathbb F_3$ and
$2^8\cdot3=5$ in $\mathbb F_7$, proving the table exactly.

Let $v$ be the relevant middle-column integer $2,4$, or $8$, and
$c$ the last-column residue. Raising the retained equality to the
$k$th power proves
\begin{equation}\label{eq:plethysm-smith-orders}
 C^k=t^{vk}\varepsilon_k,\quad\varepsilon_k\in D_\ell^\times,
 \quad\varepsilon_k\bmod t=c^k,\qquad
 \operatorname{coker}F(D_C)\simeq
 \bigoplus_{k>0}(D_\ell/(t^{vk}))^{m(k)}.
\end{equation}
This is an explicit Smith form: multiply each matrix row by
the inverse unit $\varepsilon_k^{-1}$, equivalently use target basis
vector $\varepsilon_k$ times the original vector, retaining that unit and its
specified leading residue as part of the original map. Ordered by
$k$, consecutive powers of $t$ divide one another. Its length is
$v\sum_kkm(k)$, and the exact least and largest matrix-entry orders are $v$ times
the exponents in \eqref{eq:plethysm-single-floors} or
\eqref{eq:plethysm-nested-floors}.
More precisely, the $j$th filtration quotient of the cokernel, for
$j\ge0$, has residue-field dimension
\[
 \dim\big(t^j\operatorname{coker}/t^{j+1}\operatorname{coker}\big)
 =\sum_{k:vk>j}m(k).
\]
For a single cyclic summand the classes $1,t,\ldots,t^{vk-1}$ are
a basis, proving both length and filtration formulas by direct sum.

The mixed-prime arithmetic can also be retained before reduction.
In $\mathbb Z_\ell[[t]]$, with $q=1+t$ and $\ell=3$ or $7$, define
\[
 E_\ell(t)=\Phi_\ell(1+t)
 =\frac{(1+t)^\ell-1}{t}
 =\sum_{j=0}^{\ell-1}\binom\ell{j+1}t^j.
\]
Here the quotient is the displayed polynomial identity. Its constant
term is $\ell$, its leading term is $t^{\ell-1}$, and each other
coefficient is divisible by $\ell$; the constant term is not
divisible by $\ell^2$. The exact quantum factor identity
\[
 [\ell]_q=q^{1-\ell}\Phi_\ell(q)\Phi_{2\ell}(q),\qquad
 \Phi_{2\ell}(q)=\sum_{j=0}^{\ell-1}(-q)^j,
\]
follows by multiplying the two geometric sums; $\Phi_{2\ell}(1)=1$.
The other quantum factor in $h$ has value $7$ when $\ell=3$ and
$3$ when $\ell=7$. Thus, writing that factor as $h_{\ne\ell}$,
the full equality is
\begin{equation}\label{eq:plethysm-mixed-local}
 C^k=t^{2k}E_\ell(t)^k U_\ell(t)^k,
 \quad U_\ell=u^2q^{1-\ell}\Phi_{2\ell}(q)h_{\ne\ell}
           \in\mathbb Z_\ell[[t]]^\times.
\end{equation}
The constant terms of $U_3,U_7$ are respectively $28$ and $12$,
which are units in the stated coefficient rings. This retains the
horizontal factor $t^{2k}$ and the separate arithmetic factor
$E_\ell^k$; it does not replace a two-dimensional local ring by an
unspecified valuation. Modulo $\ell$, $E_\ell=t^{\ell-1}$, recovering
exactly the orders $4k$ and $8k$.

There is an additional first Tor term if the conductor sequence is
first reduced to characteristic $3$ or $7$ and then specialized at
$t=0$. Over the corresponding Laurent polynomial ring put
$\bar h=h\bmod\ell$; it is nonzero and divisible by $t$.
The three first Tor generators are now
$\bar h^k/t$, $\bar C^k/t$, and $\bar r^{2k}/t$.
Multiplication by $\bar r^{2k}$ sends the first exactly to the
second; projection sends the second to $\bar h^k$ times the third,
which is zero on the residue field. The connecting map again sends
the third generator to $1$. Therefore, for every positive weight,
the full sequence over $\mathbb F_\ell$ is
\[
 0\longrightarrow\mathbb F_\ell\xrightarrow{1}\mathbb F_\ell
 \xrightarrow{0}\mathbb F_\ell\xrightarrow{1}\mathbb F_\ell
 \xrightarrow{0}\mathbb F_\ell\xrightarrow{1}\mathbb F_\ell
 \longrightarrow0.
\]
These descriptions follow from the same length-one resolution;
in particular the extra Tor term has been calculated, not omitted.

Finally, all exponents in the original integral maps are nonnegative.
A negative valuation does occur for the actual inverse after the
specified localization: $F(D_C)^{-1}$ has entries $C^{-k}$ over
$R[C^{-1}]$ and orders $-vk$ over the fraction field of $D_\ell$.
If a positive-weight coordinate occurs, this inverse cannot be an
$R$-linear inverse on the original lattices: an inverse image of its
unit target vector would require $C^{-k}\in R$, which is impossible
since $C$ vanishes at $q=1$. Its precise failure to be integral is
the cokernel already computed in
\eqref{eq:plethysm-smith-orders}. This supplies an exact pole/defect
correspondence without identifying a negative Laurent coefficient
with a negative plethysm multiplicity.

\subsection{Exact negative inertia at the retained circle character}

Retain the circle real structure $q^*=q^{-1}$, with conjugation on
coefficients, and the original exact character from the conic calculation:
\[
 q_0=\frac{\sqrt{14}+i\sqrt2}{4},\quad
 q_0^{-1}=\frac{\sqrt{14}-i\sqrt2}{4},\quad
 p_0=\frac{\sqrt{14}}2,\quad r_0=\frac{i}{\sqrt2},\quad
 a_0=\frac12,\quad s_0=\frac14.
\]
The coefficient field is $K_0=\mathbb Q(\sqrt{14},i\sqrt2)\subset\mathbb C$
with its displayed embedding and usual conjugation. The product of the
two displayed values for $q_0$ is $(14+2)/16=1$, so evaluation is a
homomorphism respecting the circle involution. To check every factor,
put $x_0=q_0^2+q_0^{-2}=p_0^2-2=3/2$. Then
\[
 [3]_{q_0}=x_0+1=\frac52,\qquad
 [7]_{q_0}=x_0^3+x_0^2-2x_0-1=\frac{13}{8}.
\]
The second identity expands $[7]_q$ as
$(q^6+q^{-6})+(q^4+q^{-4})+(q^2+q^{-2})+1$
and uses $q^6+q^{-6}=x^3-3x$, $q^4+q^{-4}=x^2-2$.
Consequently the full factorization retains the exact signs
\begin{equation}\label{eq:plethysm-circle-coefficient}
 r_0^2=-\frac12,\qquad h(q_0)=\frac{65}{16}>0,
 \qquad C(q_0)=-\frac{65}{32}.
\end{equation}

Let $F$ be either Schur functor already constructed, and identify its
source and target \emph{as based vector spaces} by the map taking
each source tableau basis vector to the equally labelled target
vector. This specifies the endomorphism $A_F$ representing
$F(D_C)(q_0)$ on the source; it does not infer an equality of the
original source and target modules. Give this source the positive
Hermitian form with its specified tableau basis orthonormal.
Because every diagonal entry of $A_F$ is the real number
$(-65/32)^k$, the endomorphism is selfadjoint for this form.
The Hermitian form that it defines has the complete formula
\[
 B_F(u,v)=\sum_T(-65/32)^{w(T)}\overline{u_T}v_T
\]
for a single Schur functor, and the same formula with $U,W(U)$ for
a nested one. In particular its positive, negative and zero indices
are exactly
\begin{equation}\label{eq:plethysm-inertia}
 n_+=\sum_{k\text{ even}}m(k),\qquad
 n_-=\sum_{k\text{ odd}}m(k),\qquad n_0=0.
\end{equation}
Indeed the even and odd coordinate spans give a direct orthogonal
sum of positive and negative definite spaces. A subspace of dimension
larger than the even-coordinate count intersects the odd-coordinate
space nontrivially by the dimension formula, so cannot be positive
definite. Interchanging even and odd proves maximality for the
negative index. Every diagonal coefficient is nonzero, so the
radical is zero. This proves the inertia statement over the specified
embedded field and after extension to $\mathbb C$.

The signature is thus the exact character evaluation
\[
 n_+-n_-=P_\lambda(-1)=s_\lambda(1,1,1,-1)
\]
for a single Schur functor, and
\[
 n_+-n_-=s_\lambda[s_\mu](1,1,1,-1)
 =s_\lambda(\underbrace{1,\ldots,1}_{m_\mu(\mathrm{even})},
             \underbrace{-1,\ldots,-1}_{m_\mu(\mathrm{odd})})
\]
for a nested one. The last equality follows by evaluating each
inner-tableau monomial at $(1,1,1,-1)$; it is $+1$ or $-1$ according
to its weight parity. The scalar $h(q_0)^k$ is positive for every
coordinate. Therefore $F(V)(q_0)$, whose corresponding coefficients
are $(-1/2)^k$, has exactly the same inertia. The actual conductor
factor $F(H)(q_0)$ is positive definite in these matched bases and
their matrices still multiply to $F(D_C)(q_0)$. This proves the
propagation of the negative directions through the actual endpoint
factorization, without discarding the positive residual factor.

For $S_{(2)}\circ S_{(2)}$ the previously computed profile gives
$n_+=21+12+1=34$, $n_-=18+3=21$, and signature $13$.
This also shows that the signature need not stay negative under
iteration even when some negative directions survive: the exact
parity-weight character, not an unsigned dimension, determines it.
No entry of $m(k)$ has become negative. These Hermitian forms are
specified forms on the endpoint Schur modules; no positivity form
of the nonstandard RH programme or spectral membership of $s_0$
has been identified with them.

\subsection{The original shape and an actual specialized source projector}

For the original shape $\nu=(7,5,3)$, formula
\eqref{eq:plethysm-branching} is the following explicit finite sum:
\begin{gather*}
 \sum_{a=5}^{7}\sum_{b=3}^{5}\sum_{c=0}^{3}
 \frac{(a-b+1)(b-c+1)(a-c+2)}2\,z^{15-a-b-c}\\
 =27+81z+162z^2+270z^3+300z^4+252z^5+126z^6+42z^7.
\end{gather*}
The exponents of $S_\nu(D_C)$, with multiplicities, are therefore
exactly the eight displayed coefficients. Their sum is $1260$ and
their weighted sum is $4725$. In characteristic zero at $q=1$,
the map has rank $27$, kernel and cokernel dimension $1233$, and
cokernel length $9450$ over $D_0$. The corresponding lengths over
$D_3,D_7$ are $18900$ and $37800$. These follow by multiplying
$4725$ by the proved orders $2,4,8$. They describe this particular
integral Schur transport and its stated base changes.
At the distinct circle character $q_0$ of
\eqref{eq:plethysm-circle-coefficient}, its inertia is
\begin{gather*}
 (n_+,n_-,n_0)
 =(27+162+300+126,\ 81+270+252+42,\ 0)=(615,645,0),\\
 \operatorname{signature}=-30.
\end{gather*}
This is a negative signature with $645$ individually specified
negative tableau directions, and is an immediate exact instance of
\eqref{eq:plethysm-inertia}.

There is also a proved map to the full original source, over its
already justified rational specialization. Put
$X=V\otimes W$ with the unchanged alphabet
\[
 x_1=v_1\otimes w_1,\quad x_2=v_1\otimes w_2,\quad
 x_3=v_2\otimes w_1,\quad x_4=v_2\otimes w_2.
\]
Section~\ref{sec:full-source-generators} constructs the isomorphism
$s_\nu:L_{\mathbb Q}\to Q_\nu(X)$ from the rational specialized
full source lattice to the exterior-column classical relation quotient,
and proves its inverse $j_\nu$. Here $Q_\nu(X)$ denotes that precise
quotient presentation, retaining its original column order. To compare
it with the integral Schur-image convention in this section, define
\[
 \eta_\nu:Q_\nu(X)\longrightarrow S_\nu(X),\qquad
 [\text{column exterior tensor of }T]\longmapsto b_T.
\]
Every classical adjacent column relation is the alternating shuffle
relation: the free letters form an alternating list of length $a+1$
across two columns with left height $a$, with the other letters fixed.
Its row-polynomial image vanishes by
Lemma~\ref{lem:plethysm-shuffle}; the unchanged other columns multiply
this zero polynomial. Thus $\eta_\nu$ is well defined. It is onto
because column tensors generate the Schur image. Its source has
dimension $1260$ by the proved source isomorphism, and its target has
dimension $1260$ by the explicit tableau count above. Its kernel is
therefore zero. Its inverse has a basis formula: send $b_T$ for each
semistandard $T$ to its column-tensor class. This right inverse is
also a left inverse by injectivity, proving both compositions.
Alternation, row multiplication and quotient maps are equivariant,
so both maps are equivariant as well. Put
\[
 \sigma_\nu=\eta_\nu s_\nu:L_{\mathbb Q}\longrightarrow S_\nu(X),
 \qquad\sigma_\nu^{-1}=j_\nu\eta_\nu^{-1}.
\]
Use the explicit coordinate isomorphisms
$\theta_E:E_C\otimes_R\mathbb Q\to X$ and
$\theta_N:N_R\otimes_R\mathbb Q\to X$ taking their $i$th displayed
basis vector to $x_i$. Since $C(1)=0$, their square is
\[
 \theta_N(D_C)_0=P_X\theta_E,\qquad
 P_X(x_i)=x_i\ (1\le i\le3),\quad P_X(x_4)=0.
\]
Applying $S_\nu$ and composing with the proved source maps gives the
actual endomorphism
\begin{equation}\label{eq:plethysm-original-source-projector}
\Pi_\nu=\sigma_\nu^{-1} S_\nu(P_X)\sigma_\nu:
 L_{\mathbb Q}\longrightarrow L_{\mathbb Q}.
\end{equation}
It is idempotent: $\sigma_\nu\sigma_\nu^{-1}=1$, functoriality, and $P_X^2=P_X$
give $\Pi_\nu^2=\Pi_\nu$. The Schur tableau formula shows its image
is $\sigma_\nu^{-1}S_\nu(\langle x_1,x_2,x_3\rangle)$ with its $27$-element
tableau basis, and its kernel is the span
of the $1233$ tableaux containing $x_4$, pulled back by
$\sigma_\nu^{-1}$. These are explicit inverse
descriptions of both subspaces, proving the ranks independently of
an arbitrary choice of $1260$ vectors in the source.

The original diagonal $GL(V)\times GL(W)$ torus preserves each
$x_i$, so it commutes with $P_X$. The proved equivariance of
$\sigma_\nu,\sigma_\nu^{-1}$ and Schur functoriality prove that $\Pi_\nu$ commutes
with this torus. No full raising/lowering equivariance is claimed:
already $F_Vx_2=x_4$, so $P_XF_Vx_2=0$ whereas
$F_VP_Xx_2=x_4$. The construction provides the exact coordinate
and source maps under which the endpoint specialization acts.
The rational source isomorphism is used only over $\mathbb Q$;
the integral torsion and the characteristic-$3,7$ calculations
above have not been transferred through an unproved integral or
modular source isomorphism.

The complete transport thus reaches ordinary Schur and nested Schur
modules, their exact conductor quotients, their residual prime
modules and their specified source specialization. It provides no
class-variety coordinate-ring obstruction, canonical-basis
Frobenius model, or Boolean complexity reduction by itself. Those
maps must be supplied by the continuing GCT construction; no
unproved complexity conclusion is a hypothesis in any result here.


\section{The published concentration coordinates and an exact conic heat transport}
\label{sec:ns-conic-transport}

This section uses the coordinate formulas of the public manuscript
\emph{Finite Time Blowup for Navier--Stokes}, whose author line is
OpenAI. The source reading receipt records the preserved revisions and
the audited source pin. The manuscript's Theorem~1.1 claims finite-time
unbounded velocity for every viscosity $\nu>0$, zero initial velocity,
a smooth compactly supported force, fixed compact spatial support before
time one, and uniformly bounded kinetic energy. That global claim is
not independently verified in this section. All statements proved below
are explicit coordinate, differential-algebra, and scalar-germ identities.
In particular, the leading field appearing in the source is retained as
a leading field, rather than substituted for its completed construction.
The current 166-page PDF has SHA-256
\begin{center}\small\texttt{0e779481c4da40bd28d1e642e1d8ca57447d129610df28dfa5a11e9af8ae228f}.\end{center}
The earlier 165-page PDF has SHA-256
\begin{center}\small\texttt{8c8a94ad9ac824c8b605b9827cadf7beaca48bd10b380de3cfc872a2c37afa81}.\end{center}
The source audit compares the cited passages in both versions: the
coordinates are (3.2), current p.~7 and (4.1), current p.~24;
the differential identities are Lemma~4.1 and (4.2), current p.~25;
the profiles and balances are (4.3)--(4.7), current pp.~25--26;
the growth and viscosity formulas are (10.21)--(10.23), current
p.~124. The corresponding earlier passages are on pp.~7, 24--25,
and 124. The selected equations agree; complete revision equivalence
is not asserted. The preserved source commit is
\begin{center}\small\texttt{8937a8f4cbc7abaab5e9e97d1cc7f5d2319d9538}.\end{center}
The public manuscript supplies these source objects
\cite{OpenAINSBridge2026}. Its revised historical account credits
Diego C\'ordoba, Luis Mart\'inez-Zoroa, and the indicated joint work
with Fan Zheng for preceding forced-flow constructions.
Tristan Buckmaster's primary statement separately identifies his
joint work with Levent Alp\"oge on announced smooth-forcing
incompressible-porous-media, Boussinesq, and three-dimensional Euler
blowup, and describes an unfinished verification of a hypodissipative
Navier--Stokes claim \cite{BuckmasterStatementBridge2026}.
These are the stated attributions and mathematical scopes. Neither
source observation establishes priority or a technical dependency
of the positive-viscosity manuscript on that joint work. The new
conic calculations below are proved here from the displayed objects.

\subsection{The exact source coordinates and differential operators}
To prevent collisions of source notation, write $\rho$ for the source's
cylindrical radius $r$, $Q$ for its concentration variable $q$, $\xi$
for its similarity variable $\eta$, and $\sigma=\rho^2/2$ for its
squared-radius variable $s$. The arithmetic coordinates $r,s,b,q$ retain
their previous meanings. These are names for the original variables,
not changes of their values. Retain
\begin{gather*}
 A=\tfrac12+h,\qquad D=\tfrac12-h,\qquad \tau=1-t,\\
 z=Q^D\xi,\qquad \tau=Q(1-\xi^2),\qquad X=\frac{\rho^2}{2Q},\\
 d=1-\xi^2,\qquad \ell=1-2h\xi^2.
\end{gather*}
The coordinate identities hold for $0<h<1/2$, $\tau>0$,
$Q>0$, $-1<\xi<1$; the source's profile construction uses
$0<h<1/100$. Powers of $Q$ denote their positive real values.
For each $(z,\tau)\in\mathbb R\times(0,\infty)$ there is exactly one
$Q>|z|^{1/D}$ with
\[
 Q-z^2Q^{2h}=\tau.
\]
Indeed the left side is zero at the stated endpoint, tends to infinity,
and has derivative $1-2hz^2Q^{2h-1}\geq1-2h>0$ on this interval.
The implicit-function theorem and $\ell>0$ prove smoothness. Inverting
gives $\xi=zQ^{-D}$ and $\rho=\sqrt{2QX}$, the latter for $X>0$.
The time--axial Jacobian is
\[
 \det\frac{\partial(\tau,z)}{\partial(Q,\xi)}
 =Q^D(d+2D\xi^2)=Q^D\ell>0.
\]
Differentiating these exact relations at fixed physical coordinates gives
\begin{align*}
 Q_t&=-\ell^{-1},& \xi_t&=D\xi/(Q\ell),& X_t&=X/(Q\ell),\\
 Q_z&=2\xi Q^{1-D}/\ell,&
 \xi_z&=d/(Q^D\ell),& X_z&=-2\xi X/(Q^D\ell),\\
 Q_\rho&=0,&\xi_\rho&=0,&X_\rho&=\rho/Q.
\end{align*}
For example $Q_\tau=1/\ell$, $\xi_\tau=-D\xi/(Q\ell)$,
and $X_\tau=-X/(Q\ell)$, so the first line retains the time reversal.
The $\xi_z$ calculation uses $\ell-2D\xi^2=d$.

For any real $w$ and smooth profile $f(X,\xi)$ define
$\mathcal P_w f=Q^w f(X,\xi)$ and
\begin{align*}
 T_wf&=\ell^{-1}(-wf+D\xi f_\xi+Xf_X),\\
 Z_wf&=\ell^{-1}(2w\xi f+d f_\xi-2\xi Xf_X),\qquad
 Rf=2Xf_{XX}+2f_X.
\end{align*}
The chain and product rules, with the preceding six derivatives, give
\begin{equation}\label{eq:ns-profile-operators-exact}
 \partial_t\mathcal P_w=\mathcal P_{w-1}T_w,
 \qquad \partial_z\mathcal P_w=\mathcal P_{w-D}Z_w,
 \qquad
 (\partial_\rho^2+\rho^{-1}\partial_\rho)\mathcal P_w
       =\mathcal P_{w-1}R.
\end{equation}
For the last identity, the two terms from differentiating $\rho/Q$
and the cylindrical term each contribute $Q^{w-1}f_X$, and
$\rho^2Q^{w-2}f_{XX}=2XQ^{w-1}f_{XX}$. Thus all radial constants
are retained. Applying the axial identity twice gives
\begin{align}
 (\partial_t-\nu\Delta)\mathcal P_w f
 &=\mathcal P_{w-1}(T_w-\nu R)f
     -\nu\mathcal P_{w-1+2h}Z_{w-D}Z_wf,
       \label{eq:ns-profile-heat-full}\\
 -\tfrac14\partial_t^2\mathcal P_w f
 &=-\tfrac14\mathcal P_{w-2}T_{w-1}T_wf.
       \label{eq:ns-profile-arithmetic-full}
\end{align}
Here $\Delta$ is the three-dimensional scalar Laplacian on
axisymmetric functions. These equations keep both distinct similarity
weights $w-1$ and $w-1+2h$, as well as the arithmetic factor $-1/4$.

The actual source leading profiles satisfy
\[
 u^{(0)}_\theta=Q^{-A}E,\quad u^{(0)}_z=Q^{-A}U,
 \quad \rho u^{(0)}_\rho=V_0,\quad p^{(0)}=Q^{-2A}\Pi,
 \qquad E=C_0^{-1}\sqrt{2X}\,\phi,\quad C_0>1,
\]
where $C_0$ denotes the source's amplitude constant, not the GCT
coefficient. Its incompressibility and pressure identities are
\[
 V_0=\frac X\ell\left(2\xi U-2D\xi\mathcal A_XU
                   -d\partial_\xi\mathcal A_XU\right),\quad
 \mathcal A_XU=X^{-1}\int_0^XU(y,\xi)\,dy,\qquad
 \Pi_X=E^2/(2X).
\]
The first follows by integrating
$(V_0)_X=-Z_{-A}U$ from zero, using $A+D=1$ and
integrating $XU_X$ exactly. The second is the source radial-pressure
balance $\partial_\rho p^{(0)}=(u^{(0)}_\theta)^2/\rho$.
The regular axis interpretation is the one stated by the source:
$E/\sqrt{2X}$, $U$, and $V_0/X$ are smooth at $X=0$.

For this particular leading field its physical material derivative on
axisymmetric scalars has the exact representation
\begin{equation}\label{eq:ns-source-material-profile}
 \bigl(\partial_t+u^{(0)}_\rho\partial_\rho
                     +u^{(0)}_z\partial_z\bigr)\mathcal P_w f
   =\mathcal P_{w-1}M_wf,\qquad
 M_w=T_w+V_0\partial_X+U Z_w.
\end{equation}
Indeed $u^{(0)}_\rho\partial_\rho\mathcal P_wf
=Q^{w-1}V_0f_X$, while the axial term has weight
$w-A-D=w-1$. This proves the correspondence for the actual source
ansatz, including its nonlinear dependence on the chosen profiles.

For completeness this also determines its full momentum residual,
including the cylindrical frame terms. Put $f_\rho=V_0/\sqrt{2X}$
and $R_1=R-(2X)^{-1}$. On $X>0$ the three components of
$(\partial_t+u^{(0)}\cdot\nabla)u^{(0)}-\nu\Delta u^{(0)}
+\nabla p^{(0)}$ are
\begin{align*}
 \mathfrak r_\rho={}&\mathcal P_{-3/2}
       (M_{-1/2}f_\rho-\nu R_1f_\rho)
       -\nu\mathcal P_{-3/2+2h}Z_{-1/2-D}Z_{-1/2}f_\rho,\\
 \mathfrak r_\theta={}&\mathcal P_{-A-1}
       (M_{-A}E+V_0E/(2X)-\nu R_1E)
       -\nu\mathcal P_{-A-1+2h}Z_{-A-D}Z_{-A}E,\\
 \mathfrak r_z={}&\mathcal P_{-A-1}
       (M_{-A}U-\nu RU+Z_{-2A}\Pi)
       -\nu\mathcal P_{-A-1+2h}Z_{-A-D}Z_{-A}U.
\end{align*}
To verify these formulas, differentiation of the cylindrical unit
vectors gives the radial convective term $-(u_\theta^{(0)})^2/\rho$,
the azimuthal convective term $u_\rho^{(0)}u_\theta^{(0)}/\rho$,
and the radial and azimuthal vector-Laplacian terms
$-u_\rho^{(0)}/\rho^2$ and $-u_\theta^{(0)}/\rho^2$.
The radial convective term cancels $\partial_\rho p^{(0)}$ by
the exact displayed pressure balance. The azimuthal additional term
is $\mathcal P_{-A-1}(V_0E/(2X))$.
Finally $-2A-D=-A-1$ gives the axial pressure weight. Substituting
\eqref{eq:ns-profile-operators-exact} proves every remaining term.
No assertion that this residual vanishes is needed or made.

The source constructs its leading profiles at viscosity one. Its
actual passage to arbitrary $\nu>0$ is the spatial map
\[
 y=x/\sqrt\nu,\qquad
 u_\nu(x,t)=\sqrt\nu\,u(y,t),\quad
 p_\nu(x,t)=\nu p(y,t),\quad
 f_\nu(x,t)=\sqrt\nu\,f(y,t).
\]
Every momentum term acquires the factor $\sqrt\nu$: time
differentiation preserves the velocity prefactor; the two velocity
factors in advection and the spatial derivative give
$\nu\nu^{-1/2}$; viscosity and two spatial derivatives give
$\nu\sqrt\nu\nu^{-1}$; and the pressure gradient gives
$\nu\nu^{-1/2}$. The divergence is unchanged. The Jacobian
$dx=\nu^{3/2}dy$ gives
$\|u_\nu\|_2^2=\nu^{5/2}\|u\|_2^2$ and
$\nu\int\|\nabla u_\nu\|_2^2dt
=\nu^{5/2}\int\|\nabla u\|_2^2dt$ on the same time interval.
In these rescaled fields the physical coordinates satisfy
$z=\sqrt\nu Q^D\xi$, $\rho^2=2\nu QX$ and
$\tau=Q(1-\xi^2)$. Thus $\partial_z$ gains $\nu^{-1/2}$
and the radial scalar Laplacian gains $\nu^{-1}$ in
\eqref{eq:ns-profile-operators-exact}. This is the source's specified
viscosity transport, with its time coordinate unchanged.

\subsection{The material base and the complete two-variable conic connection}
Work first over $\mathbb C$. Retain the original coordinates $s,b,c,r$
and $P=cr^3-2r^2+br-4s$. The original material base derivation and
its simple-root lift are
\[
 U_B=\partial_s+6c\partial_b,\qquad
 \mathcal D=\partial_s+6c\partial_b
        +\frac{4-6cr}{3cr^2-4r+b}\partial_r.
\]
The three contributions to $\mathcal DP$ are $-4$, $6cr$, and
$4-6cr$, so this is a derivation of the localization at $P_r$.
Under the real base identification
\begin{equation}\label{eq:ns-original-base-identification}
 s=\tau=1-t,\qquad b=z,\qquad c=c,
\end{equation}
$U_B$ becomes $-\partial_t+6c\partial_z$. This identity concerns
the stated material base; the leading physical material derivative
has instead the explicitly proved profile formula
\eqref{eq:ns-source-material-profile}. On $c=0$ the original
arithmetic operator is exactly $-\partial_s^2/4=-\partial_t^2/4$.
The previously established pullback by the original polynomial $F$
composes with \eqref{eq:ns-original-base-identification}; its two
original Euclidean diffusion coefficients are unchanged.

On $c=0$, retain the complete cover
\begin{gather*}
 \beta=b/4,\qquad \eta=r-\beta,\qquad
 m=b^2/16-2s,\qquad n=m+4,\\
 \mathcal B=\mathbb C[s,b,(mn)^{-1}],\qquad
 \mathcal E=\mathcal B[\eta,p]/(\eta^2-m,p^2-n).
\end{gather*}
The original branch factors are $\delta=16m=b^2-32s$ and
$L=n=4+b^2/16-2s$. The localization removes both and only these
two factors. The ring is free of rank four, with bases
$(1,\eta,p,\eta p)$ and $(1,r,p,rp)$; monic division proves
existence and uniqueness of the four coefficients. There are inverse
maps to the localized original Laurent presentation, given by
\begin{equation}\label{eq:ns-conic-laurent-map}
 q=(p+\eta)/2,\quad q^{-1}=(p-\eta)/2,
 \quad\eta=q-q^{-1},\quad p=q+q^{-1},\quad
 s=b^2/32-\tfrac12(q-q^{-1})^2.
\end{equation}
The two displayed images of $q,q^{-1}$ have product
$(n-m)/4=1$; sums and differences prove both inverse compositions.
Thus the arithmetic $q$ in this formula is never set equal to $Q$.

There are unique commuting derivations $\nabla_s,\nabla_b$ extending
$\partial_s,\partial_b$ to $\mathcal E$. On generators they are
\begin{align*}
 \nabla_s\eta&=-1/\eta,&\nabla_sp&=-1/p,&
 \nabla_sq&=-q/(\eta p),\\
 \nabla_b\eta&=b/(16\eta),&\nabla_bp&=b/(16p),&
 \nabla_bq&=bq/(16\eta p).
\end{align*}
Differentiating the two defining quadratic relations proves these
formulas and their uniqueness, since $2\eta$ and $2p$ are units.
Their commutator is a derivation killing $s,b$; differentiation of
the same quadratic relations then forces it to kill $\eta,p$.
This proves flatness without assuming that arbitrary differentiation
descends to an unlocalized quotient.

Here are the full matrices in the retained basis $(1,r,p,rp)$.
Column vectors are coefficient vectors, and
$\nabla_i=\partial_iI_4+\Gamma_i$:
\begin{align*}
 \Gamma_s&=\begin{pmatrix}
 0&\beta/m&0&0\\0&-1/m&0&0\\
 0&0&-1/n&\beta/m\\0&0&0&-1/m-1/n
 \end{pmatrix},\\
 \Gamma_b&=\begin{pmatrix}
 0&1/4-b\beta/(16m)&0&0\\0&b/(16m)&0&0\\
 0&0&b/(16n)&1/4-b\beta/(16m)\\
 0&0&0&b/(16m)+b/(16n)
 \end{pmatrix}.
\end{align*}
For example $\nabla_br=1/4+b\eta/(16m)$ and
$\nabla_b(rp)=p/4+(b/(16m)+b/(16n))\eta p
+b\beta p/(16n)$, which give the second and fourth columns.
The other columns follow directly from the generator formulas.
The inverse change of coefficient bases is explicit:
\[
 T=\begin{pmatrix}1&\beta&0&0\\0&1&0&0\\
 0&0&1&\beta\\0&0&0&1\end{pmatrix},\qquad
 T^{-1}=\begin{pmatrix}1&-\beta&0&0\\0&1&0&0\\
 0&0&1&-\beta\\0&0&0&1\end{pmatrix}.
\]
In the $(1,\eta,p,\eta p)$ basis write
\[
 H_s=\operatorname{diag}(0,-m^{-1},-n^{-1},-m^{-1}-n^{-1}),
 \qquad H_b=-\tfrac b{16}H_s.
\]
Then $\Gamma_i=T^{-1}H_iT+T^{-1}(\partial_iT)$, which
also verifies that the $b$-dependent basis derivative has been retained.
For clarity the entire second axial derivative in this basis is
\[
 \nabla_b^2=\partial_b^2I_4+2H_b\partial_b+J_b,
\]
where $J_b$ is diagonal with entries
\begin{align*}
 (J_b)_{00}&=0,\\
 (J_b)_{11}&=\frac1{16m}-\frac{b^2}{256m^2},\qquad
 (J_b)_{22}=\frac1{16n}-\frac{b^2}{256n^2},\\
 (J_b)_{33}&=\frac1{16}\left(\frac1m+\frac1n\right)
 -\frac{b^2}{256}\left(\frac1{m^2}+\frac1{n^2}\right)
 +\frac{b^2}{128mn}.
\end{align*}
These entries are $\partial_bH_b+H_b^2$, using
$\partial_bm=\partial_bn=b/8$. Similarly
\[
 \nabla_s^2=\partial_s^2I_4+2H_s\partial_s+J_s,
 \quad J_s=\operatorname{diag}
 \left(0,-m^{-2},-n^{-2},-m^{-2}-n^{-2}+2/(mn)\right).
\]
The formulas follow from $\partial_sm=\partial_sn=-2$.
Returning to the original basis replaces each second-order constant
matrix by $\partial_i\Gamma_i+\Gamma_i^2$ and each first-order
matrix by $2\Gamma_i$, so no basis-dependent second derivative is lost.

Let the physical real base be
\[
 \Omega=\{(\rho,z,t):\rho>0,\ t<1,\ m=z^2/16-2(1-t),\ mn\ne0\}.
\]
Tensoring $\mathcal E$ over its indicated base map with complex smooth
functions on $\Omega$ gives a rank-four complex algebra bundle.
On this bundle physical $\partial_t=-\nabla_s$,
$\partial_z=\nabla_b$, and $\partial_\rho$ differentiates the smooth
coefficients and kills $\eta,p$. Consequently its complete scalar
heat and arithmetic operators are
\begin{equation}\label{eq:ns-flat-heat-connection}
 \mathcal N_\nu=-\nabla_s-\nu\bigl(
 \partial_\rho^2+\rho^{-1}\partial_\rho+\nabla_b^2\bigr),
 \qquad \mathcal H=-\tfrac14\nabla_s^2,
 \qquad [\mathcal N_\nu,\mathcal H]=0.
\end{equation}
The commutator vanishes because the two lifted derivations commute,
all transverse derivatives commute with them, and $\nu$ is constant.
The coefficient inclusion $f\mapsto f\,1$ is an algebra map and
intertwines every displayed physical derivative. Inserting the actual
profiles into it therefore carries \eqref{eq:ns-profile-heat-full},
\eqref{eq:ns-profile-arithmetic-full}, and the three exact residual
identities to this bundle. This is the explicit differential
correspondence; it does not identify the four cover coordinates with
four independently constructed physical velocities.

\subsection{The physical real structure and a finite negative coefficient}
On the open physical region $-4<m<0$, put $a_*=\sqrt{-m}$ and
$p_*=\sqrt{m+4}$ using the positive real square roots.
The four sections of the complex cover are exactly
\[
 \eta=\epsilon i a_*,\quad p=\epsilon' p_*,\quad
 r=b/4+\epsilon i a_*,\quad
 q=\tfrac12(\epsilon'p_*+\epsilon i a_*),\qquad
 \epsilon,\epsilon'\in\{1,-1\}.
\]
Their squared moduli are $(m+4-m)/4=1$, and the maps
\eqref{eq:ns-conic-laurent-map} verify every equation. The four
sections are distinct since $a_*,p_*>0$. Ordinary complex conjugation
fixes $s,b$, fixes $p$, changes $\eta$ to $-\eta$, and sends $q$ to
$q^{-1}$. This is the real structure over the physical real base.
There is no real value of $q$ on this region: a real nonzero $q$ would
give $m=(q-q^{-1})^2\geq0$. The branch boundaries $m=0$ and $m=-4$
remain excluded from the connection, even though limits of the scalar
polynomials below exist there.

On the source's actual growth path $z=0$ and
$\rho=\sqrt{2X_{\rm in}\tau}$, retain $Q=\tau$, $\xi=0$,
$X=X_{\rm in}>0$. The base map gives $b=0,s=\tau$, so
$m=-2\tau$. For $0<\tau<2$ choose the sheet
\begin{equation}\label{eq:ns-unit-circle-growth-map}
 \theta=\arcsin\sqrt{\tau/2},\quad
 q=e^{i\theta},\quad \eta=i\sqrt{2\tau},\quad
 p=\sqrt{4-2\tau},\quad r=i\sqrt{2\tau}.
\end{equation}
Here $\theta\in(0,\pi/2)$ and the formulas agree with the four-sheet
description. In particular the source concentration $Q\downarrow0$
maps to the finite arithmetic point $q\to1$, with the original branch
point retained as a limit.

For the original GCT coefficient, every factor is retained:
\begin{align}
 C(q)&=q^{10}-q^4-q^{-4}+q^{-10}
       =(q-q^{-1})^2[7]_q[3]_q,\notag\\
 C(e^{i\theta})&=-4\sin(7\theta)\sin(3\theta),\notag\\
 C\big|_{b=0,s=\tau}
 &=-2\tau(3-2\tau)(7-28\tau+28\tau^2-8\tau^3)
   =\tau B(\tau),\label{eq:ns-finite-negative-polynomial}\\
 B(\tau)&=-42+196\tau-280\tau^2+160\tau^3-32\tau^4.
       \notag
\end{align}
The trigonometric equality follows by substituting $q=e^{i\theta}$
and applying $2\cos(10\theta)-2\cos(4\theta)
=-4\sin(7\theta)\sin(3\theta)$. For the polynomial equality,
$[3]_q=p^2-1$ and $[7]_q=p^6-5p^4+6p^2-1$ follow by multiplying
the corresponding finite Laurent sums by $q-q^{-1}$. Substituting
$p^2=4-2\tau$ and $(q-q^{-1})^2=-2\tau$ gives the displayed
factored polynomial and its full expansion. Therefore
\[
 C(q)<0\quad\hbox{for}\quad
 0<\tau<2\sin^2(\pi/7),\qquad
 \lim_{\tau\downarrow0}\frac{C(q)}\tau=-42.
\]
Both sines are strictly positive on the stated interval because
$0<\theta<\pi/7$ and $0<3\theta<3\pi/7<\pi$.
This proves a finite negative value and its exact vanishing order.

The precise scalar-germ transport of this coefficient is also explicit.
Fix $0<h<1/100$, $e_0>0$, and any bounded real function
$\varepsilon$ on $(0,\tau_0)$. Define the actual scalar function
\[
 g_\varepsilon(\tau)=\tau^{-A}
          (e_0+\tau^{2h}\varepsilon(\tau)).
\]
Multiplication by the computed GCT coefficient gives the exact identity
\begin{equation}\label{eq:ns-negative-germ-transport}
 C(q)g_\varepsilon(\tau)
 =\tau^D B(\tau)(e_0+\tau^{2h}\varepsilon(\tau)),\qquad
 \lim_{\tau\downarrow0}\tau^{-D}C(q)g_\varepsilon(\tau)=-42e_0.
\end{equation}
This follows from $1-A=D$ and \eqref{eq:ns-finite-negative-polynomial}.
Since $B(\tau)+42$ is divisible by $\tau$, and $0<2h<1$, the
right side equals $-42e_0\tau^D+O(\tau^{D+2h})$.
It is negative for sufficiently small $\tau$ and tends to zero,
since $D>0$. Multiplication raises the displayed leading exponent by
exactly one. In the formal power-series ring $\mathbb R[[\tau]]$,
the same exact polynomial has order one, because $B(0)=-42\ne0$.
No identification of $\tau$ with a prime number or of this order with
a prime-adic valuation is made.

The source's equation (10.21) reports its completed velocity component
as $\tau^{-A}(e_0+O(\tau^{2h}))$ on this path. The exact map
\eqref{eq:ns-negative-germ-transport} acts on the full stated class of
such scalar germs; identifying a member with that completed solution
retains the independent-verification status of the source theorem.
The negative coefficient result above uses only the proved conic map
and polynomial identities. It gives no asymptotic circuit-size bound,
plethysm multiplicity obstruction, or nonstandard-RH eigenvalue claim.

Finally, the source pressure normalization supplies a second exact
sign statement with a different specified observable:
$\Pi(X,\xi)=-\int_X^\infty E(y,\xi)^2/(2y)\,dy\leq0$
whenever that integral is finite. It is strictly negative if the
nonnegative integrand is positive on a set of positive measure.
Replacing $p$ by $p+k(t)$ leaves its spatial gradient and the momentum
equation unchanged, but need not preserve a chosen compact-support or
infinity normalization. Thus the displayed pressure sign belongs to
that normalization. The coefficient sign
\eqref{eq:ns-finite-negative-polynomial} has instead been proved for
the fixed original Laurent polynomial on the stated unit-circle sheet.


\section{The concentration path through the actual Schur conductor map}
\label{sec:ns-schur-path}

This section composes the previously proved conic path with the integral
Schur endpoint, retaining the original source and target bases. Put
\[
 I=(0,2\sin^2(\pi/7)),\qquad
 B(\tau)=-42+196\tau-280\tau^2+160\tau^3-32\tau^4,
\]
and use exactly the sheet in \eqref{eq:ns-unit-circle-growth-map}.
The concentration parameter is $Q=\tau$, whereas the GCT parameter is
$q=e^{i\arcsin\sqrt{\tau/2}}$. These parameters remain distinct.
The endpoint factorization of Section~\ref{sec:plethysm-endpoint-transport}
pulls back to the three based maps
\begin{align}
 D_\tau&=\operatorname{diag}(1,1,1,\tau B(\tau)),\notag\\
 H_\tau&=\operatorname{diag}(1,1,1,h_\tau),\qquad
 V_\tau=\operatorname{diag}(1,1,1,-2\tau),\notag\\
 h_\tau&=-B(\tau)/2
 =21-98\tau+140\tau^2-80\tau^3+16\tau^4,\qquad
 D_\tau=V_\tau H_\tau.\label{eq:ns-schur-original-factorization}
\end{align}
Here $h_\tau$ denotes the retained conductor factor $[7]_q[3]_q$,
and is not the source profile parameter $h$. Indeed $r^2=-2\tau$
and $C=r^2[7]_q[3]_q=\tau B(\tau)$ prove every entry.
On $I$ both $C$ and $r^2$ are negative, so $h_\tau>0$.
The pointwise maps have their earlier endpoint, intermediate, and node modules
as source and target after this base change; the displayed matrices
do not identify those three original modules.

For the formal series statement below, the exact ring map must also be
specified. Set $R_m=\mathbb Z[m]$ and
\[
 h_m=(m+3)\bigl((m+4)^3-5(m+4)^2+6(m+4)-1\bigr),
 \qquad C_m=mh_m.
\]
Take three free $R_m$-modules with the respective original endpoint,
intermediate, and node basis labels. Define their maps to have matrices
$\operatorname{diag}(1,1,1,h_m)$ and
$\operatorname{diag}(1,1,1,m)$. Their composition has fourth entry
$C_m$. The homomorphism $R_m\to R=\mathbb Z[1/2][q,q^{-1}]$ sending
$m$ to $(q-q^{-1})^2$ takes $m+4$ to $(q+q^{-1})^2$ and hence
$h_m$ to $[7]_q[3]_q$. Mapping each free basis symbol to its stated
original basis vector is an isomorphism after this base change: the
inverse sends the original coordinate tuple back to the same symbols.
The three matrices show that these isomorphisms intertwine every map.
The same argument holds after the localizations used earlier.
This proves descent of this based endpoint diagram to $R_m$.
Its separate base change $R_m\to\mathbb R[[\tau]]$, $m\mapsto-2\tau$,
is the formal diagram used here. In particular, no homomorphism sending
the full Laurent variable $q$ to an element of $\mathbb R[[\tau]]$ is
asserted: the pointwise sheet has a square-root term at $\tau=0$.
The complete Schur base-change proof carries these diagram
isomorphisms through $F$ without changing any tableau labels.

Let $F=S_\lambda$, or $F=S_\lambda\circ S_\mu$, be the precise
integral Schur functor constructed above. Write $d=|\lambda|$ in
the first case and $d=|\lambda||\mu|$ in the second. Its tableau
weight multiplicities are the previously constructed $m_F(k)$.
Functoriality and the full diagonal formula give
\begin{equation}\label{eq:ns-schur-path-diagonal}
 \begin{aligned}
 F(D_\tau)&=\operatorname{diag}_{k,j}
       \bigl(\tau^k B(\tau)^k\bigr),\\
 F(H_\tau)&=\operatorname{diag}_{k,j}(h_\tau^k),\qquad
 F(V_\tau)=\operatorname{diag}_{k,j}((-2\tau)^k),\\
 &1\leq j\leq m_F(k).
 \end{aligned}
\end{equation}
For a single tableau, $k$ is its number of fourth symbols. In the
nested case, $k$ is the sum of those counts over the inner tableaux
appearing in the outer tableau, with multiplicity. Substitution of the fourth diagonal entry therefore
gives its $k$th power; the other three diagonal entries are one.
This also proves the displayed composition entry by entry.

Over $\mathbb R[[\tau]]$, $B$ is a unit because its constant term
is $-42$. An inverse is obtained recursively: if
$B=\sum_{i=0}^4b_i\tau^i$, set $a_0=-1/42$ and
$a_n=-b_0^{-1}\sum_{i=1}^{\min(4,n)}b_i a_{n-i}$.
Thus the exact image in the $(k,j)$ target coordinate is
$\tau^k\mathbb R[[\tau]]$, its kernel is zero, and its cokernel
is $\mathbb R[[\tau]]/(\tau^k)$. The latter is zero when $k=0$;
for $k>0$ its basis over $\mathbb R$ is $1,\tau,\ldots,\tau^{k-1}$,
as follows by separating the first $k$ coefficients of a series.
Consequently the cokernel length and determinant order both equal
$\sum_k k m_F(k)$. Specialization at $\tau=0$ has rank $m_F(0)$
and kernel and cokernel dimension $\sum_{k>0}m_F(k)$.
These are statements over the specified real series ring. The integer
factors $21^k$ and residual characteristics in the earlier integral
calculation are retained there; they have not been replaced by assertions
about valuations of prime numbers along a real path.

For each $\tau\in I$, identify source and target tableau labels by
their explicitly matching based-vector-space map and give the source
the Hermitian metric for which that basis is orthonormal. The form
defined by $F(D_\tau)$ then has coefficients $\tau^kB(\tau)^k$.
Since $B(\tau)<0$ on $I$, the even-weight and odd-weight coordinate
spaces are respectively positive and negative definite. They are
orthogonal and exhaust the space. A larger positive subspace would
intersect the negative coordinate space by the dimension formula,
and conversely for a larger negative subspace. Thus throughout $I$
the exact inertia is
\[
 \left(\sum_{k\ \text{even}}m_F(k),
       \sum_{k\ \text{odd}}m_F(k),0\right).
\]
The factor $F(H_\tau)$ is positive definite in the matched metric,
so the same indices hold for $F(V_\tau)$.

For the original shape $\nu=(7,5,3)$, the completely evaluated profile is
\[
 (m_\nu(0),\ldots,m_\nu(7))
 =(27,81,162,270,300,252,126,42).
\]
Its sum is $1260$ and its weighted sum is $4725$. It follows that
\begin{equation}\label{eq:ns-schur-source-path-data}
 \det S_\nu(D_\tau)=\tau^{4725}B(\tau)^{4725},\qquad
 \operatorname{inertia}=(615,645,0),\qquad
 \operatorname{rank}_{\tau=0}=27.
\end{equation}
The specialized kernel and cokernel have dimension $1233$; the series
cokernel has length $4725$. The endpoint entries extend to $\tau=0$,
although the conic connection does not extend across its excluded
branch locus. No connection value at that branch point is used here.

There is a second exact composition that tests what happens to growth.
Use the actual scalar functions of \eqref{eq:ns-negative-germ-transport}:
fix $0<h<1/100$, $A=1/2+h$, $e_0>0$, and bounded real $\varepsilon$,
and set $g(\tau)=\tau^{-A}(e_0+\tau^{2h}\varepsilon(\tau))$.
Define the explicit new linear map
\[
 L_\tau=g(\tau)D_\tau
\]
between the same two based endpoint spaces. This definition attaches
the stated scalar germ to their map; it does not identify an endpoint
vector with a fluid velocity or a quantum state. Each tableau monomial
has total degree $d$, including every inner and outer box in the nested
case, whence
\begin{equation}\label{eq:ns-schur-growth-exponents}
 F(L_\tau)_{k,j}
 =g(\tau)^d C(q)^k
 =\tau^{k-dA}B(\tau)^k
       (e_0+\tau^{2h}\varepsilon(\tau))^d.
\end{equation}
This proves homogeneity without a division by a factorial. For fixed
$k,d$, the binomial formula and boundedness give
$(e_0+\tau^{2h}\varepsilon)^d=e_0^d+O(\tau^{2h})$.
Also $B(\tau)^k=(-42)^k+O(\tau)$ for $k>0$, and the formula is
exactly one for $k=0$. Since $2h<1$, multiplication yields
\[
 F(L_\tau)_{k,j}
 =(-42)^k e_0^d\tau^{k-dA}
       +O(\tau^{k-dA+2h}).
\]
The leading constant is nonzero. Hence the absolute value tends to
zero when $k>dA$, tends to the absolute leading constant when $k=dA$,
and diverges when $k<dA$. This classification proves the outcome for
the defined map, rather than assuming that scalar multiplication and
a degree-$d$ functor have the same effect on an exponent.

In particular $d=15$ and $0\leq k\leq7$ for $\nu=(7,5,3)$, so
$k-15A\leq-1/2-15h<0$. Every one of its $1260$ diagonal entries
diverges in absolute value, with $615$ positive and $645$ negative
entries for sufficiently small $\tau$. Its determinant is exactly
\[
 \det S_\nu(L_\tau)
 =\tau^{4725-18900A}B(\tau)^{4725}
       (e_0+\tau^{2h}\varepsilon(\tau))^{18900},
\]
with leading exponent $-4725-18900h$. In contrast, the already proved
single scalar product $C(q)g(\tau)$ tends to zero with exponent
$1/2-h$. The displayed functorial computation explains this difference:
the degree-fifteen scalar contributes exponent $-15A$, and the largest
available fourth-coordinate weight contributes only $7$.

All multiplicities here are the unchanged nonnegative tableau counts.
The rational source comparison at $q=1$ is the explicit map already
proved in Section~\ref{sec:plethysm-endpoint-transport}; this path
calculation does not create a generic quantum-group intertwiner or an
unrestricted complexity lower bound. It supplies the exact Schur
composition, signs, and growth exponents on the specified original
endpoint instead.

\section{The retained conic and the interacting material-tensor state map}
\label{sec:ym-retained-composition}

The local source \texttt{FABEL\_TENSOR\_TRANSFER.md}, equations
(97)--(105), supplies the full material tensor and finite Wilson box.
The source paths and hashes accompany this fragment. The NS input
is the explicit conic coordinates and scalar polynomial proved in
Section~\ref{sec:ns-conic-transport}. The published manuscript's
global blowup theorem retains its stated independent-verification
status; it is not used as a theorem in this calculation.

\subsection{Original Hamiltonian and complex six-state map}

Fix an integer $L\ge2$, spacing $a_{\rm YM}>0$, and coupling $g>0$.
Here $a_{\rm YM}$ is the source's spacing $a$, not the arithmetic
coefficient $a=2s$. Vertices are $\{-L,\ldots,L\}^3$; include every
positive nearest-neighbour edge $\mathsf E_L$ and every elementary
face $\mathsf P_L$ whose vertices are present, including boundaries.
On $\mathcal Q_L=\operatorname{SU}(2)^{\mathsf E_L}$ use product
Haar probability $\lambda$. With $T_a=-i\sigma_a/2$ and
$X_{e,a}F=\left.\frac{d}{du}F(\ldots,e^{uT_a}U_e,\ldots)\right|_{u=0}$,
retain
\begin{gather*}
 E_e=-\sum_{a=1}^3X_{e,a}^2,\quad
 W_p=\operatorname{tr}\bigl(U_i(n)U_j(n+\mathbf e_i)
 U_i(n+\mathbf e_j)^{-1}U_j(n)^{-1}\bigr),\\
 H=\kappa\sum_eE_e+\sum_pV_p,\qquad
 \kappa=2g^2/a_{\rm YM},\qquad
 V_p=b_{\rm YM}(2-W_p),\quad b_{\rm YM}=1/(2g^2a_{\rm YM}).
\end{gather*}
The gauge group is
$\mathcal G=\operatorname{SU}(2)^{\{-L,\ldots,L\}^3}$, acting by
$(h\cdot U)_e=h_{s(e)}^{-1}U_eh_{t(e)}$.
Its Hilbert space is $\mathscr H_L=L^2(\mathcal Q_L,\lambda)$
and its domain is $H^2(\mathcal Q_L)$. In particular its original
constant $2b_{\rm YM}|\mathsf P_L|$ is retained. Minimizing the
closed elliptic form, using compact Sobolev embedding, gives a
ground vector. Taking its modulus does not increase the form;
elliptic regularity and the strong maximum principle make the
ground vector smooth and strictly positive. Denote it by $\psi$,
with $H\psi=E_0\psi$ and $\|\psi\|=1$.
Integration by parts gives, for smooth $u$,
\[
 \langle\psi u,(H-E_0)\psi u\rangle
 =\kappa\sum_{e,a}\int\psi^2|X_{e,a}u|^2\,d\lambda .
\]
This vanishes precisely for constant $u$, by connectedness, so
$\ker(H-E_0)=\mathbb C\psi$. Gauge invariance of $H$ and uniqueness
of its positive unit ground vector prove gauge invariance of $\psi$.
Write $\mathcal A=H-E_0$.

For complex edge weights define on the common smooth domain
\[
 f_p=\tfrac14\sum_{e\in\partial p}f_e,\qquad
 D_f=\kappa\sum_ef_eE_e+\sum_pf_pV_p,\qquad
 \Xi_f=D_f\psi-\langle\psi,D_f\psi\rangle\psi .
\]
Inner products are conjugate linear in the first argument.
Thus $\Xi_f$ is complex linear in $f$, smooth, gauge invariant,
and perpendicular to $\psi$. It belongs to every power domain of
$H$. For complex $f$ this uses
$D_f=D_{\operatorname{Re}f}+iD_{\operatorname{Im}f}$ on smooth
functions; no self-adjointness of this complex weighted operator
is asserted. The physical $H$ and its adjoint remain unchanged.
On this smooth domain the ground equation gives
$\mathcal A\Xi_f=[H,D_f]\psi$ by direct subtraction.

Use the original dimensionless midpoint $m(e)=n+\mathbf e_i/2$.
The exact complexification of the source's real map is
\[
 \mathfrak M:\operatorname{Sym}_3(\mathbb C)\longrightarrow
 \mathscr V_L:=C^\infty(\mathcal Q_L)^{\mathcal G}\cap\psi^\perp,
 \qquad \mathfrak M(S)=\Xi_{m^{\mathsf T}Sm}.
\]
Physical midpoints $a_{\rm YM}m(e)$ multiply this map by
$a_{\rm YM}^2$ and its raw spectral measures by $a_{\rm YM}^4$.
For $v\in\mathscr V_L$ put
$\nu_v(B)=\langle v,1_B(\mathcal A)v\rangle$ on $(0,\infty)$.
Set $u_i=\Xi_{m_i^2}$, $v_{ij}=\Xi_{2m_im_j}$ and
\[
 \nu_0=\nu_{u_1+u_2+u_3},\qquad
 \nu_{\rm d}=\nu_{u_1-u_2},\qquad \nu_{\rm o}=\nu_{v_{12}}.
\]
For every complex symmetric $S$, the full measure is
\begin{equation}\label{eq:ym-complex-measure}
 \nu_{\mathfrak M(S)}
 =\frac{|\operatorname{tr}S|^2}{9}\nu_0
 +\frac12\sum_i|S_{ii}-\operatorname{tr}S/3|^2\nu_{\rm d}
 +\sum_{i<j}|S_{ij}|^2\nu_{\rm o}.
\end{equation}
To prove this with conjugation retained, signed coordinate
permutations act on links by their endpoint permutation and by
inversion when an edge reverses. Haar measure and the Casimir are
preserved, and the SU(2) face trace is unchanged by reversal.
Their unitaries commute with $H$ and its spectral projections and
fix $\psi$. Reflections kill every mixed $u,v$ spectral Gram entry
and every different-pair $v,v$ entry; permutations equate the
remaining $v,v$ entries. The $u,u$ Gram matrix has constant
diagonal $d$ and constant off-diagonal $h$; a transposition and
Hermitian symmetry give $h=\overline h$. Its quadratic form for
complex diagonal coefficients $s_i$, with mean
$s_{\rm av}=(s_1+s_2+s_3)/3$, is
\[
 (d-h)\sum_i|s_i-s_{\rm av}|^2
 +(d+2h)|s_1+s_2+s_3|^2/3 .
\]
Since $\nu_0=3d+6h$ and $\nu_{\rm d}=2(d-h)$, this proves
\eqref{eq:ym-complex-measure} on every Borel set. Integration
against $1,E,E^2,e^{-uE}$, with $u\ge0$, gives norms, excitation forms, squared
current norms, and time correlations. No real square has been
continued through a Hilbert adjoint.

The kernel of this six-state map is also exact. Decompose
$\operatorname{Sym}_3(\mathbb C)=\mathcal S_0\oplus
\mathcal S_{\rm d}\oplus\mathcal S_{\rm o}$ into scalar,
traceless diagonal, and off-diagonal matrices, of dimensions
$1,2,3$. Set $d_j=\nu_j((0,\infty))$. Since the state is
perpendicular to the entire zero eigenspace, its norm squared
is the total mass in \eqref{eq:ym-complex-measure}. Each of the
three coefficient forms is strictly positive on every nonzero
matrix in its corresponding channel. Nonnegativity of the
three terms therefore proves
\[
 \ker\mathfrak M=\bigoplus_{j:\ d_j=0}\mathcal S_j,\qquad
 \operatorname{rank}\mathfrak M
 =\mathbf1_{d_0>0}+2\mathbf1_{d_{\rm d}>0}
                         +3\mathbf1_{d_{\rm o}>0}.
\]
This evaluates all possible kernel directions without asserting
an unknown seed mass positive.

\subsection{The complete original tensor on the common conic}

Write $z_F$ for the Fabel parameter
$z=\sqrt{1-8\tau_F}>0$, distinct from the NS axial coordinate.
The unchanged polynomial is
\begin{align*}
 F_1(x,y,w)&=(1+xy)^3w+y^2(1+xy)(4+3xy),\\
 F_2(x,y,w)&=y+3x(1+xy)^2w+3xy^2(4+3xy),\\
 F_3(x,y,w)&=2x-3x^2y-x^3w .
\end{align*}
With $J=DF$, $x_0=(1,-3/2,13/2)$ and
$\gamma(z_F)=(z_F^{-1},-3z_F/2,13z_F^2/2)$, differentiation gives
$\det J=-2$ and
\[
 J_0=\begin{pmatrix}-9/16&-3/8&-1/8\\3/8&25/4&3/4\\-17/2&-3&-1\end{pmatrix},
 \quad J_\gamma=\begin{pmatrix}
 -9z_F^3/16&-3z_F/8&-1/8\\3z_F^2/8&25/4&3/(4z_F)\\
 -17/2&-3/z_F^2&-1/z_F^3
 \end{pmatrix}.
\]
Put $\tau_F=(1-z_F^2)/8$ and $B_{\tau_F}=I+6\tau_FE_{23}$.
The full tensor is
\begin{align}
 C_F(z_F)&=J_0^{-1}B_{\tau_F}^{-1}J_\gamma,\notag\\
 &=\begin{pmatrix}
 (17-9z_F^3)/8&(3-3z_F^3)/(4z_F^2)&(1-z_F^3)/(4z_F^3)\\
 (51-24z_F^2-27z_F^3)/16&(9+8z_F^2-9z_F^3)/(8z_F^2)&(3-3z_F^3)/(8z_F^3)\\
 (-153+36z_F^2+117z_F^3)/8&(-27-12z_F^2+39z_F^3)/(4z_F^2)&(-9+13z_F^3)/(4z_F^3)
 \end{pmatrix}.\label{eq:ym-full-tensor}
\end{align}
Every entry follows by multiplying the displayed matrices, using
$B_{\tau_F}^{-1}=I-6\tau_FE_{23}$. Its inverse is
$J_\gamma^{-1}B_{\tau_F}J_0$, and its determinant is one.
Put
\[
 K_F=C_FC_F^{\mathsf T},\qquad
 v_*=(1/4,3/8,-9/4)^{\mathsf T},\qquad
 S_*=v_*v_*^{\mathsf T}
 =\frac1{64}\begin{pmatrix}4&6&-36\\6&9&-54\\-36&-54&324\end{pmatrix}.
\]
The transpose here is not an adjoint. For real $z_F\ne0$ this
$K_F$ is positive definite since $C_F$ is invertible. For complex
$z_F$ only complex symmetry is asserted. Each of the nine entries
in \eqref{eq:ym-full-tensor} gives the convergent Laurent limits
\begin{equation}\label{eq:ym-full-endpoint}
 z_F^3C_F\longrightarrow v_*e_3^{\mathsf T},\qquad
 z_F^6K_F\longrightarrow S_* .
\end{equation}

Substitution with $xy=-3/2$ gives exactly
\[
 F(\gamma(z_F))=(-z_F^2/4,0,0),\quad
 r=-z_F/2,\quad a=-r^2=-z_F^2/4,\quad s=-z_F^2/8 .
\]
Indeed $F_1=(-13+9)z_F^2/16$,
$F_2=(-12+39-27)z_F/8$ and
$F_3=(4+9-13)/(2z_F)$. This recovers the original inverse chart
$(x,y,w)=(-1/(2r),3r,26r^2)$; the separate boundary point
$(0,0,a)$ remains in the fibre. Choose the analytic germ
$p(z_F)=\sqrt{4+z_F^2/4}$ with $p(0)=2$. The exact conic maps are
\begin{equation}\label{eq:ym-germ-conic}
 q(z_F)=\frac{p(z_F)-z_F/2}{2},\qquad
 q(z_F)^{-1}=\frac{p(z_F)+z_F/2}{2}.
\end{equation}
Their product is one, their difference is $r=-z_F/2$, and their
sum is $p$. Thus $r^2=-2s$, $p^2=4-2s$ and $b=c=0$ retain the
earlier coordinate values. The connection is used only where
$rp\ne0$, in a sufficiently small punctured neighbourhood of zero.
All factors of the original GCT coefficient remain:
\begin{align}
 C(q)&=q^{10}-q^4-q^{-4}+q^{-10}\notag\\
 &=m(n^3-5n^2+6n-1)(n-1),\quad
 m=z_F^2/4=-2s,\quad n=m+4 .\label{eq:ym-coefficient-base}
\end{align}
For direct verification, multiplying the finite Laurent sums gives
$C=(q-q^{-1})^2[7]_q[3]_q$,
$[3]_q=p^2-1$ and $[7]_q=p^6-5p^4+6p^2-1$.
In particular $z_F^{-2}C(q(z_F))\to21/4$.
For real $z_F$ near zero, $q>0$ and $m>0$, so $C(q)>0$.
The NS growth conic instead has $s=\tau=1-t>0$ and
$r=i\sqrt{2\tau}$; there is no real $z_F$ with this value of $s$,
because $s=-z_F^2/8\le0$ for real $z_F$.

The strongest direct continuation on this conic is
\begin{equation}\label{eq:ym-complex-path}
 z_F=-i\sqrt{8\tau},\quad r=i\sqrt{2\tau},\quad
 p=\sqrt{4-2\tau},\quad
 q=\tfrac12\bigl(\sqrt{4-2\tau}+i\sqrt{2\tau}\bigr),
 \qquad 0<\tau\ll1 .
\end{equation}
It is exactly the NS conic point. Its Fabel time is
$\tau_F=1/8+\tau$, so it is not the original real material
trajectory before $\tau_F=1/8$. On this path,
\[
 C(q)=\tau B(\tau),\quad
 B(\tau)=-42+196\tau-280\tau^2+160\tau^3-32\tau^4 .
\]
This is the full factored polynomial
$-2\tau(3-2\tau)(7-28\tau+28\tau^2-8\tau^3)$.
It is negative for sufficiently small positive $\tau$, and
$B(0)=-42$.

\subsection{State and spectral limits with the adjoint unchanged}

Define the actual physical state
$\Psi_\tau=\mathfrak M(K_F(-i\sqrt{8\tau}))$.
The physical Hamiltonian and its real regulator have not been
continued. Since $(-i)^6=-1$ and $8^3=512$,
\eqref{eq:ym-full-endpoint} gives
\begin{equation}\label{eq:ym-state-limits}
 \tau^3\Psi_\tau\longrightarrow-\mathfrak M(S_*)/512,
 \qquad
 \tau^2C(q)\Psi_\tau\longrightarrow(21/256)\mathfrak M(S_*) .
\end{equation}
The second limit multiplies the first by $B(\tau)$.
Equivalently $z_F^4=64\tau^2$ and
$z_F^4 C(q)K_F\to(21/4)S_*$ prove it directly.
Convergence holds in every Sobolev norm and after applying
$\mathcal A$, since the states are combinations of six fixed
smooth states with converging scalar coefficients.

The real matrix $S_*$ has trace $337/64$. Its three mean-subtracted
diagonal entries are $-325/192$, $-310/192$, and $635/192$.
Its off-diagonal entries are $6/64$, $-36/64$, and $-54/64$.
Consequently \eqref{eq:ym-complex-measure} gives
\[
 \nu_*:=\nu_{\mathfrak M(S_*)}
 =\frac{113569}{36864}\nu_0
  +\frac{100825}{12288}\nu_{\rm d}
  +\frac{531}{512}\nu_{\rm o}.
\]
It follows that
\begin{equation}\label{eq:ym-spectral-limits}
 \tau^6\nu_{\Psi_\tau}\longrightarrow\frac1{262144}\nu_*,
 \qquad
 \tau^4\nu_{C(q)\Psi_\tau}\longrightarrow\frac{441}{65536}\nu_* .
\end{equation}
These are total-variation limits: in
\eqref{eq:ym-complex-measure} bound the variation of the difference
by the sum of the three coefficient errors times the finite masses
of the seed measures. Using instead their finite first or second
moments proves convergence of those moments. The identities
include zero seed states without assuming an unknown moment
nonzero or small.

More strongly, at every fixed $\tau$ with $C(q)\ne0$,
\[
 \nu_{C(q)\Psi_\tau}=|C(q)|^2\nu_{\Psi_\tau}.
\]
This follows from linearity of each spectral projection and the
fixed Hilbert adjoint. Thus spectral supports agree, and the
energy-to-norm quotients agree whenever the states are nonzero.
The negative scalar changes a phase; it scales every positive
spectral measure by a positive scalar. Since $\mathcal A\ge0$,
none of these states has negative excitation form. This proves
an obstruction for this specific proposed sign-to-energy map
while retaining its exact state and measure correspondences.

\subsection{Arithmetic differentiation and the full quantum source}

Let $\mathcal R=\mathbb C\{z_F\}[z_F^{-1}]$, the ring of convergent
meromorphic germs at zero. Equation \eqref{eq:ym-germ-conic}
defines homomorphisms
$A=\mathbb Z[q,q^{-1}]\to\mathcal R$ and
$D=\mathbb Q[q,q^{-1}]_{(q-1)}\to\mathcal R$: every denominator
nonzero at $q=1$ becomes a unit analytic germ.
This is a germ map, not simultaneous evaluation of every element
of $D$ on an entire fixed punctured interval.
The arithmetic base map $s\mapsto-z_F^2/8$ factors through its
original rank-four conic algebra, with $r=-z_F/2$ and $p=p(z_F)$.
In particular the exact common coefficient is
\[
 C(q)=-2s(7-28s+28s^2-8s^3)(3-2s).
\]
This is the same polynomial acting on the earlier arithmetic
operator $\widehat S$ with $s$ replaced there by $\widehat S$.
The coefficient-ring map asserts neither a representation of
that topological operator on $\mathscr H_L$ nor membership of
these $s$ in its spectrum.

The lifted base derivation on this slice is
\[
 \nabla_s=-4z_F^{-1}\frac{d}{dz_F},\qquad
 -\tfrac14\nabla_s^2=-4z_F^{-2}\frac{d^2}{dz_F^2}
                         +4z_F^{-3}\frac{d}{dz_F}.
\]
Differentiating $s=-z_F^2/8$, $r=-z_F/2$, and $p^2=4-2s$
gives $\nabla_ss=1$, $\nabla_sr=-1/r$, $\nabla_sp=-1/p$.
Also $q'(z_F)=-q/(2p)$, hence $\nabla_sq=-q/(rp)$.
Thus the derivations agree on every generator of the original
conic ring. Applying the first-order operator twice gives both
terms of the second-order operator, with their signs.
Its action on finite-range state germs commutes with the fixed
linear map $\mathfrak M$, for example
\[
 -\tfrac14\nabla_s^2\mathfrak M(K_F)
 =\mathfrak M(-4z_F^{-2}K_F''+4z_F^{-3}K_F').
\]
Under $s=1-t$ this retains the scalar $-\partial_t^2/4$.
It does not identify that operator with $H$, or delete radial and
axial terms of the full NS heat operator.

Extend the actual honest source lattice by
$X_{\mathcal R}=\mathcal R\otimes_A X_A$.
It remains free of rank $1260$, with its original four generator
matrices, divided powers and both degree-fifteen central characters.
The already proved actual source splitting and full tensor-square
involution extend through $D\to\mathcal R$, with both idempotents
and their exact ranks. Put
$\mathscr V_{\mathcal R}=\mathcal R\otimes_{\mathbb C}\mathscr V_L$;
these are algebraic finite-sum tensors, not a Hilbert completion.
There is a completely specified external tensor map
\begin{equation}\label{eq:ym-external-tensor}
 X_{\mathcal R}\otimes_{\mathcal R}\operatorname{Sym}_3(\mathcal R)
 \longrightarrow X_{\mathcal R}\otimes_{\mathcal R}\mathscr V_{\mathcal R},
 \qquad x\otimes S\longmapsto x\otimes\mathfrak M_{\mathcal R}(S).
\end{equation}
Give both second factors the trivial quantum action. The source
coproduct
$\Delta(E)=E\otimes K^{-1}+1\otimes E$ and
$\Delta(F)=F\otimes1+K\otimes F$ then acts only on $x$.
For each generator, both compositions in
\eqref{eq:ym-external-tensor} equal its original action on $x$
tensored with $\mathfrak M_{\mathcal R}(S)$. This proves the
intertwiner, its divided-power identities and its central labels;
tensor balancing follows from $\mathcal R$-linearity.
Its kernel is exactly
$X_{\mathcal R}\otimes_{\mathbb C}\ker\mathfrak M$, included by
$x\otimes S\mapsto x\otimes(1\otimes S)$:
the complex six-dimensional domain splits as the kernel plus
a complementary vector space on which $\mathfrak M$ is
injective; extending scalars preserves this splitting and its
inverse on the image, and $X_{\mathcal R}$ is free.
Inserting $K_F$ and the actual scalar $C(q)$ composes this map
with the state and measure calculations above. Replacing
$X_{\mathcal R}$ by either extended tensor-square idempotent
image gives the same proved map. No quantum action on the YM
Hilbert factor, identification $X\to\mathscr H_L$, mass gap,
canonical-basis construction or complexity separation is inferred.

\section{Research provenance, failed constructions, and public workbenches}
\label{sec:research-provenance}

This manuscript records a human-directed Codex investigation. The human
participant proposed the cross-problem directions and specifically suggested
looking for negativity at a finite endpoint. Their initial wording included:
\begin{quote}
``So we're looking for something that is under the floor, not necessarily
at infinity. Maybe, I don't know.''
\end{quote}
This was an exploratory suggestion, not a theorem or a supplied proof.
Codex developed the displayed calculations, wrote the proofs and executable
checks, and used separate AI agents to review them. Those reviews are
independent calculation routes within an AI workflow; they are not external
human peer review or a formal verification of every theorem. The complete
arguments, domains, constants, and unsuccessful constructions are retained
so that readers can examine the mathematical content directly.

The original retained-fibre, conductor, arithmetic-operator, and GCT
calculations preceded this task's introduction of the public
Navier--Stokes manuscript on 8 September 2026. On that date, another
mathematics task supplied the public manuscript and the pinned source.
The 165-page revision was retained; a subsequent source notification
reported the 166-page revision downloaded at 19:17 UTC, with a server
last-modified header of 19:09:35 GMT. The two hashes, exact source formulas,
selected-passage comparison, and attribution are recorded in
Section~\ref{sec:ns-conic-transport}. The date and HTTP observations record
source introduction and retrieval; they do not certify the external
manuscript's global theorem. The new concentration-coordinate bridge was
derived after that introduction from the explicitly displayed source
formulas and the previously retained conic.

Several unsuccessful approaches are part of the result. The printed
mixed-sign GCT endpoint required a weight correction. The unsaturated
integral presentation retains forty integer-$2$ torsion summands after
inverting the specified quantum integer. Ordinary tensor permutations
fail to preserve the generic source action. The replacement involution
constructed here is equivariant in degree two but has a nonzero degree-three
braid defect. The positive signed-sheet cover recovers the original
generators while its quantum relations have nonzero defects in the
quotient kernel. Each assertion and the exact repairing map or retained
obstruction is proved in the corresponding section; none is erased from
the provenance because an initial construction failed.

The finite negative coefficient and the negative signature of the specified
Schur form are established calculations. Their derivations do not turn
plethysm multiplicities into negative integers. The full source-compatible
canonical geometric construction and the unrestricted asymptotic
complexity obstruction are still unfinished. This preprint presents
proved intermediate maps and the concrete places where the attempted
extensions stop, without asserting a resolution of a Millennium problem.

The following read-only public workbenches were supplied by the human
participant for public audit as places where Codex worked. They expose
the broader exploratory projects, rather than serving as substitutes
for the complete proofs in this manuscript. Labels follow the supplied
identifications; the fourth link was supplied without a project label.
\begin{itemize}
\item \href{https://www.overleaf.com/read/rtmyqxyrzprn\#fa24eb}{S6 Workbench}.
\item \href{https://www.overleaf.com/read/rnwvwkhfhztp\#8f63da}{Erd\H{o}s--Straus Workbench}.
\item \href{https://www.overleaf.com/read/hzthvczhdyxc\#a60fc2}{Navier--Stokes Workbench}.
\item \href{https://www.overleaf.com/read/sybvppjwxmxw\#b17484}{Additional public workbench}.
\end{itemize}
The initial DOI edition is intended to make the mathematics and these
workbenches available for scrutiny. A later publication of the contributing
transcripts, with profanity filtered and the human suggestions preserved
in context, is a separately stated future stage. No such transcript
publication is represented as already completed here.

